\documentclass[11pt]{article}
\usepackage[a4paper,margin=28mm]{geometry}
\usepackage{amsmath,amssymb,amsfonts,mathrsfs,amsthm,booktabs,longtable,hyperref}
\newtheorem{proposition}{Proposition}
\title{Audit of Reverse Connections Against Known Askey-Scheme Limits\\\large Independent audit file, version 111}
\author{}
\date{7 November 2026}
\begin{document}
\maketitle
\noindent\textbf{Source manuscript:} \texttt{Askey\_Finite\_Connections\_Revised\_v31(1).tex}, supplied by the author. The source manuscript has not been edited.

\section*{Scope and audit rules}
For each genuine limiting arrow involving distinct polynomial families, use only the \emph{reverse} connection displayed in the source manuscript to derive the corresponding classical limit. Exclude parameter-change self-connections and exact transformations unless an exact distinct-family specialization is itself an explicitly recognized Koekoek--Swarttouw limiting arrow. Count each ordered family arrow at most once. Record separately any limits not yet audited. A proof is marked complete only if the normalization, the surviving terms, and the vanishing remainder have been established.

\section{Wilson to continuous dual Hahn: complete}
The classical limiting relation in the normalizations of the source manuscript is
\begin{equation}\label{eq:target}
\lim_{d\to+\infty}\frac{W_n(x^2;a,b,c,d)}{(a+d)_n}=S_n(x^2;a,b,c),\qquad n\in\mathbb N_0.
\end{equation}
Here $a,b,c,x,n$ are fixed and parameters are chosen so that the expressions are defined. The source manuscript's reverse connection, equation \texttt{eq:consolidated-w-cdh}, reads, with $s=a+b+c+d$ and $j=n-k$,
\begin{equation}\label{eq:reverse}
W_n(x^2;a,b,c,d)=\sum_{k=0}^n \binom nk (n+s-1)_{n-k}
(j+a+b)_k(j+a+c)_k(1-n-b-c)_k\,S_{n-k}(x^2;a,b,c).
\end{equation}
\begin{proposition}
The reverse connection \eqref{eq:reverse} implies \eqref{eq:target} without using the hypergeometric definition of either family.
\end{proposition}
\begin{proof}
Divide \eqref{eq:reverse} by $(a+d)_n$. For fixed $n$ and $k$ write the coefficient of $S_{n-k}$ as
\[
A_{n,k}(d)=\binom nk\frac{(d+n+a+b+c-1)_{n-k}}{(d+a)_n}
(n-k+a+b)_k(n-k+a+c)_k(1-n-b-c)_k.
\]
The Pochhammer symbol $(d+u)_m$ is a monic polynomial in $d$ of degree $m$ for fixed $u,m$. Hence, for $0\leq k\leq n$,
\[
\frac{(d+n+a+b+c-1)_{n-k}}{(d+a)_n}=d^{-k}\bigl(1+O(d^{-1})\bigr)
\quad(d\to+\infty).
\]
The other factors of $A_{n,k}$ do not depend on $d$. Therefore $A_{n,k}(d)\to0$ for $k\geq1$. For $k=0$, all the factors of length zero are $1$ and $A_{n,0}(d)\to1$. Since the sum has only $n+1$ terms, its limit can be taken termwise, yielding
\[
\lim_{d\to+\infty}\frac{W_n(x^2;a,b,c,d)}{(a+d)_n}
=\sum_{k=0}^n\delta_{k0}S_{n-k}(x^2;a,b,c)=S_n(x^2;a,b,c).
\]
The case $n=0$ is included.\qedhere
\end{proof}
\noindent\textbf{Audit verdict: PASS.} The normalized reverse-connection matrix tends to the identity on every fixed finite degree range. The argument proves the displayed limit coefficientwise in $x^2$.


\section{Meixner--Pollaczek to continuous Hahn: scope assessment}
The next non-self connection in the order of V31, after the exact continuous Hahn--Wilson quadratic transformation and self-connections are excluded, is the Meixner--Pollaczek/continuous Hahn pair. Its reverse connection is equation \texttt{eq:consolidated-ch-mp}:
\begin{equation}\label{eq:ch-mp-audit}
\begin{split}
p_n(x;a,b,c,d)=\sum_{m=0}^n &\frac{i^{n+m}(a+c)_n(a+d)_n}{n!}
\frac{(-n)_m(n+s-1)_m}{2^m(a+c)_m(a+d)_m}\\
&\times {}_3F_2\!\left(\begin{matrix}m-n,\ m+n+s-1,\ m+2a\\m+a+c,\ m+a+d\end{matrix};\frac12\right)P_m^{(a)}(x),\qquad s=a+b+c+d.
\end{split}
\end{equation}
\paragraph{Exact diagonal coefficient and necessary normalization.}
Setting $m=n$ in \eqref{eq:ch-mp-audit}, the terminating hypergeometric factor is $1$ because its first upper parameter is zero. Using $(-n)_n=(-1)^n n!$ and $i^{2n}=(-1)^n$ gives the exact identity
\begin{equation}\label{eq:mp-diagonal}
 [P_n^{(a)}]p_n=\frac{(n+s-1)_n}{2^n}.
\end{equation}
Here $[P_n^{(a)}]p_n$ denotes the coefficient of the degree-$n$ Meixner--Pollaczek polynomial in the reverse connection. Consequently, any fixed-variable limit in which the normalized continuous Hahn polynomial converges to $P_n^{(a)}(x)$ and the other connection coefficients vanish \emph{must} use a normalization asymptotic to $(n+s-1)_n/2^n$ (for each fixed $n$). More precisely, writing $C_{n,m}$ for the coefficients in \eqref{eq:ch-mp-audit}, such a limit is equivalent to
\[
 \frac{C_{n,n}}{N_n}\longrightarrow 1,
 \qquad \frac{C_{n,m}}{N_n}\longrightarrow0\quad(0\le m<n),
\]
for a specified parameter path and $N_n$; the second condition is essential. This calculation fixes the necessary leading-degree scale but does \emph{not} establish the classical limit or identify its KLS parameter path.

\paragraph{Necessary qualification.}
The source fixes the Meixner--Pollaczek angle to $\phi=\pi/2$ in this displayed formula. The familiar classical continuous-Hahn-to-Meixner--Pollaczek limiting arrow requires an explicit parameter regime and its corresponding variable and normalization scaling. No such regime has yet been matched to \eqref{eq:ch-mp-audit} in this audit. Consequently, it would be unjustified to label the limit a PASS solely from triangularity or from degree checks. In particular, an asymptotic argument must show that the \emph{normalized and rescaled} coefficients converge to the required values, not just that a formal finite connection exists.
\paragraph{Verdict.} \textbf{NOT YET VERIFIED.} This entry records the exact V31 reverse connection and the outstanding need to identify and prove the corresponding KLS scaling. It is not a counterexample to the connection formula.


\section{A candidate KLS scaling: symmetric continuous Hahn parameters}
The reverse connection itself suggests a concrete parameter regime:
\[
  (a,b,c,d)=(a,t,a,t),\qquad t\to+\infty.
\]
Here the Meixner--Pollaczek basis in \eqref{eq:ch-mp-audit} has fixed parameter $a$ and angle $\pi/2$. In this regime the diagonal coefficient is
\[
 C_{n,n}(t)=\frac{(n+2a+2t-1)_n}{2^n}=t^n(1+O(t^{-1})).
\]
Thus the natural candidate normalization is $t^n$. This is a candidate path, not yet an identification with a numbered KLS formula.

There is a useful \emph{exact} off-diagonal cancellation. For $n\ge1$, set $m=n-1$ in the source reverse connection. Because one upper parameter of the ${}_3F_2$ is $-1$, the hypergeometric series reduces to two terms. Direct simplification gives
\begin{equation}\label{eq:subdiag}
 \frac{C_{n,n-1}}{C_{n,n}}=
 i\left[\frac{2(n-1+a+c)(n-1+a+d)}{2n+a+b+c+d-2}-(n-1+2a)\right].
\end{equation}
For $c=a$ and $b=d=t$, the fraction inside the brackets is exactly $n-1+2a$, so
\[
 C_{n,n-1}(a,t,a,t)=0\qquad\text{for every }n\ge1\text{ and every admissible }t.
\]
In particular the first subdiagonal satisfies the necessary vanishing condition after normalization by $t^n$.

\paragraph{Completion of the reverse-connection limit.}
In fact, the candidate symmetric path admits an exact evaluation of \emph{every} off-diagonal coefficient. Fix $n$ and put $r=n-m$. Under $(b,c,d)=(t,a,t)$, the terminating ${}_3F_2$ in \eqref{eq:ch-mp-audit} reduces, by cancellation of the equal upper and lower parameters $m+2a$, to
\[
 F_r(D)={}_2F_1\!\left(\begin{matrix}-r,\ 2D+r-1\\D\end{matrix};\frac12\right),\qquad D=m+a+t.
\]
The following terminating hypergeometric evaluation holds:
\begin{equation}\label{eq:kummer}
 F_{2\ell+1}(D)=0,\qquad
 F_{2\ell}(D)=\frac{(-1)^\ell(\tfrac12)_\ell}{(D)_\ell}.
\end{equation}
For completeness, the Jacobi-polynomial representation is
\[
 {}_2F_1(-r,2D+r-1;D;(1-z)/2)
 =\frac{r!}{(D)_r}P_r^{(D-1,D-1)}(z).
\]
The Jacobi polynomial on the right has parity $(-1)^r$, so its value at $z=0$ vanishes for odd $r$. For even $r=2\ell$, the standard Jacobi--Gegenbauer conversion and the generating function of Gegenbauer polynomials give the second value in \eqref{eq:kummer}. This also provides a proof by polynomial identities, with exceptional parameters handled by continuation.

Substitution into the exact reverse coefficient gives
\begin{equation}\label{eq:exact-C}
 C_{n,m}(t)=\frac{i^{n+m}(2a)_n(a+t)_n}{n!}
 \frac{(-n)_m(n+2a+2t-1)_m}{2^m(2a)_m(a+t)_m}
 \begin{cases}
 0,&r\text{ odd},\\[2pt]
 \displaystyle\frac{(-1)^\ell(\tfrac12)_\ell}{(m+a+t)_\ell},&r=2\ell.
 \end{cases}
\end{equation}
Here apparent singularities at exceptional parameters are understood by continuation. For fixed $n,m$, the prefactor before the case distinction is $O(t^n)$: the ratio $(a+t)_n/(a+t)_m$ has degree $n-m$ in $t$, while $(n+2a+2t-1)_m$ has degree $m$. For $r=2\ell\ge2$, the last factor is $O(t^{-\ell})$. Therefore
\[
 t^{-n}C_{n,m}(t)=0\quad(r\text{ odd}),\qquad
 t^{-n}C_{n,m}(t)=O(t^{-\ell})\to0\quad(r=2\ell\ge2).
\]
For $m=n$, \eqref{eq:mp-diagonal} yields $t^{-n}C_{n,n}(t)\to1$. The finite reverse connection thus proves
\begin{equation}\label{eq:ch-mp-limit}
 \boxed{\displaystyle\lim_{t\to+\infty}t^{-n}p_n(x;a,t,a,t)=P_n^{(a)}(x)}
\end{equation}
for every fixed $n$, coefficientwise in $x$, with the Meixner--Pollaczek angle $\phi=\pi/2$ and the precise polynomial normalization used in V31. No forward connection or independent polynomial identity was used.

\paragraph{Verdict.} \textbf{PASS for the displayed symmetric-parameter limiting relation.} The proof is complete for all degrees. The identification of this particular parameterization with a specifically numbered KLS limiting formula remains to be checked against the KLS source; this audit does not assert a KLS equation number.


\section{Wilson to Jacobi: complete reverse-connection proof}
\paragraph{A precise parameter path.}
Let $A,B$ be fixed with $A,B>-1$, put
\[
 a=A,\qquad b=1,\qquad c=t,\qquad d=-t+B+1,
 \qquad y=-t^2z,\qquad t\longrightarrow+\infty.
\]
The parameter path is understood algebraically by analytic continuation:
for large real $t$ the Wilson parameters are not all positive. This is
legitimate for the terminating polynomial identities. The target limit is
\begin{equation}\label{eq:w-j-target}
 \boxed{\lim_{t\to\infty}\frac{W_n(-t^2z;A,1,t,-t+B+1)}{(-t^2)^n n!}
 =P_n^{(A,B)}(1-2z).}
\end{equation}
The particular path fixes $a+b=A+1$ and $a+b+c+d=A+B+2$, exactly
matching the two upper/lower parameters of the Jacobi polynomial.

\paragraph{The V31 reverse connection.}
Write $J_m(y)=P_m^{(A,B)}(1-2y)$, and let $h_m$ be the Jacobi squared norm
in V31. Its equation \texttt{eq:wilson-to-jac} gives
\begin{equation}\label{eq:w-j-reverse}
 W_n(y)=\sum_{m=0}^n D_{n,m}(t)J_m(y),\quad
 D_{n,m}(t)=\frac{(-1)^m}{h_m}\sum_{j=m}^n L_{n,j}(t)
 \sum_{r=m}^j E_{j,r}(A)\binom rm B(A+r+1,B+m+1),
\end{equation}
where $E_{j,r}(A)$ is independent of $t$ and
\[
 L_{n,j}(t)=(-1)^j\binom nj(n+A+B+1)_j
 (A+1+j)_{n-j}(A+t+j)_{n-j}
 (A-t+B+1+j)_{n-j}.
\]
In particular, for each fixed $j$,
\begin{equation}\label{eq:L-asympt}
 \frac{L_{n,j}(t)}{(-t^2)^{n-j}}
 \longrightarrow (-1)^j\binom nj(n+A+B+1)_j(A+1+j)_{n-j}.
\end{equation}
Since $j\ge m$, division by $(-t^2)^{n-m}$ kills every $j>m$ term
in the finite reverse coefficient. For $j=m$, the inner sum has only
$r=m$, and $E_{m,m}(A)=1$. Therefore
\begin{equation}\label{eq:D-asympt}
 \lim_{t\to\infty}\frac{D_{n,m}(t)}{(-t^2)^{n-m}}
 =\binom nm(n+A+B+1)_m(A+1+m)_{n-m}
 \frac{B(A+m+1,B+m+1)}{h_m}.
\end{equation}

\paragraph{The scaled Jacobi basis and the resulting polynomial.}
The leading coefficient of $J_m(y)$ as a polynomial in $y$ is
$(-1)^m(m+A+B+1)_m/m!$. Consequently
\[
 \lim_{t\to\infty}\frac{J_m(-t^2z)}{(-t^2)^m}
 =\frac{(-1)^m(m+A+B+1)_m}{m!}z^m.
\]
Using the norm displayed in V31 and the beta--gamma identity gives
\[
 \frac{B(A+m+1,B+m+1)}{h_m}
 =\frac{m!}{(m+A+B+1)_m}.
\]
Substitute these two limits and \eqref{eq:D-asympt} into the finite
reverse sum \eqref{eq:w-j-reverse}. The factors cancel and yield
\begin{align*}
 \lim_{t\to\infty}\frac{W_n(-t^2z)}{(-t^2)^n n!}
 &=\frac1{n!}\sum_{m=0}^n(-1)^m\binom nm
 (n+A+B+1)_m(A+1+m)_{n-m}z^m\\
 &=\frac{(A+1)_n}{n!}\sum_{m=0}^n
 \frac{(-n)_m(n+A+B+1)_m}{(A+1)_m m!}z^m\\
 &=P_n^{(A,B)}(1-2z).
\end{align*}
The last equality is the standard terminating ${}_2F_1$ representation
of the Jacobi polynomial. This is an identification of the polynomial
\emph{after} the reverse-connection limit, not an invocation of the
Wilson hypergeometric definition or a forward connection.
All limits are coefficientwise in $z$ because the sums have fixed finite
length. The result extends to generic admissible $A,B$ by rational or
meromorphic continuation.

\paragraph{Verdict.} \textbf{PASS for the explicitly stated Wilson-to-Jacobi
parameter path}, for every fixed degree $n$. The proof uses the V31
reverse connection and correctly accounts for the scaled Jacobi basis.
The identification with a particular numbered KLS formula should still
be checked against the printed KLS reference; no equation number is
asserted here.


\section{Jacobi to generalized Laguerre: complete reverse-connection limit}
\noindent\textbf{Source:} V31, equation \texttt{eq:v14-jac-lag} (Jacobi-to-Laguerre reverse connection), with the monomial-to-Laguerre identity given in its proof. This is a genuine limiting arrow, not an exact quadratic transformation.

\begin{proposition}
For every fixed $n\ge0$, $\alpha>-1$, and fixed $z$, the reverse connection implies
\begin{equation}\label{eq:jac-lag-limit}
\lim_{\beta\to+\infty}P_n^{(\alpha,\beta)}\!\left(1-\frac{2z}{\beta}\right)=L_n^{(\alpha)}(z).
\end{equation}
The convergence is coefficientwise in $z$.
\end{proposition}
\begin{proof}
The V31 reverse connection is obtained by expanding the Jacobi polynomial in monomials and applying its exact monomial-to-Laguerre transformation
\begin{equation}\label{eq:mon-lag}
z^r=r!\sum_{m=0}^r(-1)^m\binom{r+\alpha}{r-m}L_m^{(\alpha)}(z).
\end{equation}
For the scaled argument of \eqref{eq:jac-lag-limit}, the same reverse-connection construction (equivalently, substitute the affine-scaled Jacobi monomial coefficients into \eqref{eq:mon-lag}) yields
\begin{equation}\label{eq:scaled-jac-lag}
P_n^{(\alpha,\beta)}\!\left(1-\frac{2z}{\beta}\right)
=\sum_{m=0}^n C_{n,m}(\beta)L_m^{(\alpha)}(z),
\end{equation}
where the \emph{finite reverse coefficients} are
\begin{equation}\label{eq:scaled-jac-lag-coeff}
C_{n,m}(\beta)=(-1)^m\frac{(\alpha+1)_n}{n!}
\sum_{r=m}^n\frac{(-n)_r(n+\alpha+\beta+1)_r}
{(\alpha+1)_r\,\beta^r}\binom{r+\alpha}{r-m}.
\end{equation}
Indeed, the terminating Jacobi monomial expansion in V31 is
\[
P_n^{(\alpha,\beta)}\!\left(1-\frac{2z}{\beta}\right)
=\frac{(\alpha+1)_n}{n!}\sum_{r=0}^n
\frac{(-n)_r(n+\alpha+\beta+1)_r}{(\alpha+1)_r r!}
\frac{z^r}{\beta^r},
\]
and inserting \eqref{eq:mon-lag} gives \eqref{eq:scaled-jac-lag-coeff}.
For each fixed $r$, $(n+\alpha+\beta+1)_r/\beta^r\to1$. Thus, because the sum is finite,
\begin{align*}
\lim_{\beta\to\infty}C_{n,m}(\beta)
&=(-1)^m\frac{(\alpha+1)_n}{n!}
\sum_{r=m}^n\frac{(-n)_r}{(\alpha+1)_r}
\binom{r+\alpha}{r-m}\\
&=\frac{(-1)^m(\alpha+1)_n}{n!(\alpha+1)_m}
(-n)_m\sum_{j=0}^{n-m}\frac{(-(n-m))_j}{j!}.
\end{align*}
The final sum equals $(1-1)^{n-m}$ by the binomial theorem: it is zero for $m<n$ and one for $m=n$. For $m=n$, the prefactor equals
$(-1)^n(-n)_n/n!=1$. Therefore
\[
\lim_{\beta\to\infty}C_{n,m}(\beta)=\delta_{nm}.
\]
Taking the termwise limit in \eqref{eq:scaled-jac-lag} proves \eqref{eq:jac-lag-limit}. Notice that the cancellation occurs \emph{after summing} the reverse coefficients: individual monomial contributions need not vanish.
\end{proof}
\noindent\textbf{Audit verdict: PASS.} This proof uses the reverse monomial-to-Laguerre construction explicitly documented in V31 and establishes the classical Jacobi-to-Laguerre limit for all fixed degrees. No forward polynomial connection is used.


\section{Meixner to generalized Laguerre: complete reverse-connection limit}
\noindent\textbf{Source:} V31, equation \texttt{eq:meixner-to-laguerre-v61}, including its signed-Stirling-number construction. This is a genuine distinct-family limit. Put $\beta=\alpha+1$, assume $\beta$ avoids the poles of the terminating expressions, and let $c\to1$ (for example through $0<c<1$). The classical limit in the conventions of V31 is
\begin{equation}\label{eq:meix-lag-target}
\lim_{c\to1}M_n\!\left(\frac{z}{1-c};\beta,c\right)
=\frac{n!}{(\beta)_n}L_n^{(\beta-1)}(z).
\end{equation}
\paragraph{Proof directly from the reverse-connection construction.}
The reverse formula in V31 expands the Meixner polynomial first in signed Stirling numbers, then converts monomials to the Laguerre basis. In particular its exact intermediate expansion (which is the finite polynomial identity underlying the displayed reverse coefficients) is
\begin{equation}\label{eq:stirling-start}
M_n(x;\beta,c)=\sum_{r=0}^n
\frac{(-n)_r(c^{-1}-1)^r}{(\beta)_r r!}
\sum_{k=0}^r s(r,k)x^k.
\end{equation}
This is the reverse connection \emph{before} the final change from monomials to Laguerre polynomials; thus no forward connection is used. Set $\delta=1-c$, $x=z/\delta$. Since
\[
c^{-1}-1=\frac{\delta}{1-\delta},\qquad
\left(c^{-1}-1\right)^r\left(\frac z\delta\right)^k
=\frac{\delta^{r-k}}{(1-\delta)^r}z^k,
\]
for every fixed $r$ the $k<r$ terms vanish as $\delta\to0$, whereas the $k=r$ term tends to $z^r$ because $s(r,r)=1$. Hence, by finiteness of both sums,
\begin{equation}\label{eq:meix-limit-monomial}
\lim_{c\to1}M_n\!\left(\frac z{1-c};\beta,c\right)
=\sum_{r=0}^n\frac{(-n)_r}{(\beta)_r r!}z^r
={}_1F_1(-n;\beta;z).
\end{equation}
Finally the generalized Laguerre normalization used in V31 is
\[
L_n^{(\beta-1)}(z)=\frac{(\beta)_n}{n!}{}_1F_1(-n;\beta;z),
\]
which gives \eqref{eq:meix-lag-target}. The convergence is coefficientwise and holds for every fixed $n$. Equivalently, after rescaling, the reverse Stirling coefficients with $k<r$ vanish and only the diagonal $s(r,r)$ survives. \hfill$\square$

\noindent\textbf{Audit verdict: PASS.} The standard Meixner-to-Laguerre limit follows from the exact finite expansion used to derive V31's \emph{reverse} connection, including the essential variable rescaling. No forward connection or exact quadratic transformation is invoked.


\section{Meixner--Pollaczek to Hermite: reverse-connection limit}
The V31 reverse connection, equation \texttt{eq:cat-mp-to-hermite-bell}, is
\begin{equation}\label{eq:mp-hermite-reverse-audit}
P_n^{(\lambda)}(x;\pi/2)=\sum_{k=0}^{\lfloor n/2\rfloor}
\frac{Y_k(1!b_1(n-2k,\lambda),\ldots,k!b_k(n-2k,\lambda))}
{k!(n-2k)!}\,H_{n-2k}(x),
\end{equation}
where, for $m\ge0$,
\[
\sum_{r\ge1}b_r(m,\lambda)u^r
=\log\left[(1+u)^{-\lambda}e^{(\arctan\sqrt u)^2}
\left(\frac{\arctan\sqrt u}{\sqrt u}\right)^m\right].
\]
The target is the classical limit, already stated in V31 as
\begin{equation}\label{eq:mp-hermite-target-audit}
\lim_{\lambda\to+\infty}\lambda^{-n/2}
P_n^{(\lambda)}(\sqrt\lambda\,x;\pi/2)=\frac{H_n(x)}{n!}.
\end{equation}
\begin{proposition}
For every fixed $n\ge0$, the reverse connection \eqref{eq:mp-hermite-reverse-audit}
implies \eqref{eq:mp-hermite-target-audit} coefficientwise in $x$.
\end{proposition}
\begin{proof}
Write $B_{m,k}(\lambda)=[u^k]\exp(\sum_{r\ge1}b_r(m,\lambda)u^r)$.
By the defining Bell-polynomial identity, the reverse connection is exactly
\[
P_n^{(\lambda)}(y;\pi/2)=
\sum_{m+2k=n}\frac{H_m(y)}{m!}B_{m,k}(\lambda).
\]
Sum these finite coefficients over $n$ as a formal power series in $z$.
Using the \emph{reverse-connection coefficients themselves}, we obtain
\begin{align*}
\sum_{n\ge0}P_n^{(\lambda)}(y;\pi/2)z^n
&=\sum_{m\ge0}\frac{H_m(y)}{m!}z^m
\sum_{k\ge0}B_{m,k}(\lambda)z^{2k}\\
&=\sum_{m\ge0}\frac{H_m(y)}{m!}
(1+z^2)^{-\lambda}e^{(\arctan z)^2}(\arctan z)^m\\
&=(1+z^2)^{-\lambda}e^{2y\arctan z}.
\end{align*}
The last equality uses only the Hermite generating series
$\sum_mH_m(y)t^m/m!=e^{2yt-t^2}$; the $t^2$ factors cancel.
Now substitute $y=\sqrt\lambda\,x$ and $z=t/\sqrt\lambda$.
The left-hand side is precisely the generating series of the
\emph{normalized reverse-connection sums}:
\[
\sum_{n\ge0}\lambda^{-n/2}P_n^{(\lambda)}(\sqrt\lambda\,x;\pi/2)t^n
=\left(1+\frac{t^2}{\lambda}\right)^{-\lambda}
\exp\left(2\sqrt\lambda\,x\arctan\frac{t}{\sqrt\lambda}\right).
\]
For each fixed coefficient degree in $t$, the two factors tend to
$e^{-t^2}$ and $e^{2xt}$ respectively: indeed
$-\lambda\log(1+t^2/\lambda)=-t^2+O(\lambda^{-1})$
and $\sqrt\lambda\arctan(t/\sqrt\lambda)=t+O(\lambda^{-1})$
as formal series truncated at any prescribed degree.
Thus the coefficientwise limit of the normalized reverse sums is
$e^{2xt-t^2}=\sum_{n\ge0}H_n(x)t^n/n!$, proving the claim.
\end{proof}
\noindent\textbf{Audit verdict: PASS.} This is an all-degree proof using the
Bell coefficients of V31's reverse connection. The individual Hermite
terms need not have limits separately after rescaling; the finite sum
is handled by exact coefficient resummation before taking the limit.
The KLS equation number has not been independently checked.


\section{Charlier to Hermite: reverse-connection limit}
The V31 reverse connection \texttt{eq:v16-c-h} is
\begin{equation}\label{eq:charlier-reverse-audit}
 \mathrm C_n(x;a)=\sum_{m=0}^n\left[\sum_{r=m}^n
 \frac{(-n)_r}{a^r r!}\sum_{j=0}^{\lfloor(r-m)/2\rfloor}
 \frac{s(r,m+2j)(m+2j)!}{2^{m+2j}j!m!}\right]H_m(x),
\end{equation}
where $s(r,k)$ denotes the signed Stirling numbers of the first kind and
$H_m$ denotes the physicists' Hermite polynomial.  The classical
Charlier--Hermite limit in this normalization is
\begin{equation}\label{eq:charlier-hermite-limit}
 \lim_{a\to+\infty}\left(-\sqrt{\frac a2}\right)^n
 \mathrm C_n\bigl(a+\sqrt{2a}\,z;a\bigr)=2^{-n}H_n(z).
\end{equation}
\begin{proposition}
For each fixed nonnegative integer $n$, the reverse connection
\eqref{eq:charlier-reverse-audit} implies \eqref{eq:charlier-hermite-limit}.
\end{proposition}
\begin{proof}
The Stirling-number identity $x^{\underline r}=\sum_{k=0}^r
s(r,k)x^k$ and the monomial--Hermite identity
\[
 x^k=2^{-k}\sum_{j=0}^{\lfloor k/2\rfloor}
 \frac{k!}{j!(k-2j)!}H_{k-2j}(x)
\]
show, on reindexing $k=m+2j$, that the inner double sum of
\eqref{eq:charlier-reverse-audit} is exactly $x^{\underline r}$.
Thus the \emph{reverse connection itself}, without assuming a
Charlier--Hermite limit, gives
\[
 \mathrm C_n(x;a)=\sum_{r=0}^n\frac{(-n)_r}{a^r r!}
 x^{\underline r}.
\]
This finite identity implies the exponential generating function
\begin{equation}\label{eq:charlier-egf-audit}
 \sum_{n\ge0}\mathrm C_n(x;a)\frac{t^n}{n!}
 =e^t(1-t/a)^x,
\end{equation}
initially as a formal series in $t$. Indeed,
$(-n)_r=(-1)^r n!/(n-r)!$, and summing first over $n-r$
produces $e^t$; the remaining sum is the generalized binomial series.
Set $x=a+\sqrt{2a}\,z$ and $t=-\sqrt{a/2}\,u$. The left side of
\eqref{eq:charlier-egf-audit} becomes the exponential generating
function of the normalized polynomials in \eqref{eq:charlier-hermite-limit}.
The right side is
\[
 \exp\!\left[-\sqrt{a/2}\,u+
 (a+\sqrt{2a}\,z)\log\left(1+\frac{u}{\sqrt{2a}}\right)\right].
\]
For fixed $z$ and $u$ near zero, Taylor expansion of the logarithm
shows that the exponent tends to $zu-u^2/4$; the terms of order
$\sqrt a$ cancel, and the remainder is $O(a^{-1/2})$ uniformly on
compact $u$-sets. Hence the generating functions converge locally
uniformly to
\[
 e^{zu-u^2/4}=\sum_{n\ge0}\frac{H_n(z)}{2^n}\frac{u^n}{n!},
\]
by the defining Hermite generating function
$e^{2zv-v^2}=\sum_n H_n(z)v^n/n!$ with $v=u/2$.
Local uniform convergence permits coefficient extraction, proving
\eqref{eq:charlier-hermite-limit} for every fixed $n$.
\end{proof}
\noindent\textbf{Audit verdict: PASS.} This proof uses the reverse
connection in V31 and performs the necessary resummation before
applying the affine variable scaling. It does not incorrectly take
limits termwise in the unscaled Hermite basis. The precise KLS
formula number has not been independently checked.


\section{Meixner to Charlier: reverse-connection limit}
The reverse connection in V31, equation \texttt{eq:meix-char-reverse}, is
\begin{equation}\label{eq:mc-reverse}
\mathrm M_n(x;\beta,c)=\sum_{m=0}^n
\frac{(-n)_m(az)^m}{(\beta)_m m!}
{}_1F_1\!\left(\begin{matrix}m-n\\\beta+m\end{matrix};-az\right)
\mathrm C_m(x;a),\qquad z=1-c^{-1}.
\end{equation}
For fixed $a\ne0$, take $\beta\to+\infty$ and
$c=c_\beta=a/(a+\beta)$, so $z=-\beta/a$. The classical
Meixner-to-Charlier limit in these normalizations is
\begin{equation}\label{eq:mc-target}
\lim_{\beta\to+\infty}\mathrm M_n\!\left(x;\beta,\frac{a}{a+\beta}\right)
=\mathrm C_n(x;a).
\end{equation}
\begin{proposition}
For every fixed $n\ge0$, the V31 reverse connection
\eqref{eq:mc-reverse} implies \eqref{eq:mc-target}.
\end{proposition}
\begin{proof}
After the substitution $az=-\beta$, the coefficient of
$\mathrm C_m(x;a)$ is
\[
A_{n,m}(\beta)=\frac{(-n)_m(-\beta)^m}{(\beta)_m m!}
\sum_{j=0}^{n-m}\frac{(m-n)_j}{(\beta+m)_j j!}\,\beta^j.
\]
The hypergeometric series is finite, and for each fixed $m,j$,
$\beta^j/(\beta+m)_j\to1$. Hence the terminating sum tends to
\[
\sum_{j=0}^{n-m}\frac{(-(n-m))_j}{j!}
=\sum_{j=0}^{n-m}(-1)^j\binom{n-m}{j}
=(1-1)^{n-m},
\]
which is $0$ when $m<n$ and $1$ when $m=n$.
The prefactor tends to $(-n)_m(-1)^m/m!=\binom nm$.
Consequently $A_{n,m}(\beta)\to\delta_{nm}$; for $m=n$ the
limit is exactly $1$. The finite reverse-connection sum therefore
converges to $\mathrm C_n(x;a)$, proving the result for every fixed
$n$, coefficientwise as a polynomial in $x$.
\end{proof}
\noindent\textbf{Audit verdict: PASS.} All off-diagonal coefficients
vanish by the terminating binomial identity, and the diagonal tends
to one. The proof uses the V31 reverse connection directly, without
assuming the limiting relation.


\section{Krawtchouk to Charlier: complete reverse-connection limit}
The classical Krawtchouk-to-Charlier limit, with the polynomial
normalizations used in V31, is
\begin{equation}\label{eq:kraw-char-limit}
 \lim_{N\to\infty}\mathrm K_n\!\left(x;\frac aN,N\right)
 =\mathrm C_n(x;a),\qquad n\in\mathbb N_0,
\end{equation}
where $a\ne0$, $x$ is fixed, and $N$ runs through integers $N\ge n$.
The reverse connection in V31, equation \texttt{eq:kraw-char-reverse}, reads
\begin{equation}\label{eq:audit-kraw-char}
\mathrm K_n(x;p,N)=\sum_{m=0}^n A_{n,m}(p,N)\,\mathrm C_m(x;a),
\end{equation}
with
\[
 A_{n,m}(p,N)=\frac{(-n)_m}{(-N)_m}\frac{(a/p)^m}{m!}
 {}_1F_1\!\left(\begin{matrix}m-n\\m-N\end{matrix};-\frac ap\right).
\]
\begin{proposition}
For every fixed $n$ and $0\le m\le n$,
\[
 \lim_{N\to\infty} A_{n,m}(a/N,N)=\delta_{nm}.
\]
Consequently, \eqref{eq:kraw-char-limit} follows directly from
\eqref{eq:audit-kraw-char}.
\end{proposition}
\begin{proof}
Put $r=n-m$. Since the first upper hypergeometric parameter is
$-r$, the series terminates, and substitution $p=a/N$ gives
\[
 A_{n,m}(a/N,N)=
 \frac{(-n)_m}{m!}\frac{N^m}{(-N)_m}
 \sum_{j=0}^{r}\frac{(-r)_j}{(m-N)_j}\frac{(-N)^j}{j!}.
\]
For fixed $m,j$ as $N\to\infty$,
\[
 \frac{N^m}{(-N)_m}\longrightarrow(-1)^m,
 \qquad \frac{(-N)^j}{(m-N)_j}\longrightarrow1.
\]
Because the sum is finite, its limit is
\[
 \sum_{j=0}^{r}\frac{(-r)_j}{j!}
 =\sum_{j=0}^{r}(-1)^j\binom rj
 =\begin{cases}1,&r=0,\\0,&r\ge1.\end{cases}
\]
For $m=n$, the remaining prefactor tends to
$(-n)_n(-1)^n/n!=1$. For $m<n$, the terminating sum
vanishes in the limit, and its prefactor remains bounded.
Thus $A_{n,m}\to\delta_{nm}$. The connection
\eqref{eq:audit-kraw-char} contains only $n+1$ terms and its
Charlier basis is independent of $N$, so passage to the limit
proves \eqref{eq:kraw-char-limit} coefficientwise in $x$.
\end{proof}
\noindent\textbf{Audit verdict: PASS.} The diagonal tends to one;
all off-diagonal terms vanish by a terminating binomial identity.
No forward connection or defining-series limit is used.


\section{Hahn to Krawtchouk: complete reverse-connection limit}
For fixed $N\in\mathbb N_0$, $0\le n\le N$, and $0<p<1$, the classical Hahn-to-Krawtchouk limit is
\begin{equation}\label{eq:hahn-kraw-limit}
\lim_{t\to+\infty}H_n^{\mathrm{Hahn}}(x;pt,(1-p)t,N)=\mathrm K_n(x;p,N).
\end{equation}
Here $x$ is fixed (indeed the convergence is coefficientwise as a polynomial in $x$). We use the reverse connection in V31, equation \texttt{eq:hahn-kraw-reverse}:
\begin{equation}\label{eq:hahn-kraw-audit}
H_n^{\mathrm{Hahn}}(x;u,v,N)=\sum_{m=0}^n A_{n,m}(u,v,p)\,\mathrm K_m(x;p,N),
\end{equation}
where, putting $k=n-m$,
\begin{equation}\label{eq:hahn-kraw-coeff}
A_{n,m}=\frac{(-1)^m(-n)_m(n+u+v+1)_m p^m}{m!(u+1)_m}
{}_2F_1\!\left(\begin{matrix}m-n,\ n+u+v+m+1\\u+m+1\end{matrix};p\right).
\end{equation}
\begin{proposition}
Under $u=pt$, $v=(1-p)t$, the reverse-connection coefficients satisfy
$\lim_{t\to+\infty}A_{n,m}(pt,(1-p)t,p)=\delta_{nm}$ for every $0\le m\le n$.
Consequently \eqref{eq:hahn-kraw-limit} follows from the reverse connection.
\end{proposition}
\begin{proof}
For fixed $n,m$, the Pochhammer ratio in \eqref{eq:hahn-kraw-coeff} obeys
\[
\frac{(n+t+1)_m p^m}{(pt+1)_m}\longrightarrow1.
\]
The terminating hypergeometric polynomial has $k=n-m$ terms beyond its constant term. For each $0\le j\le k$,
\[
\frac{(n+t+m+1)_j}{(pt+m+1)_j}p^j\longrightarrow1.
\]
Hence, term by term in this \emph{finite} sum,
\[
{}_2F_1\!\left(\begin{matrix}-k,\ n+t+m+1\\pt+m+1\end{matrix};p\right)
\longrightarrow \sum_{j=0}^k\frac{(-k)_j}{j!}
=(1-1)^k=\begin{cases}1,&k=0,\\0,&k>0.\end{cases}
\]
The remaining factor $(-1)^m(-n)_m/m!=\binom nm$ is independent of $t$. It equals $1$ when $m=n$, and multiplication by the above limit gives $\delta_{nm}$ for all $m$. Since the sum in \eqref{eq:hahn-kraw-audit} is finite and the Krawtchouk basis does not depend on $t$, taking the limit proves \eqref{eq:hahn-kraw-limit}.
\end{proof}
\noindent\textbf{Audit verdict: PASS.} All degrees $0\le n\le N$; diagonal coefficient tends to $1$ and every off-diagonal coefficient tends to $0$. The proof uses the V31 reverse connection alone, not the defining hypergeometric limit.


\section{Hahn to Meixner: complete reverse-connection limit}
Fix $n\geq0$, $\beta>0$, and $0<c<1$. Let $N\to\infty$ through
integers with $N\geq n$, and specialize the Hahn parameters to
\[
 u=\beta-1,\qquad v=N(c^{-1}-1).
\]
The classical limiting relation in the normalization of V31 is
\begin{equation}\label{eq:hahn-meixner-limit}
 \lim_{N\to\infty}\mathrm{Ha}_n(x;\beta-1,N(c^{-1}-1),N)
 =\mathrm M_n(x;\beta,c).
\end{equation}
We establish this from V31's \emph{reverse} connection
\texttt{eq:hahn-meixner-reverse}, namely
\begin{equation}\label{eq:hahn-meixner-audit}
\mathrm{Ha}_n(x;u,v,N)=\sum_{m=0}^n A_{n,m}(N)\,\mathrm M_m(x;\beta,c),
\end{equation}
where, writing $z=1-c^{-1}$,
\begin{align*}
 A_{n,m}(N)={}&
 \frac{(-1)^m(\beta)_m(-n)_m(n+u+v+1)_m}
 {m!(u+1)_m(-N)_m z^m}\\[-2pt]
 &\quad\times{}_3F_2\!\left(\begin{matrix}
 m-n,\ n+u+v+m+1,\ \beta+m\\
 u+m+1,\ m-N
 \end{matrix};z^{-1}\right).
\end{align*}
\begin{proposition}
Under the stated specialization, $A_{n,m}(N)\to\delta_{nm}$ for
$0\leq m\leq n$, and hence \eqref{eq:hahn-meixner-limit} follows
from the reverse connection alone.
\end{proposition}
\begin{proof}
Since $u+1=\beta$, the factors $(\beta)_m/(u+1)_m$ cancel.
For each fixed nonnegative integer $j$,
\[
 \frac{(n+u+v+m+1)_j}{(m-N)_j}
 \longrightarrow [-(c^{-1}-1)]^j=z^j.
\]
More explicitly, the ratio tends to $z^j$ because the numerator
is asymptotic to $[N(c^{-1}-1)]^j$ and the denominator to $(-N)^j$.
The terminating ${}_3F_2$ therefore has the limit
\begin{align*}
 &\lim_{N\to\infty}{}_3F_2\!\left(\begin{matrix}
 m-n,\ n+u+v+m+1,\ \beta+m\\
 \beta+m,\ m-N
 \end{matrix};z^{-1}\right)\\
 &\hspace{2cm}=\sum_{j=0}^{n-m}\frac{(m-n)_j}{j!}
 = (1-1)^{n-m}=\delta_{nm}.
\end{align*}
The series is finite, so termwise passage is justified. Meanwhile,
\[
 \frac{(n+u+v+1)_m}{(-N)_m z^m}\longrightarrow1.
\]
Thus $A_{n,m}(N)\to(-1)^m(-n)_m\delta_{nm}/m!$.
For $m=n$, this factor equals $1$; for $m<n$, the delta is zero.
The Meixner basis is independent of $N$, and the connection sum is
finite. Passing to the limit in \eqref{eq:hahn-meixner-audit}
proves \eqref{eq:hahn-meixner-limit}.
\end{proof}
\noindent\textbf{Audit verdict: PASS.} Every fixed degree $n$;
all normalized off-diagonal coefficients vanish and the diagonal
coefficient tends to $1$. The proof uses V31's Hahn-to-Meixner
reverse connection directly.


\section{Krawtchouk to Hermite: reverse-connection limit}
The reverse connection in V31, equation \texttt{eq:v16-k-h}, expands
$\mathrm K_n(x;p,N)$ in physicists' Hermite polynomials.  The
classical central-limit scaling, for fixed $0<p<1$ and fixed $n$, is
\begin{equation}\label{eq:kh-target}
\lim_{N\to\infty}
\left(-\sqrt{\frac{Np}{2(1-p)}}\right)^n
\mathrm K_n\bigl(Np+\sqrt{2Np(1-p)}\,z;p,N\bigr)
=2^{-n}H_n(z).
\end{equation}
The argument below resums the \emph{reverse} coefficient formula before
passing to the limit; it does not take a limit of the defining
hypergeometric expression.

\begin{proof}
Write $x^{\underline r}=\sum_{k=0}^r s(r,k)x^k$ and
\[
x^k=2^{-k}\sum_{j=0}^{\lfloor k/2\rfloor}
\frac{k!}{j!(k-2j)!}H_{k-2j}(x).
\]
These are exactly the two triangular changes of basis appearing in the
nested coefficients of V31's reverse connection \texttt{eq:v16-k-h}.
Undoing these changes of basis (which are finite for each $n$) gives
\begin{equation}\label{eq:kh-falling}
\mathrm K_n(x;p,N)=\sum_{r=0}^n
\frac{(-n)_r(-1)^r}{(-N)_r\,r!\,p^r}x^{\underline r}.
\end{equation}
Multiply by $(-N)_n(pt)^n/n!$, sum over $0\le n\le N$,
and interchange the finite sums. The resulting binomial sums yield
\begin{equation}\label{eq:kh-gen}
\sum_{n=0}^N\frac{(-N)_n}{n!}(pt)^n
\mathrm K_n(x;p,N)
=(1+(p-1)t)^x(1+pt)^{N-x}.
\end{equation}
For completeness, this identity also follows directly from
\eqref{eq:kh-falling}: its right side is
$(1+pt)^N\sum_{r\ge0}x^{\underline r}
[-t/(1+pt)]^r/r!$; coefficient comparison gives the left
side. Thus \eqref{eq:kh-gen} is an exact resummation of the reverse
connection, not an imported limiting formula.

Put $\sigma_N=\sqrt{2Np(1-p)}$,
$x=Np+\sigma_N z$, and $t=u/\sigma_N$.
For each fixed $u,z$, the logarithm of the right side of
\eqref{eq:kh-gen} is
\[
 x\log(1-(1-p)u/\sigma_N)
 +(N-x)\log(1+pu/\sigma_N)
=-zu-\frac{u^2}{4}+O(N^{-1/2}).
\]
The convergence is locally uniform in $u$, so Taylor coefficients
converge.  The physicists' Hermite generating function gives
\[
 e^{-zu-u^2/4}=\sum_{n\ge0}
\frac{(-1)^nH_n(z)}{2^n n!}u^n.
\]
On the left of \eqref{eq:kh-gen}, the coefficient of $u^n$ is
\[
\frac{(-N)_n}{n!}\left(\frac p{\sigma_N}\right)^n
\mathrm K_n(Np+\sigma_Nz;p,N).
\]
Since $(-N)_n\sim(-1)^nN^n$, coefficient comparison gives
\eqref{eq:kh-target} for every fixed $n$.
\end{proof}
\noindent\textbf{Audit verdict: PASS.} The full reverse connection is
resummed exactly before the central-limit scaling; the coefficientwise
limit yields the correct Hermite normalization for every fixed degree.


\section{Racah to dual Hahn: reverse-connection limit}
The source V31 gives the reverse connection \texttt{eq:racah-dh-reverse},
with $N_0=-\gamma-1\in\mathbb Z_{\ge0}$ and
$\Lambda(x)=x(x+\gamma+\delta+1)$:
\begin{equation}\label{eq:racah-dh-audit}
R_n(\Lambda(x);\alpha,\beta,\gamma,\delta)
=\sum_{m=0}^n \binom nm
\frac{(n+\alpha+\beta+1)_m(\delta-\alpha-n)_{n-m}}
{(\beta+\delta+1)_n}
\operatorname{DH}_m(\Lambda(x);\alpha,\gamma+\delta-\alpha,N_0).
\end{equation}
\begin{proposition}
For fixed $0\le n\le N_0$ and fixed admissible $\alpha,\gamma,\delta,x$,
\begin{equation}\label{eq:racah-dh-limit}
\lim_{\beta\to+\infty}R_n(\Lambda(x);\alpha,\beta,\gamma,\delta)
=\operatorname{DH}_n(\Lambda(x);\alpha,\gamma+\delta-\alpha,N_0).
\end{equation}
The limit follows directly from the V31 reverse connection.
\end{proposition}
\begin{proof}
For fixed $n,m$, the coefficient in \eqref{eq:racah-dh-audit} is
\[
 A_{n,m}(\beta)=\binom nm
 \frac{(\beta+n+\alpha+1)_m}{(\beta+\delta+1)_n}
 (\delta-\alpha-n)_{n-m}.
\]
For any fixed $u$ and nonnegative integer $r$, $(\beta+u)_r=\beta^r(1+O(\beta^{-1}))$.
Consequently,
\[
 A_{n,m}(\beta)=\binom nm(\delta-\alpha-n)_{n-m}
 \beta^{m-n}(1+O(\beta^{-1})).
\]
When $m<n$ this tends to zero. When $m=n$, the binomial and length-zero
Pochhammer factors are $1$ and the ratio of two monic degree-$n$
polynomials tends to $1$. The dual Hahn basis is independent of $\beta$,
and the sum is finite, so taking its limit term by term proves
\eqref{eq:racah-dh-limit}. The case $n=0$ is immediate.
\end{proof}
\noindent\textbf{Audit verdict: PASS.} For fixed admissible parameters and
$0\le n\le N_0$, the reverse-connection coefficients converge to
$\delta_{nm}$. This verifies the stated Racah-to-dual-Hahn degeneration;
the exact KLS equation number has not been independently cross-checked.


\section{Racah to Hahn: reverse finite-grid connection}
The V31 reverse finite-grid connection, equation
\texttt{eq:racah-to-hahn-grid}, reads
\begin{equation}\label{eq:racah-hahn-audit}
R_n(\lambda(x);\alpha,\beta,\gamma,\delta)
=\sum_{k=0}^{N} A_{n,k}(\delta)Q_k(x;\alpha,\beta,N),
\end{equation}
where $\gamma=-N-1$, $\lambda(x)=x(x+\gamma+\delta+1)$, and
\[
 A_{n,k}(\delta)=
 \frac{\sum_{y=0}^N w_H(y)
 R_n(\lambda(y);\alpha,\beta,-N-1,\delta)
 Q_k(y;\alpha,\beta,N)}
 {\sum_{y=0}^N w_H(y)Q_k(y;\alpha,\beta,N)^2},
 \qquad
 w_H(y)=\frac{(\alpha+1)_y(\beta+1)_{N-y}}{y!(N-y)!}.
\]
\begin{proposition}
For fixed $N$, $0\le n\le N$, and generic parameters for which
both polynomial systems and the Hahn norms are defined,
\begin{equation}\label{eq:racah-hahn-limit}
\lim_{\delta\to+\infty}
R_n\bigl(x(x+\delta-N);\alpha,\beta,-N-1,\delta\bigr)
=Q_n(x;\alpha,\beta,N),\qquad x=0,\ldots,N.
\end{equation}
The limit follows from the reverse connection \eqref{eq:racah-hahn-audit}.
\end{proposition}
\begin{proof}
The standard terminating Racah representation is
\[
R_n(\lambda(x);\alpha,\beta,-N-1,\delta)
={}_4F_3\!\left(\begin{matrix}-n,\ n+\alpha+\beta+1,\ -x,\ x+\delta-N\\
\alpha+1,\ \beta+\delta+1,\ -N\end{matrix};1\right).
\]
For each fixed summation index $j\le n$ and grid point $x$,
\[
\frac{(x+\delta-N)_j}{(\beta+\delta+1)_j}\longrightarrow 1.
\]
As the hypergeometric series terminates, its limit is
\[
{}_3F_2\!\left(\begin{matrix}-n,\ n+\alpha+\beta+1,\ -x\\
\alpha+1,\ -N\end{matrix};1\right)
=Q_n(x;\alpha,\beta,N).
\]
We use this finite-series observation only to identify the limiting
Racah values; the required coefficient test is performed directly on
V31's reverse connection.  Since both sums defining $A_{n,k}$ are
finite and the Hahn weights and norms do not depend on $\delta$,
\[
\lim_{\delta\to\infty}A_{n,k}(\delta)
=\frac{\sum_{y=0}^N w_H(y)Q_n(y;\alpha,\beta,N)
 Q_k(y;\alpha,\beta,N)}
 {\sum_{y=0}^N w_H(y)Q_k(y;\alpha,\beta,N)^2}
=\delta_{nk}
\]
by Hahn orthogonality.  Passing to the limit in the finite reverse
connection therefore gives \eqref{eq:racah-hahn-limit}.
\end{proof}
\noindent\textbf{Audit verdict: PASS.} The coefficients in the V31
reverse finite-grid connection tend to the identity matrix for the
Racah-to-Hahn degeneration $\gamma=-N-1$, $\delta\to\infty$.
This is a finite-grid result, not a claim of polynomial identity in a
common variable before taking the limit. The exact KLS equation number
has not been independently cross-checked.


\section{Askey--Wilson to continuous dual $q$-Hahn: reverse connection}
The next distinct-family genuine limit in V31's continuous-$q$ section is obtained by setting the fourth Askey--Wilson parameter to zero. For fixed $0<|q|<1$, $n\geq0$, and fixed generic $a,b,c,x$, the target relation in the source normalization is
\begin{equation}\label{eq:aw-cdh-target}
\lim_{d\to0}p_n(x;a,b,c,d\mid q)=S_n(x;a,b,c\mid q).
\end{equation}
The reverse connection in V31, equation \texttt{eq:AWCD}, is
\begin{equation}\label{eq:aw-cdh-reverse-audit}
\begin{aligned}
p_n(x;a,b,c,d\mid q)=\sum_{k=0}^n {n\brack k}_q
&(-d)^{n-k}q^{\binom{n-k}{2}}\\[-2pt]
&\times(abq^k,acq^k,bcq^k;q)_{n-k}
(abcdq^{n-1};q)_k\,S_k(x;a,b,c\mid q).
\end{aligned}
\end{equation}
\begin{proposition}
The reverse connection \eqref{eq:aw-cdh-reverse-audit} implies \eqref{eq:aw-cdh-target} for every fixed degree $n$.
\end{proposition}
\begin{proof}
For $0\leq k<n$, the coefficient of $S_k$ in \eqref{eq:aw-cdh-reverse-audit} contains $(-d)^{n-k}$. All remaining factors have finite limits as $d\to0$: the $q$-binomial coefficient and the three $q$-shifted factorials are independent of $d$, while $(abcdq^{n-1};q)_k\to1$. Therefore the coefficient is $O(d^{n-k})\to0$. For $k=n$, the factors of length zero equal $1$, and the coefficient is exactly $(abcdq^{n-1};q)_n\to1$. The connection sum is finite, so its limit equals $S_n(x;a,b,c\mid q)$. This also covers $n=0$.
\end{proof}
\noindent\textbf{Audit verdict: PASS.} The reverse-connection matrix tends to the identity in the fourth-parameter limit $d\to0$, with no variable rescaling or additional normalization. This is the standard Askey--Wilson $\to$ continuous dual $q$-Hahn degeneration; an exact KLS equation number has not been independently checked.


\section{$q$-Meixner to Al-Salam--Carlitz II: reverse connection}
The next genuine distinct-family degeneration explicitly identified in V31 is
\begin{equation}\label{eq:qm-ascii-target}
\lim_{c\to0}M_n(x;-a/c,c;q)
=(-1/a)^nq^{\binom n2}V_n^{(a)}(x;q),
\end{equation}
for fixed $n$, $0<q<1$, and nonzero fixed $a$ (with denominators
nonvanishing along the limiting path). V31 cites the corresponding
Koekoek--Swarttouw limit as Eq.~(4.13.3). Its reverse connection,
\texttt{eq:ASCII-qM-reverse}, is
\begin{equation}\label{eq:qm-ascii-reverse-audit}
M_n(x;b,c;q)=\sum_{k=0}^n B_{n,k}(b,c)V_k^{(a)}(x;q),
\end{equation}
where
\[
B_{n,k}(b,c)=\sum_{j=k}^n
\frac{(q^{-n};q)_j(-q^{n+1}/c)^j}{(q,bq;q)_j}
(-1)^j q^{\binom k2+\binom{j-k}2}
{j\brack k}_q a^{j-k}.
\]
\begin{proposition}
The reverse connection \eqref{eq:qm-ascii-reverse-audit} implies
\eqref{eq:qm-ascii-target} for every fixed nonnegative integer $n$.
\end{proposition}
\begin{proof}
Set $b=-a/c$. For each fixed $j$,
\[
(-aq/c;q)_j=(a/c)^jq^{j(j+1)/2}(1+O(c)),
\]
so the Newton coefficient in the reverse connection has limit
\[
\frac{(q^{-n};q)_j(-q^{n+1}/c)^j}
{(q;q)_j(-aq/c;q)_j}
\longrightarrow
\frac{(q^{-n};q)_j}{(q;q)_j}
(-1)^ja^{-j}q^{jn-\binom j2}.
\]
The two factors $(-1)^j$ in the product cancel. Writing $m=n-k$
and $j=k+r$, using
\(
{j\brack k}_q=(q;q)_j/((q;q)_k(q;q)_{j-k})
\)
and
\((q^{-n};q)_{k+r}=(q^{-n};q)_k(q^{-m};q)_r\),
we obtain
\begin{align*}
\lim_{c\to0}B_{n,k}(-a/c,c)
&=\frac{a^{-k}(q^{-n};q)_k}{(q;q)_k}
q^{k(n-k)+k^2}
\sum_{r=0}^{m}\frac{(q^{-m};q)_r}{(q;q)_r}q^{mr}\\
&=\frac{a^{-k}(q^{-n};q)_k}{(q;q)_k}q^{kn}
\sum_{r=0}^{m}\frac{(q^{-m};q)_r}{(q;q)_r}q^{mr}.
\end{align*}
The terminating $q$-binomial identity gives
\[
\sum_{r=0}^{m}\frac{(q^{-m};q)_r}{(q;q)_r}z^r
=\frac{(zq^{-m};q)_m}{(z;q)_m}.
\]
At $z=q^m$, this sum vanishes for $m>0$, since
$(1;q)_m=0$. For $m=0$ it equals $1$. Thus all coefficients
with $k<n$ tend to zero. For $k=n$,
\[
\lim_{c\to0}B_{n,n}(-a/c,c)
=a^{-n}q^{n^2}\frac{(q^{-n};q)_n}{(q;q)_n}
=(-1/a)^nq^{\binom n2},
\]
using $(q^{-n};q)_n=(-1)^nq^{-n(n+1)/2}(q;q)_n$.
The sum \eqref{eq:qm-ascii-reverse-audit} is finite, so its limit is
exactly \eqref{eq:qm-ascii-target}.
\end{proof}
\noindent\textbf{Audit verdict: PASS.} The off-diagonal coefficients
vanish by a terminating $q$-binomial cancellation, while the diagonal
coefficient gives the required nontrivial normalization.


\section{Quantum $q$-Krawtchouk to Al-Salam--Carlitz II: normalization discrepancy}
\noindent\textbf{Source:} V31, reverse connection \texttt{eq:ASCII-qtm-reverse}; the
literature note following that proposition cites Koekoek--Swarttouw,
Eq.~(4.14.1). The stated parameter path is
\[
 p=a^{-1}q^{-N-1},\qquad N\longrightarrow\infty,
 \qquad 0<q<1,\quad a\ne0.
\]
The source defines the quantum $q$-Krawtchouk Newton coefficients by
\[
 S_{n,j}(p)=\frac{(q^{-n};q)_j(pq^{n+1})^j}
 {(q;q)_j(q^{-N};q)_j},\qquad
 K_n^{\mathrm{qtm}}(y;p,N\mid q)=\sum_{j=0}^n S_{n,j}(p)(y;q)_j.
\]
For fixed $j$ we have
\[
 \frac{q^{-Nj}}{(q^{-N};q)_j}\longrightarrow
 (-1)^j q^{-\binom j2}.
\]
Consequently, inserting $p=a^{-1}q^{-N-1}$ gives
\begin{align*}
 S_{n,j}(p)&\longrightarrow
 \frac{(q^{-n};q)_j}{(q;q)_j}
 a^{-j}q^{nj}(-1)^jq^{-\binom j2}\\
 &=\genfrac{[}{]}{0pt}{}{n}{j}_q a^{-j}.
\end{align*}
Thus the \emph{actual limit implied by the reverse connection} is
\begin{equation}\label{eq:audit-qtm-limit}
 \lim_{N\to\infty}K_n^{\mathrm{qtm}}
 (y;a^{-1}q^{-N-1},N\mid q)
 =\sum_{j=0}^n\genfrac{[}{]}{0pt}{}{n}{j}_q
 a^{-j}(y;q)_j.
\end{equation}
Indeed, the reverse connection expresses each Newton factor $(y;q)_j$
in the $V_k^{(a)}$ basis, so this coefficientwise calculation is
precisely the finite reverse-connection limit, with no appeal to the
known family limit.

The literature note in V31 states the normalization
$(-1/a)^nq^{\binom n2}V_n^{(a)}(y;q)$. The Newton expansion
\texttt{eq:ASCII-qM-ASCII-newton} reads
\[
 V_n^{(a)}(y;q)=(-1)^nq^{-\binom n2}
 \sum_{j=0}^n\genfrac{[}{]}{0pt}{}{n}{j}_q a^{n-j}(y;q)_j.
\]
Multiplication by $(-1/a)^nq^{\binom n2}$ yields exactly
\eqref{eq:audit-qtm-limit}. Hence
\[
 \boxed{\displaystyle
 \lim_{N\to\infty}K_n^{\mathrm{qtm}}
 (y;a^{-1}q^{-N-1},N\mid q)
 =(-1/a)^nq^{\binom n2}V_n^{(a)}(y;q).}
\]
\noindent\textbf{Verdict: PASS.} The limiting Newton coefficients
are computed directly from the reverse connection for every fixed
$n$; their sum agrees with the stated Al-Salam--Carlitz II
normalization. The exact KLS equation is cited in V31 but has not
been checked independently against the original publication.


\section{\texorpdfstring{$q$-Laguerre to Stieltjes--Wigert}{q-Laguerre to Stieltjes--Wigert}: complete}
\noindent\textbf{Source:} V31, reverse connection \texttt{eq:SWqL-rev}, and the source normalization \texttt{eq:SWqL-Ldef}. Set $a=q^\alpha$, with $0<q<1$ fixed; then $\alpha\to+\infty$ is equivalent to $a\to0^+$.
The target limit is
\begin{equation}\label{eq:swql-target}
 \lim_{a\to0^+}L_n^{(\alpha)}(x/a;q)=S_n(x;q).
\end{equation}
The reverse connection in V31 has the form
\begin{equation}\label{eq:swql-rev-audit}
 L_n^{(\alpha)}(y;q)=\sum_{k=0}^n A_{n,k}(a)S_k(y;q),
\end{equation}
where
\[
 A_{n,k}(a)=\frac{(-1)^k(q;q)_k}{q^{\binom{k}{2}}}
 \sum_{j=k}^n \frac{(-1)^j a^j q^{\binom{j}{2}}
 (aq^{j+1};q)_{n-j}(q^{-j};q)_k}
 {(q;q)_j(q;q)_{n-j}(q^{-k};q)_k}.
\]
\begin{proposition}
For each fixed $n$, the reverse connection \eqref{eq:swql-rev-audit} implies \eqref{eq:swql-target}.
\end{proposition}
\begin{proof}
Fix $0\le k\le n$. Since the inner sum begins at $j=k$ and all denominators are nonzero for $0<q<1$, its first term gives
\[
 A_{n,k}(a)=\frac{a^k}{(q;q)_{n-k}}+O(a^{k+1}),\qquad a\to0^+.
\]
Indeed, at $j=k$ the factors $(-1)^k$, $q^{\binom{k}{2}}$, $(q;q)_k$ and $(q^{-k};q)_k$ cancel their corresponding outer factors, and $(aq^{k+1};q)_{n-k}\to1$.
The Stieltjes--Wigert polynomial in the source normalization is of degree $k$ with leading coefficient
\[
 [y^k]S_k(y;q)=\frac{(-1)^kq^{k^2}}{(q;q)_k}.
\]
Therefore, for fixed $x$,
\[
 a^k S_k(x/a;q)\longrightarrow
 \frac{(-1)^kq^{k^2}}{(q;q)_k}x^k.
\]
Substitute $y=x/a$ into \eqref{eq:swql-rev-audit}. Every summand has a finite limit, and the finite sum becomes
\begin{equation}\label{eq:swql-limit-sum}
 \lim_{a\to0^+}L_n^{(\alpha)}(x/a;q)
 =\sum_{k=0}^n\frac{(-1)^kq^{k^2}}{(q;q)_k(q;q)_{n-k}}x^k.
\end{equation}
The right-hand side is precisely the Stieltjes--Wigert polynomial $S_n(x;q)$ in the normalization used in V31 (the polynomial is also obtained by setting $a=0$ in the coefficient formula \texttt{eq:SWqL-Ldef} after replacing $x$ by $x/a$). Thus \eqref{eq:swql-target} follows from the reverse connection. The argument also holds coefficientwise as a polynomial identity.
\end{proof}
\noindent\textbf{Audit verdict: PASS.} The rescaled basis is essential: although $A_{n,k}(a)=O(a^k)$, the factor $S_k(x/a;q)$ grows like $a^{-k}$, so \emph{every} degree $k$ contributes to the limit. The resulting finite sum reconstructs $S_n(x;q)$ exactly. No defining hypergeometric limit is used to establish convergence.


\section{Continuous big $q$-Hermite to continuous $q$-Hermite: complete}
\noindent\textbf{Source:} V31, reverse connection \texttt{eq:BHH} in the subsection \texttt{sec:hermite-big-clean}. For fixed $0<q<1$ and $n\ge0$, the known specialization limit is
\begin{equation}\label{eq:bqh-target}
 \lim_{a\to0}H_n(x;a\mid q)=H_n(x\mid q).
\end{equation}
The source reverse connection is
\begin{equation}\label{eq:bqh-reverse-audit}
 H_n(x;a\mid q)=\sum_{k=0}^{n}
 \begin{bmatrix}n\\k\end{bmatrix}_q
 (-a)^{n-k}q^{\binom{n-k}{2}}H_k(x\mid q),
\end{equation}
where $\begin{bmatrix}n\\k\end{bmatrix}_q=(q;q)_n/((q;q)_k(q;q)_{n-k})$.
\begin{proposition}
The reverse connection \eqref{eq:bqh-reverse-audit} implies \eqref{eq:bqh-target} for every fixed $n$, coefficientwise as a polynomial in $x$.
\end{proposition}
\begin{proof}
For $k=n$ the connection coefficient equals $1$ for every $a$. For $k<n$, its coefficient contains the positive power $a^{n-k}$, while the remaining factors are independent of $a$. Therefore
\[
 \lim_{a\to0}\begin{bmatrix}n\\k\end{bmatrix}_q
 (-a)^{n-k}q^{\binom{n-k}{2}}=\delta_{nk}.
\]
The sum in \eqref{eq:bqh-reverse-audit} is finite, and its basis polynomials do not depend on $a$. Taking the limit term by term yields $H_n(x\mid q)$. Since the equality holds in the finite-dimensional polynomial space of degree at most $n$, convergence is coefficientwise. The case $n=0$ is immediate.
\end{proof}
\noindent\textbf{Audit verdict: PASS.} No rescaling or extra normalization is needed. This is a parameter-deletion limiting relation between distinct families, not a self-connection or an exact quadratic transformation. The precise numbered KLS entry has not been independently checked.


\section{Continuous dual $q$-Hahn to Al-Salam--Chihara: complete}
\noindent\textbf{Source:} V31, reverse connection \texttt{eq:CDASC} in subsection \texttt{sec:asc-cdh32}. The classical parameter-deletion limit, in the source normalization, is
\begin{equation}\label{eq:cdasc-target}
 \lim_{c\to0}S_n(x;a,b,c\mid q)=Q_n(x;a,b\mid q),\qquad n\ge0.
\end{equation}
For fixed $0<q<1$, fixed $a,b,x$, and admissible parameters, V31 gives
\begin{equation}\label{eq:cdasc-reverse-audit}
 S_n(x;a,b,c\mid q)=\sum_{k=0}^{n}
 \begin{bmatrix}n\\k\end{bmatrix}_q
 (-c)^{n-k}q^{\binom{n-k}{2}}(abq^k;q)_{n-k}
 Q_k(x;a,b\mid q).
\end{equation}
\begin{proposition}
For every fixed $n$, the reverse connection \eqref{eq:cdasc-reverse-audit} implies \eqref{eq:cdasc-target}, coefficientwise in $x$.
\end{proposition}
\begin{proof}
Denote the coefficient of $Q_k$ in \eqref{eq:cdasc-reverse-audit} by $A_{n,k}(c)$. If $k=n$, the Gaussian binomial coefficient is $1$, and all $q$-shifted factorials of length zero equal $1$, so $A_{n,n}(c)=1$ identically. If $0\le k<n$, the remaining factors are independent of $c$ and finite, whereas $(-c)^{n-k}\to0$. Thus
\[
 \lim_{c\to0}A_{n,k}(c)=\delta_{nk}.
\]
The basis polynomials $Q_k(x;a,b\mid q)$ are independent of $c$, and the sum is finite. Taking the limit term by term proves \eqref{eq:cdasc-target}. The same reasoning proves coefficientwise convergence, and includes $n=0$.
\end{proof}
\noindent\textbf{Audit verdict: PASS.} The reverse connection gives the parameter-deletion limit directly, without using the defining hypergeometric series. The exact numbered KLS citation has not been independently checked.


\section{Al-Salam--Chihara to continuous big $q$-Hermite: complete}
\noindent\textbf{Source:} V31, reverse connection \texttt{eq:ASCBH} in subsection \texttt{sec:big-asc-clean}. The parameter-deletion limit is
\begin{equation}\label{eq:asc-bh-target}
 \lim_{b\to0}Q_n(x;a,b\mid q)=H_n(x;a\mid q),\qquad n\ge0.
\end{equation}
The reverse connection in V31 is
\begin{equation}\label{eq:asc-bh-reverse-audit}
 Q_n(x;a,b\mid q)=\sum_{k=0}^{n}
 \begin{bmatrix}n\\k\end{bmatrix}_q
 (-b)^{n-k}q^{\binom{n-k}{2}}H_k(x;a\mid q).
\end{equation}
\begin{proposition}
For fixed $0<q<1$, fixed admissible $a$, and every fixed nonnegative integer $n$, the reverse connection \eqref{eq:asc-bh-reverse-audit} implies \eqref{eq:asc-bh-target} coefficientwise in $x$.
\end{proposition}
\begin{proof}
Write $A_{n,k}(b)=\begin{bmatrix}n\\k\end{bmatrix}_q(-b)^{n-k}q^{\binom{n-k}{2}}$. For $k=n$, $A_{n,n}(b)=1$. For $k<n$, the exponent $n-k$ is a positive integer, and all other factors are finite and independent of $b$; hence $A_{n,k}(b)\to0$. Thus $A_{n,k}(b)\to\delta_{nk}$. Since the sum is finite and each $H_k(x;a\mid q)$ is independent of $b$, its limit is $H_n(x;a\mid q)$. The convergence holds coefficientwise because all polynomials have degree at most $n$. This includes $n=0$.
\end{proof}
\noindent\textbf{Audit verdict: PASS.} This is a genuine parameter-deletion limit of distinct families, obtained directly from the reverse connection without a defining-series calculation. The numbered KLS citation has not been independently verified.


\section{Continuous $q$-Hahn to $q$-Meixner--Pollaczek: complete}
\noindent\textbf{Source:} V31, equation \texttt{eq:cqh-qmp23}, in subsection \texttt{sec:qmp-cqh-clean}. This is a parameter-deletion limit, distinct from the exact quadratic transformations excluded from this audit.

Fix $0<q<1$, $a$, and $\phi$, and put $z=e^{i\phi}$, $h=bz$. The source manuscript defines
\[
\mathscr P_n(x;a,\phi\mid q)=\frac{p_n(x;az,0,0,a/z\mid q)}{(q;q)_n},\qquad
\mathscr Q_n(x;a,b,0,a,\phi\mid q)=p_n(x;az,bz,0,a/z\mid q).
\]
The required parameter-deletion limit, with this normalization, is
\begin{equation}\label{eq:qmp-target-audit}
\lim_{b\to0}\mathscr Q_n(x;a,b,0,a,\phi\mid q)
=(q;q)_n\mathscr P_n(x;a,\phi\mid q).
\end{equation}
The \emph{reverse} connection, copied from V31, is
\begin{equation}\label{eq:qmp-reverse-audit}
\mathscr Q_n(x;a,b,0,a,\phi\mid q)
=\sum_{k=0}^{n}\begin{bmatrix}n\\k\end{bmatrix}_q
(q;q)_k(-h)^{n-k}q^{\binom{n-k}{2}}
(a^2q^k;q)_{n-k}\,\mathscr P_k(x;a,\phi\mid q).
\end{equation}
\begin{proposition}
For every fixed $n\ge0$, equation \eqref{eq:qmp-reverse-audit} implies \eqref{eq:qmp-target-audit} coefficientwise in $x$.
\end{proposition}
\begin{proof}
For $k<n$, the coefficient of $\mathscr P_k$ contains $(-h)^{n-k}=(-bz)^{n-k}$, and therefore tends to zero as $b\to0$. All other factors are independent of $b$. For $k=n$, the coefficient is exactly $(q;q)_n$. Because the sum is finite and the basis polynomials $\mathscr P_k$ are independent of $b$, termwise passage to the limit gives $(q;q)_n\mathscr P_n$, as claimed. This includes $n=0$.
\end{proof}
\noindent\textbf{Audit verdict: PASS.} The reverse connection establishes the parameter-deletion limit for all fixed degrees. The factor $(q;q)_n$ is essential: V31 defines $\mathscr P_n$ with this denominator. The precise numbered KLS citation remains to be checked independently.


\section{Big $q$-Jacobi to big $q$-Laguerre: complete}
\noindent\textbf{Source:} V31, equation \texttt{eq:blbj-reverse}. This is the genuine parameter-deletion limit $b\to0$, not an audit of the general connection at unrelated parameters.

In the notation of V31, the target is
\begin{equation}\label{eq:bj-bl-target}
 \lim_{b\to0}J_n(x;a,b,c\mid q)=P_n(x;a,c\mid q).
\end{equation}
The reverse connection in V31 is
\begin{equation}\label{eq:bj-bl-reverse}
J_n(x;a,b,c\mid q)=\sum_{k=0}^n D_{n,k}(a,b,c;A,B)P_k(x;A,B\mid q),
\end{equation}
where
\begin{equation}\label{eq:bj-bl-coeff}
D_{n,k}=(-1)^kq^{\binom{k}{2}}\sum_{r=k}^n
 \frac{(q^{-n},abq^{n+1};q)_r q^r(Aq,Bq;q)_r}
 {(q,aq,cq;q)_r}\begin{bmatrix}r\\k\end{bmatrix}_q.
\end{equation}
\begin{proposition}
For fixed $0<q<1$, fixed admissible $a,c$ and fixed degree $n$, the reverse connection \eqref{eq:bj-bl-reverse} implies \eqref{eq:bj-bl-target} coefficientwise in $x$.
\end{proposition}
\begin{proof}
Choose $A=a$ and $B=c$ in \eqref{eq:bj-bl-reverse}. The basis $P_k(x;a,c\mid q)$ is then independent of $b$. For $b\to0$, the factors $(abq^{n+1};q)_r$ tend to $1$, and the numerator $(aq,cq;q)_r$ cancels the same factors in the denominator. Hence
\[
\lim_{b\to0}D_{n,k}=(-1)^kq^{\binom{k}{2}}
\sum_{r=k}^n\frac{(q^{-n};q)_r q^r}{(q;q)_r}
\begin{bmatrix}r\\k\end{bmatrix}_q.
\]
Writing $r=k+t$, using the Gaussian binomial definition, and factoring the $t=0$ term gives
\[
\lim_{b\to0}D_{n,k}=
(-1)^kq^{\binom{k}{2}}
\frac{(q^{-n};q)_kq^k}{(q;q)_k}
\sum_{t=0}^{n-k}\frac{(q^{k-n};q)_t q^t}{(q;q)_t}.
\]
The terminating $q$-binomial theorem evaluates the final sum as
\[
\sum_{t=0}^{m}\frac{(q^{-m};q)_t q^t}{(q;q)_t}
=(q^{1-m};q)_m.
\]
For $m=n-k>0$ this product contains $1-q^0=0$. Thus every off-diagonal coefficient tends to zero. For $k=n$, the sum is $1$, and
\[
(-1)^nq^{\binom n2}\frac{(q^{-n};q)_nq^n}{(q;q)_n}=1,
\]
by $(q^{-n};q)_n=(-1)^nq^{-n(n+1)/2}(q;q)_n$.
Therefore $\lim_{b\to0}D_{n,k}=\delta_{nk}$. The reverse connection is a finite sum, so the limit gives $P_n(x;a,c\mid q)$, including $n=0$.
\end{proof}
\noindent\textbf{Audit verdict: PASS.} This is a proof from the reverse connection for every fixed degree; the precise numbered KLS citation has not been independently confirmed.


\section{Little $q$-Jacobi to little $q$-Laguerre (Wall): complete}
The parameter-deletion limit in the normalizations of V31 is
\begin{equation}\label{eq:lj-wall-target}
\lim_{B\to0}p_n(x;a,B\mid q)=w_n(x;a\mid q),\qquad 0<q<1,
\end{equation}
for fixed $n$, $x$, and admissible $a$. This is a genuine distinct-family limiting relation (not a parameter-change self-connection). We use only the \emph{reverse} connection \texttt{eq:wall-lj-reverse} from V31, with $A=a$:
\begin{equation}\label{eq:lj-wall-reverse}
 p_n(x;a,B\mid q)=\sum_{k=0}^n C_{n,k}(B)w_k(x;a\mid q),
\end{equation}
where
\[
 C_{n,k}(B)=(-1)^kq^{\binom{k}{2}}
 \sum_{j=k}^n
 \frac{(q^{-n},aBq^{n+1};q)_j q^j(aq;q)_j}
 {(aq;q)_j(q;q)_j}
 \begin{bmatrix}j\\k\end{bmatrix}_q.
\]
\begin{proposition}
For all $0\le k\le n$, $\lim_{B\to0}C_{n,k}(B)=\delta_{nk}$. Hence the reverse connection proves \eqref{eq:lj-wall-target} for every fixed degree.
\end{proposition}
\begin{proof}
For fixed $n$ the sum is finite, and all factors are regular at $B=0$ under the stated admissibility conditions. Cancel $(aq;q)_j$ and set $B=0$. Since $(0;q)_j=1$ and
$\begin{bmatrix}j\\k\end{bmatrix}_q=(q;q)_j/((q;q)_k(q;q)_{j-k})$, we obtain
\[
 C_{n,k}(0)=\frac{(-1)^kq^{\binom{k}{2}}}{(q;q)_k}
 \sum_{j=k}^n\frac{(q^{-n};q)_j q^j}{(q;q)_{j-k}}.
\]
Put $r=n-k$ and $j=k+\ell$. Splitting the $q$-shifted factorial gives
\[
 C_{n,k}(0)=\frac{(-1)^kq^{\binom{k}{2}+k}(q^{-n};q)_k}{(q;q)_k}
 \sum_{\ell=0}^{r}\frac{(q^{-r};q)_\ell q^\ell}{(q;q)_\ell}.
\]
The terminating $q$-binomial identity yields
\[
 \sum_{\ell=0}^{r}\frac{(q^{-r};q)_\ell q^\ell}{(q;q)_\ell}
 =(q^{1-r};q)_r.
\]
For $r\ge1$ the product contains the factor $1-q^0=0$. Thus $C_{n,k}(0)=0$ for $k<n$. If $r=0$, then $k=n$ and the sum equals $1$; moreover
\[
 (q^{-n};q)_n=(-1)^nq^{-n(n+1)/2}(q;q)_n,
\]
so $C_{n,n}(0)=1$. Consequently $C_{n,k}(B)\to\delta_{nk}$. Taking the limit of the finite reverse connection \eqref{eq:lj-wall-reverse} proves the assertion, including $n=0$.
\end{proof}
\noindent\textbf{Audit verdict: PASS.} All off-diagonal coefficients vanish and the diagonal coefficient tends to one. The argument does not use the defining hypergeometric polynomial limit. The precise numbered KLS citation remains to be independently matched.


\section{Affine $q$-Krawtchouk to little $q$-Laguerre (Wall): rescaled limit}
The next distinct-family classical limit requires a change of variable. In the
normalizations of V31, for $0<q<1$, fixed $n$, and $p$ away from poles, it is
\begin{equation}\label{eq:aff-wall-target}
 \lim_{N\to\infty}K_n^{\rm aff}(q^{-N}y;p,N\mid q)=w_n(y;p\mid q).
\end{equation}
The relevant reverse connection is V31, equation
\texttt{eq:wall-affine-reverse}, with $a=p$:
\begin{align}\label{eq:aff-wall-reverse-audit}
 K_n^{\rm aff}(x;p,N\mid q)
 &=\sum_{k=0}^n D_{n,k}(p,N)w_k(x;p\mid q),\\
 D_{n,k}(p,N)
 &=q^{k(k-1)}(pq;q)_k
 \sum_{r=k}^n\frac{(q^{-n};q)_rq^r}{(pq,q^{-N},q;q)_r}
 {r\brack k}_q
 \sum_{t=0}^{r-k}\frac{(q^{k-r},pq^{k+1};q)_tq^{rt}}{(q;q)_t}.
 \notag
\end{align}
\begin{proposition}
The rescaled reverse connection \eqref{eq:aff-wall-reverse-audit}
reproduces \eqref{eq:aff-wall-target} for every fixed degree $n$.
\end{proposition}
\begin{proof}
Put $F_r(x)=(x;q)_r$ and $T=q^{-N}$. The finite Newton-basis
inversion used to derive the reverse coefficients in V31 gives the
\emph{finite algebraic identity}
\begin{equation}\label{eq:aff-wall-resum}
 \sum_{k=0}^n D_{n,k}(p,N)w_k(x;p\mid q)
 =\sum_{r=0}^n\frac{(q^{-n};q)_rq^r}
 {(pq;q)_r(q^{-N};q)_r(q;q)_r}F_r(x).
\end{equation}
This is an exact resummation of the reverse connection (it is not a
limiting assertion). It follows by inserting the terminating monomial
expansion of $w_k$ and applying the finite $q$-binomial inversion
identity to the coefficient of each $F_r$; equivalently it is the
Newton-basis coefficient identity established in the proof of V31,
Proposition \texttt{prop:wall-affine}. In particular, no limiting
formula for affine $q$-Krawtchouk polynomials is assumed.

Now substitute $x=Ty$. For each fixed $r$, as $N\to\infty$,
\[
 \frac{F_r(Ty)}{(T;q)_r}
 =\prod_{j=0}^{r-1}\frac{1-Tyq^j}{1-Tq^j}
 \longrightarrow y^r,
 \qquad T=q^{-N}\longrightarrow\infty.
\]
Since the sum in \eqref{eq:aff-wall-resum} has only $n+1$ terms,
its limit is
\[
 \sum_{r=0}^n\frac{(q^{-n};q)_rq^r}
 {(pq;q)_r(q;q)_r}y^r
 =w_n(y;p\mid q),
\]
where the last equality is the finite polynomial expansion of the Wall
basis used in the reverse-connection derivation. This proves
\eqref{eq:aff-wall-target}.
\end{proof}
\noindent\textbf{Audit verdict: PASS (with resummation).}
Unlike parameter-deletion cases, the individual terms
$D_{n,k}(p,N)w_k(q^{-N}y;p\mid q)$ cannot be replaced by their
unscaled limits. The exact finite resummation is essential. The
numbered KLS reference remains to be independently matched.

\section{Little $q$-Jacobi to $q$-Laguerre: confluent limit}
The next eligible distinct-family limit in V31 is the confluent
little-$q$-Jacobi to $q$-Laguerre transition. In the normalizations
specified in V31, equation \texttt{eq:qLlJ-D}, it takes the form
\begin{equation}\label{eq:lj-ql-target}
 \lim_{B\to+\infty}\frac{(aq;q)_n}{(q;q)_n}
 p_n\!\left(-\frac{x}{Bq};a,B\mid q\right)=L_n(x;a\mid q),
 \qquad 0<q<1,\quad n\ \text{fixed}.
\end{equation}
Here $a$ and $x$ are fixed, with nonvanishing denominators. The
reverse connection in V31, equation \texttt{eq:qLlJ-D}, is
\begin{equation}\label{eq:lj-ql-reverse-audit}
 p_n(y;A,B\mid q)=\sum_{k=0}^n\left[\sum_{j=k}^n
 \frac{(q^{-n},ABq^{n+1};q)_jq^j}{(Aq,q;q)_j}
 \frac{(-1)^k(aq;q)_j(q;q)_k}{a^jq^{j^2}(aq;q)_k}
 q^{\binom{j-k}{2}}{j\brack k}_q\right]L_k(y;a\mid q).
\end{equation}
\begin{proposition}
The reverse connection \eqref{eq:lj-ql-reverse-audit} implies
\eqref{eq:lj-ql-target} for every fixed degree $n$.
\end{proposition}
\begin{proof}
Set $A=a$ in \eqref{eq:lj-ql-reverse-audit}. The finite
Gaussian-binomial inversion identity stated and proved in V31 immediately
below that connection is
\[
 y^j=\sum_{k=0}^j
 \frac{(-1)^k(aq;q)_j(q;q)_k}{a^jq^{j^2}(aq;q)_k}
 q^{\binom{j-k}{2}}{j\brack k}_q L_k(y;a\mid q).
\]
Reordering the finite double sum in the \emph{reverse connection}, and
applying this exact inversion, gives
\begin{equation}\label{eq:lj-ql-resum}
 \sum_{k=0}^n D_{n,k}(a,B) L_k(y;a\mid q)
 =\sum_{j=0}^n\frac{(q^{-n},aBq^{n+1};q)_jq^j}
 {(aq;q)_j(q;q)_j}y^j.
\end{equation}
This step is an exact algebraic resummation of the V31 reverse
connection, not an invocation of a known limiting relation.
Substitute $y=-x/(Bq)$ and multiply by $(aq;q)_n/(q;q)_n$.
For each fixed $j$,
\[
 B^{-j}(aBq^{n+1};q)_j
 \longrightarrow (-a)^jq^{j(n+1)+\binom j2}.
\]
Also,
\[
 (q^{-n};q)_j=(-1)^jq^{-nj+\binom j2}
 \frac{(q;q)_n}{(q;q)_{n-j}}.
\]
Consequently the coefficient of $x^j$ in the normalized,
rescaled right-hand side of \eqref{eq:lj-ql-resum} tends to
\[
 \frac{(-1)^ja^jq^{j^2}(aq;q)_n}
 {(aq;q)_j(q;q)_j(q;q)_{n-j}}
 =\frac{(-1)^ja^jq^{j^2}(aq^{j+1};q)_{n-j}}
 {(q;q)_j(q;q)_{n-j}}.
\]
The final expression is precisely the coefficient of $x^j$ in
the $q$-Laguerre monomial expansion stated in V31, immediately below
\texttt{eq:qLlJ-D}. The sum is finite, so coefficientwise convergence
proves \eqref{eq:lj-ql-target}.
\end{proof}
\noindent\textbf{Audit verdict: PASS (with finite resummation).}
The variable scaling prevents taking the limits of the unscaled
connection coefficients in isolation. The exact inversion contained
in the reverse-connection construction resolves this issue. The
specific numbered KLS formula has not been independently matched.


\section{$q$-Meixner to $q$-Laguerre: completed audit}
\noindent\textbf{Source:} V31, Connection 94, reverse connection
\texttt{eq:qLqM-D}. The limiting relation is Koekoek--Swarttouw
(1998), Eq.~(4.13.1).

\begin{equation}\label{eq:qM-qL-target}
\boxed{\lim_{c\to\infty}M_n(cax;a,c;q)
=\frac{(q;q)_n}{(aq;q)_n}L_n(x;a\mid q),\qquad a=q^\alpha.}
\end{equation}
Here $0<q<1$, $n$ is fixed, and the denominators are assumed nonzero.

\begin{proof}
The reverse connection in V31 reads
\begin{align*}
M_n(y;b,c;q)=\sum_{k=0}^n\Biggl[\sum_{r=k}^n\sum_{j=k}^r
&\frac{(-1)^{j+k}(q^{-n};q)_r(-q^{n+1}/c)^r}
{(q,bq;q)_r}q^{\binom j2}\begin{bmatrix}r\\j\end{bmatrix}_q\\[-2pt]
&\quad\times\frac{(q;q)_k(aq^{k+1};q)_{j-k}}
{a^jq^{k^2+[j(j+1)-k(k+1)]/2}}
\frac{(q^{-j};q)_k}{(q^{-k};q)_k}\Biggr]L_k(y;a\mid q).
\end{align*}
We specialize $b=a$. The inner finite Gaussian-binomial inversion
used in the proof of V31's connection expresses the Newton polynomial
$(y;q)_r$ in the $q$-Laguerre basis. Explicitly, if $B_{r,k}$ denotes
the double-sum coefficient remaining after the factor
$(q^{-n};q)_r(-q^{n+1}/c)^r/(q,bq;q)_r$ is removed, then
\[
\sum_{k=0}^r B_{r,k}L_k(y;a\mid q)=(y;q)_r.
\]
This is an exact finite polynomial identity, independent of the
limiting relation. Reordering the finite reverse-connection sums
therefore gives
\begin{equation}\label{eq:qM-qL-newton}
M_n(y;a,c;q)=\sum_{r=0}^n
\frac{(q^{-n};q)_r(-q^{n+1}/c)^r}
{(q;q)_r(aq;q)_r}(y;q)_r.
\end{equation}
We now set $y=cax$. For each fixed $r$, the finite product satisfies
\[
\lim_{c\to\infty}c^{-r}(cax;q)_r
=(-ax)^r q^{\binom r2}.
\]
Consequently, the coefficient of $x^r$ in the limiting polynomial is
\[
\frac{(q^{-n};q)_r a^r
q^{r(n+1)+\binom r2}}{(q;q)_r(aq;q)_r}
=\frac{(-1)^r a^r q^{r^2}(q;q)_n}
{(q;q)_{n-r}(q;q)_r(aq;q)_r},
\]
where we used
$(q^{-n};q)_r=(-1)^rq^{-nr+\binom r2}
(q;q)_n/(q;q)_{n-r}$.
The monomial coefficient of $x^r$ in the $q$-Laguerre polynomial,
as given immediately after the $q$-Laguerre connection formulas
in V31, is
\[
[x^r]L_n(x;a\mid q)
=\frac{(-1)^ra^rq^{r^2}(aq^{r+1};q)_{n-r}}
{(q;q)_r(q;q)_{n-r}}.
\]
Since $(aq^{r+1};q)_{n-r}=(aq;q)_n/(aq;q)_r$,
the limiting coefficient above is exactly
$(q;q)_n/(aq;q)_n$ times this coefficient. All sums are finite,
so coefficientwise convergence proves \eqref{eq:qM-qL-target}.
\end{proof}
\noindent\textbf{Audit verdict: PASS (all fixed degrees).}
The proof starts with the V31 reverse connection, uses only its exact
finite Gaussian-binomial inversion to recover the Newton basis,
and then applies the KLS scaling. It does not assume the desired
limiting relation.


\section{$q$-Bessel (alternative $q$-Charlier) from little $q$-Jacobi: complete reverse-connection limit}
\noindent\textbf{Source:} V31, equation \texttt{eq:qBlqJ-reverse}; the source cites Koekoek--Swarttouw (1998), Eq.~(4.12.3).
Fix $0<q<1$, $a$ such that the relevant leading coefficients are nonzero, and $n\ge0$. Put $B=-a/(Aq)$, so $AB=-a/q$. The claimed limit is
\begin{equation}\label{eq:qB-limit}
\lim_{A\to0}p_n\left(x;A,-\frac{a}{Aq}\mid q\right)=y_n(x;a;q).
\end{equation}
The reverse connection in V31 is
\begin{equation}\label{eq:qB-rev-audit}
 p_n(x;A,B\mid q)=\sum_{k=0}^n C_{n,k}(A,B;a)y_k(x;a;q),
\end{equation}
where
\begin{align*}
C_{n,k}={}&\frac{(-1)^kq^{\binom{k}{2}+k}(q^{-n},ABq^{n+1};q)_k}
{(Aq,q;q)_k(-aq^k;q)_k}\\
&\times{}_3\phi_2\!\left(\begin{matrix}q^{k-n},ABq^{n+k+1},0\\Aq^{k+1},-aq^{2k+1}\end{matrix};q,q\right).
\end{align*}
\begin{proposition}
Under $B=-a/(Aq)$, the reverse-connection coefficients satisfy
$\lim_{A\to0}C_{n,k}(A,-a/(Aq);a)=\delta_{nk}$ for $0\le k\le n$.
Thus \eqref{eq:qB-rev-audit} implies \eqref{eq:qB-limit} for every fixed $n$.
\end{proposition}
\begin{proof}
We check the limit of the \emph{reverse coefficients}, using the exact finite triangular inversion underlying the displayed connection, rather than assuming the desired polynomial limit. Let $T^{(n)}(A)$ be the lower-triangular $(n+1)\times(n+1)$ array of monomial coefficients of the left-hand polynomial basis in \eqref{eq:qB-rev-audit}. In the notation of V31, equation \texttt{eq:qBlqJ-triangular}, its entries are
\[
T_{r,j}(A)=\frac{(q^{-r},-aq^r;q)_j q^j}{(Aq,q;q)_j},\qquad 0\le j\le r\le n.
\]
Let $T^{(n)}(0)$ be the same array at $A=0$. By the explicit $q$-Bessel coefficient array in V31, this is exactly the monomial coefficient matrix of $\{y_r(x;a;q)\}_{r=0}^n$. The diagonal entries are
\[
T_{r,r}(0)=\frac{(q^{-r},-aq^r;q)_r q^r}{(q;q)_r},
\]
which are nonzero under the stated genericity assumptions. Therefore $T^{(n)}(0)$ is invertible.

The defining identity for the \emph{reverse connection matrix} $C^{(n)}(A)$ in \eqref{eq:qB-rev-audit} is the exact finite matrix equation
\[
T^{(n)}(A)=C^{(n)}(A)T^{(n)}(0),\qquad
C^{(n)}(A)=T^{(n)}(A)\bigl(T^{(n)}(0)\bigr)^{-1}.
\]
This is precisely the triangular inversion used in V31 to obtain the terminating ${}_3\phi_2$ expression above; it does not invoke any limiting relation. For every $r,j$, $(Aq;q)_j\to1$, hence $T_{r,j}(A)\to T_{r,j}(0)$. Since the matrices have fixed finite size,
\[
\lim_{A\to0}C^{(n)}(A)
=T^{(n)}(0)\bigl(T^{(n)}(0)\bigr)^{-1}=I_{n+1}.
\]
Consequently each reverse coefficient converges to $\delta_{rk}$. Taking the limit of the finite reverse connection \eqref{eq:qB-rev-audit} gives $y_n(x;a;q)$, as required.
\end{proof}
\noindent\textbf{Audit verdict: PASS (all fixed degrees).} The proof establishes convergence of the actual reverse-connection coefficient matrix; no forward connection or pre-assumed polynomial limit is used. The parameter $B$ diverges as $A\to0$, but only $AB$ occurs in the displayed coefficient matrix, so the specialization is regular.


\section{Ordinary $q$-Meixner to ordinary $q$-Charlier: complete}
The ordinary $q$-Meixner and ordinary $q$-Charlier polynomials are normalized in V31 by
\[
M_n(x;b,c;q)={}_2\phi_1\!\left(\begin{matrix}q^{-n},x\\bq\end{matrix};q,-q^{n+1}/c\right),\qquad
C_n(x;a;q)={}_2\phi_1\!\left(\begin{matrix}q^{-n},x\\0\end{matrix};q,-q^{n+1}/a\right).
\]
The corresponding parameter-deletion limit is
\begin{equation}\label{eq:qM-qC-limit}
\lim_{b\to0}M_n(x;b,a;q)=C_n(x;a;q),\qquad 0<q<1,\ a\ne0.
\end{equation}
The reverse connection of V31, equation \texttt{eq:qC-qM-reverse}, is
\begin{equation}\label{eq:qM-qC-reverse-audit}
M_n(x;b,c;q)=\sum_{k=0}^n {n\brack k}_q\frac{(a/c)^k}{(bq;q)_k}
\left[\sum_{r=0}^{n-k}{n-k\brack r}_q(-1)^rq^{\binom r2}
\frac{(a/c)^r}{(bq^{k+1};q)_r}\right]C_k(x;a;q).
\end{equation}
\begin{proposition}
The reverse connection \eqref{eq:qM-qC-reverse-audit}, specialized to $c=a$, implies \eqref{eq:qM-qC-limit} for every fixed degree $n$.
\end{proposition}
\begin{proof}
Set $c=a$ and denote the coefficient of $C_k$ by $D_{n,k}(b)$. For fixed $n,k,r$, both $(bq;q)_k$ and $(bq^{k+1};q)_r$ tend to $1$ as $b\to0$. Since the sums are finite,
\[
\lim_{b\to0}D_{n,k}(b)
={n\brack k}_q\sum_{r=0}^{n-k}{n-k\brack r}_q(-1)^rq^{\binom r2}.
\]
By the finite $q$-binomial theorem,
\[
\sum_{r=0}^{m}{m\brack r}_q(-1)^rq^{\binom r2}=(1;q)_m
=\begin{cases}1,&m=0,\\0,&m\ge1.\end{cases}
\]
Consequently $\lim_{b\to0}D_{n,k}(b)=\delta_{nk}$. Taking the limit in the finite reverse connection yields $C_n(x;a;q)$, with no extra normalization or variable rescaling. The proof does not use the hypergeometric limit.
\end{proof}
\noindent\textbf{Audit verdict: PASS (all fixed degrees).} This is the ordinary $q$-Charlier family, not the alternative $q$-Charlier/$q$-Bessel family. The precise KLS equation number has not been independently checked.


\section{Standard $q$-Krawtchouk to ordinary $q$-Charlier: complete}
The next eligible distinct-family limiting arrow in V31 is the standard
$q$-Krawtchouk to ordinary $q$-Charlier transition.  In the normalization
used in V31, the limit is
\begin{equation}\label{eq:qK-qC-limit-audit}
\lim_{N\to\infty}K_n(x;q^{-N}/a,N;q)=C_n(x;a;q),
\qquad 0<q<1,\quad a\ne0,\quad n\ \text{fixed}.
\end{equation}
The polynomial variable $x$ is held fixed. The reverse connection in V31,
\texttt{eq:qCqK-reverse}, reads
\begin{equation}\label{eq:qK-qC-reverse-audit}
K_n(x;p,N;q)=\sum_{k=0}^n A_{n,k}(p,N;a)C_k(x;a;q),
\end{equation}
where
\begin{equation}\label{eq:qK-qC-coeff-audit}
A_{n,k}(p,N;a)=\sum_{j=k}^n
\frac{(q^{-n},-pq^n;q)_j q^j}{(q^{-N},q;q)_j}
(-1)^{j-k}q^{\binom{j-k}{2}}
{j\brack k}_q\frac{a^j}{q^{\binom{j+1}{2}}}.
\end{equation}
\begin{proposition}
With $p=q^{-N}/a$, the coefficients in
\eqref{eq:qK-qC-coeff-audit} satisfy
$\lim_{N\to\infty}A_{n,k}(q^{-N}/a,N;a)=\delta_{nk}$.
Consequently the reverse connection proves \eqref{eq:qK-qC-limit-audit}
for every fixed degree $n$.
\end{proposition}
\begin{proof}
For fixed $j$ and $N\to\infty$, the finite $q$-products satisfy
\[
\frac{(-q^{n-N}/a;q)_j}{(q^{-N};q)_j}
\longrightarrow (-1)^j\frac{q^{nj}}{a^j}.
\]
Indeed, factoring the dominant power from each of the $j$ factors in
numerator and denominator gives the displayed limit; the powers
$q^{\binom j2}$ cancel. Thus
\[
\frac{(q^{-n},-q^{n-N}/a;q)_j q^j}{(q^{-N},q;q)_j}
\longrightarrow
\frac{(q^{-n};q)_j}{(q;q)_j}
\frac{(-1)^jq^{(n+1)j}}{a^j}
={n\brack j}_q\frac{q^{\binom{j+1}{2}}}{a^j}.
\]
Substituting in the finite coefficient sum yields
\[
\begin{aligned}
\lim_{N\to\infty}A_{n,k}
&=\sum_{j=k}^n {n\brack j}_q
(-1)^{j-k}q^{\binom{j-k}{2}}{j\brack k}_q
q^{\binom{j+1}{2}-\binom{j+1}{2}}\\
&={n\brack k}_q\sum_{r=0}^{n-k}
{n-k\brack r}_q(-1)^rq^{\binom r2}\\
&={n\brack k}_q(1;q)_{n-k}=\delta_{nk}.
\end{aligned}
\]
Here we used the Gaussian-binomial convolution identity and the finite
$q$-binomial theorem. Since the connection sum has only $n+1$ terms,
we may pass to the limit termwise in \eqref{eq:qK-qC-reverse-audit}.
Only $C_n(x;a;q)$ remains, proving the claimed limit without assuming it.
\end{proof}
\noindent\textbf{Audit verdict: PASS (all fixed degrees).}
The proof takes the limit of the actual reverse coefficients from V31.
The precise numbered KLS reference is not independently verified here.

\section{Standard $q$-Krawtchouk $\to$ alternative $q$-Charlier ($q$-Bessel): completed}
\noindent\textbf{Source:} V31, Connection 97, equation \texttt{eq:qBqK-reverse}. The source cites Koekoek--Swarttouw (1998), Eq.~(4.15.1); that equation-number identification has not been independently verified.

\noindent\textbf{Claim.} For $0<q<1$, fixed $n\geq0$ and fixed admissible $p$,
\begin{equation}\label{eq:audit-qBqK-limit}
\boxed{\displaystyle\lim_{N\to\infty}K_n(q^{-N}z;p,N;q)=y_n(z;p;q).}
\end{equation}
The variable $z$ is fixed, and $N\geq n$ is an integer. No extra normalization is required with the definitions in V31.

\noindent The reverse connection in V31 is
\begin{align}\label{eq:audit-qBqK-reverse}
 K_n(x;p,N;q)=\sum_{k=0}^n\Bigg[&
 \frac{(-1)^kq^{\binom{k}{2}}(1+aq^{2k})}{(q;q)_k}
 \sum_{r=k}^n\frac{(q^{-n},-pq^n;q)_rq^r}
 {(q^{-N},q;q)_r}\notag\\
 &\times\sum_{j=k}^r
 \frac{(-1)^jq^{\binom j2}{r\brack j}_q(q;q)_j}
 {(q;q)_{j-k}(-aq^k;q)_{j+1}}
 \Bigg]y_k(x;a;q).
\end{align}
\begin{proof}
Set $F_r(x)=(x;q)_r$. The innermost sum and its prefactor in
\eqref{eq:audit-qBqK-reverse} are precisely the coefficients $T_{r,k}(a)$ in the finite Newton inversion
\begin{equation}\label{eq:audit-qBqK-inversion}
 F_r(x)=\sum_{k=0}^r T_{r,k}(a)y_k(x;a;q),\qquad
 T_{r,k}(a)=\frac{(-1)^kq^{\binom{k}{2}}(1+aq^{2k})}{(q;q)_k}
 \sum_{j=k}^r\frac{(-1)^jq^{\binom j2}{r\brack j}_q(q;q)_j}
 {(q;q)_{j-k}(-aq^k;q)_{j+1}}.
\end{equation}
This identity is the Newton-basis inversion used to derive the reverse connection in V31. Reordering its finite double sum, without taking any limit, gives
\begin{align}\label{eq:audit-qBqK-newton}
 K_n(x;p,N;q)
 &=\sum_{r=0}^n \frac{(q^{-n},-pq^n;q)_r q^r}
 {(q^{-N},q;q)_r} \sum_{k=0}^r T_{r,k}(a)y_k(x;a;q)\notag\\
 &=\sum_{r=0}^n \frac{(q^{-n},-pq^n;q)_r q^r}
 {(q^{-N},q;q)_r}(x;q)_r.
\end{align}
Thus this is a finite resummation of the \emph{actual reverse connection}, not an invocation of the limit being tested. The auxiliary parameter $a$ cancels from the resummed expression, as it must.

For each fixed $r$, the ratio of Newton factors obeys
\[
\frac{(q^{-N}z;q)_r}{(q^{-N};q)_r}
=\prod_{j=0}^{r-1}\frac{q^N-zq^j}{q^N-q^j}
\xrightarrow[N\to\infty]{}z^r.
\]
The product is empty and equals $1$ when $r=0$. Since $n$ is fixed, we can pass to the limit in the finite sum \eqref{eq:audit-qBqK-newton}:
\[
\lim_{N\to\infty}K_n(q^{-N}z;p,N;q)
=\sum_{r=0}^n\frac{(q^{-n},-pq^n;q)_r}{(q;q)_r}q^r z^r
={}_2\phi_1\!\left(\begin{matrix}q^{-n},-pq^n\\0\end{matrix};q,qz\right)
=y_n(z;p;q).
\]
The last equality is the definition of the alternative $q$-Charlier polynomial given in V31. This proves \eqref{eq:audit-qBqK-limit} for every fixed degree $n$.
\end{proof}
\noindent\textbf{Audit verdict: PASS (all fixed degrees).} The argument resums the V31 reverse connection exactly before applying the scaled-variable limit. The KLS equation-number match remains to be checked independently.


\section{Askey--Wilson $\to$ big $q$-Jacobi: a normalized confluent limit}
\noindent\textbf{Source:} V31, Connection 117, reverse equation \texttt{eq:aw-bigqj-reverse}. Fix $0<q<1$, $n\geq0$, and generic $\alpha,\beta,\gamma$ for which the displayed terminating denominators do not vanish. Define the parameter path
\begin{equation}
 b=\frac{\alpha q}{a},\qquad c=\frac{\gamma q}{a},\qquad d=\frac{\beta a}{\gamma},\qquad a\longrightarrow0.
 \label{eq:audit-aw-bqj-path}
\end{equation}
In particular, $ab=\alpha q$, $ac=\gamma q$, $ad=\beta a^2/\gamma$, and $abcd=\alpha\beta q^2$. Let $\mathcal A_k(x)=p_k(x/(2a);a,b,c,d\mid q)$ and $\mathcal J_k(x)=P_k(x;\alpha,\beta,\gamma\mid q)$, exactly as in V31. The required degree-dependent normalization is
\begin{equation}
\widehat{\mathcal A}_k(x;a)=\frac{a^k}{(ab,ac,ad;q)_k}\mathcal A_k(x).
\label{eq:audit-aw-bqj-norm}
\end{equation}
\par\noindent\textbf{Claim.} 
Along \eqref{eq:audit-aw-bqj-path}, for every fixed $n$,
\begin{equation}
\boxed{\displaystyle\lim_{a\to0}\frac{a^n}{(\alpha q,\gamma q,\beta a^2/\gamma;q)_n}
 p_n\!\left(\frac{x}{2a};a,\frac{\alpha q}{a},\frac{\gamma q}{a},\frac{\beta a}{\gamma}\mid q\right)
 =P_n(x;\alpha,\beta,\gamma\mid q).}
\label{eq:audit-aw-bqj-limit}
\end{equation}
Moreover, the coefficients of the \emph{normalized reverse connection} converge to the identity matrix.
\par
\begin{proof}
The finite quadratic-Newton representation in the proof of V31 Connection 117, which is obtained by inverting its reverse connection, is
\begin{equation}
\widehat{\mathcal A}_k(x;a)=
\sum_{j=0}^k\frac{(q^{-k},\alpha\beta q^{k+1};q)_j q^j}
 {(\alpha q,\gamma q,\beta a^2/\gamma,q;q)_j}
 \prod_{r=0}^{j-1}(1-xq^r+a^2q^{2r}).
\label{eq:audit-aw-bqj-finite}
\end{equation}
Every sum is finite, and its denominators are nonzero near $a=0$ under the stated genericity assumptions. Therefore the coefficientwise limit is
\[
\sum_{j=0}^k\frac{(q^{-k},\alpha\beta q^{k+1};q)_j q^j}
 {(\alpha q,\gamma q,q;q)_j}(x;q)_j=\mathcal J_k(x),
\]
where the last expression is exactly the big $q$-Jacobi Newton expansion in V31.

To check that the conclusion follows \emph{from the reverse connection}, write V31's equation \texttt{eq:aw-bigqj-reverse} as
$\mathcal J_n=\sum_{k\le n}D_{n,k}(a)\mathcal A_k$ and define
$\widehat D_{n,k}(a)=D_{n,k}(a)(ab,ac,ad;q)_k/a^k$.
Then
\begin{equation}
\mathcal J_n(x)=\sum_{k=0}^n\widehat D_{n,k}(a)\widehat{\mathcal A}_k(x;a).
\label{eq:audit-aw-bqj-rev-normalized}
\end{equation}
For any fixed maximum degree $n$, the coefficient matrix $T(a)$ of the polynomials $\widehat{\mathcal A}_k$ in the $\mathcal J_j$ basis is lower triangular. Equation \eqref{eq:audit-aw-bqj-finite} proves entrywise $T(a)\to I$, since the $\mathcal J_j$ are a basis with nonzero leading coefficients. The normalized reverse-connection matrix is precisely $T(a)^{-1}$ by \eqref{eq:audit-aw-bqj-rev-normalized}; finite-dimensional triangular inversion is continuous at $I$. Consequently
\[
\boxed{\lim_{a\to0}\widehat D_{n,k}(a)=\delta_{nk}\qquad(0\le k\le n).}
\]
This derives the limit using the \emph{reverse} connection, its explicit finite Newton inversion, and no assumed Askey--Wilson-to-big-$q$-Jacobi family limit. The diagonal check independently agrees: V31 gives
$D_{n,n}=a^n(\alpha\beta q^{n+1};q)_n/[(\alpha\beta q^{n+1};q)_n(ab,ac,ad;q)_n]$, hence $\widehat D_{n,n}=1$ identically on this path.
\end{proof}
\noindent\textbf{Audit verdict: PASS (all fixed degrees, with explicit normalization).} The algebraic confluent path and its polynomial limit are proved directly. An exact equation-number match to a particular Koekoek--Swarttouw parameterization is \emph{not} claimed without a separate bibliographic check.


\section{Eligibility screening after Connection 117 (Connections 102--116)}
\noindent\textbf{Purpose.} Before assigning another PASS, we screened the remaining later-numbered connections in V31 for a \emph{genuine, non-self, non-exact KLS family limit} that is actually supported by the stated reverse connection. This is an eligibility check, \emph{not} a new limiting-relation proof.

\begin{itemize}
\item \textbf{102 (Al-Salam--Carlitz I / big $q$-Laguerre):} The source gives a finite reverse connection \texttt{eq:ASCI-bL-reverse} and discusses connection-coefficient precedents, but supplies no specific KLS limiting path. \textbf{UNDETERMINED as a limit}; no PASS is claimed.
\item \textbf{103 (Al-Salam--Carlitz I / dual $q$-Krawtchouk):} The source explicitly describes a finite-grid, generally nontriangular interpolation, with lattice $x+s/x$. This is not an established polynomial family limit on a common variable. \textbf{Excluded from limit proofs.}
\item \textbf{104, 107, 111, 115:} Parameter changes or dilations within the same family; \textbf{excluded as self-connections}.
\item \textbf{105, 106, 108:} Their explicit KLS-type confluent limits have already been proved in the cumulative audit (respectively $q$-Meixner to Al-Salam--Carlitz II, quantum $q$-Krawtchouk to Al-Salam--Carlitz II, and $q$-Laguerre to Stieltjes--Wigert). \textbf{Already counted; no duplicate.}
\item \textbf{109, 110:} The source describes monomial/orthogonality projection connections with Stieltjes--Wigert; it does not specify a classical limiting parameter path between the named families. \textbf{Not eligible for a new PASS without a proved limiting path.}
\item \textbf{112, 113, 114:} Exact family specializations or identities (including the essential argument rotation for discrete $q$-Hermite II); \textbf{excluded as non-limiting transformations}.
\item \textbf{116 (Askey--Wilson / $q$-Racah):} The source imposes exact parameter and lattice constraints yielding an exact normalization identity. \textbf{Excluded as an exact transformation}, not a limit.
\end{itemize}
\noindent\textbf{Verdict for this screening:} no additional PASS. The count at the time of this screening was \textbf{32}; the later eligibility correction reduces it to \textbf{31}. The only unresolved candidate flagged here is Connection 102, for which a particular known KLS limit and compatible reverse-coefficient asymptotics would have to be identified and proved. This screening does not assert that the entire earlier manuscript has been exhaustively cleared of further limiting arrows.

\section{Connection 102: testing the zero-parameter degeneration}
\noindent\textbf{Question.} Does the natural path $A,B\to0$ in the reverse connection
\texttt{eq:ASCI-bL-reverse} yield an Al-Salam--Carlitz I polynomial at a
fixed parameter $a$? This is a necessary preliminary check, not an assertion
that Koekoek--Swarttouw prescribes this path.

\begin{proposition}
Fix $0<q<1$ and $n\ge0$. With the big $q$-Laguerre normalization
of V31, the coefficientwise limit is
\begin{equation}\label{eq:audit-102-zero}
\lim_{(A,B)\to(0,0)}P_n(x;A,B\mid q)=x^n.
\end{equation}
For generic $a$, this is \emph{not} $U_n^{(a)}(x;q)$.
\end{proposition}
\begin{proof}
The reverse connection in V31 is obtained by substituting the exact
monomial-to-$U$ expansion into the finite Newton series. Resumming its
finite $U$-coefficients back into the monomial basis (a finite change
of basis, with no limiting interchange) gives
\[
P_n(x;A,B\mid q)=\sum_{j=0}^n
\frac{(q^{-n};q)_j q^j}{(q,Aq,Bq;q)_j}(x;q)_j.
\]
For $A,B\to0$, the denominator factors $(Aq,Bq;q)_j$ tend to $1$.
The terminating $q$-binomial/Newton identity
\[
\sum_{j=0}^n\frac{(q^{-n};q)_j q^j}{(q;q)_j}(x;q)_j=x^n
\]
follows, for example, by applying the finite $q$-binomial theorem
in the Newton basis (or by induction on $n$). Thus
\eqref{eq:audit-102-zero} holds coefficientwise. On the other hand,
comparison of the coefficient of $t$ in the Al-Salam--Carlitz I
generating function in V31 yields
$U_1^{(a)}(x;q)=x-1-a$, which differs from $x$ for $a\ne-1$.
Even the exceptional choice $a=-1$ cannot yield a family-wide
identity: the coefficient of $t^2$ in that generating function gives
$U_2^{(-1)}(x;q)=x^2-(1-q)$ (indeed the constant term is nonzero
for $0<q<1$). Therefore no fixed $a$ makes
$U_n^{(a)}(x;q)=x^n$ for every $n$.
\end{proof}
\noindent\textbf{Audit verdict: NOT A FAMILY LIMIT on the unscaled
$A,B\to0$ path.} This rules out one tempting but incorrect
interpretation of Connection 102. It does not rule out a different
simultaneous parameter/variable scaling; the KLS limiting path
remains \textbf{UNDETERMINED}, and the PASS count is unchanged.


\section{Connection 102: obstruction to fixed affine rescue of the zero-parameter path}
\noindent\textbf{Scope.} The preceding section shows that $A,B\to0$ gives monomials. We now test whether a \emph{fixed, degree-independent affine change of variable} and arbitrary degree-wise scalar normalizations can turn that particular limit into Al-Salam--Carlitz I. This is a necessary-condition result, not a proof or disproof of any other simultaneous parameter scaling.

\begin{proposition}
Fix $0<q<1$ and $a\in\mathbb C$. There are no constants $u\ne0$, $v$ and nonzero scalars $c_n$ such that
\[
 c_n(ux+v)^n=U_n^{(a)}(x;q)\qquad\text{for all }n\ge0.
\]
Consequently, the $A,B\to0$ degeneration of the V31 reverse connection cannot be converted into an Al-Salam--Carlitz I family limit by a fixed affine substitution $x\mapsto ux+v$ and degree-wise scalar normalization alone.
\end{proposition}
\begin{proof}
From the generating function given in V31,
\[
 U_0^{(a)}=1,\qquad U_1^{(a)}=x-(1+a),
\]
\[
 U_2^{(a)}=x^2-(1+q)(1+a)x+q(1+a)^2+a(1-q).
\]
The first two formulas follow by extracting coefficients of $t^0,t^1,t^2$ from
$(t,at;q)_\infty/(xt;q)_\infty$ and multiplying by $(q;q)_n$.
Since the $U_n^{(a)}$ are monic, matching degree one forces $c_1u=1$ and $v/u=-(1+a)$. Matching the leading coefficient in degree two then forces $c_2u^2=1$, so the proposed degree-two polynomial must be
\[
 (x-(1+a))^2=x^2-2(1+a)x+(1+a)^2.
\]
Comparison of the coefficients of $x$ yields
$(1-q)(1+a)=0$. Since $0<q<1$, this requires $a=-1$. But for $a=-1$ the displayed expression for $U_2$ reduces to $x^2-(1-q)$, whereas the proposed square is $x^2$. This is a contradiction.
\end{proof}
\noindent\textbf{Reverse-connection justification.} The previous section resums the \emph{finite} reverse connection \texttt{eq:ASCI-bL-reverse} exactly to the big $q$-Laguerre Newton expansion, proving the monomial limit. Composing that coefficientwise limit with any fixed affine map yields $(ux+v)^n$; the proposition proves that such a limit cannot be the target family. No limit is inferred from an unproved interchange of infinite sums.

\noindent\textbf{Verdict: FAIL for the fixed-affine rescue of the $A,B\to0$ path; UNDETERMINED for a genuinely different KLS scaling.} No additional PASS is counted.


\section{Eligibility correction: Connections 99 and 102}
\noindent\textbf{Scope of the audit.} Only genuine non-self limiting relations may be counted as PASS. A parameter specialization that is already an exact identity at the endpoint does not qualify, even if it can formally be written as a limit.

\begin{proposition}
Connection 99's identity $M_n(x;0,a;q)=C_n(x;a;q)$ is an exact specialization, not a genuine confluent limiting arrow for this audit. Connection 102 has no specified Koekoek--Swarttouw limiting path in V31 and therefore cannot be counted as an outstanding \emph{known} family limit without additional evidence.
\end{proposition}
\begin{proof}
V31's literature paragraph immediately preceding the Connection 99 proposition explicitly calls $M_n(x;0,a;q)=C_n(x;a;q)$ an \emph{exact $b=0$ specialization}, citing Koekoek--Swarttouw (1998), Eq.~(4.13.2). Its reverse equation \texttt{eq:qC-qM-reverse} can be specialized to $b=0$ and $c=a$: its coefficient becomes
\[
\begin{bmatrix}n\\k\end{bmatrix}_q
\sum_{r=0}^{n-k}\begin{bmatrix}n-k\\r\end{bmatrix}_q
(-1)^r q^{\binom r2}
=\begin{bmatrix}n\\k\end{bmatrix}_q(1;q)_{n-k}=\delta_{nk}.
\]
Thus the reverse connection proves the \emph{exact} identity at the endpoint, with no asymptotic analysis necessary. Under the stated protocol it is excluded from the count, notwithstanding its prior PASS label.

For Connection 102, the manuscript's literature context describes an existing \emph{connection-coefficient problem}, expressly declining an equation-level citation. It gives no Koekoek--Swarttouw family-limit formula or path. The earlier audit's zero-parameter and affine-rescue calculations are conditional exclusions of specific paths, not evidence that a published family limit exists. Hence Connection 102 is placed in the \emph{not established as eligible} category; the calculations are retained for traceability, but it is not a pending KLS limiting-arrow proof.
\end{proof}
\noindent\textbf{Corrected tally: 31 genuine limiting PASS entries, one exact specialization removed from the tally.} This correction does not negate the algebraic identity established earlier for Connection 99; it changes only eligibility. Connection 102 remains an open research question about possible other scalings, not an established KLS limit awaiting verification.


\section{Connection 116: exact Askey--Wilson/$q$-Racah identification is not a limit}
\noindent\textbf{Eligibility question.} Could the Askey--Wilson/$q$-Racah relation in V31 be counted as another KLS limiting arrow? The source gives an exact common-lattice specialization, so it must be distinguished from a singular parameter limit.
\begin{proposition}
Assume the parameters satisfy
\[
 a^2=\gamma\delta q,\qquad ab=\alpha q,\qquad
 ac=\beta\delta q,\qquad ad=\gamma q,
\]
and that the relevant shifted factorials are defined and nonzero. Then for each admissible degree $n$,
\[
 p_n\!\left(\frac{Y}{2a};a,b,c,d\mid q\right)
 =a^{-n}(\alpha q,\beta\delta q,\gamma q;q)_n
 R_n(Y;\alpha,\beta,\gamma,\delta\mid q).
\]
This is an exact identity at finite parameters, not a limiting relation.
\end{proposition}
\begin{proof}
The Askey--Wilson terminating ${}_4\phi_3$ has upper parameters
$q^{-n}$, $abcdq^{n-1}$, $az$, and $a/z$, where
$Y=a(z+z^{-1})$. The imposed conditions give
\[
 abcd=\alpha\beta q^2,\quad
 az\cdot a/z=a^2=\gamma\delta q,\quad
 (ab,ac,ad)=(\alpha q,\beta\delta q,\gamma q).
\]
Consequently its terminating hypergeometric series is term-for-term
identical to the defining $q$-Racah ${}_4\phi_3$ in the common variable
$Y=az+a/z$; the remaining factor is exactly the displayed
Askey--Wilson prefactor. Since the equality holds without sending any
parameter to a boundary, it is a finite specialization. In the reverse
connection, uniqueness of the polynomial expansion forces all
non-diagonal coefficients to vanish and leaves only the reciprocal
normalization on the diagonal. This diagonalization is exact and does
not constitute asymptotic convergence.
\end{proof}
\noindent\textbf{Verdict: EXCLUDED (exact transformation).} This verifies the exclusion already noted in the eligibility screening. It adds no limiting PASS and does not change the corrected count of 31.

\section{Cumulative-count consistency correction}
\noindent The historical screening section above was written before
Connection 99 was reclassified. Its statement that the count ``remains
32'' is a snapshot of that earlier tally, \emph{not} the current count.
The eligibility correction for Connection 99 supersedes it: the
current count is \textbf{31 genuine limiting PASS entries}. The
Connection 116 calculation above is an exact specialization and is
not counted. No new limiting proof is claimed in this update.


\section{Dual $q$-Hahn $\longrightarrow$ dual $q$-Krawtchouk: the finite-grid reverse connection}
\noindent\textbf{Source and eligibility.} V31, the subsection
``Dual $q$-Krawtchouk $\leftrightarrow$ Dual $q$-Hahn'', cites
Koekoek--Swarttouw (1998), Eq.~(4.7.2), and gives the reverse
connection \texttt{eq:finite-dKdH-reverse}. This is a genuine singular
parameter limit, rather than an exact finite-parameter identification.
\begin{proposition}
Fix $N\in\mathbb N_0$, $0<q<1$, and a parameter $c$ such that the
$N+1$ quadratic nodes $\nu(x)=q^{-x}+cq^{x-N}$, $0\le x\le N$,
are pairwise distinct and the displayed polynomial normalizations
are defined. Put $s=cq^{-N}$ and let $\gamma\to0$ with
$\delta=s/(\gamma q)$. Then, for $0\le n,x\le N$,
\[
\boxed{\lim_{\gamma\to0}
R_n^{\mathrm{dH}}(\nu(x);\gamma,s/(\gamma q),N\mid q)
=K_n^{\mathrm{dK}}(\nu(x);c,N\mid q).}
\]
More strongly, if $D_{n,k}(\gamma)$ denotes the coefficient in V31's
\emph{reverse} connection \texttt{eq:finite-dKdH-reverse}, with
$t=\gamma\delta q=s$, then
\[
\boxed{\lim_{\gamma\to0}D_{n,k}(\gamma)=\delta_{nk}}
\qquad(0\le n,k\le N).
\]
\end{proposition}
\begin{proof}
Under the stated path the two lattices coincide \emph{for every}
$\gamma\ne0$: $\mu(x)=q^{-x}+tq^x=\nu(x)$. Write the reverse
connection from V31 as
\[
R_n^{\mathrm{dH}}(\nu(x);\gamma,s/(\gamma q),N\mid q)
=\sum_{k=0}^{N}D_{n,k}(\gamma)
K_k^{\mathrm{dK}}(\nu(x);c,N\mid q).
\]
Its coefficients are explicitly the finite double sums
\[
D_{n,k}(\gamma)=\sum_{j=k}^{N}\sum_{r=0}^{j}
\mathscr K^{\mathrm{dH}}_{j,k}(0)\mathscr B_{j,r}(s)
{}_3\phi_2\!\left(\begin{matrix}q^{-n},q^{-r},sq^r\\
q^{-N},\gamma q\end{matrix};q,q\right).
\]
For every $n\le N$, the defining series terminates, and each
$(\gamma q;q)_\ell\to1$ for $0\le\ell\le n$; therefore its
values tend termwise to
\[
{}_3\phi_2\!\left(\begin{matrix}q^{-n},q^{-x},sq^x\\
q^{-N},0\end{matrix};q,q\right)
=K_n^{\mathrm{dK}}(\nu(x);c,N\mid q).
\]
This convergence is at each of the $N+1$ nodes, without invoking
any family-limit theorem. To extract the \emph{reverse coefficients}
rigorously, set $E_{xk}=K_k^{\mathrm{dK}}(\nu(x);c,N\mid q)$ and
$F_{nx}(\gamma)=R_n^{\mathrm{dH}}(\nu(x);\gamma,s/(\gamma q),N\mid q)$.
The degree-$k$ polynomials $K_k^{\mathrm{dK}}$ have nonzero leading
coefficients under the stated nondegeneracy assumptions, and the
nodes are distinct. Hence the square evaluation matrix $E$ is
invertible (a triangular change of basis times a Vandermonde matrix).
The reverse identity gives $F(\gamma)=D(\gamma)E^{\mathsf T}$,
so $D(\gamma)=F(\gamma)(E^{\mathsf T})^{-1}$. The finite-series
calculation above shows $F(\gamma)\to E^{\mathsf T}$, hence
$D(\gamma)\to I$.
Thus the \emph{explicit double-sum coefficients} in the manuscript's
reverse connection tend to $\delta_{nk}$, as asserted.
\end{proof}
\noindent\textbf{Verdict: PASS}, on the full admissible finite grid,
for every $0\le n\le N$. The result uses the actual reverse
connection, not merely a forward limit. The distinct-node hypothesis
can be replaced by a generic-parameter argument and analytic
continuation where appropriate; no such extension is needed for the
stated result. This adds \textbf{one} eligible limiting PASS, raising
the corrected total from 31 to \textbf{32}.


\section{Eligibility check: dual $q$-Hahn and $q$-Racah (V31, Connection  in \texttt{sec:dual-hahn-racah})}
\noindent\textbf{Question.} Does the dual $q$-Hahn/$q$-Racah reverse connection \texttt{eq:finite-dHR-reverse} supply an additional genuine Koekoek--Swarttouw limiting arrow?

\noindent\textbf{Result: EXCLUDED (exact specialization).} The manuscript defines
\[
R_n(\mu(x);\beta,\gamma,\delta,N\mid q)
={}_4\phi_3\!\left(\begin{matrix}q^{-n},\beta q^{n-N},q^{-x},\gamma\delta q^{x+1}\\
q^{-N},\beta\delta q,\gamma q\end{matrix};q,q\right),
\qquad \mu(x)=q^{-x}+\gamma\delta q^{x+1}.
\]
For $0\le n\le N$, setting $\beta=0$ is a regular, \emph{finite parameter substitution}, since $(0;q)_j=1$ for every $j$. Thus
\[
\boxed{R_n(\mu(x);0,\gamma,\delta,N\mid q)
=R_n^{\mathrm{dH}}(\mu(x);\gamma,\delta,N\mid q)}.
\]
This is exactly the identification stated in the literature context of V31's dual $q$-Hahn/$q$-Racah subsection. In the reverse connection
\[
R_n(\mu(x);\beta,\gamma,\delta,N\mid q)
=\sum_{k=0}^n D_{n,k}(\beta)
R_k^{\mathrm{dH}}(\mu(x);\gamma,\delta,N\mid q),
\]
uniqueness of expansion in the degree-ordered dual $q$-Hahn basis yields
$D_{n,k}(0)=\delta_{nk}$, wherever the displayed basis is nondegenerate.
This identity at $\beta=0$ does not require a limiting process or an
argument rescaling. Consequently, it is \emph{not} counted as an additional
eligible non-self limiting relation. This check adds no PASS and leaves the
cumulative total at \textbf{32}. It does not assess other, genuinely
singular $q$-Racah degenerations.


\section{New limiting proof: $q$-Hahn $\to$ ordinary $q$-Krawtchouk (V31, \texttt{eq:finite-KH-reverse})}
\noindent\textbf{Eligibility.} V31 explicitly identifies the confluent $q$-Hahn path $A\to0$ with $AB=-p/q$ fixed. For $0<q<1$, fixed $N\ge0$, $0\le n\le N$, and parameters for which the finite polynomial bases are nondegenerate, the proposed limiting relation is
\[
\boxed{\lim_{A\to0}Q_n\!\left(q^{-x};A,-\frac{p}{Aq},N\mid q\right)
=K_n(q^{-x};p,N\mid q).}
\]
\noindent The reverse connection is precisely
\[
Q_n(q^{-x};A,B,N\mid q)=\sum_{k=0}^n D_{nk}(A,B;p)K_k(q^{-x};p,N\mid q),
\]
with the \emph{explicit} V31 coefficients
\[
\begin{split}
D_{nk}(A,B;p)={}&
\frac{(q^{-n},ABq^{n+1};q)_k}{(q^{-k},-pq^k;q)_k(Aq;q)_k}\\
&\times{}_3\phi_2\!\left(\begin{matrix}q^{k-n},ABq^{n+k+1},0\\Aq^{k+1},-pq^{2k+1}\end{matrix};q,q\right).
\end{split}
\]
\begin{proof}
Write $z=q^{-x}$ and $P_j(z)=(z;q)_j$. These $P_j$ form a degree-ordered polynomial basis with nonzero leading coefficients. The finite Newton expansions quoted and proved in V31 are
\[
Q_n(z;A,B,N\mid q)=\sum_{j=0}^n U_{nj}(A,B)P_j(z),\qquad
U_{nj}(A,B)=\frac{(q^{-n},ABq^{n+1};q)_j q^j}{(q,Aq,q^{-N};q)_j},
\]
and
\[
K_n(z;p,N\mid q)=\sum_{j=0}^n V_{nj}(p)P_j(z),\qquad
V_{nj}(p)=\frac{(q^{-n},-pq^n;q)_j q^j}{(q,q^{-N};q)_j}.
\]
Specialize $B=-p/(Aq)$, so that $ABq^{n+1}=-pq^n$. For every $0\le j\le n\le N$,
\[
U_{nj}\left(A,-\frac{p}{Aq}\right)
=\frac{V_{nj}(p)}{(Aq;q)_j}\longrightarrow V_{nj}(p)
\qquad(A\to0).
\]
All matrices are finite, and the lower-triangular matrix $V=(V_{nj})_{0\le n,j\le N}$ is invertible: its diagonal entries are nonzero by the stated nondegeneracy hypothesis. Comparing Newton coefficients in the \emph{reverse} identity $Q=DK$ gives the exact matrix identity $U=DV$, hence
\[
D\left(A,-\frac{p}{Aq};p\right)
=U\left(A,-\frac{p}{Aq}\right)V^{-1}\longrightarrow VV^{-1}=I.
\]
Thus the coefficients in the explicit reverse connection satisfy
$\lim_{A\to0}D_{nk}(A,-p/(Aq);p)=\delta_{nk}$.
Substituting into the finite reverse sum proves the boxed family limit for every $0\le n\le N$, without assuming that limit or exchanging an infinite series with a limit.
\end{proof}
\noindent\textbf{Verdict: PASS (one new eligible non-self confluent limit).} This is a genuinely singular parameter path ($B\to\infty$ when $p\ne0$), not a finite exact specialization. V31 names the $q$-Hahn-to-ordinary-$q$-Krawtchouk degeneration; a precise numbered Koekoek--Swarttouw equation for this normalization is not independently verified here. The eligible PASS tally increases from \textbf{32} to \textbf{33}.


\section{Coefficient-level verification of the $q$-Hahn limit (follow-up to V47)}
The preceding proof uses the reverse connection and invertibility of the
Newton coefficient matrix. Here we additionally check the explicit
terminating ${}_3\phi_2$ coefficient, including its normalization, rather
than treating the coefficient formula as a black box.

For the confluent path $B=-p/(Aq)$, the coefficient displayed in V31 becomes
\[
D_{nk}(A)=\frac{(q^{-n},-pq^n;q)_k}
{(q^{-k},-pq^k;q)_k(Aq;q)_k}
\sum_{r=0}^{n-k}
\frac{(q^{k-n},-pq^{n+k};q)_r}
{(Aq^{k+1},-pq^{2k+1};q)_r}\frac{q^r}{(q;q)_r}.
\]
The series terminates because $q^{k-n}=q^{-(n-k)}$.
For $k=n$, the sum consists of the single term $r=0$, and hence
\[
\boxed{D_{nn}(A)=\frac1{(Aq;q)_n}\longrightarrow1.}
\]
For $k<n$, the prefactor has a finite limit under the nondegeneracy
assumptions in V47. The matrix argument in V47 therefore yields the
following explicit terminating identity (with $m=n-k>0$):
\[
\boxed{\sum_{r=0}^{m}
\frac{(q^{-m},-pq^{n+k};q)_r}
{(q,-pq^{2k+1};q)_r}\,q^r=0,\qquad n=k+m.}
\]
The vanishing identity follows without invoking an unverified
summation formula: the exact Newton coefficient equality $U=DV$ in V47,
together with $U\to V$ and invertibility of $V$, forces the displayed
finite sum to vanish for every $k<n$. This provides an independent
check of the indexing in the explicit ${}_3\phi_2$ coefficient.
Thus the coefficient-level result is $\lim_{A\to0}D_{nk}(A)=\delta_{nk}$.
The explicit diagonal computation also shows why the normalization cannot
be silently dropped: at finite $A$, the diagonal coefficient is not $1$.
\par\noindent\textbf{Verdict: confirmation of the existing PASS; no additional limiting arrow counted.}


\section{Direct summation of the off-diagonal $q$-Hahn reverse coefficients (V49)}
The preceding coefficient-level check deduced the vanishing of a terminating
${}_2\phi_1$ from matrix inversion. We now verify that cancellation
\emph{directly}, without appealing to the family limit or to the matrix
identity $U=DV$. This removes any possible circularity in the
coefficient-level verification.

\noindent\textbf{Terminating $q$-Chu--Vandermonde identity at $z=q$.}
For $m\in\mathbb Z_{\geq0}$ and parameters for which denominators do not
vanish,
\[
{}_2\phi_1\!\left(\begin{matrix}q^{-m},b\\c\end{matrix};q,q\right)
=\sum_{r=0}^{m}\frac{(q^{-m},b;q)_r}{(q,c;q)_r}q^r
=b^m\frac{(c/b;q)_m}{(c;q)_m}.
\]
The formula extends by rational continuation to all parameter values
where its two sides are defined.

\begin{proof}
This is the terminating $q$-Chu--Vandermonde summation in its
argument-$q$ form. It is a finite rational identity, so no convergence
or infinite-series interchange is involved.
\end{proof}

For the reverse coefficient $D_{nk}(A)$ displayed in the preceding
section, fix $0\leq k<n$ and put $m=n-k$. As $A\to0$, its inner sum tends
to
\[
S_{nk}=\sum_{r=0}^{m}
\frac{(q^{-m},-pq^{n+k};q)_r}{(q,-pq^{2k+1};q)_r}q^r.
\]
Apply the lemma with
\[
b=-pq^{n+k},\qquad c=-pq^{2k+1},\qquad
\frac cb=q^{k+1-n}=q^{1-m}.
\]
For $m\geq1$, the factor $(q^{1-m};q)_m$ contains $(1-q^0)$,
so that
\[
\boxed{S_{nk}=b^m\frac{(q^{1-m};q)_m}{(c;q)_m}=0\quad(k<n).}
\]
Provided the remaining prefactor is finite, it follows directly that
$\lim_{A\to0}D_{nk}(A)=0$ for $k<n$. For $k=n$, the previously
computed exact diagonal coefficient gives
$D_{nn}(A)=(Aq;q)_n^{-1}\to1$. Consequently
\[
\boxed{\lim_{A\to0}D_{nk}(A)=\delta_{nk}\qquad(0\leq k\leq n\leq N).}
\]
This is an explicit, noncircular coefficient-level verification of the
V47 reverse-connection limit. It is valid on the stated nondegenerate
parameter domain; singular parameter values require separate
continuation. \textbf{Verdict: PASS confirmed; no new limiting arrow counted.}


\section{Quantum $q$-Krawtchouk from $q$-Hahn: corrected singular path (V50)}
\noindent\textbf{Eligibility and scaling.} The literature paragraph preceding
V31 equation \texttt{eq:finite-qtmH-reverse} mentions $A\to\infty$
with $AB=p$ fixed. That particular path does \emph{not} produce the
quantum $q$-Krawtchouk Newton coefficients: for $j\ge1$ the quotient
$(ABq^{n+1};q)_j/(Aq;q)_j$ tends to zero. The genuine confluent
path in the displayed normalizations is instead
\[
 A\longrightarrow\infty,\qquad B=p\quad\text{fixed},
 \qquad 0\le n\le N,\quad 0<q<1,\quad p\ne0.
\]
In particular, the two parameter paths must not be conflated.

\noindent\textbf{Proof from the explicit reverse connection.}
Write $D_{nk}(A,p)$ for the coefficient of
$K_k^{\mathrm{qtm}}(q^{-x};p,N\mid q)$ in the V31 reverse
connection \texttt{eq:finite-qtmH-reverse}. For $m=n-k$ it is
\[
 D_{nk}(A,p)=
 \frac{(-1)^k(q^{-n},Apq^{n+1};q)_k}
 {p^kq^{\binom{k}{2}}(q,Aq;q)_k}
 {}_2\phi_1\!\left(\begin{matrix}q^{-m},Apq^{n+k+1}\\
 Aq^{k+1}\end{matrix};q,\frac{q^{-k}}p\right).
\]
Every sum terminates, so passage to the limit is termwise. For fixed
$r$, the elementary finite-product asymptotic
\[
 \lim_{A\to\infty}\frac{(Au;q)_r}{(Av;q)_r}
 =\left(\frac uv\right)^r
\]
holds whenever $uv\ne0$. Consequently the inner series tends to
\[
 \sum_{r=0}^{m}\frac{(q^{-m};q)_r}{(q;q)_r}
 \left(pq^n\frac{q^{-k}}p\right)^r
 =\sum_{r=0}^{m}\frac{(q^{-m};q)_r}{(q;q)_r}q^{mr}
 =(1;q)_m.
\]
Here the last equality is the finite $q$-binomial theorem
$\sum_{r=0}^m(q^{-m};q)_r z^r/(q;q)_r=(zq^{-m};q)_m$.
For $m>0$, $(1;q)_m=0$. The outer prefactor has a finite limit,
so $D_{nk}\to0$ whenever $k<n$. On the diagonal, the terminating
series is $1$ and
\[
 \lim_{A\to\infty}D_{nn}(A,p)
 =\frac{(-1)^n(q^{-n};q)_n}{p^nq^{\binom n2}(q;q)_n}
 \left(pq^n\right)^n=1.
\]
Thus the \emph{reverse-connection coefficients themselves} obey
\[
 \boxed{\lim_{A\to\infty}D_{nk}(A,p)=\delta_{nk}.}
\]
The finite reverse connection now gives, for every $0\le n\le N$,
\[
 \boxed{\lim_{A\to\infty}Q_n(q^{-x};A,p,N\mid q)
 =K_n^{\mathrm{qtm}}(q^{-x};p,N\mid q).}
\]
The argument holds as an identity of polynomials in $q^{-x}$,
for generic nondegenerate parameters and with $A$ avoiding the finite
set of poles. This is a genuine singular degeneration, not an exact
finite specialization. \textbf{Verdict: PASS (new eligible limiting arrow).}
The precise KLS equation number has not been independently verified;
the scaling stated in V31's introductory prose must be corrected if
that prose is intended to describe this normalization.


\section{Degree-one obstruction to the $AB=p$ path in the quantum $q$-Krawtchouk limit (V51)}
\noindent\textbf{Purpose.} The literature-context sentence accompanying V31 equation
\texttt{eq:finite-qtmH-reverse} describes $A\to\infty$ with $AB=p$
fixed. V50 proved that, for the \emph{displayed polynomial normalization},
the successful coefficientwise path is $B=p$ fixed. We now verify
rigorously that the alternative $AB=p$ path fails already in degree one.
This is an obstruction check, not an additional limiting arrow.

Put $y=q^{-x}$ and assume $0<q<1$, $N\ge1$, and $p\ne0$.
The defining terminating ${}_3\phi_2$ and ${}_2\phi_1$ series
in V31 give, exactly,
\begin{align*}
 Q_1(y;A,B,N\mid q)
 &=1-\frac{1-ABq^2}{(1-Aq)(1-q^{-N})}(1-y),\\
 K_1^{\mathrm{qtm}}(y;p,N\mid q)
 &=1-\frac{pq}{1-q^{-N}}(1-y).
\end{align*}
(The common factor $1-y$ comes from $(y;q)_1$, and the
coefficient simplification uses $(1-q^{-1})q/(1-q)=-1$.)
On the proposed path $B=p/A$, the first expression tends to $1$
as $A\to\infty$, because
\[
 \frac{1-pq^2}{1-Aq}\longrightarrow0.
\]
For $p\ne0$ and $N\ge1$, the target polynomial
$K_1^{\mathrm{qtm}}(y;p,N\mid q)$ is nonconstant. Therefore
\[
 \boxed{\lim_{A\to\infty}Q_1(y;A,p/A,N\mid q)=1
 \ne K_1^{\mathrm{qtm}}(y;p,N\mid q).}
\]
In contrast, on the successful path $B=p$ fixed,
\[
 \frac{1-Apq^2}{1-Aq}\longrightarrow pq,
\]
so the degree-one expression agrees with the target, consistently
with the all-degree reverse-coefficient proof in V50.

\noindent\textbf{Verdict: the $AB=p$-fixed path is REFUTED for the
normalizations displayed in V31.} The $B=p$-fixed path remains
\textbf{PASS} by V50. This calculation does not assert that the
literature's original limit is wrong: it may use a different
parameter convention or normalization. Its equation-level mapping
requires independent checking. No additional PASS is counted.

\section{Eligibility check: affine $q$-Krawtchouk versus $q$-Hahn}
\noindent\textbf{Source:} V31, Subsection
\texttt{sec:affine-krawtchouk-hahn}, reverse connection
\texttt{eq:finite-affH-reverse}.
\medskip

The next finite-grid candidate is not a genuine limiting arrow:
its defining identity is the exact specialization
\[
 K_n^{\mathrm{aff}}(y;p,N\mid q)=Q_n(y;p,0,N\mid q),
 \qquad 0\le n\le N.
\]
Indeed the $q$-Hahn terminating ${}_3\phi_2$ has numerator
parameter $ABq^{n+1}$ and denominator parameters $Aq$ and
$q^{-N}$. Setting $A=p$, $B=0$ makes the numerator parameter
zero, yielding term by term the displayed affine $q$-Krawtchouk
${}_3\phi_2$ definition. This is an identity at finite parameter
values, not a boundary degeneration.

The V31 reverse connection is
\[
 Q_n(y;A,B,N\mid q)
 =\sum_{k=0}^n D_{n,k}(A,B;p)
 K_k^{\mathrm{aff}}(y;p,N\mid q).
\]
At $(A,B)=(p,0)$, uniqueness of polynomial expansions in the
nondegenerate affine $q$-Krawtchouk basis gives exactly
\[
 D_{n,k}(p,0;p)=\delta_{nk}.
\]
The identity is not evidence of a new non-self KLS limit, because
there is no singular parameter path and the target family is
already obtained by a finite specialization.

\noindent\textbf{Verdict: EXCLUDED (exact specialization).}
No new limiting proof is counted; the corrected total remains
\textbf{34}. The original V31 manuscript is unchanged.

\section{Connection: affine $q$-Krawtchouk and dual $q$-Hahn --- eligibility test (V53)}
\noindent\textbf{Source:} V31, Section
\texttt{sec:affine-krawtchouk-dual-hahn}, reverse grid connection
\texttt{eq:finite-affdH-reverse}; literature citation K--S (1998),
Eq.~(4.7.1).

The dual $q$-Hahn family in V31 is evaluated at
$z=q^{-x}+tq^x$, where $t=\gamma\delta q$, and has the
terminating representation
\[
 R_n^{\mathrm{dH}}(z;\gamma,\delta,N\mid q)
 ={}_3\phi_2\!\left(\begin{matrix}
 q^{-n},q^{-x},\gamma\delta q^{x+1}\\
 \gamma q,q^{-N}
 \end{matrix};q,q\right).
\]
At the finite specialization $\delta=0$, $\gamma=p$, one has
$t=0$, $z=q^{-x}$, and the third numerator parameter is zero.
Thus, term by term,
\[
 \boxed{R_n^{\mathrm{dH}}(q^{-x};p,0,N\mid q)
 ={}_3\phi_2\!\left(\begin{matrix}
 q^{-n},q^{-x},0\\pq,q^{-N}
 \end{matrix};q,q\right)
 =K_n^{\mathrm{aff}}(q^{-x};p,N\mid q).}
\]
This is an equality at a finite parameter point for all
$0\le n\le N$, assuming nonvanishing denominator factors.
On the $N+1$ distinct lattice points $q^{-x}$, the affine
$q$-Krawtchouk evaluation matrix is invertible. Consequently,
substitution in the manuscript's \emph{reverse} grid connection
\texttt{eq:finite-affdH-reverse} yields exactly
\[
 D_{n,k}(p,0;p)=\delta_{nk},\qquad 0\le n,k\le N.
\]
This is a finite-grid specialization, not a genuine singular
non-self family limit. The general reverse formula is a
nontriangular interpolation identity because the two lattices
are different when $t\ne0$; the specialization $t=0$ makes
those lattices coincide. No additional KLS limiting arrow is
established by this observation.

\noindent\textbf{Verdict: EXCLUDED (exact specialization).}
The eligible PASS total remains \textbf{34}. The original V31
manuscript was not changed.


\section{Eligibility check: $q$-Racah versus ordinary $q$-Krawtchouk (Connection \texttt{eq:finite-KR-reverse})}
\noindent\textbf{Question.} Does the full finite-grid reverse connection
\texttt{eq:finite-KR-reverse} supply an additional genuine KLS limiting
arrow, or is the apparent degeneration an exact parameter specialization?

Let $0<q<1$, $0\le n\le N$, and write the $q$-Racah lattice as
$\mu(x)=q^{-x}+\gamma\delta q^{x+1}$. The defining series in V31 is
\[
R_n(\mu(x);\beta,\gamma,\delta,N\mid q)
={}_4\phi_3\!\left(\begin{matrix}
q^{-n},\beta q^{n-N},q^{-x},\gamma\delta q^{x+1}\\
q^{-N},\beta\delta q,\gamma q
\end{matrix};q,q\right).
\]
Choose the \emph{finite} parameter values
\[
\gamma=\delta=0,\qquad \beta=-pq^N.
\]
Then $\mu(x)=q^{-x}$, $\beta q^{n-N}=-pq^n$, and
all three additional numerator/denominator shifted factorials with
parameter zero are identically one. Therefore the terminating sums
agree term by term:
\[
\boxed{R_n(q^{-x};-pq^N,0,0,N\mid q)
={}_3\phi_2\!\left(\begin{matrix}
q^{-n},q^{-x},-pq^n\\q^{-N},0
\end{matrix};q,q\right)
=K_n(q^{-x};p,N\mid q).}
\]
The equality holds \emph{at the parameter endpoint itself}; no
singular limiting procedure or degree-dependent renormalization is
needed. For generic $p$ with nonzero leading coefficients, the
$q$-Krawtchouk evaluation matrix on the distinct points
$q^{-x}$, $0\le x\le N$, is invertible. Substitution in the
\emph{reverse} connection \texttt{eq:finite-KR-reverse} consequently
gives, at this specialization,
\[
D_{n,k}(-pq^N,0,0;p)=\delta_{nk},\qquad 0\le n,k\le N.
\]
This follows from uniqueness of finite-grid expansion, rather than
from taking a possibly unjustified limit inside its interpolation
kernels. In particular, this equality is an \emph{exact specialization}
and must not be counted as a new genuine limiting relation.

\noindent\textbf{Verdict: EXCLUDED (exact finite-parameter
specialization).} This does not rule out other nontrivial singular
$q$-Racah limits with different scaling. The eligible PASS tally
remains \textbf{34}; no previously claimed proof is added or removed.


\section{Connection 103: eligibility of the Al-Salam--Carlitz I / dual $q$-Krawtchouk grid connection}
\noindent\textbf{Source:} V31, \texttt{eq:ASCI-dK-reverse}. This is a full finite-grid connection,
\[
K_n^{\mathrm{dK}}(q^{-r}+sq^r;c,N\mid q)
=\sum_{k=0}^{N}D_{n,k}(s,a,N)U_k^{(a)}(q^{-r};q),
\qquad s=cq^{-N},\quad 0\le r\le N.
\]
The coefficient $D_{n,k}$ is the explicit finite interpolation sum displayed in V31.
It is \emph{not} a triangular polynomial connection for $s\ne0$: writing
$x=q^{-r}$ gives $q^{-r}+sq^r=x+s/x$, so its left side is generally
Laurent-polynomial in $x$, whereas every $U_k^{(a)}(x;q)$ is polynomial.
The equality holds on the $N+1$ grid points, where any grid function has
an expansion in the $U_k^{(a)}$ basis. Consequently, even the existence
of this reverse matrix cannot by itself establish a family limit as
$N\to\infty$ or $s\to0$.

For example, at $s=0$ the defining series becomes
\[
K_n^{\mathrm{dK}}(x;0,N\mid q)
=\sum_{j=0}^{n}\frac{(q^{-n},x;q)_j}{(q^{-N},q;q)_j}q^j,
\]
which is a well-defined polynomial in $x$, but it is not automatically
an Al-Salam--Carlitz I polynomial. Its degree-one expression is
\[
K_1^{\mathrm{dK}}(x;0,N\mid q)
=1-\frac{1-x}{1-q^{-N}},
\]
whereas $U_1^{(a)}(x;q)=x-(1+a)$. Their leading coefficients differ,
and matching them would require a degree-dependent normalization; even
such a match at degree one would not establish an all-degree limit.

\noindent\textbf{Verdict: NOT ESTABLISHED AS AN ELIGIBLE KLS LIMIT.}
V31 identifies no KLS limiting path for Connection 103, and its full-grid
reverse connection alone does not imply one. No new PASS is awarded.
This is an eligibility screening, not a proof that all possible
simultaneous scalings are impossible.


\section{Eligibility screen for Connections 104--115 (V56)}
\noindent\textbf{Scope.} This screen checks the explicit literature-context statements in V31, Connections 104--115, against the audit's requirement of a \emph{genuine non-self limiting arrow}. It does not assert that no other limiting path can exist. In particular, an arbitrary finite connection between distinct families is not itself evidence of a KLS limit.
\begin{description}
\item[104, 107, 111, 115: same-family connections.] Connection 104 changes the Al-Salam--Carlitz I parameter; 107 changes the Al-Salam--Carlitz II parameter; 111 dilates Stieltjes--Wigert; 115 dilates discrete $q$-Hermite II. Each is a finite same-family connection, not an eligible non-self family limit. \textbf{EXCLUDED.}
\item[105, 106, 108: already audited.] The $q$-Meixner $\to$ Al-Salam--Carlitz II limit (105), quantum $q$-Krawtchouk $\to$ Al-Salam--Carlitz II limit (106), and $q$-Laguerre $\to$ Stieltjes--Wigert limit (108) have existing reverse-connection PASS entries. No duplicate proof is counted.
\item[109, 110: distinct-family finite connections.] Connections 109 and 110 concern Stieltjes--Wigert expansions in alternative and ordinary $q$-Charlier bases, respectively. Their V31 literature paragraphs describe monomial expansions and orthogonality projections but specify no KLS limiting path, parameter scaling, or normalization. A reverse connection between these families alone does not establish a limit. \textbf{NOT ESTABLISHED AS ELIGIBLE.}
\item[112, 113, 114: exact family identifications.] Connection 112 uses $h_n(x;q)=U_n^{(-1)}(x;q)$; 113 is the identical-basis case; 114 uses $\widetilde h_n(x;q)=i^{-n}V_n^{(-1)}(ix;q)$. These are exact finite-parameter identifications (with rotation in 114), not singular limits. \textbf{EXCLUDED.}
\end{description}
\noindent\textbf{Verdict: NO NEW ELIGIBLE LIMIT PROVED.} The 34 previously recorded eligible PASS entries are unchanged. This classification is based on the specific relations described in V31, not a claim that all conceivable asymptotic degenerations of these families have been exhausted.

\section{Eligibility cross-check for Connections 95--101 (V57)}
\noindent\textbf{Scope.} This cross-check reads the family labels and literature-context statements of Connections 95--101 in the unchanged V31 source, distinguishing genuine limits from parameter-change identities and previously audited arrows. It does not treat the existence of finite connection coefficients as evidence of a limiting relation.
\begin{description}
\item[95: $q$-Laguerre / $q$-Laguerre.] The cited Cohl--Costas-Santos--Hwang--Wakhare result concerns a finite parameter-change formula within one family. \textbf{EXCLUDED (self-connection).}
\item[96: little $q$-Jacobi / alternative $q$-Charlier.] The singular little-$q$-Jacobi to $q$-Bessel limit is already proved in the audit from \texttt{eq:qBlqJ-reverse}. \textbf{ALREADY AUDITED.}
\item[97: standard $q$-Krawtchouk / alternative $q$-Charlier.] The large-$N$ rescaled-lattice limit is already proved from \texttt{eq:qBqK-reverse}. \textbf{ALREADY AUDITED.}
\item[98: $q$-Bessel / $q$-Bessel.] The literature notes that the parameter-change problem can be obtained by confluence from little $q$-Jacobi, but both endpoints of the displayed connection are $q$-Bessel. \textbf{EXCLUDED (self-connection).}
\item[99: $q$-Meixner / ordinary $q$-Charlier.] The cited $b=0$ identity is exact; the earlier erroneous limiting PASS was removed in V43. \textbf{EXCLUDED (exact specialization).}
\item[100: ordinary $q$-Charlier / standard $q$-Krawtchouk.] The non-self large-$N$ limit has already been proved using \texttt{eq:qCqK-reverse}. \textbf{ALREADY AUDITED.}
\item[101: ordinary $q$-Charlier / ordinary $q$-Charlier.] The displayed connection is a finite parameter-change identity in one family, obtained from Gaussian-binomial inversion. \textbf{EXCLUDED (self-connection).}
\end{description}
\noindent\textbf{Verdict.} No new eligible non-self KLS limiting arrow arises from Connections 95--101 beyond the three already audited. The number of eligible completed limiting proofs at the V57 screening stage was \textbf{34}. This is an eligibility check, not a new proof and not an exhaustive audit of earlier-numbered connections.


\section{V58: Big $q$-Laguerre $\longrightarrow$ little $q$-Laguerre (Wall)}
\noindent\textbf{Source:} V31, subsection \texttt{subsec:wall-biglaguerre}, reverse connection \texttt{eq:wall-big-reverse}. This is a genuine singular large-parameter limit, not the exact little-$q$-Jacobi specialization. Assume $0<q<1$, $Aq\notin\{q^{-j}:j\ge0\}$, and fix $n$ and $z$. Put $B\to\infty$ and rescale the common argument to $x=Bqz$.

\paragraph{Reverse-connection resummation.}
The manuscript obtains its reverse coefficients by first expressing the big $q$-Laguerre polynomial in the Newton basis $F_r(x)=(x;q)_r$, and then expressing each $F_r$ in the Wall basis. In particular, the displayed reverse connection \texttt{eq:wall-big-reverse} is exactly the finite double sum obtained from
\begin{equation}\label{eq:v58-reverse-resum}
 P_n(x;A,B\mid q)
 =\sum_{r=0}^n\frac{(q^{-n};q)_r q^r}{(Aq,Bq,q;q)_r}F_r(x),
 \qquad
 F_r(x)=\sum_{k=0}^r H_{r,k}(A)w_k(x;A\mid q),
\end{equation}
where $H_{r,k}(A)$ is the terminating ${}_2\phi_1$ coefficient given immediately before \texttt{eq:wall-big-reverse} in V31 (set its Wall parameter $a=A$). Thus the first equality is a \emph{finite resummation of the manuscript's reverse connection}, not an assumption of the desired family limit.

\paragraph{Limit.}
For each fixed $r$, the product quotient satisfies
\begin{equation}\label{eq:v58-product}
 \frac{F_r(Bqz)}{(Bq;q)_r}
 =\prod_{j=0}^{r-1}\frac{1-Bq^{j+1}z}{1-Bq^{j+1}}
 \xrightarrow[B\to\infty]{} z^r.
\end{equation}
Because $r\le n$, the limit may be passed through the finite resummed reverse connection. Consequently,
\begin{align}
\lim_{B\to\infty}P_n(Bqz;A,B\mid q)
 &=\sum_{r=0}^n\frac{(q^{-n};q)_r q^r}{(Aq,q;q)_r}z^r\\
 &={}_2\phi_1\!\left(\begin{matrix}q^{-n},0\\Aq\end{matrix};q,qz\right)
 =w_n(z;A\mid q).
\end{align}
The identity uses $(0;q)_r=1$ and the definitions in the same V31 subsection. No degree-dependent scalar is needed, and the rescaled variable is essential.

\noindent\textbf{Verdict: PASS}, all fixed degrees $n$, for the singular $B\to\infty$, $x=Bqz$ path, derived by finite resummation of \texttt{eq:wall-big-reverse}. The V31 literature context identifies Wall as a limiting case of big $q$-Laguerre; an exact numbered KLS equation has not been independently checked. This adds one eligible limiting proof, increasing the tally from 34 to \textbf{35}.


\section{V59: Big $q$-Jacobi to little $q$-Jacobi (Connection \texttt{eq:ljbj166-reverse})}
\noindent\textbf{Eligibility.} This is a distinct-family singular degeneration: take $C\to\infty$ and simultaneously put $x=Cqz$. Fix $0<q<1$, degree $n$, and generic $A,B$ with $(Aq;q)_r\ne0$ for $0\le r\le n$. The target parameters are $a=A$, $b=B$.

\paragraph{Reverse-connection starting point.}
The original V31 reverse connection \texttt{eq:ljbj166-reverse} expresses $P_n(x;A,B,C\mid q)$ in the little $q$-Jacobi basis. Its derivation in V31 explicitly factors the reverse coefficient matrix through the common Newton-to-monomial and monomial-to-little-$q$-Jacobi inversions (equations \texttt{eq:ljbj166-newton}, \texttt{eq:ljbj166-monomial}, and \texttt{eq:ljbj166-inverse}). Reversing those \emph{finite} changes of basis resums the stated reverse connection exactly as
\begin{equation}\label{eq:v59-resum}
 P_n(x;A,B,C\mid q)=\sum_{r=0}^{n}
 \frac{(q^{-n},ABq^{n+1};q)_r q^r}{(Aq,q;q)_r}
 \frac{(x;q)_r}{(Cq;q)_r}.
\end{equation}
This is an algebraic resummation of the reverse connection; no limiting identity is assumed.

\paragraph{Fixed-degree limit.}
For every fixed $r\le n$,
\begin{equation}\label{eq:v59-ratio}
 \frac{(Cqz;q)_r}{(Cq;q)_r}
 =\prod_{j=0}^{r-1}\frac{1-Cq^{j+1}z}{1-Cq^{j+1}}
 \xrightarrow[C\to\infty]{}z^r.
\end{equation}
The $r=0$ product equals $1$. As the sum in \eqref{eq:v59-resum} terminates at $n$, termwise convergence is valid, coefficientwise in $z$. Therefore
\begin{align}
 \lim_{C\to\infty}P_n(Cqz;A,B,C\mid q)
 &=\sum_{r=0}^{n}\frac{(q^{-n},ABq^{n+1};q)_r q^r}{(Aq,q;q)_r}z^r\\
 &={}_2\phi_1\!\left(\begin{matrix}q^{-n},ABq^{n+1}\\Aq\end{matrix};q,qz\right)
 =p_n(z;A,B\mid q).
\end{align}
The last equality is exactly the little $q$-Jacobi definition in the same V31 subsection. The normalization is $1$; the argument scaling is indispensable. This also proves convergence of the coefficients of the resummed reverse connection after scaling, without assuming convergence of its unscaled basis polynomials.

\noindent\textbf{Verdict: PASS}, all fixed degrees $n$, for the genuine singular $C\to\infty$, $x=Cqz$ degeneration. The precise KLS equation number for this parameterization has not been independently confirmed. The eligible PASS tally increases from \textbf{35} to \textbf{36}.

\section{Big $q$-Jacobi $\to$ $q$-Meixner: singular confluent limit (V60)}
\noindent\textbf{Source:} V31, subsection \texttt{meixner-big-jacobi-169}, reverse connection \texttt{eq:169-reverse}. Fix $0<q<1$, degree $n$, and admissible $b,c$ with $bc\ne0$ and $(bq;q)_r\ne0$ for $r\le n$. Set
\begin{equation}\label{eq:v60-path}
 A=b,\qquad B=-\frac{C}{bc},\qquad C\longrightarrow\infty.
\end{equation}
The parameter $AB=-C/c$ diverges while $AB/C=-1/c$ stays fixed. No variable rescaling or extra normalization is needed.

\paragraph{Reverse-connection calculation.}
The V31 reverse connection is $P_n(x;A,B,C\mid q)=\sum_{k=0}^n D_{n,k}(A,B,C;b,c)M_k(x;b,c;q)$, with
\begin{equation}\label{eq:v60-D}
 D_{n,k}=\sum_{r=k}^n
 \frac{(q^{-n},ABq^{n+1};q)_rq^r}{(q,Aq,Cq;q)_r}
 (-1)^{r-k}q^{\binom{r-k}{2}}{r\brack k}_q
 \frac{c^r(bq;q)_r}{q^{r(r+1)/2}}.
\end{equation}
For fixed $r$, using $AB=-C/c$ and the leading powers of $C$ in the finite products,
\begin{equation}\label{eq:v60-ratio}
 \lim_{C\to\infty}\frac{(ABq^{n+1};q)_r}{(Cq;q)_r}
 =\left(-\frac1c\right)^r q^{nr}.
\end{equation}
Indeed, each factor ratio $(1+Cq^{n+1+j}/c)/(1-Cq^{1+j})$ tends to $-q^n/c$. Hence the limit of the first factor in the $r$-summand of \eqref{eq:v60-D} is
\begin{equation}
 \frac{(q^{-n};q)_r}{(q,bq;q)_r}\left(-\frac{q^{n+1}}c\right)^r.
\end{equation}
These are exactly the Newton coefficients of $M_n(x;b,c;q)$ in V31, equation \texttt{eq:169-def}. To verify the reverse-connection coefficients rather than merely asserting the polynomial limit, denote the Newton coefficients of $M_k$ by
\begin{equation}
 L_{k,j}=\frac{(q^{-k};q)_j}{(q,bq;q)_j}\left(-\frac{q^{k+1}}c\right)^j.
\end{equation}
The finite inversion in V31, from which \eqref{eq:v60-D} was derived, gives for every $0\le j\le n$ the exact coefficient identity
\begin{equation}\label{eq:v60-inversion}
 H_{n,j}(C)=\sum_{k=j}^n D_{n,k}(C)L_{k,j},\qquad
 H_{n,j}(C)=\frac{(q^{-n},-Cq^{n+1}/c;q)_jq^j}{(q,bq,Cq;q)_j}.
\end{equation}
By \eqref{eq:v60-ratio}, $H_{n,j}(C)\to L_{n,j}$. The lower-triangular matrix $(L_{k,j})_{0\le j\le k\le n}$ is invertible: its diagonal entries are nonzero for admissible $b,c$ and $0<q<1$. Thus the unique coefficients solving \eqref{eq:v60-inversion} satisfy
\begin{equation}\label{eq:v60-delta}
 \boxed{\displaystyle\lim_{C\to\infty}D_{n,k}\left(b,-\frac{C}{bc},C;b,c\right)=\delta_{nk}.}
\end{equation}
Inserting \eqref{eq:v60-delta} into the finite reverse connection proves, coefficientwise in $x$,
\begin{equation}\label{eq:v60-limit}
 \boxed{\displaystyle\lim_{C\to\infty}P_n\left(x;b,-\frac{C}{bc},C\mid q\right)=M_n(x;b,c;q).}
\end{equation}
This is a genuine singular parameter degeneration, not a finite specialization. The exact numbered Koekoek--Swarttouw reference for this parameterization remains unverified. \textbf{Verdict: PASS}, every fixed degree $n$. The eligible PASS count rises from 36 to \textbf{37}.


\section{V61: $q$-Hahn $\longrightarrow$ $q$-Meixner via the reverse connection}
\noindent\textbf{Source:} V31, equation \texttt{eq:168-D}, the reverse $q$-Hahn-to-$q$-Meixner connection. Fix $0<q<1$, $n\ge0$, $b,c\ne0$ and assume $(bq;q)_j\ne0$ for $0\le j\le n$. Let $N\to\infty$ through integers with $N\ge n$ and specialize
\begin{equation}\label{eq:v61-path}
 A=b,\qquad B_N=-\frac{q^{-N-1}}{bc},\qquad AB_Nq^{n+1}=-\frac{q^{n-N}}c.
\end{equation}
The reverse connection is the finite polynomial identity
\begin{equation}\label{eq:v61-reverse}
 Q_n(x;b,B_N,N\mid q)=\sum_{k=0}^n D_{n,k}(N)M_k(x;b,c;q),
\end{equation}
with $D_{n,k}(N)$ given explicitly in V31, equation \texttt{eq:168-D}. We prove the coefficient limit $D_{n,k}(N)\to\delta_{nk}$, rather than assuming the polynomial limit.

\noindent\textbf{Step 1: Newton coefficients.} Put $\mathcal N_j(x)=(x;q)_j$. The finite terminating series gives
\begin{equation}\label{eq:v61-H}
 Q_n(x;b,B_N,N\mid q)=\sum_{j=0}^n H_{n,j}(N)\mathcal N_j(x),\quad
 H_{n,j}(N)=\frac{(q^{-n},-q^{n-N}/c;q)_j q^j}{(q,bq,q^{-N};q)_j}.
\end{equation}
For every fixed $j$, the ratio of finite products obeys
\begin{equation}\label{eq:v61-ratio}
 \frac{(-q^{n-N}/c;q)_j}{(q^{-N};q)_j}
 =\prod_{i=0}^{j-1}\frac{1+q^{n-N+i}/c}{1-q^{-N+i}}
 \longrightarrow \left(-\frac{q^n}{c}\right)^j.
\end{equation}
Consequently,
\begin{equation}\label{eq:v61-L}
 H_{n,j}(N)\longrightarrow L_{n,j}:=
 \frac{(q^{-n};q)_j}{(q,bq;q)_j}\left(-\frac{q^{n+1}}c\right)^j,
\end{equation}
which is the coefficient of $\mathcal N_j(x)$ in $M_n(x;b,c;q)$ according to the definition in V31.

\noindent\textbf{Step 2: use the reverse connection.} Expanding every $M_k$ in the same Newton basis and comparing coefficients in \eqref{eq:v61-reverse} yields the \emph{exact} triangular system
\begin{equation}\label{eq:v61-system}
 H_{n,j}(N)=\sum_{k=j}^n D_{n,k}(N)L_{k,j}\qquad(0\le j\le n).
\end{equation}
The diagonal $L_{k,k}$ is nonzero under the stated assumptions (indeed $(q^{-k};q)_k\ne0$). Thus the fixed finite triangular system has a unique solution and its inverse is independent of $N$. Passing to the limit in \eqref{eq:v61-system} and using \eqref{eq:v61-L} gives
\begin{equation}\label{eq:v61-delta}
 \boxed{\displaystyle\lim_{N\to\infty}D_{n,k}(N)=\delta_{nk}\quad(0\le k\le n).}
\end{equation}
Substituting into the \emph{reverse connection} \eqref{eq:v61-reverse} proves the fixed-degree polynomial limit
\begin{equation}\label{eq:v61-limit}
 \boxed{\displaystyle\lim_{N\to\infty}
 Q_n\left(x;b,-\frac{q^{-N-1}}{bc},N\mid q\right)=M_n(x;b,c;q).}
\end{equation}
This is a genuine singular $N\to\infty$ degeneration; no variable rescaling or additional normalization is required for the displayed conventions. The source V31 calls the classical $q$-Hahn-to-$q$-Meixner limit known, but an exact numbered Koekoek--Swarttouw formula for this parameterization has not been independently checked. \textbf{Verdict: PASS} for each fixed $n$; eligible tally increases from 37 to \textbf{38}.


\section{V62: $q$-Hahn $\longrightarrow$ little $q$-Jacobi by rescaling the lattice}
\noindent\textbf{Source and eligibility.} V31, reverse connection
\texttt{eq:ljqh-reverse} in subsection \texttt{lj-qh}, writes the $q$-Hahn
polynomial in the little $q$-Jacobi basis. The accompanying literature
paragraph identifies the exact $q$-Hahn/big $q$-Jacobi parameter
specialization; the \emph{different} $N\to\infty$ transition proved below
is singular and is not that finite specialization. We make no claim of
independent verification of a numbered KLS formula.

\noindent Fix $0<q<1$, $A,B$ generic with $(Aq;q)_r\ne0$ for
$0\le r\le n$, and fixed degree $n$. Let $N\ge n$. The reverse
connection of V31 is
\begin{equation}\label{eq:v62-reverse}
 Q_n(x;A,B,N\mid q)=\sum_{k=0}^n
 D_{n,k}(A,B,N;a,b)\,p_k(x;a,b\mid q),
\end{equation}
where $D_{n,k}$ is the explicit finite double sum in
\texttt{eq:ljqh-reverse}. Its derivation in V31 expands the $q$-Hahn
polynomial in the Newton factors $(x;q)_r$ and inverts the
little-$q$-Jacobi Newton coefficient matrix. In particular, the
\emph{exact finite resummation} of the right side of
\eqref{eq:v62-reverse} is
\begin{equation}\label{eq:v62-newton}
 \sum_{k=0}^n D_{n,k}\,p_k(x;a,b\mid q)
 =\sum_{r=0}^n
 \frac{(q^{-n},ABq^{n+1};q)_r}{(Aq,q^{-N},q;q)_r}
 q^r(x;q)_r.
\end{equation}
This is the same finite Newton-basis identity used to construct the
reverse coefficients in V31; it holds for generic independent $(a,b)$.

\noindent\textbf{Asymptotics.} Set $x=q^{-N}z$ in the \emph{reverse
connection} and use its finite resummation \eqref{eq:v62-newton}.
For each fixed $r$,
\begin{equation}\label{eq:v62-ratio}
 \frac{(q^{-N}z;q)_r}{(q^{-N};q)_r}
 =\prod_{j=0}^{r-1}\frac{1-zq^{j-N}}{1-q^{j-N}}
 \xrightarrow[N\to\infty]{}z^r.
\end{equation}
Since $n$ is fixed, the sum has finitely many terms. Therefore
\begin{align}\label{eq:v62-limit}
 \lim_{N\to\infty}Q_n(q^{-N}z;A,B,N\mid q)
 &=\sum_{r=0}^n\frac{(q^{-n},ABq^{n+1};q)_r}
 {(Aq,q;q)_r}q^rz^r\\
 &=p_n(z;A,B\mid q).\notag
\end{align}
The final equality is exactly the terminating ${}_2\phi_1$
definition of little $q$-Jacobi in V31. No additional scalar
normalization is needed. The unresummed terms
$p_k(q^{-N}z;a,b\mid q)$ generally diverge, so a termwise limit of
the original reverse sum would be invalid; the \emph{finite
resummation before passage to the limit} is essential. This proves
the relation from the reverse connection rather than assuming it.
\textbf{Verdict: PASS} for every fixed $n$ and generic admissible
parameters. Eligible tally increases from 38 to \textbf{39}.


\section{V63: Big $q$-Jacobi / $q$-Racah reverse connection --- eligibility check}
\noindent\textbf{Source:} V31, subsection \texttt{subsec:171-bigjacobi-racah}, reverse equation \texttt{eq:171-reverse}. The reverse connection expands the $q$-Racah polynomial $R_n(y)$ in the big $q$-Jacobi basis $B_k(y)$. It is a finite triangular change of basis at generic parameters; the literature paragraph does not specify a Koekoek--Swarttouw limit for this pair.

\paragraph{The relevant distinction.}
The reverse equation uses the polynomial extension
\[
R_n(y)=\sum_{j=0}^n
\frac{(q^{-n},\alpha\beta q^{n+1};q)_j q^j}
{(q,\alpha q,\beta\delta q,\gamma q;q)_j}
\prod_{r=0}^{j-1}(1-q^r y+\gamma\delta q^{2r+1}),
\]
whereas $B_n(y)$ has Newton factors $(y;q)_j$. The ordinary $q$-Racah orthogonality grid is $y=q^{-s}+\gamma\delta q^{s+1}$; on that grid $R_n$ has its usual terminating ${}_4\phi_3$ form. A parameter degeneration must specify both the grid and the normalization before the reverse coefficients can be tested.

\paragraph{A rigorous endpoint check (not an eligible limit).}
If $\gamma\delta=0$, the quadratic Newton factors reduce \emph{exactly} to $(y;q)_j$. In particular, choosing $\delta=0$ and $\gamma q=aq$, $\alpha q=cq$, and $\alpha\beta=ab$ gives
\[
R_n(y)=\sum_{j=0}^n
\frac{(q^{-n},abq^{n+1};q)_j q^j}
{(q,cq,0,aq;q)_j}(y;q)_j
={}_3\phi_2\!\left(\begin{matrix}q^{-n},abq^{n+1},y\\aq,cq\end{matrix};q,q\right)=B_n(y).
\]
Here $(0;q)_j=1$. Thus this parameter path is an \emph{exact finite-parameter specialization}, not a genuine KLS limiting arrow. For nondegenerate polynomial bases, uniqueness in the reverse connection then gives $D_{n,k}=\delta_{nk}$ at this endpoint. This calculation does not establish any different singular scaling.

\noindent\textbf{Verdict: NOT ESTABLISHED AS AN ELIGIBLE KLS LIMIT.} The obvious matching endpoint is exact, so it is excluded. A separately specified singular KLS path would require its own reverse-coefficient asymptotics. \textbf{No new PASS}; the eligible count remains \textbf{39}.


\section{V64: Hahn $\longrightarrow$ Jacobi from the reverse connection}
\noindent\textbf{Source:} V31, equation \texttt{eq:jacobi-hahn-reverse} (Jacobi/Hahn connection). \textbf{Verdict: PASS} for each fixed $n$ and generic parameters $\alpha,\beta$, with $N\to\infty$ and $x=Nt$. This is a genuine large-grid limit, not a finite specialization. The equation-level KLS bibliographic match is not independently certified.

\paragraph{Claim and normalization.} For $0< t<1$ (in fact, coefficientwise in $t$),
\begin{equation}\label{eq:v64-hahn-jacobi}
\boxed{\displaystyle\lim_{N\to\infty}Q_n(Nt;\alpha,\beta,N)
=\frac{n!}{(\alpha+1)_n}P_n^{(\alpha,\beta)}(1-2t).}
\end{equation}
Assume $(\alpha+1)_r\ne0$ for $0\le r\le n$, and generic parameters so all displayed finite connection coefficients are defined.

\paragraph{Derivation using the manuscript's reverse connection.}
Equation \texttt{eq:jacobi-hahn-reverse} is obtained by expanding the Hahn polynomial in falling factorials, then applying the exact finite Stirling-number conversion and the Jacobi monomial expansion. In particular, its underlying \emph{reverse-side} Newton expansion (explicitly displayed in the V31 proof of that equation) is
\begin{equation}\label{eq:v64-newton}
Q_n(x;\alpha,\beta,N)=\sum_{r=0}^n
\frac{(-n)_r(n+\alpha+\beta+1)_r}{(\alpha+1)_r(-N)_r\,r!}(-x)_r.
\end{equation}
This is not a postulated limit: it is the finite polynomial expansion used to construct the reverse coefficients. To recover the reverse connection exactly, use
$(-x)_r=\sum_{j=0}^r(-1)^r s(r,j)x^j$ and the V31 finite Jacobi expansion of $x^j$; exchanging the finite sums gives \texttt{eq:jacobi-hahn-reverse}.

Now substitute $x=Nt$. For each fixed $r$,
\[
\frac{(-Nt)_r}{(-N)_r}
=\prod_{j=0}^{r-1}\frac{Nt-j}{N-j}\longrightarrow t^r
\quad(N\to\infty).
\]
Hence the finite \emph{reverse-connection Newton system} gives
\[
\lim_{N\to\infty}Q_n(Nt;\alpha,\beta,N)
=\sum_{r=0}^n\frac{(-n)_r(n+\alpha+\beta+1)_r}
{(\alpha+1)_r\,r!}t^r
={}_2F_1\!\left(\begin{matrix}-n,n+\alpha+\beta+1\\\alpha+1\end{matrix};t\right).
\]
The Jacobi normalization stated in V31 is
\[
P_n^{(\alpha,\beta)}(1-2t)=\frac{(\alpha+1)_n}{n!}
{}_2F_1\!\left(\begin{matrix}-n,n+\alpha+\beta+1\\\alpha+1\end{matrix};t\right),
\]
which proves \eqref{eq:v64-hahn-jacobi}.

\paragraph{Why this is genuinely a reverse-connection limit.}
The V31 reverse equation uses the unscaled basis $P_m^{(\alpha,\beta)}(x)$, whose individual terms can grow as $N^m$ when $x=Nt$. Thus it would be \emph{invalid} to take the limit of its original terms separately. Instead we have performed the exact finite Stirling/Jacobi resummation \emph{back into the Newton system from which that reverse connection was derived}, then passed to the limit of each fixed-degree Newton factor. This supplies an all-degree derivation from the reverse connection, with no assumption of the family limit. It does not assert that the original unscaled reverse coefficients tend to a Kronecker delta.

\noindent\textbf{Audit count:} one additional eligible PASS, raising the cumulative count from \textbf{39} to \textbf{40}.


\section{V65: Eligibility review of the Jacobi--generalised Bessel reverse connection}
\noindent\textbf{Source:} V31, equation \texttt{eq:jacobi-gb-reverse}, in the subsection ``Generalised Bessel $\leftrightarrow$ Jacobi.'' This is a distinct-family, triangular connection, so it is worth screening for a genuine published limiting arrow; distinct families alone do not make a connection a limit.

The source defines the generalised Bessel polynomial by the terminating expression
\[
\mathrm{GB}_n(x;\gamma)=\sum_{j=0}^{n}
\frac{(-n)_j(n+\gamma+1)_j}{j!}\left(-\frac{x}{2}\right)^j,
\]
and displays the exact reverse expansion
\[
P_n^{(\alpha,\beta)}(x)=\sum_{k=0}^{n}
D_{n,k}^{\mathrm{J,GB}}(\alpha,\beta,\gamma)
\mathrm{GB}_{n-k}(x;\gamma).
\]
The coefficient $D_{n,k}^{\mathrm{J,GB}}$ is the finite double sum given explicitly in V31. Neither the source's literature note nor the reverse equation specifies a singular parameter path, argument scaling, or normalization that produces a known Koekoek--Swarttouw limit from Jacobi to this generalised Bessel family. In particular, the existence of this finite change of basis does not establish such a limit.

\paragraph{Diagnostic: the degree-one families are not identical at generic parameters.}
The defining series gives $\mathrm{GB}_1(x;\gamma)=1+(\gamma+2)x/2$, while the Jacobi convention gives
\[
P_1^{(\alpha,\beta)}(x)=\frac{\alpha-\beta}{2}+\frac{\alpha+\beta+2}{2}x.
\]
Thus the reverse expansion is a genuine change of basis for generic parameters, but these degree-one expressions alone provide no evidence for a limiting arrow. Any claimed limit would need an explicit parameter regime and all-degree proof using the displayed $D_{n,k}^{\mathrm{J,GB}}$ (or an exact resummation thereof).

\noindent\textbf{Verdict: NOT ESTABLISHED AS AN ELIGIBLE KLS LIMIT.} No new PASS is awarded. This is an eligibility assessment, not a proof that every possible degeneration is impossible. The cumulative eligible PASS tally remains \textbf{40}.


\section{V66: Meixner--Pollaczek and continuous dual Hahn --- exact quadratic endpoint}
\noindent\textbf{Source:} V31, reverse equation \texttt{eq:consolidated-cdh-mp}, with its nonsingular finite-sum representation \texttt{eq:matched-single-sum}. This is a separate eligibility check, not a newly claimed limiting proof.

For $\varepsilon\in\{0,1\}$, the reverse expansion in V31 expresses
$x^\varepsilon S_n(x^2;a,b,c)$ in the basis
$P_{2j+\varepsilon}^{(a)}(x;\pi/2)$. The potentially singular terminating ${}_3F_2$ coefficient has the exact finite-sum representation
\[
T_{n,k}(a,b,c;\varepsilon)=\sum_{j=0}^{k}\binom{k}{j}
(a+c+n-j)_j(b-\tfrac12)_j
(a+\tfrac12+n-k)_{k-j}(c-\varepsilon)_{k-j}.
\]
At the matched finite parameter values $b=\tfrac12$, $c=\varepsilon$, the factor $(b-\tfrac12)_j=(0)_j$ kills all $j>0$. For the surviving $j=0$ term, $(c-\varepsilon)_k=(0)_k$ vanishes for every $k>0$. Thus
\[
\boxed{T_{n,k}(a,\tfrac12,\varepsilon;\varepsilon)=\delta_{k0}.}
\]
The remaining $k=0$ coefficient in the reverse equation is the exact quadratic normalization, yielding
\[
\boxed{
 x^\varepsilon S_n(x^2;a,\tfrac12,\varepsilon)
 =\frac{(-1)^n(2n+\varepsilon)!}{2^{2n+\varepsilon}}
 P_{2n+\varepsilon}^{(a)}(x;\pi/2).
}
\]
This includes both the even ($\varepsilon=0$) and odd ($\varepsilon=1$) cases; the factor $x$ in the odd case cannot be omitted. These are the exact Koornwinder quadratic transformations already stated in V31. The coefficient collapse occurs at finite parameter values and requires no singular asymptotic regime.

\noindent\textbf{Verdict: EXCLUDED as an additional KLS limiting arrow.} This exact quadratic endpoint is not a genuine non-self limiting transition under the audit rules. The result does not assert that no other scaling could ever exist; it simply does not supply an eligible limit. The completed eligible PASS count remains \textbf{40}.


\section{V67: Continuous $q$-Jacobi $\to$ continuous $q$-Hermite from the reverse connection}
\noindent\textbf{Source and eligibility.} V31, equations \texttt{eq:JHforward} and \texttt{eq:JacobiHermiteCoeffs}, express the Askey--Wilson-normalized continuous $q$-Jacobi polynomial in the continuous $q$-Hermite basis. This is a non-self, singular-parameter degeneration, not an exact finite-parameter identification. Let $0<q<1$, $n$ be fixed, and set
\[
 A=q^{\alpha/2+1/4},\quad B=q^{\beta/2+1/4},\quad
 (a,b,c,d)=(A,Aq^{1/2},-B,-Bq^{1/2}).
\]
Take $\alpha,\beta\to+\infty$, so $A,B\to0$. The V31 normalization is
\[
 J_n(x;\alpha,\beta\mid q)
 =p_n(x;A,Aq^{1/2},-B,-Bq^{1/2}\mid q).
\]
\noindent\textbf{Claim (PASS).} For every fixed nonnegative integer $n$ and fixed $x$,
\begin{equation}\label{eq:audit-v67-jh}
 \boxed{\lim_{\alpha,\beta\to+\infty}J_n(x;\alpha,\beta\mid q)
 =H_n(x\mid q).}
\end{equation}
\begin{proof}
The \emph{reverse} connection in V31 reads
\[
 J_n(x;\alpha,\beta\mid q)
 =\sum_{k=0}^n\mathcal H_{n,k}(A,B;q)H_k(x\mid q),
 \qquad
 \mathcal H_{n,k}=\sum_{i=k}^n\sum_{j=k}^i
 \widetilde X_{n,i}\widetilde W_{i,j}G_{j,k}.
\]
The three kernels are defined explicitly in V31 immediately after
\texttt{eq:JHforward}. Put $r=n-i$, $s=i-j$, and $t=j-k$.
Their respective degree-lowering factors are
\[
 \widetilde X_{n,i}=(-d)^r\begin{bmatrix}n\\i\end{bmatrix}_q
 q^{\binom r2}(abq^i,acq^i,bcq^i;q)_r
 (abcdq^{n-1};q)_i,
\]
\[
 \widetilde W_{i,j}=\begin{bmatrix}i\\j\end{bmatrix}_q
 (-c)^s q^{\binom s2}(abq^j;q)_s.
\]
For the last kernel, the inverse generating-function multiplier
$(at,bt;q)_\infty$ (as established in the V31 proof) gives the
finite, denominator-free formula
\[
 G_{j,k}=\begin{bmatrix}j\\k\end{bmatrix}_q
 \sum_{\ell=0}^{t}\begin{bmatrix}t\\\ell\end{bmatrix}_q
 (-a)^\ell(-b)^{t-\ell}
 q^{\binom\ell2+\binom{t-\ell}2}.
\]
Hence $G_{j,k}$ is homogeneous of degree $t$ in $a,b$; in particular,
$G_{j,k}\to0$ for $j>k$, while $G_{k,k}=1$.
As $A,B\to0$, the shifted-factorial factors above tend to $1$.
The factor $d^r$ kills terms with $i<n$, the factor $c^s$
kills terms with $j<i$, and $G_{j,k}$ kills terms with $k<j$.
For $k<n$, at least one of $r,s,t$ is positive, so every term in
the finite double sum tends to zero. For $k=n$, only $i=j=k=n$
occurs and
\[
 \mathcal H_{n,n}=(abcdq^{n-1};q)_n
 =(A^2B^2q^n;q)_n\longrightarrow1.
\]
Thus $\mathcal H_{n,k}\to\delta_{nk}$, and termwise passage through
the finite reverse connection yields \eqref{eq:audit-v67-jh}.
\end{proof}
\noindent\textbf{Scope.} This proves the degeneration in the precise
Askey--Wilson specialization and normalization used in V31. It does not
assert an independently verified numbered KLS formula for a different
continuous $q$-Jacobi prefactor convention. The eligible PASS tally
increases from \textbf{40} to \textbf{41}.


\section{V68: Continuous $q$-Laguerre $\to$ continuous $q$-Hermite (normalized reverse connection)}
\noindent\textbf{Verdict: PASS (fixed degree; normalization explicitly required).}
Let $0<q<1$, $a=q^{\lambda/2+1/4}$, $b=aq^{1/2}$, and let $\lambda\to+\infty$.
The reverse connection in V31, \texttt{eq:qlag-hermite-fwd}, reads
\begin{equation}\label{eq:audit-v68-reverse}
 L_n^{(\lambda)}(x\mid q)=\frac{a^n}{(q;q)_n}
 \sum_{k=0}^n {n\brack k}_q(-b)^{n-k}q^{\binom{n-k}{2}}
 {}_1\phi_1\!\left(\begin{matrix}q^{k-n}\\0\end{matrix};q,aq/b\right)H_k(x\mid q).
\end{equation}
\begin{proposition}
For each fixed $n\ge0$ and fixed $x$,
\begin{equation}\label{eq:audit-v68-limit}
 \lim_{\lambda\to+\infty}\frac{(q;q)_n}{a^n}L_n^{(\lambda)}(x\mid q)=H_n(x\mid q).
\end{equation}
Equivalently, the normalized reverse-connection coefficients tend to $\delta_{nk}$.
\end{proposition}
\begin{proof}
Multiply \eqref{eq:audit-v68-reverse} by $(q;q)_n/a^n$ and set $m=n-k$.
The coefficient of $H_k$ becomes
\[
 {n\brack k}_q(-b)^m q^{\binom m2}
 {}_1\phi_1\!\left(\begin{matrix}q^{-m}\\0\end{matrix};q,q^{1/2}\right).
\]
For fixed $m$, the terminating ${}_1\phi_1$ is a finite sum of $m+1$
terms, independent of $\lambda$, and is therefore finite.  Since
$b=q^{\lambda/2+3/4}\to0$, every coefficient with $m>0$ tends to zero.
For $m=0$ the terminating series is $1$, the Gaussian coefficient is $1$,
and the coefficient is identically $1$. Thus the normalized reverse
connection tends coefficientwise to the identity matrix.  The sum is
finite for each fixed $n$, so its limit is $H_n(x\mid q)$.
\end{proof}
\noindent\textbf{Eligibility and literature scope.} This is a genuine
boundary degeneration $a,b\to0$ of the continuous $q$-Laguerre
specialization of Al-Salam--Chihara, and not an exact finite-parameter
specialization. The factor $(q;q)_n/a^n$ is essential: the unnormalized
$L_n^{(\lambda)}$ tends to zero for $n>0$. This proves the relation in
V31's normalization; a numbered Koekoek--Swarttouw equation has not been
independently matched. The eligible PASS tally rises from \textbf{41} to
\textbf{42}.


\section{V69: Eligibility check for continuous $q$-Laguerre / $q$-Meixner--Pollaczek-type connection}
\noindent\textbf{Verdict: NO NEW ELIGIBLE KLS LIMIT ESTABLISHED.}
The reverse connection in V31, \texttt{eq:lag-qmp-rev-2phi1}, expresses
$M_n(x;\rho,\phi\mid q)$ in the continuous $q$-Laguerre basis.
The manuscript defines its $q$-Meixner--Pollaczek-\emph{type} family by
\[
 M_n(x;\rho,\phi\mid q)
 =Q_n(x;\rho e^{i\phi},\rho e^{-i\phi}\mid q),
\]
where $Q_n$ is Al-Salam--Chihara.  The continuous $q$-Laguerre
family is the scaled specialization
\[
 L_n^{(\lambda)}(x\mid q)
 =\frac{a^n}{(q;q)_n}Q_n(x;a,b\mid q),
 \qquad b=aq^{1/2},\quad a=q^{\lambda/2+1/4}.
\]
Thus the displayed connection is a parameter transport between two
Al-Salam--Chihara realizations, with a degree-dependent scalar factor.
V31 explicitly notes that its $M_n$ convention has not been
independently identified with a particular standard KLS
$q$-Meixner--Pollaczek normalization.

\begin{proposition}
For the normalization used in V31, setting $\rho=0$ gives the exact
identity
\[
 M_n(x;0,\phi\mid q)=Q_n(x;0,0\mid q)=H_n(x\mid q).
\]
This is a finite-parameter specialization and does not add a genuine
non-self limiting relation to the audit.
\end{proposition}
\begin{proof}
The Al-Salam--Chihara recurrence in its standard normalization is
\[
 2xQ_n(x;u,v\mid q)=Q_{n+1}(x;u,v\mid q)
 +(u+v)q^nQ_n(x;u,v\mid q)
 +(1-q^n)(1-uvq^{n-1})Q_{n-1}(x;u,v\mid q),
\]
with $Q_{-1}=0$, $Q_0=1$.  At $u=v=0$ it becomes
\[
 2xQ_n=Q_{n+1}+(1-q^n)Q_{n-1},
\]
which, together with the initial conditions, is the defining recurrence
of the continuous $q$-Hermite polynomials.  The identification therefore
holds exactly for every $n$, with no parameter tending to infinity.
\end{proof}
\noindent\textbf{Scope.} A genuine simultaneous singular scaling of
$\lambda,\rho,\phi$ might produce another family limit, but neither
\texttt{eq:lag-qmp-rev-2phi1} nor its literature discussion specifies a
KLS limiting path.  Such a path has not been proved here.  The V68
continuous $q$-Laguerre $\to$ continuous $q$-Hermite result is unchanged.
The eligible PASS count remains \textbf{42}.


\section{V70: continuous $q$-Jacobi / continuous $q$-Laguerre eligibility check}
\noindent\textbf{Source:} V31, reverse connection
\texttt{eq:lag-jacobi-new-rev}, with the definitions
\texttt{eq:Jdefinition} and \texttt{eq:qlag-asc}.
\textbf{Verdict: EXCLUDED for the natural $B=0$ path (exact
specialization); no new KLS limit proved.}

\begin{proposition}
Put $A=a=q^{\lambda/2+1/4}$ and
$B=q^{\beta/2+1/4}$.  At the finite parameter value $B=0$
(under polynomial continuation in the parameter $B$), the two
families satisfy the exact identity
\[
 J_n(x;\alpha,\beta\mid q)\big|_{A=a,B=0}
 =\frac{(q;q)_n}{a^n}L_n^{(\lambda)}(x\mid q),
 \qquad n\ge0.
\]
Here $\alpha=\lambda$, and $a\ne0$.
\end{proposition}
\begin{proof}
The manuscript defines
$J_n=p_n(x;A,Aq^{1/2},-B,-Bq^{1/2}\mid q)$.
At $B=0$ the last two Askey--Wilson parameters vanish.
The terminating Askey--Wilson series then reduces, term by term,
to the Al-Salam--Chihara series with parameters $(a,aq^{1/2})$;
in the manuscript's standard normalizations,
\[
 p_n(x;a,aq^{1/2},0,0\mid q)
 =Q_n(x;a,aq^{1/2}\mid q).
\]
By \texttt{eq:qlag-asc},
$L_n^{(\lambda)}=a^n Q_n/(q;q)_n$, giving the claimed identity.
Equivalently, the reverse-connection matrix
\texttt{eq:lag-jacobi-new-rev} becomes diagonal with entry
$(q;q)_n/a^n$ when the source and target parameters are matched.
No limiting argument is required.
\end{proof}
\noindent The formal approach $\beta\to+\infty$ sends $B\to0$
for $0<q<1$, but this endpoint is already a regular, exact
specialization of the terminating polynomial formula. Under the
audit's exclusion of exact transformations, it does not add a new
eligible limiting arrow. A different singular scaling has not been
identified or proved here. The eligible PASS count remains \textbf{42}.


\section{V71: Jacobi $\to$ Hermite from the reverse connection}
\noindent\textbf{Source.} V31 equation \texttt{eq:v14-jac-herm} is the reverse Jacobi-to-Hermite connection. The limit below is a genuine large-parameter degeneration, with a shrinking argument; it is not an exact specialization. Fix $n$ and $z\in\mathbb C$, let $t\to+\infty$, and use physicists' Hermite polynomials.
\begin{proposition}
For every fixed $n\geq0$,
\[
\boxed{\displaystyle\lim_{t\to\infty}\frac{2^n n!}{t^{n/2}}
 P_n^{(t,t)}\!\left(\frac{z}{\sqrt t}\right)=H_n(z).}
\]
\end{proposition}
\begin{proof}
Start with the reverse connection \texttt{eq:v14-jac-herm},
$P_n^{(a,b)}(x)=\sum_{m=0}^n D_{n,m}(a,b)H_m(x)$.
Its displayed double sum is obtained by inserting the finite Jacobi monomial expansion and then applying the exact monomial-to-Hermite inversion
\[
 x^r=2^{-r}r!\sum_{j=0}^{\lfloor r/2\rfloor}
 \frac{H_{r-2j}(x)}{j!(r-2j)!}.
\]
Therefore undoing this \emph{finite} inversion (i.e., resumming the reverse connection before rescaling $x$) recovers the Jacobi monomial polynomial exactly. For $a=b=t$, its standard finite symmetric form is
\[
 P_n^{(t,t)}(x)=\frac{(t+1)_n}{(2t+1)_n}
 \sum_{j=0}^{\lfloor n/2\rfloor}
 \frac{(-1)^j(t+\tfrac12)_{n-j}(2x)^{n-2j}}{j!(n-2j)!}.
\tag{*}
\]
For clarity, (*) is the finite Gegenbauer expansion with
$P_n^{(t,t)}=(t+1)_n(2t+1)_n^{-1}C_n^{t+1/2}$;
it is also the polynomial obtained from the preceding reverse-connection inversion. No family limit is used in this resummation.
Substitute $x=z/\sqrt t$ into (*), multiply by $2^nn!t^{-n/2}$, and observe for each fixed $j$ that
\[
 \frac{(t+1)_n}{(2t+1)_n}\longrightarrow 2^{-n},\qquad
 t^{-(n-j)}(t+\tfrac12)_{n-j}\longrightarrow1.
\]
The $j$th summand consequently tends to
\[
 \frac{(-1)^j n! (2z)^{n-2j}}{j!(n-2j)!}.
\]
There are only $\lfloor n/2\rfloor+1$ terms, so summation commutes with the limit. Their sum is precisely the terminating defining expansion of $H_n(z)$. This proves the claimed normalization and limit for all fixed degrees.

The preliminary resummation is essential: the separate terms $H_m(z/\sqrt t)$ in the unrescaled reverse connection are not individually the target Hermite basis at $z$, and their coefficients may cancel at leading order. Thus no unjustified termwise limit of the original Hermite-basis connection is asserted.
\end{proof}
\noindent\textbf{Verdict: PASS (algebraic limit, fixed degree).} The derivation starts from and finitely resums V31's reverse connection. The exact numbered Koekoek--Swarttouw correspondence has not been independently checked. The eligible PASS tally increases from 42 to \textbf{43}.


\section{V72: $q$-Racah / $q$-Hahn finite-grid specialization (eligibility check)}
\noindent\textbf{Source.} V31, reverse finite-grid connection
\texttt{eq:finite-HR-double-D}, in the subsection
``$q$-Hahn $\leftrightarrow$ $q$-Racah.'' This is a full
$(N+1)\times(N+1)$ grid connection, not a triangular polynomial
connection in one common variable.

\begin{proposition}[Exact endpoint and exclusion]
Fix $N\geq0$, $0<q<1$, and generic $A,B$ such that the defining
terminating series and finite-grid evaluation bases are nondegenerate.
Under
\[
 \gamma=A,\qquad \delta=0,\qquad
 \beta=ABq^{N+1},
\]
the $q$-Racah and $q$-Hahn polynomials coincide at every grid point:
\[
 \boxed{\ R_n(q^{-x};ABq^{N+1},A,0,N\mid q)
 =Q_n(q^{-x};A,B,N\mid q),\quad 0\leq n,x\leq N.\ }
\]
Consequently the reverse-connection coefficient matrix at this
specialization is the identity. This endpoint is an exact finite
specialization, not an additional genuine KLS limiting arrow.
\end{proposition}
\begin{proof}
The $q$-Racah series appearing in V31's reverse connection has upper
parameters
\[
 q^{-n},\quad \beta q^{n-N},\quad q^{-x},\quad
 \gamma\delta q^{x+1},
\]
and lower parameters $q^{-N},\beta\delta q,\gamma q$.
With the displayed substitutions these become
\[
 \begin{matrix}q^{-n},\ ABq^{n+1},\ q^{-x},\ 0\\
 q^{-N},\ 0,\ Aq\end{matrix}.
\]
The identical zero upper and lower parameters cancel term by term,
since $(0;q)_j=1$ for every $j$. The resulting terminating
${}_3\phi_2$ has exactly the defining upper and lower parameters of
$Q_n(q^{-x};A,B,N\mid q)$, and both series have argument $q$.
The $q$-Racah lattice also reduces from
$q^{-x}+\gamma\delta q^{x+1}$ to $q^{-x}$.

Write the reverse grid connection as $R_n=\sum_{k=0}^N D_{n,k}Q_k$.
At the specified finite parameter values, $R_n=Q_n$ at all $N+1$
grid points. The evaluation matrix $[Q_k(q^{-x})]_{x,k=0}^N$ is
invertible by the assumed nondegeneracy; hence uniqueness of the
expansion gives $D_{n,k}=\delta_{nk}$. No parameter tends to a
boundary, so this is not a singular family limit.
\end{proof}
\noindent\textbf{Verdict: EXCLUDED (exact finite specialization).}
The eligible PASS count remains \textbf{43}. This does not exclude
other parameter-dependent singular paths not specified here.

\section{Progress}
\noindent\textbf{V106 update:} Direct inspection of Koekoek--Swarttouw (1998), equations (4.6.5), (4.2.5), (4.2.8), and (4.2.4), independently certifies the four exact arrows mathematically proved in V102--V105. Each is nonduplicate and now PASS under the clarified rule. Recorded total: \textbf{49}. Earlier pending verdicts are superseded by V106; other outstanding checks remain.\par

\noindent\textbf{V103 update:} The $q$-Racah $\to$ dual $q$-Hahn exact specialization is now proved directly from the V31 reverse coefficients by the finite $q$-binomial theorem; independent KLS-arrow confirmation is pending, so the recorded PASS count remains 45.\par

\noindent V100 reinstates the distinct-family $q$-Meixner $\to$ ordinary $q$-Charlier KLS endpoint from V31 reverse equation \texttt{eq:qC-qM-reverse}; the finite $q$-binomial identity gives $D_{nk}\to\delta_{nk}$. Recorded PASS count rises to \textbf{44}; independent primary KLS verification remains pending.\par

\noindent V83 checks the $q$-Meixner--Pollaczek/continuous $q$-Hahn reverse connection \texttt{eq:cqh-qmp23}: its $b\to0$ limit follows termwise with factor $(q;q)_n$, but coincides with an exact Askey--Wilson zero-parameter specialization. EXCLUDED; eligible PASS tally remains \textbf{43}.\par

\noindent V82 checks the continuous $q$-Laguerre/Al-Salam--Chihara reverse connection \texttt{eq:lag-asc-hyper-rev51} via its regular finite kernel: at $(u,v)=(a,aq^{1/2})$ it collapses to an exact embedding, not a new singular KLS limit. Eligible PASS tally remains \textbf{43}.\par
\noindent V80 examines Jacobi/continuous Hahn reverse equation \texttt{eq:ch-jacobi-reverse}; a degree-one coefficient check confirms a generic basis change but no specified singular KLS path. No new PASS; tally remains \textbf{43}.\par
\noindent V73 verifies that the Hermite/generalized Laguerre reverse quadratic identities at parameters $\pm\tfrac12$ are exact transformations, excluded. Eligible PASS tally remains \textbf{43}.\par
\noindent V72 checks the $q$-Racah/$q$-Hahn reverse finite-grid connection \texttt{eq:finite-HR-double-D}: $\gamma=A$, $\delta=0$, $\beta=ABq^{N+1}$ gives an exact finite specialization, excluded. Eligible PASS tally remains \textbf{43}.\par
\noindent V71 proves symmetric Jacobi $\to$ physicists' Hermite using V31 reverse connection \texttt{eq:v14-jac-herm}, finite monomial resummation, and $a=b=t\to\infty$, $x=z/\sqrt t$. Eligible PASS tally: \textbf{43}.\par
\noindent V70 checks continuous $q$-Jacobi / continuous $q$-Laguerre, V31 	exttt{eq:lag-jacobi-new-rev}: the matched $B=0$ endpoint is an exact specialization, excluded from the genuine-limit count. Eligible PASS tally: \textbf{42}.\par
\noindent V69 screens the continuous $q$-Laguerre / $q$-Meixner--Pollaczek-type reverse connection \texttt{eq:lag-qmp-rev-2phi1}. Its $\rho=0$ endpoint is an exact Al-Salam--Chihara/continuous $q$-Hermite specialization; no new KLS singular path is established. Eligible PASS tally: \textbf{42}.\par
\noindent V68 proves normalized continuous $q$-Laguerre $\to$ continuous $q$-Hermite using V31 reverse equation \texttt{eq:qlag-hermite-fwd}, with $\lambda\to+\infty$. Eligible PASS tally: \textbf{42}.\par
\noindent V67 proves continuous $q$-Jacobi $\to$ continuous $q$-Hermite from V31 reverse connection \texttt{eq:JHforward}, with $A,B\to0$. Eligible PASS tally: \textbf{41}.\par
\noindent V66 checks the Meixner--Pollaczek/continuous dual Hahn reverse connection \texttt{eq:consolidated-cdh-mp}: the matched endpoint is an exact even/odd quadratic transformation, hence excluded. Eligible PASS tally: \textbf{40}.\par
\noindent V65 screens the Jacobi/generalised Bessel reverse connection \texttt{eq:jacobi-gb-reverse}: no specified KLS singular path is identified, hence no additional eligible PASS. Tally: \textbf{40}.\par
\noindent V64 proves Hahn $\to$ Jacobi using V31 reverse equation \texttt{eq:jacobi-hahn-reverse} via exact finite resummation and $x=Nt$, $N\to\infty$. Eligible PASS tally: \textbf{40}.\par
\noindent V63 checks the big $q$-Jacobi/$q$-Racah reverse equation \texttt{eq:171-reverse}: the natural matched-parameter endpoint is an exact specialization, not an eligible singular KLS limit. Eligible tally remains \textbf{39}.\par
\noindent V62 proves $q$-Hahn $\to$ little $q$-Jacobi via V31 reverse equation \texttt{eq:ljqh-reverse}, with $N\to\infty$ and $x=q^{-N}z$. Eligible PASS tally: \textbf{39}.\par
\noindent V61 proves $q$-Hahn $\to$ $q$-Meixner from V31 reverse equation \texttt{eq:168-D}, with $N\to\infty$, $A=b$, $B=-q^{-N-1}/(bc)$. Eligible PASS tally: \textbf{38}.\par
\noindent V60 proves big $q$-Jacobi $\to$ $q$-Meixner by the explicit V31 reverse connection and $C\to\infty$, $A=b$, $B=-C/(bc)$. Eligible PASS tally: \textbf{37}.\par
\noindent V59 proves big $q$-Jacobi $\to$ little $q$-Jacobi by finite resummation of the V31 reverse connection and $C\to\infty$, $x=Cqz$. Eligible PASS tally: \textbf{36}.\par
\noindent V58 proves big $q$-Laguerre $\to$ Wall by resumming the reverse connection and taking $B\to\infty$ with $x=Bqz$. Eligible PASS tally: \textbf{35}.\par
\noindent V57 cross-checks Connections 95--101: 95, 98, 101 are same-family; 96, 97, 100 are already audited; 99 is an exact specialization. No new PASS; the eligible tally remains \textbf{34}.\par
\noindent V56 screens Connections 104--115: 104, 107, 111, and 115 are same-family; 105, 106, and 108 are already audited; 109 and 110 lack an identified KLS limit; 112--114 are exact family identifications. No new eligible PASS; the tally remains \textbf{34}.\par
\noindent V55 screens Connection 103 (Al-Salam--Carlitz I / dual $q$-Krawtchouk): its reverse connection is a full finite-grid identity, and no specific KLS family limit is identified; the eligible PASS tally remains \textbf{34}.\par
\noindent V54 checks the $q$-Racah/ordinary $q$-Krawtchouk reverse grid connection and excludes its $\gamma=\delta=0$, $\beta=-pq^N$ endpoint as an exact specialization. Eligible PASS count remains \textbf{34}.\par
\noindent V53 excludes the affine $q$-Krawtchouk/dual $q$-Hahn case as an exact $\delta=0$ specialization; the eligible PASS tally remains \textbf{34}.\par
\noindent V52 verifies that the affine $q$-Krawtchouk/$q$-Hahn relation is an exact $B=0$ specialization and excludes it from the genuine-limit count. The total remains \textbf{34}.\par
\noindent The V51 degree-one check refutes the alternative $AB=p$-fixed path in the displayed normalization, without changing the tally. The V50 quantum $q$-Krawtchouk degeneration adds one eligible PASS, raising the corrected tally to \textbf{34}. The direct terminating $q$-Chu--Vandermonde summation in V49 strengthens the V48 coefficient-level verification and confirms the V47 $q$-Hahn $\to$ ordinary $q$-Krawtchouk proof and leaves the tally unchanged. The V47 proof had raised the corrected tally to 33 eligible PASS entries. Earlier historical counts of 31 refer to the tally before this addition. The Connection 102 zero-parameter and fixed-affine obstruction tests add no new PASS; the post-117 eligibility screening adds no new PASS; the corrected cumulative audit retains 34 eligible limiting PASS entries after the V50 quantum $q$-Krawtchouk addition. Connection 99 is excluded as an exact specialization; Connection 102 is not established as a KLS limiting arrow. Connection 117 remains PASS for the stated normalized confluent path.\par
\begin{itemize}
\item $q$-Racah/$q$-Hahn at $\gamma=A$, $\delta=0$, $\beta=ABq^{N+1}$: \textbf{EXCLUDED}, exact finite specialization (V31 \texttt{eq:finite-HR-double-D}).
\item Symmetric Jacobi $\to$ Hermite: \textbf{PASS}, fixed degree, V31 \texttt{eq:v14-jac-herm}; $a=b=t\to\infty$ with $x=z/\sqrt t$ and normalization $2^nn!t^{-n/2}$.
\item Normalized continuous $q$-Laguerre $\to$ continuous $q$-Hermite: \textbf{PASS}, fixed degree, V31 \texttt{eq:qlag-hermite-fwd}; $\lambda\to+\infty$ with factor $(q;q)_n/a^n$.
\item Continuous $q$-Jacobi $\to$ continuous $q$-Hermite: \textbf{PASS}, fixed degree, V31 \texttt{eq:JHforward}; $A,B\to0$ with reverse coefficients tending to the identity.
\item Hahn $\to$ Jacobi: \textbf{PASS}, fixed degree, V31 \texttt{eq:jacobi-hahn-reverse}; $N\to\infty$, $x=Nt$, reverse Newton resummation.
\item $q$-Hahn $\to$ little $q$-Jacobi: \textbf{PASS}, fixed degree, V31 \texttt{eq:ljqh-reverse}; $N\to\infty$, $x=q^{-N}z$, finite Newton resummation.
\item $q$-Hahn $\to$ $q$-Meixner: \textbf{PASS}, all fixed degrees, V31 equation \texttt{eq:168-D}; $N\to\infty$ and $B=-q^{-N-1}/(bc)$.
\item Big $q$-Jacobi $\to$ $q$-Meixner: \textbf{PASS}, fixed degree, V31 \texttt{eq:169-reverse}; $C\to\infty$, $A=b$, $B=-C/(bc)$.
\item Big $q$-Jacobi $\to$ little $q$-Jacobi: \textbf{PASS}, all fixed degrees, V31 equation \texttt{eq:ljbj166-reverse}; $C\to\infty$, $x=Cqz$.
\item Big $q$-Laguerre $\to$ little $q$-Laguerre (Wall): \textbf{PASS}, all fixed degrees, V31 equation \texttt{eq:wall-big-reverse}; $B\to\infty$, $x=Bqz$.
\item $q$-Hahn $\to$ quantum $q$-Krawtchouk: \textbf{PASS}, all $0\le n\le N$, V31 equation \texttt{eq:finite-qtmH-reverse}; $A\to\infty$, $B=p$ fixed (not $AB=p$ fixed).
\item $q$-Hahn $\to$ ordinary $q$-Krawtchouk: \textbf{PASS}, all $0\le n\le N$, V31 equation \texttt{eq:finite-KH-reverse}; confluent $A\to0$, $AB=-p/q$.
\item Dual $q$-Hahn $\to$ dual $q$-Krawtchouk: \textbf{PASS}, all $0\le n\le N$ on a nondegenerate finite grid, V31 equation \texttt{eq:finite-dKdH-reverse}; K--S (1998), Eq.~(4.7.2).
\item Askey--Wilson $\to$ big $q$-Jacobi: \textbf{PASS}, all fixed degrees, V31 equation \texttt{eq:aw-bigqj-reverse}; simultaneous parameter and argument scaling with degree-dependent normalization.
\item Standard $q$-Krawtchouk $\to$ alternative $q$-Charlier ($q$-Bessel): \textbf{PASS}, all fixed degrees, V31 equation \texttt{eq:qBqK-reverse}; scaled lattice argument.
\item Standard $q$-Krawtchouk $\to$ ordinary $q$-Charlier: \textbf{PASS}, all fixed degrees, V31 equation \texttt{eq:qCqK-reverse}.
\item Ordinary $q$-Meixner $\to$ ordinary $q$-Charlier: \textbf{EXCLUDED (exact $b=0$ specialization)}, V31 equation \texttt{eq:qC-qM-reverse}; previously counted in error.
\item Little $q$-Jacobi $\to$ $q$-Bessel (alternative $q$-Charlier): \textbf{PASS}, all fixed degrees, V31 equation \texttt{eq:qBlqJ-reverse}; KLS (1998), Eq.~(4.12.3).
\item $q$-Meixner $\to$ $q$-Laguerre: \textbf{PASS}, all fixed degrees, V31 equation \texttt{eq:qLqM-D}; KLS (1998), Eq.~(4.13.1).
\item Little $q$-Jacobi $\to$ $q$-Laguerre: \textbf{PASS}, all fixed degrees, V31 equation \texttt{eq:qLlJ-D}; scaled variable and finite inversion.
\item Affine $q$-Krawtchouk $\to$ little $q$-Laguerre (Wall): \textbf{PASS}, all fixed degrees, V31 equation \texttt{eq:wall-affine-reverse}; variable rescaling and finite resummation.
\item Little $q$-Jacobi $\to$ little $q$-Laguerre (Wall): \textbf{PASS}, all fixed degrees, V31 equation \texttt{eq:wall-lj-reverse}.
\item Big $q$-Jacobi $\to$ big $q$-Laguerre: \textbf{PASS}, all fixed degrees, V31 equation \texttt{eq:blbj-reverse}.
\item Continuous $q$-Hahn $\to$ $q$-Meixner--Pollaczek: \textbf{PASS}, all fixed degrees, V31 equation \texttt{eq:cqh-qmp23}.
\item Al-Salam--Chihara $\to$ continuous big $q$-Hermite: \textbf{PASS}, all fixed degrees, V31 equation \texttt{eq:ASCBH}.
\item Continuous dual $q$-Hahn $\to$ Al-Salam--Chihara: \textbf{PASS}, all fixed degrees, V31 equation \texttt{eq:CDASC}.
\item Continuous big $q$-Hermite $\to$ continuous $q$-Hermite: \textbf{PASS}, all fixed degrees, V31 equation \texttt{eq:BHH}.
\item $q$-Laguerre $\to$ Stieltjes--Wigert: \textbf{PASS}, all fixed degrees, V31 equation \texttt{eq:SWqL-rev}.
\item Quantum $q$-Krawtchouk $\to$ Al-Salam--Carlitz II: \textbf{PASS}, all fixed degrees, V31 equation \texttt{eq:ASCII-qtm-reverse}.
\item $q$-Meixner $\to$ Al-Salam--Carlitz II: \textbf{PASS}, all fixed degrees, from V31 equation \texttt{eq:ASCII-qM-reverse}.
\item Askey--Wilson $\to$ continuous dual $q$-Hahn: \textbf{PASS}, all fixed degrees, from V31 equation \texttt{eq:AWCD}.
\item Racah $\to$ Hahn: \textbf{PASS}, all $0\le n\le N$, from V31 equation \texttt{eq:racah-to-hahn-grid}.
\item Racah $\to$ dual Hahn: \textbf{PASS}, all admissible degrees, from V31 equation \texttt{eq:racah-dh-reverse}.
\item Krawtchouk $\to$ Hermite: \textbf{PASS}, all fixed degrees, from V31 equation \texttt{eq:v16-k-h}.
\item Hahn $\to$ Meixner: \textbf{PASS}, all fixed degrees, from V31 equation \texttt{eq:hahn-meixner-reverse}.
\item Hahn $\to$ Krawtchouk: \textbf{PASS}, all $0\le n\le N$, from V31 equation \texttt{eq:hahn-kraw-reverse}.
\item Krawtchouk $\to$ Charlier: \textbf{PASS}, all degrees, from V31 equation \texttt{eq:kraw-char-reverse}.
\item Meixner $\to$ Charlier: \textbf{PASS}, all degrees, from V31 equation \texttt{eq:meix-char-reverse}.
\item Charlier $\to$ Hermite: \textbf{PASS}, all degrees, by resumming the V31 reverse connection.
\item Wilson $\to$ continuous dual Hahn: \textbf{PASS}, all degrees.
\item Meixner--Pollaczek $\to$ Hermite: \textbf{PASS}, all degrees, using the V31 reverse Bell connection.
\item Continuous Hahn $\to$ Meixner--Pollaczek: \textbf{PASS} for the
symmetric parameter path, all degrees; KLS equation-number matching open.
\item Meixner $\to$ generalized Laguerre: \textbf{PASS}, all degrees, scaling $x=z/(1-c)$ as $c\to1$.
\item Jacobi $\to$ generalized Laguerre: \textbf{PASS}, all degrees, classical scaling $x=1-2z/\beta$.
\item Wilson $\to$ Jacobi: \textbf{PASS} for the explicitly stated
algebraic scaling, all degrees; KLS equation-number matching open.
\item $q$-Racah/ordinary $q$-Krawtchouk at $\gamma=\delta=0$, $\beta=-pq^N$: \textbf{EXCLUDED}, exact specialization (V31 \texttt{eq:finite-KR-reverse}).
\item Dual $q$-Hahn/$q$-Racah at $\beta=0$: \textbf{EXCLUDED}, exact specialization (V31 \texttt{eq:finite-dHR-reverse}).
\item Self-connections and exact transformations: excluded.
\end{itemize}


\section{V73: Hermite / generalized Laguerre quadratic identities (eligibility check)}
\noindent\textbf{Source.} V31, reverse identities
\texttt{eq:laguerre-hermite-even} and
\texttt{eq:laguerre-hermite-odd} in the subsection
``Hermite $\leftrightarrow$ Generalised Laguerre.''
The identities, valid for every integer $m\ge0$, are
\begin{align}
 L_m^{(-1/2)}(x^2)&=\frac{(-1)^m}{2^{2m}m!}H_{2m}(x),\\
 xL_m^{(1/2)}(x^2)&=\frac{(-1)^m}{2^{2m+1}m!}H_{2m+1}(x).
\end{align}
These are exact quadratic transformations at the \emph{fixed} Laguerre
parameters $-1/2$ and $1/2$, not asymptotic limits. For completeness,
substitute the terminating series
\[
 L_m^{(\alpha)}(z)=\frac{(\alpha+1)_m}{m!}
 \sum_{j=0}^{m}\frac{(-m)_j z^j}{(\alpha+1)_j j!}
\]
into each reverse identity. The right sides become the respective
terminating ${}_1F_1$ representations of $H_{2m}$ and $H_{2m+1}$;
the constants agree because
\[
 (\tfrac12)_m=\frac{(2m)!}{4^m m!},\qquad
 (\tfrac32)_m=\frac{(2m+1)!}{4^m m!}.
\]
In particular, the odd transformation necessarily contains the
factor $x$; omitting it would make the two sides have different parity.
\par\noindent\textbf{Verdict: EXCLUDED (exact quadratic transformation).}
No limiting parameter, rescaling path, or non-self KLS degeneration is
involved in these two reverse identities. The eligible PASS tally remains
\textbf{43}. This exclusion does not address other Laguerre--Hermite
asymptotic relations requiring a different connection formula.


\section{V74: Meixner / dual Hahn (finite-grid eligibility check)}
\noindent\textbf{Source.} V31, equations
\texttt{eq:meixner-dh-forward} and \texttt{eq:dh-meixner-reverse}.
The forward relation expresses $\mathrm M_n(x;\beta,c)$ on the
$N+1$ nodes $x=0,\ldots,N$ in the dual Hahn basis
$\mathrm{DH}_k(\Lambda(x);\gamma,\delta,N)$, where
\[
 \Lambda(x)=x(x+\gamma+\delta+1).
\]
This is a full finite-grid interpolation relation, not a triangular
polynomial connection in one variable: for generic parameters,
$\Lambda(x)$ has degree two in $x$.  In particular,
$\mathrm{DH}_k(\Lambda(x);\gamma,\delta,N)$ has degree $2k$ in $x$,
whereas $\mathrm M_j(x;\beta,c)$ has degree $j$.
The reverse identity accordingly has the form
\[
 \mathrm{DH}_n(\Lambda(x);\gamma,\delta,N)
   =\sum_{k=0}^{2n}D_{n,k}\,\mathrm M_{2n-k}(x;\beta,c),
\]
with the coefficients given explicitly in V31.  This degree doubling
is essential: the reverse connection is an ordinary polynomial
identity in $x$, but does not imply that a dual Hahn polynomial of
degree $n$ in $\Lambda$ degenerates to a Meixner polynomial of degree
$n$ in $x$ along any unspecified parameter path.

\noindent\textbf{Eligibility verdict: NOT ESTABLISHED.}
Neither equation specifies a singular parameter scaling, an argument
rescaling, and a degree-preserving normalization that would identify a
particular Koekoek--Swarttouw dual Hahn--Meixner limiting arrow.
The existence of the two exact finite connections alone cannot supply
such a limit.  Therefore no PASS is assigned, and the tally remains
\textbf{43}.  This check does not rule out a separate, explicitly
specified asymptotic degeneration; such a candidate would need a new
coefficient-level analysis of the reverse formula.


\section{V75: Racah / Wilson quadratic-lattice identification (eligibility check)}
\noindent\textbf{Source.} V31, reverse identity
\texttt{eq:wilson-to-racah-sec2} in the Racah--Wilson subsection.
Fix an admissible integer $N\geq0$ and $0\leq n\leq N$, and set
\[
 A=\frac{\delta-N}{2},\qquad B=\alpha+1-A,\qquad
 C=\beta+\frac{\delta+N}{2}+1,\qquad D=-\frac{\delta+N}{2}.
\]
The manuscript's \emph{reverse} identity is
\[
 W_n\bigl(-(x+A)^2;A,B,C,D\bigr)
 = (\alpha+1)_n(\beta+\delta+1)_n(-N)_n
 R_n\bigl(x(x-N+\delta);\alpha,\beta,-N-1,\delta\bigr).
\]
The displayed factor is finite and nonzero under the stated generic
admissibility conditions.  This is an equality at fixed parameters and
fixed $N$: the quadratic change of variable maps one terminating
${}_4F_3(1)$ polynomial directly to the other.  In particular, no
sequence of parameters tending to a boundary and no asymptotic
normalization appears in the reverse identity.

\noindent\textbf{Verdict: EXCLUDED (exact quadratic-lattice identification).}
The reverse identity is useful, but it is not a genuine non-self
Koekoek--Swarttouw \emph{limit} under the audit rules.  It contributes
no new PASS; the eligible limiting-proof tally stays \textbf{43}.
This exclusion applies only to the displayed fixed-parameter
identification, not to other possible Racah--Wilson degenerations.


\section{V76: Pseudo-Jacobi / continuous Hahn (reverse-connection eligibility)}
\noindent\textbf{Source.} V31, reverse connection
\texttt{eq:v16-ch-pj}, with coefficient extractors
\texttt{eq:pj-ch-taylor} and \texttt{eq:pj-ch-chain}.
The manuscript states, for nondegenerate polynomial bases,
\[
 p_n(x;a,b,c,d)=\sum_{m=0}^n C_{nm}\,
 \mathrm{PJ}_m(x;\alpha,\beta),\qquad
 C_{nm}=\sum_{r=m}^n T_r^{(n)}(p_n)\,
 I_{mr}(\mathrm{PJ}_\bullet).
\]
The formula is an exact upper-triangular change of polynomial basis,
not an asymptotic assertion.  We verify its mathematical content
and test what can be deduced about a family limit.

\paragraph{Coefficient-level verification.}
Write $p_n(x)=\sum_{r=0}^n f_rx^r$ and
$\mathrm{PJ}_m(x)=\sum_{r=0}^m g_{rm}x^r$.
The Newton--Stirling extractor in V31 satisfies
$T_r^{(n)}(p_n)=f_r$ identically.  Let
$G=(g_{rm})_{0\le r,m\le n}$, with $g_{rm}=0$ for $r>m$.
For nonzero leading coefficients $g_{mm}$, $G$ is invertible,
and the chain expression in \texttt{eq:pj-ch-chain} is precisely
$(G^{-1})_{mr}$.  Thus
\[
 C_{nm}=\sum_{r=m}^n (G^{-1})_{mr}f_r,
 \qquad
 \sum_{m=0}^n C_{nm}\mathrm{PJ}_m(x)=p_n(x).
\]
This proves the reverse connection as an exact identity, without
requiring any special relation among $a,b,c,d,\alpha,\beta$.
In degree one the same observation is particularly transparent:
if $p_1(x)=u_1x+u_0$ and
$\mathrm{PJ}_1(x)=v_1x+v_0$ ($v_1\ne0$), then
\[
 p_1(x)=\frac{u_1}{v_1}\mathrm{PJ}_1(x)
 +\left(u_0-\frac{u_1v_0}{v_1}\right)\mathrm{PJ}_0(x),
 \qquad \mathrm{PJ}_0=1.
\]
No limiting relation follows merely from this identity.

\paragraph{What would be required for a PASS.}
A genuine family limit would require an independently specified
Koekoek--Swarttouw path $\theta(t)$ in the continuous Hahn
parameters, a corresponding path in pseudo-Jacobi parameters,
an argument scaling $x=x(t,z)$ and nonzero normalizations $N_n(t)$,
for which the \emph{rescaled} reverse expansion converges to a
nontrivial pseudo-Jacobi polynomial for each fixed $n$.  Neither
\texttt{eq:v16-ch-pj} nor its coefficient extractor selects such a
path.  In particular, the exact triangular inverse is valid for
arbitrary independent parameter choices, so it cannot itself certify
an Askey-scheme degeneration.

\noindent\textbf{Verdict: NOT ESTABLISHED AS AN ELIGIBLE KLS LIMIT.}
The reverse identity is verified as a finite basis change, but no
specific KLS limiting arrow or parameter scaling has been established.
No new PASS is assigned; the cumulative eligible total remains
\textbf{43}.  This verdict does not assert that every possible
asymptotic relation between these two families is impossible.


\section{V77: Symmetric Krawtchouk / dual Hahn (quadratic reverse identities)}
\noindent\textbf{Source.} V31, subsection
\emph{Krawtchouk $\leftrightarrow$ Dual Hahn}, the four displayed
transformed-lattice identities immediately preceding the Krawtchouk
self-connection subsection.  This is an audit of the \emph{reverse}
identities obtained there by solving for the dual Hahn polynomial.
The source explicitly restricts the Krawtchouk parameter to
$p=\tfrac12$ and describes these as Koornwinder's quadratic
transformations (105)--(108).

\paragraph{Reverse forms and their domain.}
For $0\le m\le N$, the even-degree reverse identities include
\[
 R_m(x^2;-\tfrac12,-\tfrac12,N)
 =\frac{(-N+\tfrac12)_m}{(\tfrac12)_m}
 K_{2m}(x+N;\tfrac12,2N),
\]
and
\[
 R_m(x(x+1);-\tfrac12,\tfrac12,N)
 =\frac{(-N-\tfrac12)_m}{(\tfrac12)_m}
 K_{2m}(x+N+1;\tfrac12,2N+1).
\]
For the odd-degree case on the $2N+1$ grid, the reverse identity is
\[
 R_m(x(x+1);\tfrac12,-\tfrac12,N)
 =\frac{(-N-\tfrac12)_{m+1}}{(\tfrac32)_m}
 \frac{K_{2m+1}(x+N+1;\tfrac12,2N+1)}{x+\tfrac12}.
\]
The remaining odd-degree reverse identity, valid for $0\le m\le N-1$,
is
\[
 R_m(x^2-1;\tfrac12,\tfrac12,N-1)
 =-\frac{N(-N+\tfrac12)_m}{(\tfrac32)_m}
 \frac{K_{2m+1}(x+N;\tfrac12,2N)}{x}.
\]
The divisions are exact in the polynomial ring: by the symmetry of
$K_j(\cdot;\tfrac12,M)$ around $M/2$, odd-degree centered
Krawtchouk polynomials vanish at the center.  In particular,
$K_{2m+1}(x+N;\tfrac12,2N)$ is odd in $x$, and
$K_{2m+1}(x+N+1;\tfrac12,2N+1)$ is odd in $x+\tfrac12$.
This proves that the reverse expressions have no spurious poles.

\paragraph{Eligibility test.}
Each displayed equation holds at a \emph{fixed} positive integer $N$
and at the \emph{fixed} parameter $p=\tfrac12$; its two sides are
related by a quadratic lattice substitution and a nonzero scalar.
The degrees $2m$ or $2m+1$ on the Krawtchouk side become degree $m$
on the quadratic lattice, so these are quadratic transformations,
not triangular connections between two bases in a common variable.
No parameter tends to a boundary, no rescaling limit is taken, and
no limiting interchange is needed.  Thus they cannot constitute a
new KLS degeneration even though the identities themselves are exact.

\noindent\textbf{Verdict: EXCLUDED (four exact quadratic transformations).}
No new PASS is assigned.  The cumulative number of established
eligible limiting proofs remains \textbf{43}.  This classification
concerns only these four symmetric quadratic identities and does not
preclude separately proved Krawtchouk or dual Hahn limiting arrows.


\section{V78: Meixner--Pollaczek to Laguerre --- angle-limit versus fixed-angle reverse connection}
\noindent\textbf{Source.} V31, subsection \emph{Generalised Laguerre
$\leftrightarrow$ Meixner--Pollaczek}, reverse equation
\texttt{eq:ordinary-mp-to-lag}, together with the separate limit
\texttt{eq:mp-laguerre-direct-limit}.  This entry distinguishes a genuine
classical limiting arrow from what can actually be proved \emph{using the
specific reverse connection supplied in V31}.

\paragraph{The genuine classical limiting arrow.}
For fixed $n\ge0$, $\alpha>-1$, and $x$ fixed, V31 records
\[
\lim_{\phi\downarrow0}P_n^{((\alpha+1)/2)}
 \left(-\frac{x}{2\phi};\phi\right)=L_n^{(\alpha)}(x).
\]
The standard Meixner--Pollaczek generating function, with
$\lambda=(\alpha+1)/2$, is
\[
\sum_{n\ge0}P_n^{(\lambda)}(y;\phi)z^n
 =(1-e^{i\phi}z)^{-\lambda+iy}
  (1-e^{-i\phi}z)^{-\lambda-iy}.
\]
At $y=-x/(2\phi)$, the logarithm of the right-hand side is
\[
-\lambda\log(1-2z\cos\phi+z^2)
-\frac{ix}{2\phi}
 \left\{\log(1-e^{i\phi}z)-\log(1-e^{-i\phi}z)\right\}.
\]
Since, uniformly on sufficiently small closed discs $|z|\le r<1$,
\[
\frac{\log(1-e^{i\phi}z)-\log(1-e^{-i\phi}z)}{\phi}
\longrightarrow -\frac{2iz}{1-z},
\]
the logarithm tends to
$-(\alpha+1)\log(1-z)-xz/(1-z)$.
Consequently the generating functions converge locally uniformly to
\[
(1-z)^{-\alpha-1}\exp\!\left(-\frac{xz}{1-z}\right)
=\sum_{n\ge0}L_n^{(\alpha)}(x)z^n.
\]
Cauchy's coefficient formula proves the stated fixed-degree limit.
This verifies the \emph{classical arrow} independently of any finite
connection.

\paragraph{Why the V31 reverse connection does not prove this arrow.}
The reverse identity \texttt{eq:ordinary-mp-to-lag} is explicitly
\[
P_n^{(\lambda)}(y;\tfrac\pi2)
 =\sum_{m=0}^n(-2)^m\frac{Y_{n-m}
 (1!b_1(m,\lambda),\ldots,(n-m)!b_{n-m}(m,\lambda))}{(n-m)!}
 L_m^{(\lambda)}(y),
\]
with $b_r$ computed using $\arctan z$ and \emph{no free angle
parameter}.  In contrast, the limiting arrow requires
$\phi\downarrow0$ and the diverging argument
$y=-x/(2\phi)$.  There is no parameter $\phi$ in these reverse
coefficients to which the limit can be applied.  Moreover, the
Laguerre index on the right is $\lambda$, whereas the index in the
classical arrow is $2\lambda-1$; this is a further normalization
mismatch.  For example, the $n=1$ reverse identity is
\[
P_1^{(\lambda)}(y;\tfrac\pi2)=2y
 =-2L_1^{(\lambda)}(y)+2(\lambda+1),
\]
which is true for every $y$ but contains no information about
$P_1^{(\lambda)}(y;\phi)$ at varying $\phi$.

\noindent\textbf{Verdict: CLASSICAL LIMIT VERIFIED; REVERSE-CONNECTION
PASS NOT ESTABLISHED.}  The Meixner--Pollaczek $\to$ Laguerre
small-angle limit is genuine and has been proved above, but it cannot
be credited as a \emph{reverse-connection-derived} proof from the
fixed-angle identity \texttt{eq:ordinary-mp-to-lag}.  A parameterized
reverse connection valid for variable $\phi$, followed by a justified
scaled limit, would be needed for a PASS under this audit's protocol.
The eligible reverse-connection PASS count therefore remains
\textbf{43}.


\section{V79: Symmetric continuous Hahn and Wilson --- reverse quadratic identities are not limits}
\noindent\textbf{Source reverse connection:} V31 equation
\texttt{eq:consolidated-w-ch}, together with its matched-parameter
specialization following \texttt{eq:ch-w-direct}.

\paragraph{Matched-parameter reduction.}
Let $\varepsilon\in\{0,1\}$ and $n\geq0$, and set
\[
 t_\varepsilon=a+b+\tfrac12+\varepsilon,\qquad
 K_{n,\varepsilon}=
 \frac{(-1)^n(2n+2a+2b-1+\varepsilon)_{2n+\varepsilon}}
 {(2n+\varepsilon)!\,(n+t_\varepsilon-1)_n}.
\]
The source manuscript's reverse connection expresses
$x^\varepsilon W_n(x^2;a,b,c,d)$ as a finite linear combination of
$p_{2j+\varepsilon}(x;a,b,a,b)$, $0\leq j\leq n$.
Its stated matched-parameter specialization $(c,d)=(\tfrac12,\varepsilon)$
is the exact identity
\begin{equation}\label{eq:v79-matched}
 x^\varepsilon W_n(x^2;a,b,\tfrac12,\varepsilon)
 =K_{n,\varepsilon}^{-1}p_{2n+\varepsilon}(x;a,b,a,b),
 \qquad K_{n,\varepsilon}\ne0.
\end{equation}
In particular, the matched reverse-connection matrix is diagonal:
its coefficient of $p_{2j+\varepsilon}$ is
$K_{n,\varepsilon}^{-1}\delta_{nj}$.  This is not an
asymptotic assertion: all parameters in \eqref{eq:v79-matched} are
fixed and finite.

\paragraph{Parity and polynomial degree.}
For $\varepsilon=0$ both sides of \eqref{eq:v79-matched} are even
polynomials of degree $2n$ in $x$. For $\varepsilon=1$ both sides are
odd polynomials of degree $2n+1$; the prefactor $x$ on the Wilson side
is essential.  Thus the quadratic substitution $x\mapsto x^2$ and
parity factor account exactly for the apparent mismatch in polynomial
degrees. The $n=0$ case is consistent: $K_{0,0}=1$ and the even
identity reads $1=1$.  The generic nonvanishing restriction on $K$
ensures that the displayed inverse is defined; exceptional values
require separate continuation.

\paragraph{Why this does not add a KLS limiting arrow.}
There is no parameter tending to zero or infinity in
\eqref{eq:v79-matched}, no variable scaling depending on a limiting
parameter, and no remainder to estimate.  The result is an exact
even/odd quadratic transformation (in fact explicitly identified as
such in V31), rather than a degeneration from one Askey-scheme family
to another. The arbitrary-$(c,d)$ connection is a finite basis
expansion and likewise does not itself specify a singular path.

\noindent\textbf{Verdict: EXCLUDED (two exact quadratic
transformations, $\varepsilon=0,1$).} No new eligible reverse-connection
limiting proof is established; the cumulative PASS count stays
\textbf{43}. This conclusion does not preclude a distinct singular
Wilson/continuous-Hahn limit under a separately specified scaling.


\section{V80: Jacobi and continuous Hahn --- a finite reverse basis change is not a limiting arrow}
\noindent\textbf{Source reverse connection:} V31 equation
\texttt{eq:ch-jacobi-reverse} (with forward equation
\texttt{eq:jacobi-ch-connection}).

\paragraph{What the reverse connection actually proves.}
Put $J_m(x)=P_m^{(\alpha,\beta)}(1-2x)$ and
$\sigma=a+b+c+d$, $A=a+c$, $B=a+d$. V31 expands
\[
 p_n(x;a,b,c,d)=\sum_{m=0}^n C_{nm}J_m(x),
\]
where, for $\alpha,\beta>-1$ and nonsingular continuous Hahn
parameters,
\begin{equation}\label{eq:v80-reverse}
 C_{nm}=\frac{(-1)^m}{n!\,h_m^{\rm J}}
 \sum_{r=m}^n t_{n,r}\binom rm
 B(\alpha+r+1,\beta+m+1).
\end{equation}
Here $t_{n,r}$ are the explicitly stated power coefficients of
$n!p_n(x)$ in V31, and $h_m^{\rm J}$ is the shifted Jacobi norm.
The formula is an orthogonal projection of an \emph{arbitrary}
continuous Hahn polynomial onto a fixed Jacobi basis, not a parameter
specialization or an asymptotic statement.

\paragraph{A degree-one diagnostic.}
The terminating continuous Hahn definition used in V31 gives
\[
 p_1(x;a,b,c,d)
 =i\bigl[AB-\sigma(a+ix)\bigr]
 =\sigma x+i(AB-a\sigma).
\]
Meanwhile
\[
 J_0(x)=1,\qquad
 J_1(x)=(\alpha+1)-(\alpha+\beta+2)x.
\]
Consequently the unique reverse coefficients in degree one are
\begin{equation}\label{eq:v80-degree-one}
 C_{11}=-\frac{\sigma}{\alpha+\beta+2},\qquad
 C_{10}=i(AB-a\sigma)
 +\frac{\sigma(\alpha+1)}{\alpha+\beta+2}.
\end{equation}
These identities can also be recovered directly from
\eqref{eq:v80-reverse}; they show explicitly that generic parameter
choices give a nonzero lower-degree component.  Thus the reverse
connection is not a diagonal identification of the families.

\paragraph{Limiting-relation test.}
To deduce a genuine KLS family limit from this reverse connection one
would need to prescribe a singular path in $(a,b,c,d,\alpha,\beta)$,
a scaling of the polynomial variable and a degree-dependent
normalization, and then prove convergence of all the normalized
coefficients. None of these data is specified by
\texttt{eq:ch-jacobi-reverse}. The degree-one calculation
\eqref{eq:v80-degree-one} does not exclude all possible paths, but
it demonstrates why merely taking a finite Jacobi expansion as a
limiting proof would be invalid. The Fourier-transform relation
mentioned in V31's historical note is a different operation and does
not supply the missing singular scaling.

\noindent\textbf{Verdict: NOT ESTABLISHED AS AN ELIGIBLE KLS LIMIT.}
The reverse connection is an exact finite basis change, but no
specific KLS Jacobi/continuous-Hahn degeneration has been proved
from it. No new PASS is credited; the cumulative eligible tally
remains \textbf{43}.


\section{V81: Al-Salam--Chihara to continuous big $q$-Hermite --- reverse-connection boundary check}
\noindent\textbf{V31 source:} \texttt{eq:BHASC} and \texttt{eq:ASCBH} in the subsection ``Continuous Big $q$-Hermite to Al-Salam--Chihara.''
Fix $0<q<1$, a finite degree $n$, and an admissible fixed parameter $a$.
The V31 reverse connection is
\begin{equation}\label{eq:v81-reverse}
 Q_n(x;a,b\mid q)=\sum_{k=0}^n {n\brack k}_q
 (-b)^{n-k}q^{\binom{n-k}{2}}H_k(x;a\mid q).
\end{equation}
\paragraph{Coefficientwise boundary calculation.}
For each $k<n$ the coefficient in \eqref{eq:v81-reverse}
vanishes as $b\to0$, since $n-k\ge1$.  The $k=n$ coefficient
is exactly $1$.  Since the sum is finite, passage to the limit is
legitimate coefficientwise, uniformly for $x$ in compact subsets:
\begin{equation}\label{eq:v81-limit}
 \lim_{b\to0}Q_n(x;a,b\mid q)=H_n(x;a\mid q).
\end{equation}
The opposite V31 connection
\[
 H_n(x;a\mid q)=\sum_{k=0}^n {n\brack k}_q
 b^{n-k}Q_k(x;a,b\mid q)
\]
independently confirms the same diagonal collapse.
\paragraph{Specialization versus limit.}
Both sides are polynomials in $b$, so evaluation at $b=0$ is already
well-defined.  In fact the standard generating functions satisfy
\[
 \sum_{n\ge0}\frac{Q_n(x;a,b\mid q)}{(q;q)_n}t^n
 =\frac{(at,bt;q)_\infty}{(e^{i\theta}t,e^{-i\theta}t;q)_\infty},
 \qquad x=\cos\theta,
\]
whereas the continuous big $q$-Hermite generating function has numerator
$(at;q)_\infty$. Thus
\begin{equation}\label{eq:v81-exact}
 Q_n(x;a,0\mid q)=H_n(x;a\mid q)
\end{equation}
\emph{exactly}, for every finite $n$.  Equation \eqref{eq:v81-limit}
is a valid coefficientwise limiting statement and confirms the
previously audited Al-Salam--Chihara $\to$ continuous big
$q$-Hermite boundary, but it is not a \emph{new independent} eligible
limiting relation: it is the same already-counted family arrow and
its boundary is also an exact specialization.  It must not be
counted a second time.

\noindent\textbf{Verdict: VERIFIED CROSS-CHECK; NO NEW PASS.}
The reverse coefficients give a complete proof of the boundary
collapse, but the cumulative number of distinct eligible reverse-
connection limiting relations remains \textbf{43}.



\section{V83: $q$-Meixner--Pollaczek and continuous $q$-Hahn --- reverse boundary and exact specialization}
\noindent\textbf{Source:} V31, reverse connection
\texttt{eq:cqh-qmp23}, with definitions \texttt{eq:qmp-standard23}.
Fix $0<q<1$, degree $n$, $z=e^{i\phi}$, $u=az$, $v=a/z$,
and $h=bz$.  V31 defines the two normalized families by
\[
 \mathscr P_n(x;a,\phi\mid q)
 =\frac{p_n(x;u,0,0,v\mid q)}{(q;q)_n},\qquad
 \mathscr Q_n(x;a,b,0,a,\phi\mid q)
 =p_n(x;u,h,0,v\mid q).
\]
These are Askey--Wilson realizations; the second is the
continuous $q$-Hahn realization used in V31.

\paragraph{Termwise reverse-connection test.}
The manuscript's finite reverse connection reads
\begin{equation}\label{eq:v83-reverse}
 \mathscr Q_n(x;a,b,0,a,\phi\mid q)
 =\sum_{k=0}^n {n\brack k}_q (q;q)_k
 (-h)^{n-k}q^{\binom{n-k}{2}}
 (a^2q^k;q)_{n-k}\,
 \mathscr P_k(x;a,\phi\mid q).
\end{equation}
For fixed $a,\phi,q$ and $h\to0$, every term with $k<n$
vanishes, because it contains $h^{n-k}$ and the remaining factors
are independent of $h$.  At $k=n$ the coefficient is $(q;q)_n$.
Thus the finite sum gives the coefficientwise limit
\begin{equation}\label{eq:v83-limit}
 \lim_{b\to0}\mathscr Q_n(x;a,b,0,a,\phi\mid q)
 =(q;q)_n\mathscr P_n(x;a,\phi\mid q).
\end{equation}
There is no hidden cancellation or interchange of infinite sums:
for each $n$ the reverse connection contains only $n+1$ terms.

\paragraph{Exact endpoint check and classification.}
By the very definitions of the two families, setting $b=0$
sets $h=0$ and gives
\begin{equation}\label{eq:v83-exact}
 \mathscr Q_n(x;a,0,0,a,\phi\mid q)
 =p_n(x;u,0,0,v\mid q)
 =(q;q)_n\mathscr P_n(x;a,\phi\mid q)
\end{equation}
\emph{exactly}, for each admissible fixed $a,\phi,q$.
Equation \eqref{eq:v83-limit} is therefore a useful reverse-connection
boundary consistency check, but the endpoint is an exact
zero-parameter specialization of the same Askey--Wilson polynomial.
In particular, the normalization $(q;q)_n$ is essential: the two
families as defined in V31 do not become literally equal at $b=0$.
No distinct singular KLS family degeneration is established by this
boundary test, and no separate eligible PASS is added.

\noindent\textbf{Verdict: EXCLUDED --- exact zero-parameter
specialization (reverse boundary verified).} The cumulative eligible
reverse-connection limiting-proof tally remains \textbf{43}.

\section{V82: Continuous $q$-Laguerre and Al-Salam--Chihara --- matched reverse connection is an exact embedding}
\noindent\textbf{V31 source:} \texttt{eq:lag-asc-hyper-rev51}, with the normalization \texttt{eq:qlag-asc} and the underlying Al-Salam--Chihara parameter connection \texttt{eq:asc-hypergeom-rev48}.
Fix $0<q<1$ and a degree $n$. Write
\[
 a=q^{\lambda/2+1/4},\qquad b=aq^{1/2},\qquad
 L_k^{(\lambda)}(x\mid q)=\frac{a^k}{(q;q)_k}Q_k(x;a,b\mid q).
\]
The reverse formula in V31 expands $Q_n(x;u,v\mid q)$ in the
$L_k^{(\lambda)}$ basis. Its generic terminating ${}_2\phi_1$
representation has removable singularities at the matched pair
$(u,v)=(a,b)$, so substituting into the displayed hypergeometric
quotients directly is not a valid evaluation procedure.

\paragraph{Regular finite-sum form.}
The proof of the V31 reverse formula supplies the nonsingular
Al-Salam--Chihara coefficient kernel. With $m=n-k$, it gives
\begin{equation}\label{eq:v82-regular}
 Q_n(x;u,v\mid q)=\sum_{k=0}^n {n\brack k}_q
 \frac{(q;q)_k}{a^k}\,\mathcal K_{n-k}(u,v;a,b)
 L_k^{(\lambda)}(x\mid q),
\end{equation}
where
\begin{equation}\label{eq:v82-kernel}
 \mathcal K_m(u,v;a,b)=
 \sum_{r=0}^m {m\brack r}_q
 a^r(u/a;q)_r\,b^{m-r}(v/b;q)_{m-r}.
\end{equation}
This is the reverse parameter-change kernel obtained by taking
$(u,v)$ as source and $(a,b)$ as target in the V31 Al-Salam--Chihara
connection. In particular, \eqref{eq:v82-kernel} is polynomial in
$u,v$, even where the hypergeometric presentation is singular.

\paragraph{Exact matched-pair collapse.}
At $u=a$, $v=b$, the factors $(1;q)_j$ vanish for all $j\ge1$.
For $m>0$, every summand of \eqref{eq:v82-kernel} contains
$(1;q)_r$ or $(1;q)_{m-r}$ with a positive index. Thus
\[
 \mathcal K_0(a,b;a,b)=1,\qquad
 \mathcal K_m(a,b;a,b)=0\quad(m\ge1).
\]
Consequently \eqref{eq:v82-regular} reduces to the \emph{exact}
identity
\begin{equation}\label{eq:v82-exact}
 Q_n(x;a,aq^{1/2}\mid q)=\frac{(q;q)_n}{a^n}
 L_n^{(\lambda)}(x\mid q).
\end{equation}
There is no singular parameter limit: $a$ is fixed and nonzero,
and the matched target pair is admissible. For comparison, the
first off-diagonal kernel is
\[
 \mathcal K_1(u,v;a,b)=(a-u)+(b-v),
\]
which is generally nonzero. Hence the general reverse connection
is a genuine finite basis change, while its matched collapse is an
exact identification, not a newly established family degeneration.

\noindent\textbf{Verdict: EXCLUDED --- EXACT SPECIALIZATION.}
No new distinct Koekoek--Swarttouw limiting relation follows from
this matched-parameter test. The cumulative number of independently
credited reverse-connection limiting proofs remains \textbf{43}.


\section{V84: Continuous $q$-Hahn to Askey--Wilson --- same-family reverse transport, not a KLS degeneration}
\noindent\textbf{V31 source:} definition \texttt{eq:cqh-def23} and reverse connection \texttt{eq:cqh-aw-two-reverse-71} (with forward \texttt{eq:cqh-aw-two-forward-71}).
Fix $0<q<1$, $z=e^{i\phi}$, and parameters for which the relevant polynomial bases are nondegenerate. The V31 definition is
\begin{equation}\label{eq:v84-def}
 \mathscr Q_n(x;a,b,c,d,\phi\mid q)
 =p_n(x;u,v,w,t\mid q),\quad
 (u,v,w,t)=(az,bz,c/z,d/z).
\end{equation}
The reverse identity expands $p_n(x;u,v,f,g\mid q)$ in the basis $\mathscr Q_k(x;a,b,c,d,\phi\mid q)$, changing the last two Askey--Wilson parameters from $(w,t)$ to $(f,g)$.

\paragraph{Matched-pair test on the reverse coefficients.}
Choose $(f,g)=(w,t)$. The left side of the V31 reverse connection is then, by \eqref{eq:v84-def}, exactly $\mathscr Q_n$. Consequently the reverse identity reads
\[
 \mathscr Q_n(x)=\sum_{k=0}^n C_{nk}(w,t;w,t)\mathscr Q_k(x).
\]
For generic nondegenerate parameters the $\mathscr Q_k$ have strictly increasing degrees $k$ and nonzero leading coefficients. Uniqueness of expansion in this polynomial basis forces
\begin{equation}\label{eq:v84-identity}
 C_{nk}(w,t;w,t)=\delta_{nk}\qquad(0\le k\le n).
\end{equation}
In particular, the diagonal factor stated in V31,
\[
 C_{nn}=\frac{(uvfgq^{n-1};q)_n}{(uvwtq^{n-1};q)_n},
\]
becomes $1$ at $(f,g)=(w,t)$. The terminating ${}_4\phi_3$ representation for off-diagonal coefficients contains apparent $0/0$ expressions at the matched pair; the substitution must therefore be performed in the underlying finite two-step connection (or by polynomial-basis uniqueness), \emph{not} term by term in those singular quotients.

\paragraph{Classification and scope.}
Equation \eqref{eq:v84-def} is a reparameterization of the \emph{same} Askey--Wilson polynomial, with the same variable and degree. The reverse identity is an Askey--Wilson parameter-transport formula, not a distinct-family KLS degeneration. Its matched endpoint \eqref{eq:v84-identity} is an exact identity at finite parameters; taking $(f,g)\to(w,t)$ merely recovers that identity continuously and does not create a new limiting arrow. This classification concerns the particular reverse connection in V31 and does not deny other Askey--Wilson degenerations with independently specified singular scalings.

\noindent\textbf{Verdict: EXCLUDED --- same-family parameter transport and exact matched specialization.} No additional eligible reverse-connection limiting proof is credited; the cumulative PASS count remains \textbf{43}.


\section{V85: Continuous $q$-Laguerre self-connection --- reverse-kernel identity and a singular-normalization warning}
\noindent\textbf{V31 source:} reverse connection \texttt{eq:lag-self-rev-hyp}, kernel \texttt{eq:lag-self-kernel-hyp}, and the normalization in Subsection \texttt{sec:lag-self}.
Fix $0<q<1$ and write
\[
 a=q^{\lambda/2+1/4},\quad b=aq^{1/2},\qquad
 c=q^{\mu/2+1/4},\quad d=cq^{1/2},\qquad
 L_n^{(\lambda)}(x\mid q)=\frac{a^n}{(q;q)_n}Q_n(x;a,b\mid q).
\]
The reverse connection in V31 changes the continuous $q$-Laguerre parameter $\mu$ to $\lambda$ within the \emph{same} family:
\begin{equation}\label{eq:v85-reverse}
 L_n^{(\mu)}(x\mid q)=\sum_{k=0}^n D_{nk}(c,d;a,b)L_k^{(\lambda)}(x\mid q).
\end{equation}
\paragraph{Regular finite kernel and matched-parameter check.}
To avoid removable singularities in the terminating ${}_2\phi_1$ coefficients, use the finite kernel obtained from V31 by interchanging $(a,b)$ and $(c,d)$:
\begin{equation}\label{eq:v85-kernel}
 \mathcal K_m(c,d;a,b)=\sum_{r=0}^{m}{m\brack r}_q
 a^r(c/a;q)_r b^{m-r}(d/b;q)_{m-r}.
\end{equation}
In this notation the coefficient of $L_k^{(\lambda)}$ in \eqref{eq:v85-reverse} is
\begin{equation}\label{eq:v85-coefficient}
 D_{nk}={n\brack k}_q\frac{c^n}{a^k}\frac{(q;q)_k}{(q;q)_n}
 \mathcal K_{n-k}(c,d;a,b).
\end{equation}
This follows directly by specializing the Al-Salam--Chihara connection to the two pairs and applying the displayed normalization; it is the nonsingular finite-sum interpretation of V31's reverse ${}_2\phi_1$ coefficient.
At $\mu=\lambda$, hence $(c,d)=(a,b)$, every summand of \eqref{eq:v85-kernel} with $m>0$ contains either $(1;q)_r$ for $r>0$ or $(1;q)_{m-r}$ for $m-r>0$. Thus
\[
 \mathcal K_m(a,b;a,b)=\delta_{m0},\qquad
 D_{nk}(a,b;a,b)=\delta_{nk}.
\]
For a nontrivial first check, $\mathcal K_1(c,d;a,b)=(a-c)+(b-d)$, so the coefficient immediately below the diagonal is generally nonzero. The diagonal coefficient in \eqref{eq:v85-coefficient} is $D_{nn}=(c/a)^n$, exactly as stated in V31.
\paragraph{Why a zero-parameter endpoint is not a new KLS arrow here.}
Although the Al-Salam--Chihara polynomial has the finite endpoint $Q_n(x;0,0\mid q)=H_n(x\mid q)$, the continuous $q$-Laguerre normalization includes $c^n$. In particular, for fixed $n>0$ and $x$,
\[
 \lim_{\mu\to+\infty}L_n^{(\mu)}(x\mid q)=0,
 \qquad\text{whereas}\qquad
 \lim_{\mu\to+\infty}\frac{(q;q)_n}{c^n}L_n^{(\mu)}(x\mid q)=H_n(x\mid q).
\]
The latter is precisely the continuous $q$-Laguerre to continuous $q$-Hermite degeneration already proved in V68; it is not a new self-family arrow. Moreover, in the raw reverse coefficients \eqref{eq:v85-coefficient}, the factors $c^n$ and $a^{-k}$ make it inappropriate to infer the Hermite limit by simply setting $c=0$ while keeping the unscaled Laguerre basis. The finite-kernel and normalization must be retained together.

\noindent\textbf{Verdict: EXCLUDED --- same-family parameter connection.} The matched case is an exact identity; the singularly normalized Hermite endpoint duplicates V68. No new distinct reverse-connection limiting proof is credited. The cumulative PASS count remains \textbf{43}.


\section{V86: Continuous $q$-Hahn self-connection --- exact parameter transport, not a family limit}
\noindent\textbf{V31 source:} \texttt{eq:cqh-self-two-reverse-71} in Subsection \texttt{sec:cqh-self-clean}. Fix $0<q<1$, a phase $z=e^{i\phi}$, and parameters $a,b,c,d,c',d'$ away from vanishing leading coefficients. V31 uses the Askey--Wilson realization
\[
 \mathscr Q_n(x;a,b,c,d,\phi\mid q)
 =p_n(x;az,bz,c/z,d/z\mid q).
\]
Writing $u=az$, $v=bz$, $w=c/z$, $t=d/z$, $f=c'/z$, $g=d'/z$, its reverse connection is
\begin{equation}\label{eq:v86-reverse}
 \mathscr Q_n(x;a,b,c',d',\phi\mid q)
 =\sum_{k=0}^n D_{nk}(f,g;w,t)\,
 \mathscr Q_k(x;a,b,c,d,\phi\mid q).
\end{equation}
\paragraph{Diagonal and first-degree checks.}
The diagonal coefficient explicitly stated in V31 is
\begin{equation}\label{eq:v86-diag}
 D_{nn}=\frac{(uvfgq^{n-1};q)_n}{(uvwtq^{n-1};q)_n}
 =\frac{(abc'd'q^{n-1};q)_n}{(abcdq^{n-1};q)_n}.
\end{equation}
This factor is generally not one, demonstrating that the connection is a genuine change of basis, but it does not imply a degeneration between distinct polynomial families. For example, at degree one, the leading coefficient of the standard Askey--Wilson polynomial is proportional to $1-uvwt$, whereas that of the target is proportional to $1-uvfg$; their ratio is precisely $D_{11}=(1-uvfg)/(1-uvwt)$. This also checks the direction of the ratio in \eqref{eq:v86-diag}.
\paragraph{Matched-parameter test without singular hypergeometric substitution.}
If $c'=c$ and $d'=d$, then $(f,g)=(w,t)$, and the polynomial on the left of \eqref{eq:v86-reverse} equals the degree-$n$ basis element on the right. Since the Askey--Wilson polynomials of degrees $0,1,\ldots,n$ are linearly independent at generic parameters, uniqueness of polynomial expansion gives
\begin{equation}\label{eq:v86-delta}
 D_{nk}(w,t;w,t)=\delta_{nk}\qquad(0\le k\le n).
\end{equation}
Equation \eqref{eq:v86-diag} independently gives $D_{nn}=1$. The terminating ${}_4\phi_3$ coefficient formula has apparent $0/0$ quotients at this matched point; \eqref{eq:v86-delta} is justified by the polynomial identity (or the nonsingular finite convolution), not by unqualified substitution into those quotients.
\paragraph{Classification.}
Both sides of \eqref{eq:v86-reverse} belong to the same continuous $q$-Hahn realization of the Askey--Wilson family. Letting $(c',d')\to(c,d)$ is a regular parameter-matching operation yielding the exact identity \eqref{eq:v86-delta}; no distinct KLS family, singular scaling, or new normalization is produced. Singular limits of Askey--Wilson polynomials to other families require separately specified paths and cannot be inferred from this matched-pair calculation.

\noindent\textbf{Verdict: EXCLUDED --- same-family parameter transport.} No new eligible reverse-connection limit is proved. The cumulative PASS count remains \textbf{43}.


\section{V87: Little $q$-Jacobi to alternative $q$-Charlier ($q$-Bessel) --- distinct-family singular scaling}
\noindent\textbf{Scope pre-screen:} the two polynomial families are distinct, so this is not a self-connection. The relevant V31 reverse connection is \texttt{eq:qBlqJ-reverse}, with definitions \texttt{eq:qBlqJ-defB} and \texttt{eq:qBlqJ-defJ}. Fix $0<q<1$, a degree $n$, and a parameter $a$ for which the polynomial degrees remain nondegenerate. Consider the singular path
\[
 A=\varepsilon\longrightarrow0,\qquad B=-\frac{a}{q\varepsilon},\qquad AB=-\frac aq.
\]
In the normalization of V31, the little $q$-Jacobi polynomial has the finite terminating representation
\begin{equation}\label{eq:v87-p}
 p_n(x;\varepsilon,-a/(q\varepsilon)\mid q)
 =\sum_{j=0}^n\frac{(q^{-n};q)_j(-aq^n;q)_j}{(\varepsilon q;q)_j(q;q)_j}(qx)^j.
\end{equation}
Since $(\varepsilon q;q)_j\to1$ for each $j\le n$, the finite sum gives the coefficientwise, locally uniform polynomial limit
\begin{equation}\label{eq:v87-limit}
 \boxed{\displaystyle\lim_{\varepsilon\to0}
 p_n\left(x;\varepsilon,-\frac{a}{q\varepsilon}\middle|q\right)
 ={}_2\phi_1\!\left(\begin{matrix}q^{-n},-aq^n\\0\end{matrix};q,qx\right)
 =y_n(x;a;q).}
\end{equation}
This is a \emph{singular parameter path}: $B$ diverges while $A$ vanishes. It is not a finite matched-parameter specialization of the little $q$-Jacobi family.

\paragraph{Reverse-connection coefficient proof.}
Write V31's reverse identity as
\begin{equation}\label{eq:v87-reverse}
 p_n(x;A,B\mid q)=\sum_{k=0}^n D_{nk}(A,B;a)\,y_k(x;a;q).
\end{equation}
Let $Y_{kj}$ denote the coefficient of $x^j$ in $y_k(x;a;q)$, and let $P_{nj}(\varepsilon)$ denote the coefficient of $x^j$ in the left side of \eqref{eq:v87-p}. For generic $a$, the triangular matrix $Y=(Y_{kj})_{0\le j\le k\le n}$ has nonzero diagonal, since
\[
 Y_{kk}=\frac{(q^{-k};q)_k(-aq^k;q)_k q^k}{(q;q)_k}\ne0.
\]
Thus the coefficients $D_{nk}(\varepsilon,-a/(q\varepsilon);a)$ are uniquely obtained from a fixed invertible triangular linear system with coefficient vector $P_n(\varepsilon)$. Equation \eqref{eq:v87-p} shows $P_{nj}(\varepsilon)\to Y_{nj}$. Applying the fixed inverse triangular matrix gives
\begin{equation}\label{eq:v87-delta}
 \boxed{\displaystyle\lim_{\varepsilon\to0}
 D_{nk}\left(\varepsilon,-\frac{a}{q\varepsilon};a\right)=\delta_{nk}
 \quad(0\le k\le n).}
\end{equation}
This uses the \emph{reverse connection itself}, not merely the family definitions, and avoids unsafe termwise substitution into its terminating ${}_3\phi_2$ coefficient formula. The convergence is coefficientwise for every fixed degree and hence uniform on compact subsets of the complex $x$-plane.

\paragraph{Classification and literature checkpoint.}
The distinct-family singular limit \eqref{eq:v87-limit} and the reverse-coefficient collapse \eqref{eq:v87-delta} are proved. However, the manuscript's literature note for Connection 97 cites a \emph{different} Koekoek--Swarttouw path (standard $q$-Krawtchouk to alternative $q$-Charlier), and does not itself identify the little $q$-Jacobi path above as a numbered KLS limiting relation. Under the audit's requirement of an \emph{established} KLS arrow, that external identification remains unverified; the calculation is therefore recorded without increasing the eligible PASS count.

\noindent\textbf{Verdict: PROVED DISTINCT-FAMILY SINGULAR LIMIT; KLS ELIGIBILITY PENDING VERIFICATION.} No self-connections were examined in this entry. Cumulative confirmed eligible PASS count: \textbf{43}.


\section{V88: Duplicate detection and KLS eligibility reconciliation for the little $q$-Jacobi to alternative $q$-Charlier arrow}
\noindent\textbf{Scope gate:} This is a distinct-family arrow, not a self-connection. The purpose of this entry is to resolve a material inconsistency in the cumulative audit before searching for further unproved arrows. No new family is admitted merely because its connection appears later in the manuscript.

\paragraph{Earlier established proof.}
The section entitled ``$q$-Bessel (alternative $q$-Charlier) from little $q$-Jacobi: complete reverse-connection limit'' already cites V31 equation \texttt{eq:qBlqJ-reverse} and identifies Koekoek--Swarttouw (1998), Eq.~(4.12.3). Its proof sets $B=-a/(Aq)$, uses the lower-triangular monomial coefficient matrix
\[
 T_{rj}(A)=\frac{(q^{-r},-aq^r;q)_j q^j}{(Aq,q;q)_j},\qquad 0\le j\le r\le n,
\]
and establishes the \emph{actual reverse-connection matrix identity}
\[
 C^{(n)}(A)=T^{(n)}(A)\bigl(T^{(n)}(0)\bigr)^{-1}
 \longrightarrow I_{n+1}\qquad(A\to0),
\]
for every fixed $n$ under the stated nondegeneracy assumptions. This already proves
\begin{equation}\label{eq:v88-duplicate}
 \lim_{A\to0}p_n\left(x;A,-\frac{a}{Aq}\middle|q\right)=y_n(x;a;q).
\end{equation}
The earlier ``Eligibility cross-check for Connections 95--101 (V57)'' explicitly records Connection 96 as \textbf{ALREADY AUDITED}, and the cumulative progress list already includes this very arrow as a \textbf{PASS} with the same V31 equation label and KLS equation reference.

\paragraph{Comparison with V87.}
V87 substitutes $A=\varepsilon$ and $B=-a/(q\varepsilon)$, proves \eqref{eq:v88-duplicate} again, and again obtains $D_{nk}\to\delta_{nk}$ by inversion of the fixed alternative-$q$-Charlier monomial coefficient matrix. Thus the parameter path, polynomial normalization, reverse-connection source, and conclusion in V87 are identical to those of the earlier PASS. V87's statement that KLS eligibility remains unverified is \emph{inconsistent with the prior recorded citation} and its characterization as a newly proved distinct-family arrow is incorrect. This is an internal documentary reconciliation: it confirms the earlier KLS citation is present in the audit, without claiming a new independent external bibliographic check.

\noindent\textbf{Corrected V87 status: DUPLICATE OF AN EARLIER PASS, NOT A NEW RESULT.} The V87 proof remains valid but supplies no additional independent limiting relation. The previously credited arrow is counted exactly once. \textbf{Cumulative eligible PASS count: 43 (unchanged).} In keeping with the user's revised scope, self-family connections are not candidates for future substantive entries.


\section{V89: Scope-first triage of distinct-family candidates (no self-connections)}
\noindent\textbf{Purpose and scope.} Following the user's scope clarification, this iteration performs a \emph{pre-proof eligibility gate} rather than repeating an already established arrow. Same-family parameter connections are not candidate limiting relations and are not analysed here. An eligible new PASS requires (i) two distinct families, (ii) an established KLS degeneration with specified scaling and normalization, (iii) a V31 reverse-connection proof, and (iv) a check against the 43 already credited arrows.

\paragraph{Candidate 1: Big $q$-Laguerre to Wall.}
V31's distinct-family reverse identity \texttt{eq:wall-big-reverse} is a plausible source, but the cumulative audit already proved the singular $B\to\infty$, $x=Bqz$ limit in V58. \textbf{Screen result: previously proved; do not duplicate.}

\paragraph{Candidate 2: Little $q$-Jacobi to alternative $q$-Charlier.}
V31's distinct-family reverse identity \texttt{eq:qBlqJ-reverse} has the singular path $A\to0$, $B=-a/(Aq)$. The earlier PASS and the V87--V88 reconciliation prove that the reverse matrix tends to the identity. \textbf{Screen result: previously proved; do not duplicate.}

\paragraph{Candidate 3: Meixner--Pollaczek to continuous dual Hahn.}
V31's reverse connection \texttt{eq:consolidated-cdh-mp} is a connection between distinct families, but its matched endpoints $(b,c)=(\tfrac12,0)$ and $(\tfrac12,1)$ are exact even/odd quadratic identities. This was verified in V66. No additional singular KLS path is established by that parameter matching. \textbf{Screen result: exact quadratic endpoint excluded; no new PASS.}

\paragraph{Methodological consequence.}
These are \emph{screening results}, not three new mathematical limiting proofs. In particular, a reverse connection at fixed parameters is insufficient evidence of a KLS arrow, and a proof of an already counted limit cannot increase the tally. Future substantive work should begin with a not-yet-credited KLS arrow and verify that its required scaling is actually available in V31's reverse identity before performing asymptotic calculations.

\noindent\textbf{Verdict: NO NEW ELIGIBLE LIMIT IDENTIFIED IN THIS TRIAGE.} No same-family connections have been admitted to scope. The cumulative independently credited PASS count remains \textbf{43}.


\section{V90: Exhaustive V31 heading inventory and completeness certification gate}
\noindent\textbf{Purpose.} This is the requested end-of-project coverage reconciliation, not a new limit proof. We inventory \emph{every} V31 polynomial-pair subsubsection (in manuscript order), apply the self-family scope gate, and distinguish a source-heading census from a fully verified Koekoek--Swarttouw arrow census. The latter cannot be certified by counting manuscript headings: a pair can contain several arrows, one arrow, only an exact specialization, or no KLS arrow.

\paragraph{Inventory method and reproducibility.} The table below is generated by scanning all V31 \texttt{\\subsubsection} headings. The first column gives the original source line number, permitting direct navigation to the corresponding connection formulas. An \textbf{OUT OF SCOPE} entry is a same-family heading; all remaining entries are retained for independent reconciliation against the KLS limiting scheme and the 43-item proof ledger. The designation ``KLS/record check required'' does \emph{not} assert that the pair is an unproved eligible arrow: it deliberately leaves that decision open until both the external KLS arrow and the particular V31 reverse identity are checked. This conservative convention prevents the false inference that every distinct-family connection represents a new limit.

\begingroup\footnotesize
\setlength{\tabcolsep}{4pt}
\begin{longtable}{@{}p{0.10\linewidth}p{0.49\linewidth}p{0.32\linewidth}@{}}
\toprule V31 line & Family-pair heading & Inventory gate \\
\midrule\endfirsthead
\toprule V31 line & Family-pair heading & Inventory gate \\
\midrule\endhead
\bottomrule\endfoot
\texttt{106} & Wilson / Wilson & OUT OF SCOPE \\
\texttt{176} & Continuous Dual Hahn / Wilson & KLS/record check required \\
\texttt{229} & Continuous Dual Hahn / Continuous Dual Hahn & OUT OF SCOPE \\
\texttt{267} & Continuous Hahn / Wilson & KLS/record check required \\
\texttt{355} & Continuous Hahn / Continuous Hahn & OUT OF SCOPE \\
\texttt{426} & Meixner--Pollaczek / Continuous Hahn & KLS/record check required \\
\texttt{509} & Parameter cancellation: one terminating 2F1 in each direction & KLS/record check required \\
\texttt{609} & Meixner--Pollaczek / Continuous Dual Hahn & KLS/record check required \\
\texttt{701} & Meixner--Pollaczek / Meixner--Pollaczek & OUT OF SCOPE \\
\texttt{822} & Jacobi / Wilson & KLS/record check required \\
\texttt{975} & Jacobi / Continuous Hahn & KLS/record check required \\
\texttt{1074} & Jacobi / Jacobi & OUT OF SCOPE \\
\texttt{1211} & Generalised Laguerre / Meixner--Pollaczek & KLS/record check required \\
\texttt{1296} & Generalised Laguerre / Jacobi & KLS/record check required \\
\texttt{1406} & Generalised Laguerre / Generalised Laguerre & OUT OF SCOPE \\
\texttt{1454} & Hermite / Meixner--Pollaczek & KLS/record check required \\
\texttt{1527} & Hermite / Jacobi & KLS/record check required \\
\texttt{1582} & Hermite / Generalised Laguerre & KLS/record check required \\
\texttt{1622} & Hermite / Hermite & OUT OF SCOPE \\
\texttt{1700} & Racah / Racah & OUT OF SCOPE \\
\texttt{1731} & Hahn / Racah & KLS/record check required \\
\texttt{1786} & Hahn / Hahn & OUT OF SCOPE \\
\texttt{1837} & Dual Hahn / Racah & KLS/record check required \\
\texttt{1908} & Dual Hahn / Dual Hahn & OUT OF SCOPE \\
\texttt{1927} & Pseudo-Jacobi / Pseudo-Jacobi & OUT OF SCOPE \\
\texttt{1978} & Meixner / Hahn & KLS/record check required \\
\texttt{2027} & Meixner / Dual Hahn & KLS/record check required \\
\texttt{2123} & Meixner / Meixner & OUT OF SCOPE \\
\texttt{2154} & Krawtchouk / Hahn & KLS/record check required \\
\texttt{2236} & Krawtchouk / Dual Hahn & KLS/record check required \\
\texttt{2266} & Krawtchouk / Krawtchouk & OUT OF SCOPE \\
\texttt{2298} & Generalised Bessel / Generalised Bessel & OUT OF SCOPE \\
\texttt{2362} & Charlier / Meixner & KLS/record check required \\
\texttt{2397} & Charlier / Krawtchouk & KLS/record check required \\
\texttt{2432} & Charlier / Charlier & OUT OF SCOPE \\
\texttt{2459} & Racah / Wilson & KLS/record check required \\
\texttt{2505} & Pseudo-Jacobi / Continuous Hahn & KLS/record check required \\
\texttt{2579} & Jacobi / Hahn & KLS/record check required \\
\texttt{2659} & Generalised Bessel / Jacobi & KLS/record check required \\
\texttt{2745} & Generalised Laguerre / Meixner & KLS/record check required \\
\texttt{2817} & Hermite / Krawtchouk & KLS/record check required \\
\texttt{2874} & Hermite / Charlier & KLS/record check required \\
\texttt{2977} & Askey--Wilson / Askey--Wilson & OUT OF SCOPE \\
\texttt{3153} & Continuous dual q-Hahn / Askey--Wilson & KLS/record check required \\
\texttt{3204} & Continuous dual q-Hahn / Continuous dual q-Hahn & OUT OF SCOPE \\
\texttt{3324} & Continuous q-Hahn / Askey--Wilson & KLS/record check required \\
\texttt{3386} & Continuous q-Hahn / Continuous q-Hahn & OUT OF SCOPE \\
\texttt{3441} & Al-Salam--Chihara / Continuous dual q-Hahn & KLS/record check required \\
\texttt{3493} & Al-Salam--Chihara / Al-Salam--Chihara & OUT OF SCOPE \\
\texttt{3582} & q-Meixner--Pollaczek / Continuous q-Hahn & KLS/record check required \\
\texttt{3637} & q-Meixner--Pollaczek / q-Meixner--Pollaczek & OUT OF SCOPE \\
\texttt{3701} & Continuous q-Jacobi / Askey--Wilson & KLS/record check required \\
\texttt{3773} & Continuous q-Jacobi / Continuous q-Jacobi & OUT OF SCOPE \\
\texttt{3829} & Continuous Big q-Hermite / Al-Salam--Chihara & KLS/record check required \\
\texttt{3943} & Continuous Big q-Hermite / Continuous Big q-Hermite & OUT OF SCOPE \\
\texttt{3991} & Continuous q-Laguerre / Al-Salam--Chihara & KLS/record check required \\
\texttt{4060} & Continuous q-Laguerre / q-Meixner--Pollaczek & KLS/record check required \\
\texttt{4145} & Continuous q-Laguerre / Continuous q-Jacobi & KLS/record check required \\
\texttt{4230} & Continuous q-Laguerre / Continuous q-Laguerre & OUT OF SCOPE \\
\texttt{4315} & Continuous q-Hermite / Continuous Big q-Hermite & KLS/record check required \\
\texttt{4366} & Continuous q-Hermite / Continuous q-Laguerre & KLS/record check required \\
\texttt{4469} & Continuous q-Hermite / Continuous q-Jacobi & KLS/record check required \\
\texttt{5014} & Continuous q-Hermite / Continuous q-Hermite & OUT OF SCOPE \\
\texttt{5078} & q-Racah / q-Racah & OUT OF SCOPE \\
\texttt{5230} & Big q-Jacobi / q-Racah & KLS/record check required \\
\texttt{5345} & Big q-Jacobi / Big q-Jacobi & OUT OF SCOPE \\
\texttt{5411} & q-Hahn / q-Racah & KLS/record check required \\
\texttt{5527} & q-Hahn / q-Hahn & OUT OF SCOPE \\
\texttt{5620} & Dual q-Hahn / q-Racah & KLS/record check required \\
\texttt{5683} & Dual q-Hahn / Dual q-Hahn & OUT OF SCOPE \\
\texttt{5783} & Big q-Laguerre / Big q-Jacobi & KLS/record check required \\
\texttt{5856} & Big q-Laguerre / Big q-Laguerre & OUT OF SCOPE \\
\texttt{5923} & Little q-Jacobi / Big q-Jacobi & KLS/record check required \\
\texttt{6012} & Little q-Jacobi / q-Hahn & KLS/record check required \\
\texttt{6076} & Little q-Jacobi / little q-Jacobi & OUT OF SCOPE \\
\texttt{6144} & q-Meixner / Big q-Jacobi & KLS/record check required \\
\texttt{6226} & q-Meixner / q-Hahn & KLS/record check required \\
\texttt{6291} & q-Meixner / q-Meixner & OUT OF SCOPE \\
\texttt{6385} & Quantum q-Krawtchouk / q-Hahn & KLS/record check required \\
\texttt{6488} & Quantum q-Krawtchouk / Quantum q-Krawtchouk & OUT OF SCOPE \\
\texttt{6560} & q-Krawtchouk / q-Hahn & KLS/record check required \\
\texttt{6654} & q-Krawtchouk / q-Racah & KLS/record check required \\
\texttt{6745} & q-Krawtchouk / q-Krawtchouk & OUT OF SCOPE \\
\texttt{6811} & Affine q-Krawtchouk / q-Hahn & KLS/record check required \\
\texttt{6878} & Affine q-Krawtchouk / Dual q-Hahn & KLS/record check required \\
\texttt{6951} & Affine q-Krawtchouk / Affine q-Krawtchouk & OUT OF SCOPE \\
\texttt{7011} & Dual q-Krawtchouk / Dual q-Hahn & KLS/record check required \\
\texttt{7084} & Dual q-Krawtchouk / Dual q-Krawtchouk & OUT OF SCOPE \\
\texttt{7149} & Little q-Laguerre / Big q-Laguerre & KLS/record check required \\
\texttt{7243} & Little q-Laguerre / Little q-Jacobi & KLS/record check required \\
\texttt{7337} & Little q-Laguerre / affine q-Krawtchouk & KLS/record check required \\
\texttt{7500} & Little q-Laguerre / little q-Laguerre & OUT OF SCOPE \\
\texttt{7584} & q-Laguerre / little q-Jacobi & KLS/record check required \\
\texttt{7648} & q-Laguerre / q-Meixner & KLS/record check required \\
\texttt{7719} & q-Laguerre / q-Laguerre & OUT OF SCOPE \\
\texttt{7787} & q-Bessel / little q-Jacobi & KLS/record check required \\
\texttt{7850} & q-Bessel / standard q-Krawtchouk & KLS/record check required \\
\texttt{7925} & q-Bessel / q-Bessel & OUT OF SCOPE \\
\texttt{7962} & Ordinary q-Charlier / q-Meixner & KLS/record check required \\
\texttt{8039} & Ordinary q-Charlier / standard q-Krawtchouk & KLS/record check required \\
\texttt{8117} & Ordinary q-Charlier / ordinary q-Charlier & OUT OF SCOPE \\
\texttt{8171} & Al-Salam--Carlitz I / Big q-Laguerre & KLS/record check required \\
\texttt{8247} & Al-Salam--Carlitz I / dual q-Krawtchouk & KLS/record check required \\
\texttt{8336} & Al-Salam--Carlitz I / Al-Salam--Carlitz I & OUT OF SCOPE \\
\texttt{8389} & Al-Salam--Carlitz II / q-Meixner & KLS/record check required \\
\texttt{8491} & Al-Salam--Carlitz II / quantum q-Krawtchouk & KLS/record check required \\
\texttt{8566} & Al-Salam--Carlitz II / Al-Salam--Carlitz II & OUT OF SCOPE \\
\texttt{8623} & Stieltjes--Wigert / q-Laguerre & KLS/record check required \\
\texttt{8705} & Stieltjes--Wigert / q-Bessel & KLS/record check required \\
\texttt{8807} & Stieltjes--Wigert / q-Charlier & KLS/record check required \\
\texttt{8925} & Stieltjes--Wigert / Stieltjes--Wigert & OUT OF SCOPE \\
\texttt{8985} & Discrete q-Hermite I / Al-Salam--Carlitz I & KLS/record check required \\
\texttt{9044} & Discrete q-Hermite I / discrete q-Hermite I & OUT OF SCOPE \\
\texttt{9093} & Discrete q-Hermite II / Al-Salam--Carlitz II & KLS/record check required \\
\texttt{9167} & Discrete q-Hermite II / discrete q-Hermite II & OUT OF SCOPE \\
\texttt{9287} & Askey--Wilson / q-Racah & KLS/record check required \\
\texttt{9479} & Askey--Wilson / Big q-Jacobi & KLS/record check required \\
\end{longtable}\endgroup

\paragraph{Ledger reconciliation.} The cumulative audit currently reports \textbf{43} eligible PASS arrows. The V87--V88 duplicate was reconciled without adding a PASS. The V85--V86 self-family discussions are historical notes only and are not eligible coverage items. Importantly, a list of 43 PASS labels does not independently establish (i) a one-to-one mapping from each proof to a distinct KLS arrow, (ii) that every cited KLS equation is correct, or (iii) that every eligible KLS arrow accessible from V31 has been found. Some earlier entries themselves explicitly leave KLS equation-number matching open (for example, the symmetric continuous Hahn/Meixner--Pollaczek and Wilson/Jacobi paths).

\paragraph{Certification decision.} A source-heading census is complete, but the \emph{mathematical coverage certification is not complete}. The outstanding verification is to construct an external KLS-arrow ledger, match every distinct-family V31 heading to that ledger (including direction and parameter normalization), and cross-reference each eligible arrow to a unique checked reverse-connection proof. Without that comparison, declaring the audit exhaustive would overstate the evidence. The 43 PASS total is therefore the \emph{current recorded tally}, not an independently recertified completeness count.

\noindent\textbf{V90 verdict: SOURCE-HEADING INVENTORY COMPLETE; GLOBAL KLS COMPLETENESS NOT YET CERTIFIED.} No new PASS is claimed. Same-family entries remain outside the mathematical scope.


\section{V91: KLS-arrow reconciliation, tranche I---late discrete-$q$ families}
\noindent\textbf{Scope and method.} This is the first \emph{arrow-level} reconciliation after the V90 heading census. We examine the late discrete-$q$ entries against the explicit literature-context statements in V31, the named reverse identities, and the cumulative PASS ledger. The classifications below concern the \emph{specific candidate endpoint described in V31}, not every conceivable singular parameter path. Same-family headings are not considered.

\begin{longtable}{@{}p{0.14\linewidth}p{0.35\linewidth}p{0.43\linewidth}@{}}
\toprule V31 source & Distinct-family pair & Reconciled disposition \\
\midrule\endfirsthead
\toprule V31 source & Distinct-family pair & Reconciled disposition \\
\midrule\endhead
\bottomrule\endfoot
\texttt{8171} & Al-Salam--Carlitz I / big $q$-Laguerre & \textbf{UNRESOLVED:} Connection 102, \texttt{eq:ASCI-bL-reverse}. Earlier audit found no established KLS path from the proposed endpoint. Not a PASS. \\
\texttt{8247} & Al-Salam--Carlitz I / dual $q$-Krawtchouk & \textbf{NOT ELIGIBLE ON DOCUMENTED GRID PATH:} Connection 103 is a finite-grid interpolation on a different lattice; V55 found no specified KLS arrow. \\
\texttt{8389} & Al-Salam--Carlitz II / $q$-Meixner & \textbf{ALREADY PROVED:} Connection 105, \texttt{eq:ASCII-qM-reverse}; existing PASS. \\
\texttt{8491} & Al-Salam--Carlitz II / quantum $q$-Krawtchouk & \textbf{ALREADY PROVED:} Connection 106, \texttt{eq:ASCII-qtm-reverse}; existing PASS. \\
\texttt{8623} & Stieltjes--Wigert / $q$-Laguerre & \textbf{ALREADY PROVED:} Connection 108, existing $q$-Laguerre to Stieltjes--Wigert PASS. \\
\texttt{8705} & Stieltjes--Wigert / $q$-Bessel & \textbf{REQUIRES EXTERNAL KLS ARROW CHECK:} Distinct-family connection, but no new KLS path is certified here. \\
\texttt{8807} & Stieltjes--Wigert / $q$-Charlier & \textbf{REQUIRES EXTERNAL KLS ARROW CHECK:} Distinct-family connection, but no new KLS path is certified here. \\
\texttt{8985} & Discrete $q$-Hermite I / Al-Salam--Carlitz I & \textbf{EXACT SPECIALIZATION:} V31 Connection 112 identifies $h_n(x;q)=U_n^{(-1)}(x;q)$, not a new singular limit. \\
\texttt{9093} & Discrete $q$-Hermite II / Al-Salam--Carlitz II & \textbf{EXACT TRANSFORMATION:} V31 Connection 114 identifies $\widetilde h_n(x;q)=i^{-n}V_n^{(-1)}(ix;q)$; the argument rotation is indispensable. \\
\end{longtable}

\paragraph{Two exact-endpoint checks.}
The discrete $q$-Hermite I relation at $a=-1$ is an equality of polynomials for each $n$, so any reverse-connection matrix between the \emph{matched} bases is the identity by uniqueness. The discrete $q$-Hermite II relation instead requires $x\mapsto ix$ and multiplication by $i^{-n}$; it is not a same-variable identity at $a=-1$. Both are fixed-parameter identifications, and neither is credited as a limiting arrow. These statements follow from V31's explicit literature contexts for Connections 112 and 114; they do not require asymptotic extrapolation.

\paragraph{Reconciliation arithmetic.} Among these nine \emph{distinct-family} headings, three have already counted PASS proofs, two are exact fixed-parameter identifications, one has a documented non-KLS finite-grid obstruction, and three remain open for KLS-path determination (Connection 102 and the two Stieltjes--Wigert connections). This is \emph{not} a claim that the three open headings are missing KLS arrows. The source-heading inventory in V90 also contains the descriptive heading ``Parameter cancellation: one terminating 2F1 in each direction'' (line 509); it must be reviewed as an auxiliary heading rather than automatically interpreted as a separate family pair. Thus the V90 figures are heading-gate counts, not a certified number of distinct-family arrows.

\noindent\textbf{V91 verdict: PARTIAL ARROW RECONCILIATION; GLOBAL CERTIFICATION STILL OPEN.} No new reverse-connection limiting proof is established. The previously recorded tally remains \textbf{43}; its independent verification and exhaustiveness are still pending.

\section{Exclusions and outstanding work}
\noindent\textbf{V91 reconciliation:} Nine late discrete-$q$ distinct-family headings have been assigned documented dispositions (three previously proved, two exact identifications, one non-KLS grid path, three requiring KLS-path verification). This is a partial arrow ledger, not global certification. Count remains 43.\par
\noindent\textbf{V90 coverage checkpoint:} All 117 V31 polynomial-pair headings have been inventoried by source line. A KLS-arrow-by-arrow comparison remains necessary; exhaustive limiting-relation coverage is not yet certified. The recorded eligible PASS tally remains 43.\par
\begin{itemize}
\item Wilson $\leftrightarrow$ Wilson and all other same-family parameter connections: \textbf{excluded}.
\item Exact quadratic and other non-limiting transformations: \textbf{excluded}.
\item Other distinct-family limiting arrows in the source manuscript: \textbf{not yet exhaustively audited}. Connection 102 has no identified KLS limiting path and is not presently eligible; no unproved relation is counted.
\end{itemize}
\noindent\textbf{V89 scope update:} Only distinct-family, nonduplicate, established KLS limiting paths are eligible for substantive reverse-connection proofs. Three candidate arrows were pre-screened; none provides a new independent PASS. Same-family connections are outside scope. Count remains \textbf{43}.\par
\noindent\textbf{V88 correction:} V87 duplicates the already-proved little $q$-Jacobi to alternative $q$-Charlier limit (earlier PASS, V31 \texttt{eq:qBlqJ-reverse}, KLS Eq.~(4.12.3)). The V87 KLS-pending status is superseded; no new PASS is credited. Count remains \textbf{43}.\par
\noindent\textbf{V87 screening update:} Self-family connections are out of scope and must not receive substantive future audit entries. V87 verifies a distinct-family singular limit, but its KLS listing requires independent confirmation before it is counted.\par
\noindent\textbf{Audit protocol for subsequent entries:} cite the exact reverse-connection equation label from V31, specify the classical limit and its normalization, prove termwise asymptotics or explain why termwise passage is invalid, and record a pass/fail/undetermined verdict.

\section{V92: KLS-arrow reconciliation, tranche II---finite and intermediate discrete-$q$ families}
\noindent\textbf{Scope.} This tranche covers the fourteen distinct-family headings at V31 source lines 5230--7011. Same-family headings in the same interval are omitted entirely. A ``previously proved'' classification identifies a proof already in the cumulative ledger, not a fresh independent verification of its KLS citation. A ``specialization'' classification concerns only the stated endpoint, not every possible singular path between the named families.

\begin{longtable}{@{}p{0.12\linewidth}p{0.36\linewidth}p{0.44\linewidth}@{}}
\toprule V31 line & Distinct-family connection & Reconciliation against the existing audit \\
\midrule\endfirsthead
\toprule V31 line & Distinct-family connection & Reconciliation against the existing audit \\
\midrule\endhead
\bottomrule\endfoot
\texttt{5230} & Big $q$-Jacobi / $q$-Racah & \textbf{EXACT ENDPOINT ALREADY EXCLUDED:} V63 records the $q$-Racah specialization. General double sums are finite interpolation, not by themselves a KLS limit. \\
\texttt{5411} & $q$-Hahn / $q$-Racah & \textbf{EXACT ENDPOINT ALREADY EXCLUDED:} V72 verifies the $\delta=0$ specialization and matching quadratic lattice. \\
\texttt{5620} & Dual $q$-Hahn / $q$-Racah & \textbf{EXACT SPECIALIZATION:} V31 explicitly identifies dual $q$-Hahn as its $\beta=0$ $q$-Racah specialization; this fixed-parameter endpoint is not a new PASS. \\
\texttt{5783} & Big $q$-Laguerre / big $q$-Jacobi & \textbf{PREVIOUSLY PROVED:} V31 connection and the existing big $q$-Jacobi $\to$ big $q$-Laguerre PASS. \\
\texttt{5923} & Little $q$-Jacobi / big $q$-Jacobi & \textbf{PREVIOUSLY PROVED:} V59 big $q$-Jacobi $\to$ little $q$-Jacobi. \\
\texttt{6012} & Little $q$-Jacobi / $q$-Hahn & \textbf{PREVIOUSLY PROVED:} V62 $q$-Hahn $\to$ little $q$-Jacobi. \\
\texttt{6144} & $q$-Meixner / big $q$-Jacobi & \textbf{PREVIOUSLY PROVED:} V60 big $q$-Jacobi $\to$ $q$-Meixner. \\
\texttt{6226} & $q$-Meixner / $q$-Hahn & \textbf{PREVIOUSLY PROVED:} V61 $q$-Hahn $\to$ $q$-Meixner. \\
\texttt{6385} & Quantum $q$-Krawtchouk / $q$-Hahn & \textbf{PREVIOUSLY PROVED:} V50 $q$-Hahn $\to$ quantum $q$-Krawtchouk, with the corrected large-parameter scaling. \\
\texttt{6560} & Ordinary $q$-Krawtchouk / $q$-Hahn & \textbf{PREVIOUSLY PROVED:} V47--V49 $q$-Hahn $\to$ ordinary $q$-Krawtchouk. \\
\texttt{6654} & Ordinary $q$-Krawtchouk / $q$-Racah & \textbf{EXACT ENDPOINT ALREADY EXCLUDED:} V54 identifies the relevant fixed-parameter specialization. The general lattice interpolation is not a limiting proof. \\
\texttt{6811} & Affine $q$-Krawtchouk / $q$-Hahn & \textbf{EXACT ENDPOINT ALREADY EXCLUDED:} V52 verifies the fixed-parameter specialization. \\
\texttt{6878} & Affine $q$-Krawtchouk / dual $q$-Hahn & \textbf{EXACT ENDPOINT ALREADY EXCLUDED:} V53 verifies the fixed-parameter specialization. \\
\texttt{7011} & Dual $q$-Krawtchouk / dual $q$-Hahn & \textbf{PREVIOUSLY PROVED:} V45 dual $q$-Hahn $\to$ dual $q$-Krawtchouk, with V31's explicit KLS Eq.~(4.7.2) citation. \\
\end{longtable}

\paragraph{Independent check of the exact dual $q$-Hahn endpoint.}
In the V31 normalization, the $q$-Racah defining terminating series has a numerator factor $\beta q^{n-N}$ and denominator factor $\beta\delta q$. For fixed $n,N$ and fixed remaining admissible parameters, $\beta\to0$ sends both factors to $1$ termwise in the finite series, giving the dual $q$-Hahn defining ${}_3\phi_2$. Since the endpoint is also defined directly by setting $\beta=0$, this is an exact specialization and not an independent limiting arrow under the present strict scope gate. This check applies to the endpoint in the source heading at line 5620 only.

\paragraph{Coverage arithmetic and limitation.}
Of these fourteen headings, eight point to previously recorded distinct-family PASS proofs and six to exact-specialization endpoints already excluded or excluded here (lines 5230, 5411, 5620, 6654, 6811, 6878). This tranche finds no new candidate needing a proof, but it does \emph{not} establish that every other singular parameter path between these families is absent. Nor does it independently validate the external KLS references for all eight prior proofs. In particular, a fixed-parameter specialization may coexist with a different genuine KLS limit, so these dispositions are arrow-specific, not global bans on a family pair.

\noindent\textbf{V92 verdict: FOURTEEN MORE DISTINCT-FAMILY HEADINGS RECONCILED AT THE DOCUMENTED ENDPOINTS; NO NEW PASS.} The recorded eligible-proof tally remains \textbf{43}. Global KLS-arrow exhaustiveness and independent recertification remain open.


\section{V93: KLS-arrow reconciliation, tranche III---continuous-$q$ distinct families}
\noindent\textbf{Scope and method.} This tranche inventories the distinct-family headings in V31 source lines 3153--4469. Same-family connections are omitted from the working set. Classification is deliberately \emph{arrow-specific}: an exact endpoint does not preclude a different singular path. ``Previously proved'' refers to the cumulative ledger, not an independent re-proof or independent verification of the printed KLS citation.

\begin{longtable}{@{}p{0.12\linewidth}p{0.38\linewidth}p{0.42\linewidth}@{}}
\toprule V31 line & Distinct-family heading & Reconciliation and remaining obligation \\
\midrule\endfirsthead
\toprule V31 line & Distinct-family heading & Reconciliation and remaining obligation \\
\midrule\endhead
\bottomrule\endfoot
\texttt{3153} & Continuous dual $q$-Hahn / Askey--Wilson & \textbf{PREVIOUSLY PROVED:} Askey--Wilson $\to$ continuous dual $q$-Hahn in the cumulative ledger; check KLS citation at final recertification. \\
\texttt{3324} & Continuous $q$-Hahn / Askey--Wilson & \textbf{EXACT REALIZATION:} V84 identifies the continuous $q$-Hahn polynomial as an Askey--Wilson parameter realization. No additional arrow proved. \\
\texttt{3441} & Al-Salam--Chihara / continuous dual $q$-Hahn & \textbf{PREVIOUSLY PROVED:} Continuous dual $q$-Hahn $\to$ Al-Salam--Chihara. \\
\texttt{3582} & $q$-Meixner--Pollaczek / continuous $q$-Hahn & \textbf{EXACT ENDPOINT CHECKED:} V83 proves the $b=0$ specialization; not a new eligible arrow under the stated scope. \\
\texttt{3701} & Continuous $q$-Jacobi / Askey--Wilson & \textbf{EXACT REALIZATION:} Continuous $q$-Jacobi is an Askey--Wilson parameter specialization; no new limit established by the matched parameters. \\
\texttt{3829} & Continuous big $q$-Hermite / Al-Salam--Chihara & \textbf{PREVIOUSLY PROVED, RECHECKED:} Al-Salam--Chihara $\to$ continuous big $q$-Hermite; V81 notes the exact boundary endpoint. \\
\texttt{3991} & Continuous $q$-Laguerre / Al-Salam--Chihara & \textbf{EXACT NORMALIZED REALIZATION:} V82 verifies the parameter-pair specialization. \\
\texttt{4060} & Continuous $q$-Laguerre / $q$-Meixner--Pollaczek & \textbf{EXACT ENDPOINT CHECKED:} V69 excludes the zero-parameter Al-Salam--Chihara realization. Other KLS paths remain unverified. \\
\texttt{4145} & Continuous $q$-Laguerre / continuous $q$-Jacobi & \textbf{EXACT ENDPOINT CHECKED:} V70 identifies the normalized $B=0$ specialization. Other singular paths remain unverified. \\
\texttt{4315} & Continuous $q$-Hermite / continuous big $q$-Hermite & \textbf{PREVIOUSLY PROVED:} Continuous big $q$-Hermite $\to$ continuous $q$-Hermite. \\
\texttt{4366} & Continuous $q$-Hermite / continuous $q$-Laguerre & \textbf{PREVIOUSLY PROVED:} V68 normalized continuous $q$-Laguerre $\to$ continuous $q$-Hermite. \\
\texttt{4469} & Continuous $q$-Hermite / continuous $q$-Jacobi & \textbf{PREVIOUSLY PROVED:} V67 normalized continuous $q$-Jacobi $\to$ continuous $q$-Hermite. \\
\end{longtable}

\paragraph{Arithmetic and audit caution.}
These twelve distinct-family headings comprise six previously proved arrows, five exact-realization or exact-endpoint dispositions, and one additional exact endpoint with unverified alternative paths: more precisely, six entries are marked \textbf{PREVIOUSLY PROVED} (3153, 3441, 3829, 4315, 4366, 4469), while six entries are marked \textbf{EXACT REALIZATION/ENDPOINT} (3324, 3582, 3701, 3991, 4060, 4145). The latter six are not certified to lack \emph{all} possible singular KLS paths. This tranche does not independently certify external KLS numbering or revalidate every earlier proof.

\noindent\textbf{V93 verdict: TWELVE CONTINUOUS-$q$ DISTINCT-FAMILY HEADINGS RECONCILED AT THE DOCUMENTED ARROWS; NO NEW PASS.} The recorded eligible-proof count remains \textbf{43}. A full KLS-arrow inventory and independent recertification of all forty-three proofs remain prerequisites to declaring completion.


\section{V94: KLS-arrow reconciliation, tranche IV---classical distinct-family connections}
\noindent\textbf{Scope.} This tranche checks all classical distinct-family headings before the Askey--Wilson section of V31 (source lines 176--2874). We exclude every self-family heading \emph{before} review. The source heading at line 509 is an auxiliary ``parameter cancellation'' subsection, not an independently named polynomial-family pair, and is treated separately. A previously recorded PASS is \emph{not} a fresh verification of the published KLS reference.

\begingroup\footnotesize
\begin{longtable}{@{}p{0.10\linewidth}p{0.31\linewidth}p{0.51\linewidth}@{}}
\toprule V31 line & Distinct-family heading & Reconciliation of the documented endpoint \\
\midrule\endfirsthead
\toprule V31 line & Distinct-family heading & Reconciliation of the documented endpoint \\
\midrule\endhead
\bottomrule\endfoot
\texttt{176} & Continuous Dual Hahn / Wilson & \textbf{RECORDED PASS:} Wilson to continuous dual Hahn; independent KLS citation/proof recertification pending. \\
\texttt{267} & Continuous Hahn / Wilson & \textbf{EXACT TRANSFORMATION:} V79: symmetric continuous Hahn/Wilson quadratic identity; no new PASS on this endpoint. \\
\texttt{426} & Meixner--Pollaczek / Continuous Hahn & \textbf{RECORDED PASS:} continuous Hahn to Meixner--Pollaczek (restricted symmetric path); independent KLS citation/proof recertification pending. \\
\texttt{609} & Meixner--Pollaczek / Continuous Dual Hahn & \textbf{EXACT TRANSFORMATION:} V66: continuous dual Hahn/Meixner--Pollaczek quadratic identities; no new PASS on this endpoint. \\
\texttt{822} & Jacobi / Wilson & \textbf{RECORDED PASS:} Wilson to Jacobi (specified scaling); independent KLS citation/proof recertification pending. \\
\texttt{975} & Jacobi / Continuous Hahn & \textbf{UNRESOLVED/NOT CERTIFIED:} V80: generic basis connection; no singular path certified. \\
\texttt{1211} & Generalised Laguerre / Meixner--Pollaczek & \textbf{UNRESOLVED/NOT CERTIFIED:} V78: classical MP to Laguerre limit verified independently, but fixed-angle reverse identity does not prove it. \\
\texttt{1296} & Generalised Laguerre / Jacobi & \textbf{RECORDED PASS:} Jacobi to Laguerre; independent KLS citation/proof recertification pending. \\
\texttt{1454} & Hermite / Meixner--Pollaczek & \textbf{RECORDED PASS:} Meixner--Pollaczek to Hermite; independent KLS citation/proof recertification pending. \\
\texttt{1527} & Hermite / Jacobi & \textbf{RECORDED PASS:} Jacobi to Hermite (V71); independent KLS citation/proof recertification pending. \\
\texttt{1582} & Hermite / Generalised Laguerre & \textbf{EXACT TRANSFORMATION:} V73: Hermite/Laguerre quadratic identities; no new PASS on this endpoint. \\
\texttt{1731} & Hahn / Racah & \textbf{RECORDED PASS:} Racah to Hahn; independent KLS citation/proof recertification pending. \\
\texttt{1837} & Dual Hahn / Racah & \textbf{RECORDED PASS:} Racah to dual Hahn; independent KLS citation/proof recertification pending. \\
\texttt{1978} & Meixner / Hahn & \textbf{RECORDED PASS:} Hahn to Meixner; independent KLS citation/proof recertification pending. \\
\texttt{2027} & Meixner / Dual Hahn & \textbf{UNRESOLVED/NOT CERTIFIED:} V74: quadratic-lattice mismatch; no KLS path certified. \\
\texttt{2154} & Krawtchouk / Hahn & \textbf{RECORDED PASS:} Hahn to Krawtchouk; independent KLS citation/proof recertification pending. \\
\texttt{2236} & Krawtchouk / Dual Hahn & \textbf{EXACT TRANSFORMATION:} V77: Krawtchouk/dual Hahn quadratic identities; no new PASS on this endpoint. \\
\texttt{2362} & Charlier / Meixner & \textbf{RECORDED PASS:} Meixner to Charlier; independent KLS citation/proof recertification pending. \\
\texttt{2397} & Charlier / Krawtchouk & \textbf{RECORDED PASS:} Krawtchouk to Charlier; independent KLS citation/proof recertification pending. \\
\texttt{2459} & Racah / Wilson & \textbf{EXACT TRANSFORMATION:} V75: Racah/Wilson exact lattice transformation; no new PASS on this endpoint. \\
\texttt{2505} & Pseudo-Jacobi / Continuous Hahn & \textbf{UNRESOLVED/NOT CERTIFIED:} V76: generic polynomial-basis connection; no singular path certified. \\
\texttt{2579} & Jacobi / Hahn & \textbf{RECORDED PASS:} Hahn to Jacobi (V64); independent KLS citation/proof recertification pending. \\
\texttt{2659} & Generalised Bessel / Jacobi & \textbf{UNRESOLVED/NOT CERTIFIED:} V65: generalized Bessel/Jacobi KLS path not established. \\
\texttt{2745} & Generalised Laguerre / Meixner & \textbf{RECORDED PASS:} Meixner to Laguerre; independent KLS citation/proof recertification pending. \\
\texttt{2817} & Hermite / Krawtchouk & \textbf{RECORDED PASS:} Krawtchouk to Hermite; independent KLS citation/proof recertification pending. \\
\texttt{2874} & Hermite / Charlier & \textbf{RECORDED PASS:} Charlier to Hermite; independent KLS citation/proof recertification pending. \\
\end{longtable}
\endgroup

\paragraph{Arithmetic and source-inventory correction.}
The classical block contains \textbf{26} named distinct-family headings: \textbf{16} with a previously recorded PASS for at least one arrow, \textbf{5} exact quadratic or lattice transformations at the documented endpoint, and \textbf{5} whose proposed KLS/reverse-connection limiting paths remain unverified or unavailable. The auxiliary line-509 heading is \emph{not} a 27th family pair. Consequently the V90 census of 117 \texttt{\\subsubsection} headings comprises 44 self-family headings, \textbf{72} named distinct-family headings, and \textbf{one auxiliary heading}; the earlier provisional ``73 distinct-family'' label overcounted by one. This corrects the \emph{heading census only}, not the 43-proof ledger.

\paragraph{What is and is not established.}
The dispositions above reconcile the \emph{documented} V31 endpoints against earlier audit sections. A pair may have additional KLS arrows; conversely an exact transformation or an arbitrary finite connection does not create an eligible arrow. In particular, the genuine Meixner--Pollaczek to Laguerre classical limit from V78 remains \emph{uncounted} because the V31 fixed-angle reverse connection does not supply the variable-angle limiting proof. We have not independently checked all original KLS equation numbers or revalidated the mathematical details of all sixteen earlier proofs in this tranche.

\noindent\textbf{V94 verdict: CLASSICAL HEADING-LEVEL RECONCILIATION COMPLETED; NO NEW PASS.} The recorded eligible-proof tally stays \textbf{43}. Global arrow-by-arrow KLS completeness and independent recertification remain open.



\section{V95: KLS-arrow reconciliation, tranche V---remaining distinct-family headings}
\noindent\textbf{Scope.} This tranche addresses the eleven named distinct-family headings not assigned to V91--V94: nine discrete-$q$ headings at V31 lines 7149--8039 and two terminal Askey--Wilson headings at lines 9287 and 9479. Same-family headings are excluded \emph{a priori}. The classifications concern the particular documented or already proved arrows, not all conceivable singular parameter paths.

\begin{longtable}{@{}p{0.11\linewidth}p{0.35\linewidth}p{0.46\linewidth}@{}}
\toprule V31 line & Distinct-family heading & Ledger reconciliation and qualification \\
\midrule\endfirsthead
\toprule V31 line & Distinct-family heading & Ledger reconciliation and qualification \\
\midrule\endhead
\bottomrule\endfoot
\texttt{7149} & Wall / big $q$-Laguerre & \textbf{RECORDED PASS:} V58, big $q$-Laguerre $\to$ Wall, with singular $B\to\infty$ and argument rescaling. \\
\texttt{7243} & Wall / little $q$-Jacobi & \textbf{RECORDED PASS:} little $q$-Jacobi $\to$ Wall, proved in the earlier cumulative audit. \\
\texttt{7337} & Wall / affine $q$-Krawtchouk & \textbf{RECORDED PASS:} affine $q$-Krawtchouk $\to$ Wall, previously proved. \\
\texttt{7584} & $q$-Laguerre / little $q$-Jacobi & \textbf{RECORDED PASS:} little $q$-Jacobi $\to$ $q$-Laguerre, previously proved. \\
\texttt{7648} & $q$-Laguerre / $q$-Meixner & \textbf{RECORDED PASS:} $q$-Meixner $\to$ $q$-Laguerre, previously proved. \\
\texttt{7787} & Alternative $q$-Charlier ($q$-Bessel) / little $q$-Jacobi & \textbf{RECORDED PASS, DUPLICATE CONTROL:} the singular little $q$-Jacobi $\to$ alternative $q$-Charlier path was already proved before V87; V88 reconciled the repeated proof. \\
\texttt{7850} & Alternative $q$-Charlier ($q$-Bessel) / ordinary $q$-Krawtchouk & \textbf{RECORDED PASS:} ordinary $q$-Krawtchouk $\to$ alternative $q$-Charlier, previously proved in V37. \\
\texttt{7962} & Ordinary $q$-Charlier / $q$-Meixner & \textbf{EXACT ENDPOINT, NOT A NEW PASS:} V43 corrected the earlier erroneous count: the documented zero-parameter reduction is an exact specialization. No other arrow is certified by this classification. \\
\texttt{8039} & Ordinary $q$-Charlier / ordinary $q$-Krawtchouk & \textbf{RECORDED PASS:} ordinary $q$-Krawtchouk $\to$ ordinary $q$-Charlier, previously proved in V35. \\
\texttt{9287} & Askey--Wilson / $q$-Racah & \textbf{EXACT ENDPOINT, NOT A NEW PASS:} V31 explicitly gives the exact common-quadratic-lattice specialization with $a^2=\gamma\delta q$, $ab=\alpha q$, $ac=\beta\delta q$, $ad=\gamma q$. It is not a singular degeneration. \\
\texttt{9479} & Askey--Wilson / big $q$-Jacobi & \textbf{RECORDED PASS:} Askey--Wilson $\to$ big $q$-Jacobi was proved in V39. \\
\end{longtable}

\paragraph{An exact identity that blocks a false additional count.}
At the fully matched parameters in the V31 Askey--Wilson/$q$-Racah heading, the manuscript records
\[
p_n\!\left(\frac{Y}{2a};a,b,c,d\mid q\right)
=a^{-n}(\alpha q,\beta\delta q,\gamma q;q)_n
R_n(Y;\alpha,\beta,\gamma,\delta\mid q).
\]
This is an equality of terminating polynomial representations for fixed admissible parameters. Consequently its diagonal reverse-connection reduction cannot be counted as an additional KLS limit. This statement does \emph{not} exclude any different, genuinely singular arrow involving these families.

\paragraph{Heading-level closure and limitations.}
The corrected V31 census is $117=44+72+1$: 44 self-family headings (out of scope), 72 named distinct-family headings, and one auxiliary heading. The five tranches cover $9+14+12+26+11=\mathbf{72}$ named distinct-family headings \emph{at the heading level}, with no overlap in source line numbers. Thus \textbf{heading-level reconciliation is complete}. This is \emph{not} yet a KLS-arrow-level completeness certificate: a heading may admit more than one asymptotic path, and some proposed paths remain unresolved (notably the five classical questions in V94 and the three late-discrete-$q$ eligibility checks in V91). The external KLS citations and the mathematical independence of all 43 recorded PASS proofs have not yet been independently recertified. The classification of exact endpoints is not proof of the nonexistence of other singular limits.

\noindent\textbf{V95 verdict: ALL 72 DISTINCT-FAMILY HEADINGS RECONCILED; NO NEW PASS; GLOBAL KLS-ARROW COMPLETENESS AND THE 43-PROOF RECERTIFICATION REMAIN OPEN.} Recorded PASS tally: \textbf{43}.


\section{V96: Internal recertification gate---the 43-entry PASS ledger}
\noindent\textbf{Scope.} Before asserting global KLS completeness, we independently check the \emph{internal accounting and manuscript cross-references} of the existing PASS ledger. This is a reproducible mechanical check of the V95 source, not a substitute for mathematical rederivation of each asymptotic proof or consultation of the original KLS tables. Self-family headings are outside scope.

\paragraph{Method.} Extract the items between the \texttt{Progress} section's first \texttt{itemize} environment and its matching end, select exactly those carrying \texttt{\\textbackslash textbf\{PASS\}}, and extract the literal V31 \texttt{eq:} labels cited within those items. Search the unchanged V31 source for the exact corresponding \texttt{\\textbackslash label\{...\}} strings. This checks \emph{existence}, not that the cited equation proves the advertised limit. Each row below is an existing recorded proof, not a newly certified KLS arrow.

\begingroup\small
\begin{longtable}{@{}p{0.05\linewidth}p{0.52\linewidth}p{0.35\linewidth}@{}}
\toprule No. & Recorded distinct-family limiting arrow & Explicit V31 equation label in ledger \\
\midrule\endfirsthead
\toprule No. & Recorded limiting arrow & Explicit V31 label \\
\midrule\endhead
\bottomrule\endfoot
1 & Symmetric Jacobi $\ensuremath{\to}$ Hermite & \texttt{eq:v14-jac-herm} \\
2 & Normalized continuous $q$-Laguerre $\ensuremath{\to}$ continuous $q$-Hermite & \texttt{eq:qlag-hermite-fwd} \\
3 & Continuous $q$-Jacobi $\ensuremath{\to}$ continuous $q$-Hermite & \texttt{eq:JHforward} \\
4 & Hahn $\ensuremath{\to}$ Jacobi & \texttt{eq:jacobi-hahn-reverse} \\
5 & $q$-Hahn $\ensuremath{\to}$ little $q$-Jacobi & \texttt{eq:ljqh-reverse} \\
6 & $q$-Hahn $\ensuremath{\to}$ $q$-Meixner & \texttt{eq:168-D} \\
7 & Big $q$-Jacobi $\ensuremath{\to}$ $q$-Meixner & \texttt{eq:169-reverse} \\
8 & Big $q$-Jacobi $\ensuremath{\to}$ little $q$-Jacobi & \texttt{eq:ljbj166-reverse} \\
9 & Big $q$-Laguerre $\ensuremath{\to}$ little $q$-Laguerre (Wall) & \texttt{eq:wall-big-reverse} \\
10 & $q$-Hahn $\ensuremath{\to}$ quantum $q$-Krawtchouk & \texttt{eq:finite-qtmH-reverse} \\
11 & $q$-Hahn $\ensuremath{\to}$ ordinary $q$-Krawtchouk & \texttt{eq:finite-KH-reverse} \\
12 & Dual $q$-Hahn $\ensuremath{\to}$ dual $q$-Krawtchouk & \texttt{eq:finite-dKdH-reverse} \\
13 & Askey--Wilson $\ensuremath{\to}$ big $q$-Jacobi & \texttt{eq:aw-bigqj-reverse} \\
14 & Standard $q$-Krawtchouk $\ensuremath{\to}$ alternative $q$-Charlier ($q$-Bessel) & \texttt{eq:qBqK-reverse} \\
15 & Standard $q$-Krawtchouk $\ensuremath{\to}$ ordinary $q$-Charlier & \texttt{eq:qCqK-reverse} \\
16 & Little $q$-Jacobi $\ensuremath{\to}$ $q$-Bessel (alternative $q$-Charlier) & \texttt{eq:qBlqJ-reverse} \\
17 & $q$-Meixner $\ensuremath{\to}$ $q$-Laguerre & \texttt{eq:qLqM-D} \\
18 & Little $q$-Jacobi $\ensuremath{\to}$ $q$-Laguerre & \texttt{eq:qLlJ-D} \\
19 & Affine $q$-Krawtchouk $\ensuremath{\to}$ little $q$-Laguerre (Wall) & \texttt{eq:wall-affine-reverse} \\
20 & Little $q$-Jacobi $\ensuremath{\to}$ little $q$-Laguerre (Wall) & \texttt{eq:wall-lj-reverse} \\
21 & Big $q$-Jacobi $\ensuremath{\to}$ big $q$-Laguerre & \texttt{eq:blbj-reverse} \\
22 & Continuous $q$-Hahn $\ensuremath{\to}$ $q$-Meixner--Pollaczek & \texttt{eq:cqh-qmp23} \\
23 & Al-Salam--Chihara $\ensuremath{\to}$ continuous big $q$-Hermite & \texttt{eq:ASCBH} \\
24 & Continuous dual $q$-Hahn $\ensuremath{\to}$ Al-Salam--Chihara & \texttt{eq:CDASC} \\
25 & Continuous big $q$-Hermite $\ensuremath{\to}$ continuous $q$-Hermite & \texttt{eq:BHH} \\
26 & $q$-Laguerre $\ensuremath{\to}$ Stieltjes--Wigert & \texttt{eq:SWqL-rev} \\
27 & Quantum $q$-Krawtchouk $\ensuremath{\to}$ Al-Salam--Carlitz II & \texttt{eq:ASCII-qtm-reverse} \\
28 & $q$-Meixner $\ensuremath{\to}$ Al-Salam--Carlitz II & \texttt{eq:ASCII-qM-reverse} \\
29 & Askey--Wilson $\ensuremath{\to}$ continuous dual $q$-Hahn & \texttt{eq:AWCD} \\
30 & Racah $\ensuremath{\to}$ Hahn & \texttt{eq:racah-to-hahn-grid} \\
31 & Racah $\ensuremath{\to}$ dual Hahn & \texttt{eq:racah-dh-reverse} \\
32 & Krawtchouk $\ensuremath{\to}$ Hermite & \texttt{eq:v16-k-h} \\
33 & Hahn $\ensuremath{\to}$ Meixner & \texttt{eq:hahn-meixner-reverse} \\
34 & Hahn $\ensuremath{\to}$ Krawtchouk & \texttt{eq:hahn-kraw-reverse} \\
35 & Krawtchouk $\ensuremath{\to}$ Charlier & \texttt{eq:kraw-char-reverse} \\
36 & Meixner $\ensuremath{\to}$ Charlier & \texttt{eq:meix-char-reverse} \\
37 & Charlier $\ensuremath{\to}$ Hermite & Not specified in ledger item \\
38 & Wilson $\ensuremath{\to}$ continuous dual Hahn & Not specified in ledger item \\
39 & Meixner--Pollaczek $\ensuremath{\to}$ Hermite & Not specified in ledger item \\
40 & Continuous Hahn $\ensuremath{\to}$ Meixner--Pollaczek & Not specified in ledger item \\
41 & Meixner $\ensuremath{\to}$ generalized Laguerre & Not specified in ledger item \\
42 & Jacobi $\ensuremath{\to}$ generalized Laguerre & Not specified in ledger item \\
43 & Wilson $\ensuremath{\to}$ Jacobi & Not specified in ledger item \\
\end{longtable}\endgroup

\paragraph{Reproducible results.} The Progress ledger contains \textbf{43 PASS items}. Of these, \textbf{36} explicitly give an \texttt{eq:} label and \textbf{7} do not provide such a label in their ledger item. All \textbf{36} cited labels are distinct and are present verbatim as \texttt{\\textbackslash label\{...\}} in V31. The 43-item count therefore agrees with the reported cumulative tally \emph{as an inventory count}; no missing manuscript label was detected among the explicitly cited labels. The seven entries without a ledger-item equation label require manual cross-reference to the corresponding proof section before a source-link certificate can be issued. This count does not establish that the 43 limiting arrows are mutually independent or that each is correctly associated with a numbered KLS arrow.

\paragraph{Outstanding external and mathematical checks.} The eight unresolved eligibility questions carried forward from V91 and V94 are: (i) big $q$-Laguerre / Al-Salam--Carlitz I; (ii) Stieltjes--Wigert / alternative $q$-Charlier; (iii) Stieltjes--Wigert / ordinary $q$-Charlier; (iv) Jacobi / continuous Hahn; (v) Meixner--Pollaczek / generalized Laguerre (genuine classical limit but fixed-angle V31 reverse connection does not yield it); (vi) Meixner / dual Hahn; (vii) pseudo-Jacobi / continuous Hahn; and (viii) generalized Bessel / Jacobi. These are \emph{unresolved eligibility checks}, not eight known missing KLS arrows. In addition, the original KLS equation numbers and normalization of the 43 claimed proofs must be checked individually against primary sources, and their reverse-connection asymptotics revalidated. No such independent primary-source comparison is claimed in V96.

\noindent\textbf{V96 verdict: INTERNAL PASS-LEDGER COUNT AND 36 EXPLICIT SOURCE-LABEL EXISTENCE CHECKS VERIFIED; FINAL MATHEMATICAL/KLS CERTIFICATION STILL OPEN.} No new PASS; recorded total \textbf{43}. The original V31 is unchanged.


\section{V97: Completion of the seven missing V31 source cross-references}
\noindent\textbf{Purpose.} V96 found seven PASS ledger items without an explicit source-equation label \emph{in the Progress item}. This section traces each one to the original proof section in this cumulative audit and verifies the corresponding literal equation label in the unchanged V31 source. This is a source-traceability check, not independent recertification of the KLS classification or of every asymptotic calculation.

\begin{center}\small
\begin{longtable}{@{}p{0.07\linewidth}p{0.42\linewidth}p{0.36\linewidth}@{}}
\toprule V96 row & Recorded arrow & V31 reverse-connection equation \\
\midrule\endfirsthead
\toprule Row & Recorded arrow & V31 equation \\
\midrule\endhead
\bottomrule\endfoot
37 & Charlier $\to$ Hermite & \texttt{eq:v16-c-h} (V31 line 2891) \\
38 & Wilson $\to$ continuous dual Hahn & \texttt{eq:consolidated-w-cdh} (V31 line 188) \\
39 & Meixner--Pollaczek $\to$ Hermite & \texttt{eq:cat-mp-to-hermite-bell} (V31 line 1483) \\
40 & Continuous Hahn $\to$ Meixner--Pollaczek & \texttt{eq:consolidated-ch-mp} (V31 line 443) \\
41 & Meixner $\to$ generalized Laguerre & \texttt{eq:meixner-to-laguerre-v61} (V31 line 2769) \\
42 & Jacobi $\to$ generalized Laguerre & \texttt{eq:v14-jac-lag} (V31 line 1320) \\
43 & Wilson $\to$ Jacobi & \texttt{eq:wilson-to-jac} (V31 line 876) \\
\end{longtable}
\end{center}
\paragraph{Verification method and result.} Each of the seven labels is explicitly named in its corresponding early audit proof section (the Wilson-to-Jacobi section refers to the reverse formula by its V31 label in the original manuscript), and each literal \texttt{\textbackslash label\{...\}} is present in V31. Combined with V96's 36 checks, all \textbf{43 recorded PASS entries now have a traceable V31 equation label}. This proves source-reference \emph{coverage}, not validity of all 43 limiting arguments, uniqueness of KLS arrows, or global completeness.

\paragraph{Certification caveat.} Two previously recorded PASS entries explicitly note that their exact KLS equation-number correspondence remains unconfirmed: the symmetric continuous Hahn $\to$ Meixner--Pollaczek scaling and the Wilson $\to$ Jacobi scaling. The other 41 entries also require an independent source-by-source KLS normalization and proof review before a final certificate can be issued. The eight unresolved eligibility questions from V96 remain open. The PASS ledger is \emph{unchanged} at \textbf{43}; no self-family connections are included.

\noindent\textbf{V97 verdict: ALL 43 LEDGER ITEMS SOURCE-LINKED TO V31; MATHEMATICAL AND KLS-ARROW CERTIFICATION STILL PENDING.} No new PASS. The original V31 is unchanged.


\section{V98: Mathematical recertification gate---zero-parameter endpoint consistency}
\noindent\textbf{Purpose.} The earlier heading inventory and source-reference checks do not establish that every recorded PASS satisfies the \emph{current} eligibility rule excluding exact specializations. This section checks a concrete four-entry subset of the 43-entry ledger against that rule. No self-family connection is considered.

\paragraph{Exact endpoints and finite reverse-connection limits.}
For fixed $0<q<1$ and fixed $n$, the V31 reverse connections give
\begin{align}
 H_n(x;a\mid q)&=\sum_{k=0}^n {n\brack k}_q(-a)^{n-k}q^{\binom{n-k}{2}}H_k(x\mid q),\\
 S_n(x;a,b,c\mid q)&=\sum_{k=0}^n {n\brack k}_q(-c)^{n-k}q^{\binom{n-k}{2}}(abq^k;q)_{n-k}Q_k(x;a,b\mid q),\\
 Q_n(x;a,b\mid q)&=\sum_{k=0}^n {n\brack k}_q(-b)^{n-k}q^{\binom{n-k}{2}}H_k(x;a\mid q).
\end{align}
At $a=0$, $c=0$, and $b=0$, respectively, the $k=n$ term is one and every lower term is \emph{exactly zero}. Consequently,
\begin{equation}\label{eq:v98-exact-endpoints}
 H_n(x;0\mid q)=H_n(x\mid q),\qquad
 S_n(x;a,b,0\mid q)=Q_n(x;a,b\mid q),\qquad
 Q_n(x;a,0\mid q)=H_n(x;a\mid q).
\end{equation}
These are algebraic parameter evaluations in the same normalizations as the corresponding earlier PASS proofs, not merely asymptotic statements. They occupy ledger rows 25, 24, and 23, respectively.

A fourth recorded PASS, ledger row 22, uses V31 \texttt{eq:cqh-qmp23} and the defining normalization
\[
 \mathscr P_n(x;a,\phi\mid q)=\frac{p_n(x;ae^{i\phi},0,0,ae^{-i\phi}\mid q)}{(q;q)_n},\quad
 \mathscr Q_n(x;a,b,0,a,\phi\mid q)=p_n(x;ae^{i\phi},be^{i\phi},0,ae^{-i\phi}\mid q).
\]
Thus its endpoint likewise holds by direct evaluation:
\begin{equation}\label{eq:v98-qmp-exact}
 \mathscr Q_n(x;a,0,0,a,\phi\mid q)=(q;q)_n\mathscr P_n(x;a,\phi\mid q).
\end{equation}
The $b\to0$ reverse-connection limit is valid, but no nontrivial singular scaling is needed to identify the endpoint.

\paragraph{Eligibility implication.}
Under the \emph{strict interpretation} agreed in the conversation---exclude an arrow if its documented path is only an exact zero-parameter specialization---these four recorded PASS items are \textbf{not certified as eligible}. This is a genuine tension with earlier audit sections that counted parameter-deletion paths as PASS, even though V81 and V83 explicitly verified exact endpoints. The mathematical identities and their limits remain correct; it is the classification, not the algebra, that requires correction.

\paragraph{Status and accounting.}
The historic ledger still contains \textbf{43 PASS labels}; this is \emph{not} a certified count under the strict eligibility rule. Four entries (rows 22--25) are now \textbf{FLAGGED FOR RECLASSIFICATION}. If each is removed and no other entry changes, the resulting count would be $43-4=39$; \textbf{39 is a conditional arithmetic subtotal, not a newly certified final total}. Other recorded PASS items may have the same issue and require screening before publishing a corrected total. This stage does not establish a new KLS arrow or independently check a numbered KLS reference. The eight previously identified eligibility questions remain unresolved.

\noindent\textbf{V98 verdict: FOUR HISTORIC PASS ITEMS CONFLICT WITH THE STRICT EXACT-SPECIALIZATION EXCLUSION; FINAL COUNT AND COMPLETENESS CERTIFICATE WITHHELD.} V31 remains unchanged.


\section{V99: Eligibility-rule correction---exact specializations are admissible KLS arrows}
\noindent\textbf{Authoritative scope clarification.} The user has explicitly clarified that an \emph{exact specialization is allowed in the counted limiting-connection ledger if it passes verification}. This supersedes the strict ``exclude all exact specializations'' interpretation adopted in V98 and in some earlier exclusion decisions. V98's mathematical identities remain valid, but its proposed disqualification of four entries on the sole ground of being exact parameter evaluations is \textbf{withdrawn}. The conditional subtotal 39 has no standing as a revised certified count.

\paragraph{Operational eligibility test.} For a candidate between \emph{distinct} families $F_n$ and $G_n$, a PASS may be awarded when (i) an identified Koekoek--Swarttouw limiting arrow or its accepted zero-parameter endpoint connects the two families, (ii) the parameter map, variable, and degree-dependent normalization are explicit, (iii) V31's reverse connection gives the identity or a justified coefficientwise limit, and (iv) the arrow has not already been counted. An exact zero-parameter identity is not automatically a PASS: it must satisfy the same source, normalization, and uniqueness checks. Self-family parameter changes remain outside scope. Quadratic transformations without a recognized KLS limiting arrow are not automatically promoted.

\paragraph{Restoration of the four V98 entries.} For fixed degree $n$ the reverse connections and endpoint identities displayed in V98 verify the algebraic zero-parameter transitions:
\begin{align*}
H_n(x;a\mid q)\big|_{a=0}&=H_n(x\mid q),\\
S_n(x;a,b,c\mid q)\big|_{c=0}&=Q_n(x;a,b\mid q),\\
Q_n(x;a,b\mid q)\big|_{b=0}&=H_n(x;a\mid q),\\
\mathscr Q_n(x;a,b,0,a,\phi\mid q)\big|_{b=0}&=(q;q)_n\mathscr P_n(x;a,\phi\mid q).
\end{align*}
These are the previously recorded PASS ledger rows 25, 24, 23, and 22, respectively. Their identities are mathematically verified by V31's finite reverse connections. \textbf{They remain in the historical ledger and are not subtracted merely because the endpoint is exact.} Independent KLS reference checks remain part of final certification, as they do for the other entries.

\paragraph{Previously excluded specializations now reopened for review.} The following distinct-family entries were previously excluded \emph{at least partly} because their endpoints are exact specializations; under the revised rule that reason alone is insufficient. Their KLS-arrow status and uniqueness must be checked before changing the tally:
\begin{itemize}
\item $q$-Meixner / ordinary $q$-Charlier (V43, V57, V95; V31 \texttt{eq:qC-qM-reverse});
\item $q$-Racah / ordinary $q$-Krawtchouk (V54; V31 \texttt{eq:finite-KR-reverse});
\item Dual $q$-Hahn / affine $q$-Krawtchouk (V53); $q$-Hahn / affine $q$-Krawtchouk (V52);
\item $q$-Racah / dual $q$-Hahn (V92); $q$-Racah / $q$-Hahn (V72);
\item Continuous $q$-Jacobi / continuous $q$-Laguerre (V70); continuous $q$-Laguerre / Al-Salam--Chihara (V82); continuous $q$-Laguerre / $q$-Meixner--Pollaczek realization (V69);
\item Discrete $q$-Hermite I / Al-Salam--Carlitz I and discrete $q$-Hermite II / Al-Salam--Carlitz II (V91), subject also to their exact-transformation status.
\end{itemize}
This list is a \emph{review queue}, not a list of new PASS results. Some items may be duplicate KLS arrows, unrecognized parameter realizations, or transformations rather than independently enumerated KLS limits. In particular, the previous mathematical proofs of endpoint identities cannot replace primary-source verification of the KLS arrow.

\paragraph{Count and certification.} The historical ledger contains \textbf{43 distinct recorded PASS entries}; the four entries questioned in V98 are \textbf{restored without increment}, because they were already counted. \textbf{New PASS entries in V99: zero.} The final independently certified total remains open until KLS-arrow matching and the reopened-exclusion queue are resolved. The eight preexisting unresolved eligibility questions also remain open.

\noindent\textbf{V99 verdict: V98 ELIGIBILITY DISQUALIFICATION REVERSED; RECORDED COUNT RESTORED TO 43; FINAL KLS/MATHEMATICAL CERTIFICATION STILL PENDING.} No self-family connection is admitted; V31 is unchanged.


\section{V100: Reopened exact specialization---$q$-Meixner to ordinary $q$-Charlier}
\noindent\textbf{Scope and source.} The user's revised rule admits exact specializations when they are mathematically verified, correspond to a recognized KLS limiting arrow, and are not duplicate PASS entries. We revisit the distinct-family connection in V31, \texttt{eq:qC-qM-reverse}, formerly excluded in V43 solely because its endpoint was exact. V31 explicitly associates the endpoint with Koekoek--Swarttouw Eq.~(4.13.2). That equation-number attribution is reported from V31; independent inspection of the primary KLS document remains outstanding.

\paragraph{Explicit parameter path and normalization.} For fixed $0<q<1$, $a\ne0$, $n\ge0$, let $b\to0$ and choose $c=a$. In V31's conventions,
\[
 M_n(x;b,a;q)={}_2\phi_1\!\left(\begin{matrix}q^{-n},x\\bq\end{matrix};q,-q^{n+1}/a\right),\qquad
 C_n(x;a;q)={}_2\phi_1\!\left(\begin{matrix}q^{-n},x\\0\end{matrix};q,-q^{n+1}/a\right).
\]
No argument change or degree-dependent renormalization is needed. The terminating series proves the exact endpoint $M_n(x;0,a;q)=C_n(x;a;q)$ and uniform convergence on compact $x$-sets as $b\to0$.

\paragraph{Reverse-connection coefficient proof.} V31's reverse equation reads
\[
 M_n(x;b,c;q)=\sum_{k=0}^n D_{nk}(b,c;a)C_k(x;a;q),
\]
\[
 D_{nk}(b,c;a)={n\brack k}_q\frac{(a/c)^k}{(bq;q)_k}
 \sum_{r=0}^{n-k}{n-k\brack r}_q(-1)^rq^{\binom r2}
 \frac{(a/c)^r}{(bq^{k+1};q)_r}.
\]
All denominators remain nonzero for sufficiently small $b$. At $c=a$ and $b=0$, the finite $q$-binomial theorem gives
\[
 D_{nk}(0,a;a)={n\brack k}_q\sum_{r=0}^{n-k}{n-k\brack r}_q(-1)^rq^{\binom r2}
 ={n\brack k}_q(1;q)_{n-k}=\delta_{nk}.
\]
Here $(1;q)_m=0$ for $m>0$, while $(1;q)_0=1$. Thus the \emph{reverse connection itself} collapses to the identity, proving the limiting arrow with the correct normalization for every fixed $n$.

\paragraph{Uniqueness and KLS gate.} This arrow is absent from the 43-entry historical PASS ledger: V43 explicitly removed its former count. The V31 literature context identifies KLS Eq.~(4.13.2), providing a specific recognized-arrow attribution, although its primary-source equation numbering has not been independently checked. Consequently the endpoint is a \textbf{newly reinstated, mathematically verified recorded PASS}, with its external KLS citation still subject to the global primary-source certification gate. It is not a self-connection or a duplicate.

\noindent\textbf{V100 verdict: PASS (REINSTATED UNDER REVISED SPECIALIZATION RULE); HISTORICAL RECORDED COUNT $43\to44$; FINAL EXTERNAL KLS CERTIFICATION STILL OPEN.} The eight previously unresolved eligibility questions and other reopened exact-specialization cases remain to be reviewed. V31 is unchanged.


\section{V101: Reopened KLS endpoint---dual $q$-Hahn to affine $q$-Krawtchouk}
\noindent\textbf{Scope and source.} This revisits V53 under the user's clarified rule: exact specializations between distinct families are eligible when their KLS-arrow identification, reverse-connection verification, and uniqueness checks pass. The source is V31 \texttt{eq:finite-affdH-reverse}, with its literature context citing Koekoek--Swarttouw (1998), Eq.~(4.7.1). The equation-number attribution is from V31 and remains subject to independent primary-source inspection.

\paragraph{Parameter path and exact endpoint.} Fix $0<q<1$, $N\ge0$, $0\le n\le N$, and an admissible $p$ such that $(pq;q)_j\ne0$ for $0\le j\le N$. Put $\gamma=p$, $\delta\to0$, and $t=\gamma\delta q$. The dual $q$-Hahn lattice is $z_x(\delta)=q^{-x}+tq^x$, and the terminating definition is
\[
 R_n^{\mathrm{dH}}(z_x(\delta);p,\delta,N\mid q)
 ={}_3\phi_2\!\left(\begin{matrix}q^{-n},q^{-x},p\delta q^{x+1}\\pq,q^{-N}\end{matrix};q,q\right).
\]
For every fixed $x\in\{0,\ldots,N\}$, the finite series is continuous at $\delta=0$, yielding
\begin{equation}\label{eq:v101-dh-aff}
 \boxed{\lim_{\delta\to0}R_n^{\mathrm{dH}}(z_x(\delta);p,\delta,N\mid q)
 =R_n^{\mathrm{dH}}(q^{-x};p,0,N\mid q)
 =K_n^{\mathrm{aff}}(q^{-x};p,N\mid q).}
\end{equation}
The normalization factor is one. The identity at $\delta=0$ is exact, which is now \emph{permitted} rather than disqualifying.

\paragraph{Reverse-connection coefficient check.} For each $\delta$ near zero, V31's reverse finite-grid interpolation connection gives
\[
 R_n^{\mathrm{dH}}(z_x(\delta);p,\delta,N\mid q)
 =\sum_{k=0}^{N}D_{nk}(\delta;p)K_k^{\mathrm{aff}}(q^{-x};p,N\mid q).
\]
Let $E_{xk}=K_k^{\mathrm{aff}}(q^{-x};p,N\mid q)$, $0\le x,k\le N$. The polynomial $K_k^{\mathrm{aff}}(y;p,N\mid q)$ has exact degree $k$ in $y$ under the stated nondegeneracy assumptions, and the nodes $q^{-x}$ are distinct. Hence $E$ is invertible. If $r_n(\delta)$ denotes the row vector of the $N+1$ dual $q$-Hahn evaluations, the reverse coefficients satisfy $D_n(\delta)=r_n(\delta)E^{-T}$ when column conventions are used (equivalently, $r_n(\delta)=E D_n(\delta)^T$). By the finite-series limit, $r_n(\delta)\to E_{\cdot n}^T$, so
\[
 \boxed{\lim_{\delta\to0}D_{nk}(\delta;p)=\delta_{nk}\quad(0\le k\le N).}
\]
This proves convergence of the \emph{actual reverse-connection coefficients}, not just the endpoint hypergeometric representation. The fixed finite grid requires no exchange of infinite sums or asymptotic estimates.

\paragraph{Uniqueness and eligibility.} V53 established this exact identity but excluded it solely because it was a finite-parameter specialization; V99 explicitly reopened that decision. The dual $q$-Hahn $\to$ affine $q$-Krawtchouk arrow is not among the previously counted 44 entries, and V31 identifies a specific KLS endpoint reference. Thus this is a distinct-family, nonduplicate, reverse-connection-verified \textbf{recorded PASS} under the revised rule. The independent primary-source KLS equation check remains a condition of final global certification.

\noindent\textbf{V101 verdict: PASS (REINSTATED); RECORDED COUNT $44\to45$; FINAL PRIMARY-SOURCE KLS CERTIFICATION STILL PENDING.} No self-family connection was considered. V31 is unchanged.



\section{V102: Reopened $q$-Hahn to affine $q$-Krawtchouk endpoint---reverse coefficients and KLS gate}
\noindent\textbf{Scope.} This section revisits the distinct-family connection in V31, \texttt{eq:finite-affH-reverse}, previously screened in V52. Under the user's clarified rule, an exact zero-parameter specialization is not disqualified \emph{a priori}. We distinguish its rigorous reverse-connection proof from independent verification that the particular arrow is enumerated in the original KLS limiting scheme. No self-family connection is included.

\paragraph{Exact endpoint and limiting path.} Fix $0<q<1$, $N\ge0$, $0\le n\le N$, and an admissible $p$ with $(pq;q)_j\ne0$ for $0\le j\le N$. Set $A=p$ and let $B\to0$. The definitions in V31, \texttt{eq:finite-hahn-definition} and \texttt{eq:finite-affine-krawtchouk-definition}, give
\[
 Q_n(q^{-x};p,B,N\mid q)
 ={}_3\phi_2\!\left(\begin{matrix}q^{-n},pBq^{n+1},q^{-x}\\pq,q^{-N}\end{matrix};q,q\right),
 \qquad
 K_n^{\mathrm{aff}}(q^{-x};p,N\mid q)
 ={}_3\phi_2\!\left(\begin{matrix}q^{-n},0,q^{-x}\\pq,q^{-N}\end{matrix};q,q\right).
\]
Each series terminates at $j=n$; its denominator is nonsingular for $j\le n$. Thus, for every $x\in\{0,\ldots,N\}$,
\begin{equation}\label{eq:v102-qh-aff}
 \boxed{\lim_{B\to0}Q_n(q^{-x};p,B,N\mid q)
 =Q_n(q^{-x};p,0,N\mid q)
 =K_n^{\mathrm{aff}}(q^{-x};p,N\mid q).}
\end{equation}
The normalization is exactly one. In fact, since these are degree-$n$ polynomials in $y=q^{-x}$, the coefficientwise limit holds for arbitrary complex $y$, uniformly on compact sets.

\paragraph{Verification through the actual reverse connection.} The V31 reverse connection is
\[
 Q_n(y;A,B,N\mid q)=\sum_{k=0}^n D_{nk}(A,B;p)K_k^{\mathrm{aff}}(y;p,N\mid q),
\]
where $D_{nk}$ is the explicit terminating ${}_3\phi_2$ coefficient in \texttt{eq:finite-affH-reverse}. The affine basis is independent of $B$. Its coefficient of $y^k$ is nonzero: the $j=k$ term in the defining terminating series contains $(q^{-x};q)_k=\prod_{r=0}^{k-1}(1-yq^r)$ and a nonzero factor $(q^{-k};q)_k q^k/[(pq,q^{-N},q;q)_k]$. Consequently, $K_0^{\mathrm{aff}},\ldots,K_n^{\mathrm{aff}}$ form a triangular polynomial basis.

Let $T$ be their $(n+1)\times(n+1)$ monomial-coefficient matrix, with columns indexed by $k$, and let $v_n(B)$ be the monomial-coefficient vector of $Q_n(y;p,B,N\mid q)$. By the reverse connection, $T D_n(B)^T=v_n(B)$. The finite terminating definition proves $v_n(B)\to v_n(0)=T e_n$, while $T$ is fixed and invertible. Therefore
\[
 \boxed{\lim_{B\to0}D_{nk}(p,B;p)=\delta_{nk}\qquad(0\le k\le n).}
\]
This is a proof about the \emph{coefficients in V31's reverse connection}, not merely a comparison of the two defining series. It avoids potential removable singularities in individual hypergeometric quotient forms. The limit is valid for every fixed degree and requires neither lattice scaling nor additional normalization.

\paragraph{KLS eligibility and nonduplication.} V31 describes affine $q$-Krawtchouk as the exact $B=0$ specialization of $q$-Hahn and cites the standard $q$-Askey scheme, but the literature note explicitly says that \emph{no precise equation-level published match is claimed} for the displayed connection coefficients. The present audit has not independently inspected the original KLS limiting table for an explicit, separately countable $q$-Hahn $\to$ affine $q$-Krawtchouk arrow. Thus the \emph{mathematical} endpoint and reverse-coefficient limit are verified, but the required primary-source KLS arrow gate is still open. This pair is absent from the existing 45 recorded PASS entries, so there is no identified duplication if the KLS gate is later satisfied.

\noindent\textbf{V102 verdict: MATHEMATICALLY VERIFIED EXACT ENDPOINT; KLS-ARROW ELIGIBILITY PENDING. NEW RECORDED PASS: ZERO; TOTAL REMAINS 45.} This does not exclude the endpoint on the ground that it is exact. V31 is unchanged.


\clearpage
\section{V103: Reopened $q$-Racah to dual $q$-Hahn specialization---explicit reverse-coefficient collapse}
\noindent\textbf{Scope and criterion.} Under the clarified audit rule, an exact specialization may count if it is a recognized Koekoek--Swarttouw limiting arrow, is established from the V31 reverse connection, and is not a duplicate. We examine the distinct-family reverse connection \texttt{eq:finite-dHR-reverse}; self-family connections are not considered.

\paragraph{Terminating polynomial limit.} For $0<q<1$, fixed $N$, $0\le n,x\le N$, and admissible $\gamma,\delta$, V31 defines
\[
R_n(\mu(x);\beta,\gamma,\delta,N\mid q)
={}_4\phi_3\!\left(\begin{matrix}q^{-n},\beta q^{n-N},q^{-x},\gamma\delta q^{x+1}\\q^{-N},\beta\delta q,\gamma q\end{matrix};q,q\right),
\qquad \mu(x)=q^{-x}+\gamma\delta q^{x+1}.
\]
The finite sum has no singularity at $\beta=0$ when the remaining denominator factors are nonzero. Hence
\begin{equation}\label{eq:v103-r-dh}
\boxed{\lim_{\beta\to0}R_n(\mu(x);\beta,\gamma,\delta,N\mid q)
=R_n^{\mathrm{dH}}(\mu(x);\gamma,\delta,N\mid q).}
\end{equation}
The lattice and normalization remain unchanged.

\paragraph{Direct check of the actual reverse coefficients.} V31, \texttt{eq:finite-dHR-reverse}, states
\[
R_n(\mu(x);\beta,\gamma,\delta,N\mid q)
=\sum_{k=0}^{n}D_{nk}(\beta)R_k^{\mathrm{dH}}(\mu(x);\gamma,\delta,N\mid q),
\]
with
\[
D_{nk}(\beta)=\frac{T^R_{n,k}(\beta)}{T^{\mathrm{dH}}_{k,k}}
{}_2\phi_1\!\left(\begin{matrix}q^{k-n},\beta q^{n-N+k}\\\beta\delta q^{k+1}\end{matrix};q,q\right).
\]
At $\beta=0$, the $T$ ratio is finite, with $T^R_{n,n}(0)=T^{\mathrm{dH}}_{n,n}$, and the terminating hypergeometric factor becomes, for $m=n-k$,
\[
{}_2\phi_1\!\left(\begin{matrix}q^{-m},0\\0\end{matrix};q,q\right)
=\sum_{j=0}^{m}\frac{(q^{-m};q)_j}{(q;q)_j}q^j
=(q^{1-m};q)_m
=\begin{cases}1,&m=0,\\0,&m>0.\end{cases}
\]
The second equality is the finite $q$-binomial theorem. Therefore the V31 reverse coefficients themselves obey
\begin{equation}\label{eq:v103-reverse-delta}
\boxed{\lim_{\beta\to0}D_{nk}(\beta)=\delta_{nk}\qquad(0\le k\le n\le N).}
\end{equation}
This independently verifies the full reverse-connection collapse without relying only on polynomial-basis uniqueness.

\paragraph{KLS certification gate.} V31 explicitly cites Koekoek--Swarttouw and KLS for the $\beta=0$ specialization, but does not identify a numbered limiting equation. The original KLS arrow has not been independently checked in this stage. Thus the endpoint and reverse-coefficient proof are \emph{verified}, but the audit's primary-source KLS gate is still pending. This arrow is not among the 45 currently recorded PASS entries; it must not be added without that check.

\noindent\textbf{V103 verdict: MATHEMATICALLY VERIFIED EXACT SPECIALIZATION; INDEPENDENT KLS-ARROW CONFIRMATION PENDING. NEW PASS: ZERO; RECORDED TOTAL: 45.} The V31 source remains unchanged.


\clearpage
\section{V104: Reopened $q$-Racah to ordinary $q$-Krawtchouk endpoint---finite-grid reverse-connection verification}
\noindent\textbf{Scope.} We reassess the distinct-family reverse connection \texttt{eq:finite-KR-reverse} under the clarified rule allowing exact specializations when independently recognized as KLS limiting arrows. This is a mathematical endpoint certification, not a claim of independent primary-source KLS verification.

\paragraph{Matched parameter path.} Fix $0<q<1$, a nonnegative integer $N$, and an admissible ordinary $q$-Krawtchouk parameter $p$. In V31's $q$-Racah notation choose
\[
 \gamma=\delta=0,\qquad \beta=-p q^N,\qquad t=\gamma\delta q=0.
\]
Then the $q$-Racah lattice $\mu(x)=q^{-x}+tq^x$ equals $q^{-x}$ and its defining terminating series becomes
\[
 \begin{aligned}
 R_n(q^{-x};-pq^N,0,0,N\mid q)
 &= {}_4\phi_3\!\left(\begin{matrix}q^{-n},-pq^n,q^{-x},0\\q^{-N},0,0\end{matrix};q,q\right)\\
 &= {}_3\phi_2\!\left(\begin{matrix}q^{-n},q^{-x},-pq^n\\q^{-N},0\end{matrix};q,q\right)
 =K_n(q^{-x};p,N\mid q).
 \end{aligned}
\]
The cancellation of one upper and one lower zero parameter is legitimate because $(0;q)_j=1$ for every $j$. Both series terminate at $j=n\le N$ and have no vanishing $q^{-N}$ denominator before termination. Hence the equality holds for all $0\le n,x\le N$ without normalization.

\paragraph{Reverse-connection coefficient check.} The actual V31 reverse connection is the full finite-grid identity
\[
 R_n(\mu(x);\theta,N\mid q)=\sum_{k=0}^N D_{nk}(\theta;p)K_k(q^{-x};p,N\mid q),
\]
with $D_{nk}$ given by the explicit double sum in \texttt{eq:finite-KR-reverse}. Define the evaluation matrix $E_{xk}=K_k(q^{-x};p,N\mid q)$ and the vector $r_n(\theta)_x=R_n(\mu(x);\theta,N\mid q)$. Under the nondegeneracy assumptions of the V31 proposition, $E$ is invertible. Thus the displayed reverse-connection coefficients satisfy the \emph{unique} linear system
\[
 E D_n(\theta)^T=r_n(\theta).
\]
At the matched parameters, the terminating-series calculation gives $r_n=E_{\cdot n}$, and therefore
\begin{equation}\label{eq:v104-r-k-collapse}
 \boxed{D_{nk}(-pq^N,0,0;p)=\delta_{nk}\quad(0\le n,k\le N).}
\end{equation}
This proves collapse of V31's \emph{actual full-grid reverse connection}, not merely equality of two series. For any admissible parameter path approaching this point, continuity of the finite terminating series and the fixed inverse $E^{-1}$ yields the same coefficient limit. No claim is made about paths encountering singular denominators.

\paragraph{KLS gate and count.} V54 excluded this matched endpoint solely as an exact specialization; that reason is no longer dispositive. However, the source subsection in V31 does not identify a numbered KLS limiting arrow for this particular $q$-Racah $\to$ ordinary $q$-Krawtchouk parameter path. A fixed-parameter identity and its coefficient collapse do not by themselves prove that KLS lists it as a separately countable limiting relation. Pending independent primary-source verification, we record \emph{mathematical verification only} and do not increment the count.

\noindent\textbf{V104 verdict: EXACT ENDPOINT AND REVERSE CONNECTION VERIFIED; KLS ELIGIBILITY PENDING. NEW PASS: ZERO; RECORDED TOTAL: 45.} No self-family connection was considered; original V31 remains unchanged.



\clearpage
\section{V105: Reopened $q$-Racah to $q$-Hahn exact endpoint---full-grid reverse-connection certification}
\noindent\textbf{Scope.} We examine V31's distinct-family $q$-Hahn--$q$-Racah reverse connection, \texttt{eq:finite-HR-double-D}, under the revised rule that exact specializations are admissible \emph{provided} the KLS limiting-arrow gate is independently satisfied. This section establishes the mathematical endpoint and the actual finite-grid coefficient collapse; it does not assume that the KLS gate has been met.

\paragraph{Matched specialization.} Fix $0<q<1$, $N\in\mathbb N_0$, $0\le n,x\le N$, and parameters $A,B$ for which the $q$-Hahn evaluation basis is nondegenerate and all terminating-series denominators are nonzero. In V31's $q$-Racah notation put
\[
 \gamma=A,\qquad \delta=0,\qquad \beta=ABq^{N+1},\qquad t=\gamma\delta q=0.
\]
Then $\mu(x)=q^{-x}$ and the V31 terminating $q$-Racah definition reduces term by term to
\begin{align}
 R_n(\mu(x);ABq^{N+1},A,0,N\mid q)
 &= {}_4\phi_3\!\left(\begin{matrix}q^{-n},ABq^{n+1},q^{-x},0\\q^{-N},0,Aq\end{matrix};q,q\right)\\
 &= {}_3\phi_2\!\left(\begin{matrix}q^{-n},ABq^{n+1},q^{-x}\\Aq,q^{-N}\end{matrix};q,q\right)
 =Q_n(q^{-x};A,B,N\mid q).
 \label{eq:v105-r-h-endpoint}
\end{align}
The zero upper/lower parameters cancel because $(0;q)_j=1$ for every $j$. This is an exact identity on the grid, and the terminating expressions vary continuously along any admissible path approaching these parameters. No normalization is needed.

\paragraph{Actual V31 reverse-connection coefficients.} Equation \texttt{eq:finite-HR-double-D} in V31 has the form
\[
 R_n(\mu(x);\beta,\gamma,\delta,N\mid q)
 =\sum_{k=0}^{N}D_{nk}(\beta,\gamma,\delta;A,B)Q_k(q^{-x};A,B,N\mid q),
\]
where $D_{nk}$ is the explicit finite double sum involving $\mathscr K^H_{j,k}$ and $\mathscr B_{j,s}(0)$. Let $E^H_{xk}=Q_k(q^{-x};A,B,N\mid q)$. The $q$-Hahn polynomials have exact degree $k$ in $y=q^{-x}$ with nonzero leading coefficients under the stated assumptions, and the $N+1$ nodes $q^{-x}$ are distinct. Hence $E^H$ is invertible. The V31 reverse coefficients are therefore uniquely determined by
\[
 E^H D_n^T=r_n,\qquad
 (r_n)_x=R_n(\mu(x);\beta,\gamma,\delta,N\mid q).
\]
At the matched specialization, equation \eqref{eq:v105-r-h-endpoint} gives $r_n=E^H_{\cdot n}$. Thus
\begin{equation}\label{eq:v105-r-h-delta}
 \boxed{D_{nk}(ABq^{N+1},A,0;A,B)=\delta_{nk}\quad(0\le n,k\le N).}
\end{equation}
For any admissible approaching path, continuity of the terminating evaluations and multiplication by the \emph{fixed} inverse $(E^H)^{-1}$ proves $D_{nk}\to\delta_{nk}$. This certifies the full-grid reverse coefficients, rather than inferring the limit solely from a family-name identification. It does not assert a triangular connection in a common variable away from the grid.

\paragraph{KLS and count gate.} V31's literature context for this subsection cites classical $q$-Racah/$q$-Hahn specializations but explicitly cautions that the finite-grid changes of basis are not themselves KLS limiting relations. The present audit has not independently established a numbered KLS limiting arrow for the \emph{particular} matched path above. Accordingly, mathematical verification is complete, but an additional counted PASS is withheld. This relation is not counted as a duplicate, nor rejected for being exact: the sole unresolved issue is independent KLS-arrow eligibility.

\noindent\textbf{V105 verdict: EXACT ENDPOINT AND V31 REVERSE-COEFFICIENT COLLAPSE VERIFIED; INDEPENDENT KLS-ARROW ELIGIBILITY PENDING. NEW PASS: ZERO; RECORDED TOTAL: 45.} The original V31 is unchanged.



\clearpage
\section{V106: Primary-source Koekoek--Swarttouw certification of four reopened finite-$q$ arrows}
\noindent\textbf{Purpose and source.} The previous V102--V105 mathematical verifications withheld PASS solely pending inspection of the original Koekoek--Swarttouw (1998) limit-relations table. That primary source has now been examined directly: R.~Koekoek and R.~F.~Swarttouw, \emph{The Askey-scheme of hypergeometric orthogonal polynomials and its $q$-analogue}, TU Delft Report 98-17, Chapter 4, available at
\texttt{https://fa.ewi.tudelft.nl/\string~koekoek/documents/as98.pdf}.
The equation numbers below refer to that report, \emph{not} to the later KLS book's numbering. The user's clarified eligibility rule admits exact specializations explicitly listed as inter-family limit relations in this source.

\paragraph{(A) $q$-Hahn $\to$ affine $q$-Krawtchouk: NEW PASS 46.}
The original report's \textbf{(4.6.5)} states, verbatim in mathematical notation,
\[
 Q_n(q^{-x};p,0,N\mid q)=K_n^{\mathrm{aff}}(q^{-x};p,N;q).
\]
This matches V102's $A=p$, $B\to0$ endpoint and its proven reverse-coefficient limit for \texttt{eq:finite-affH-reverse}. The pair was absent from the preceding 45-entry ledger. Hence the previously pending KLS gate is \textbf{PASSED}; the new count is 46.

\paragraph{(B) $q$-Racah $\to$ dual $q$-Hahn: NEW PASS 47.}
The report's \textbf{(4.2.5)} states
\[
 R_n(\mu(x);q^{-N-1},0,\gamma,\delta\mid q)
 =R_n^{\mathrm{dH}}(\mu(x);\gamma,\delta,N\mid q),
 \qquad \mu(x)=q^{-x}+\gamma\delta q^{x+1}.
\]
This is the $\beta=0$ path analyzed in V103, with the fixed $q$-Racah truncation parameter $\alpha=q^{-N-1}$ restored explicitly. V103 proved collapse of the V31 reverse coefficients \texttt{eq:finite-dHR-reverse}. The exact arrow is separately listed in K--S and was not previously counted. \textbf{PASS}; new count 47.

\paragraph{(C) $q$-Racah $\to$ ordinary $q$-Krawtchouk: NEW PASS 48.}
The report's \textbf{(4.2.8)} states
\[
 R_n(q^{-x};q^{-N-1},-pq^N,0,0\mid q)
 =K_n(q^{-x};p,N;q).
\]
This exactly matches V104's $\beta=-pq^N$, $\gamma=\delta=0$ path with the implicit truncation parameter $\alpha=q^{-N-1}$ made explicit. V104 established the V31 full-grid reverse-coefficient collapse from \texttt{eq:finite-KR-reverse}. This is a distinct, separately listed K--S arrow. \textbf{PASS}; new count 48.

\paragraph{(D) $q$-Racah $\to$ $q$-Hahn: NEW PASS 49.}
The report's \textbf{(4.2.4)} states
\[
 R_n(\mu(x);q^{-N-1},\beta\gamma q^{N+1},\gamma,0\mid q)
 =Q_n(q^{-x};\gamma,\beta,N\mid q),\qquad \mu(x)=q^{-x}.
\]
Putting $\gamma=A$ and $\beta=B$ in the \emph{report's right-hand parameters} (equivalently the $q$-Racah parameter $\beta_{R}=ABq^{N+1}$) recovers precisely V105's identity. V105 verified the V31 reverse-coefficient collapse in \texttt{eq:finite-HR-double-D}. This arrow is independently enumerated in K--S and was not among the first 48. \textbf{PASS}; new count 49.

\paragraph{Source cross-check and counting discipline.}
All four exact specializations appear in the original report's \emph{Chapter 4: limit relations}, and their explicit equations agree with the parameters and normalizations proved from V31 in V102--V105. This removes the \emph{specific} primary-source K--S gate left open in those four sections; it does not imply global verification of every other historical PASS. The four targets are distinct ordered family pairs, not self-family arrows, and none duplicates the previously counted 45. The historical V102--V105 verdicts are retained as chronological records but \textbf{superseded by this V106 certification}. The original V31 manuscript remains untouched.

\paragraph{Additional primary-source finding, not yet counted.}
The report's \textbf{(4.10.1)} explicitly gives continuous $q$-Jacobi $\to$ continuous $q$-Laguerre as $\beta\to\infty$. This strengthens the K--S eligibility of V70's reopened case, but the precise matching of V70's $B\to0$ parametrization and normalization to that reported path must be checked against the V31 reverse connection before adding a PASS. Similarly, the report contains other exact arrows; no automatic credit is assigned without a matching V31 proof and duplicate check.

\noindent\textbf{V106 verdict: FOUR NEW PRIMARY-SOURCE-CERTIFIED, REVERSE-CONNECTION-VERIFIED PASS ENTRIES; RECORDED TOTAL $45\to49$.} Global completeness and the remaining historical proofs' external verification are still open.


\clearpage
\section{V107: Continuous $q$-Jacobi $\to$ continuous $q$-Laguerre---K--S normalization and reverse-connection limit}
\noindent\textbf{Reassessment of V70.} Under the clarified counting rule, an exact endpoint is admissible when the corresponding arrow occurs in Koekoek--Swarttouw's published limiting scheme. The 1998 report explicitly states the continuous $q$-Jacobi $\to$ continuous $q$-Laguerre arrow as $\beta\to+\infty$ in equation \textbf{(4.10.1)}. We check that the \emph{normalizations} and the reverse connection in V31 actually reproduce that arrow, rather than identifying two unnormalized families by name.

\paragraph{Parameters and normalizations.}
Fix $0<q<1$, $n\in\mathbb N_0$, $x$ fixed, and a real $\alpha$ in the nondegenerate range. Set
\[
 a=A=q^{\alpha/2+1/4}\ne0,\qquad b=aq^{1/2},\qquad B=q^{\beta/2+1/4}.
\]
V31 defines $J_n(x;\alpha,\beta\mid q)=p_n(x;A,Aq^{1/2},-B,-Bq^{1/2}\mid q)$ and $L_n^{(\alpha)}(x\mid q)=a^n Q_n(x;a,aq^{1/2}\mid q)/(q;q)_n$. These are \emph{not} identically the same normalization as the report's $P_n^{(\alpha,\beta)}$ and $P_n^{(\alpha)}$. Comparison of the terminating Askey--Wilson and Al-Salam--Chihara definitions gives
\begin{align}
 P_n^{(\alpha,\beta)}(x\mid q)&=\frac{a^n}{(q;q)_n}J_n(x;\alpha,\beta\mid q),\label{eq:v107-jnorm}\\
 P_n^{(\alpha)}(x\mid q)&=L_n^{(\alpha)}(x\mid q).\label{eq:v107-lnorm}
\end{align}
Indeed, the Askey--Wilson prefactor $(ab,ac,ad;q)_n/a^n$ divided by the standard continuous-$q$-Jacobi prefactor $(q^{\alpha+1};q)_n/(q;q)_n$ is $(q;q)_n/a^n$, since $ab=q^{\alpha+1}$; the analogous continuous-$q$-Laguerre comparison follows from V31's Al-Salam--Chihara normalization. These equalities fix the previously unresolved scaling.

\paragraph{Collapse of V31's actual reverse connection.}
Write equation \texttt{eq:lag-jacobi-new-rev} as
\[
 J_n(x;\alpha,\beta\mid q)=\sum_{k=0}^n D_{nk}(B) L_k^{(\alpha)}(x\mid q).
\]
V70 already established, using the terminating Askey--Wilson specialization at $B=0$, that
\[
 J_n(x;\alpha,\beta\mid q)\big|_{B=0}=\frac{(q;q)_n}{a^n}L_n^{(\alpha)}(x\mid q).
\]
The degree-$0,\ldots,n$ continuous-$q$-Laguerre polynomials are a basis of polynomials of degree at most $n$ for the stated nondegenerate parameters. Uniqueness of the expansion therefore gives the \emph{individual reverse coefficients}
\begin{equation}\label{eq:v107-delta}
 \boxed{\lim_{\beta\to+\infty}D_{nk}(q^{\beta/2+1/4})=\frac{(q;q)_n}{a^n}\,\delta_{nk}\quad(0\le k\le n).}
\end{equation}
The limit is justified coefficientwise: for fixed $n$, the terminating Askey--Wilson polynomial is regular at $B=0$, and its expansion in the fixed Laguerre basis is obtained by an invertible, parameter-independent triangular change of coordinates. Thus no unproved interchange of an infinite sum and limit is involved. Multiplying the reverse connection by $a^n/(q;q)_n$ and using \eqref{eq:v107-jnorm}--\eqref{eq:v107-lnorm} yields
\begin{equation}\label{eq:v107-limit}
 \boxed{\lim_{\beta\to+\infty}P_n^{(\alpha,\beta)}(x\mid q)=P_n^{(\alpha)}(x\mid q),}
\end{equation}
exactly the normalization and direction of Koekoek--Swarttouw (1998), equation \textbf{(4.10.1)}. Note that the \emph{unscaled} V31 reverse matrix tends to a diagonal matrix, not the identity; the identity arises only after the required source normalization.

\paragraph{Independence and verdict.}
The source and target are distinct families; this arrow is not a self-family connection and is not among the previously recorded 49 ordered-pair PASS entries. V70's earlier exclusion was based solely on the exact-specialization rule and is superseded here. The independently listed K--S arrow and the V31 reverse-connection coefficient limit jointly satisfy the revised eligibility test.

\noindent\textbf{V107 verdict: NEW PASS 50---continuous $q$-Jacobi $\to$ continuous $q$-Laguerre. RECORDED TOTAL: $49\to50$.} The original V31 source remains unchanged; global completeness certification remains open.


\clearpage
\section{V108: Al-Salam--Carlitz I $\to$ discrete $q$-Hermite I---exact K--S arrow and explicit reverse coefficients}
\noindent\textbf{Reopened case:} V91's Connection 112 was withheld because the endpoint is an exact specialization. Under the user's clarified eligibility rule, this is no longer an exclusion. Koekoek--Swarttouw (1998), Chapter 4, section 4.24, explicitly enumerates ``Al-Salam--Carlitz I $\to$ Discrete $q$-Hermite I'' and records equation \textbf{(4.24.1)}:
\[
 U_n^{(-1)}(x;q)=h_n(x;q).
\]
This is a distinct-family arrow listed in the original limiting-relations chapter, even though the endpoint is attained by parameter substitution.

\paragraph{The actual V31 reverse connection.}
Equation \texttt{eq:dHI-carlitz-rev} gives, for every $n\geq0$,
\begin{equation}\label{eq:v108-reverse}
 U_n^{(a)}(x;q)=\sum_{k=0}^n (-1)^{n-k}
 {n\brack k}_q(-a;q)_{n-k}\,h_k(x;q).
\end{equation}
For $0<q<1$ and $a\to-1$, the factors $(-a;q)_{n-k}$ are polynomials in $a$, so passage through this \emph{finite} sum is immediate. In the limit,
\[
 (-a;q)_m\big|_{a=-1}=(1;q)_m=\begin{cases}1,&m=0,\\0,&m>0.\end{cases}
\]
Thus the reverse-connection coefficients satisfy the explicit coefficientwise identity
\begin{equation}\label{eq:v108-delta}
 \boxed{\lim_{a\to-1}(-1)^{n-k}{n\brack k}_q(-a;q)_{n-k}=\delta_{nk},}
\end{equation}
which yields $\lim_{a\to-1}U_n^{(a)}(x;q)=h_n(x;q)$ for every fixed $n$ and $x$. No renormalization or variable transformation is required. This is a \emph{direct evaluation of V31's displayed reverse coefficients}, not merely an appeal to uniqueness of polynomial expansions.

\paragraph{Duplicate and orientation check.}
The arrow is Al-Salam--Carlitz I $\to$ discrete $q$-Hermite I, not the converse. V91 recorded it as an uncounted exact endpoint; no prior numbered PASS in the V107 ledger covers this ordered family pair. The identity is not a self-family connection. The published source equation (4.24.1) agrees with V31's normalization and variable, so the K--S gate and the reverse-connection gate both pass.

\noindent\textbf{V108 verdict: NEW PASS 51---Al-Salam--Carlitz I $\to$ discrete $q$-Hermite I. RECORDED TOTAL $50\to51$.}

\paragraph{Companion case not counted.}
Koekoek--Swarttouw equation \textbf{(4.25.1)} also lists Al-Salam--Carlitz II $\to$ discrete $q$-Hermite II, but its endpoint is
$\widetilde h_n(x;q)=i^{-n}V_n^{(-1)}(ix;q)$, \emph{not} $V_n^{(-1)}(x;q)=\widetilde h_n(x;q)$. V31's \texttt{eq:dHII-carlitz-rev} uses the unrotated variable $x$, so merely setting $a=-1$ in that reverse matrix does not prove an identity collapse. This arrow remains a candidate requiring an explicitly transformed reverse-connection argument and normalization check. No additional PASS is awarded for it in V108.

\paragraph{Primary-source reference.} R. Koekoek and R. F. Swarttouw, \emph{The Askey-scheme of hypergeometric orthogonal polynomials and its $q$-analogue}, Delft University of Technology Report 98-17 (1998), Chapter 4, equations (4.24.1) and (4.25.1), \texttt{https://fa.ewi.tudelft.nl/\string~koekoek/documents/as98.pdf}. The original V31 manuscript is unchanged. Global completeness certification remains open.


\clearpage
\section{V109: Al-Salam--Carlitz II $\to$ discrete $q$-Hermite II---rotated reverse connection and K--S (4.25.1)}
\noindent\textbf{Question carried over from V108.} Koekoek--Swarttouw (1998), equation (4.25.1), states
\begin{equation}\label{eq:v109-target}
 \widetilde h_n(x;q)=i^{-n}V_n^{(-1)}(ix;q).
\end{equation}
The rotation $x\mapsto ix$ and degree-dependent factor $i^{-n}$ are essential. We verify the transformed reverse connection, not the false unrotated identity $V_n^{(-1)}(x)=\widetilde h_n(x)$.

\paragraph{Reverse connection and its correct transformation.}
Write V31's \texttt{eq:dHII-carlitz-rev} as
\begin{align}\label{eq:v109-rev}
V_n^{(a)}(y;q)&=\sum_{k=0}^n C_{nk}(a)\widetilde h_k(y;q),\\
C_{nk}(a)&={n\brack k}_q(-1)^{n-k}q^{\binom{k}{2}-\binom{n}{2}}
\sum_{r=0}^{n-k}{n-k\brack r}_q(i;q)_r
\prod_{j=0}^{n-k-r-1}(a+iq^j).\notag
\end{align}
This is a finite sum with polynomial dependence on $a$, so its specialization at $a=-1$ is legitimate. Define the polynomial dilation matrix $T_{kj}(i)$ by
\[
 \widetilde h_k(ix;q)=\sum_{j=0}^k T_{kj}(i)\widetilde h_j(x;q).
\]
It exists uniquely since the $\widetilde h_j$ are monic polynomials of degree $j$. Substitution of $y=ix$ into \eqref{eq:v109-rev} followed by multiplication by $i^{-n}$ yields the \emph{actual rotated reverse-connection coefficients}
\begin{equation}\label{eq:v109-transformed}
 i^{-n}V_n^{(a)}(ix;q)=\sum_{j=0}^n F_{nj}(a)\widetilde h_j(x;q),\qquad
 F_{nj}(a)=i^{-n}\sum_{k=j}^n C_{nk}(a)T_{kj}(i).
\end{equation}
The task is to show $F_{nj}(-1)=\delta_{nj}$, not $C_{nj}(-1)=\delta_{nj}$. The latter is false already in degree two.

\paragraph{All-degree evaluation of the transformed coefficients.}
V31 proves the reverse connection using the generating functions
\begin{align*}
\sum_{n\ge0}\frac{(-1)^nq^{\binom n2}}{(q;q)_n}V_n^{(a)}(y;q)t^n
 &=\frac{(yt;q)_\infty}{(t,at;q)_\infty},\\
\sum_{n\ge0}\frac{(-1)^nq^{\binom n2}}{(q;q)_n}\widetilde h_n(x;q)t^n
 &=\frac{(xt;q)_\infty}{(-t^2;q^2)_\infty}.
\end{align*}
At $a=-1$, set $y=ix$ and replace the generating variable by $t/i$. Since $(ix)(t/i)=xt$ and
\[
 (t/i,-t/i;q)_\infty=(-it,it;q)_\infty=(-t^2;q^2)_\infty,
\]
the first generating function becomes \emph{exactly} the second. Comparing coefficients of $t^n$ proves
\begin{equation}\label{eq:v109-identity}
 i^{-n}V_n^{(-1)}(ix;q)=\widetilde h_n(x;q)\quad(n\ge0).
\end{equation}
Now specialize the finite reverse connection \eqref{eq:v109-transformed} to $a=-1$. Its left-hand side is \eqref{eq:v109-identity}; uniqueness of expansion in the fixed monic $\widetilde h$ basis forces
\begin{equation}\label{eq:v109-collapse}
 \boxed{\lim_{a\to-1}F_{nj}(a)=i^{-n}\sum_{k=j}^n C_{nk}(-1)T_{kj}(i)=\delta_{nj}.}
\end{equation}
This establishes collapse of the \emph{rotated and normalized} V31 reverse-connection matrix for every degree. The generating functions provide the independent all-degree endpoint calculation; the displayed finite reverse connection and basis uniqueness then establish the individual transformed coefficients. No claim is made that the unrotated coefficients themselves collapse.

\paragraph{Explicit degree-two guard against a false identity.}
The degree-two discrete $q$-Hermite II polynomial is
$\widetilde h_2(x;q)=x^2-(1-q)/q$. Direct substitution into V31's coefficient formula gives
$C_{22}(-1)=1$, $C_{21}(-1)=0$, and $C_{20}(-1)=2(1-q)/q$. Thus
\[
 V_2^{(-1)}(x;q)=\widetilde h_2(x;q)+2(1-q)/q=x^2+(1-q)/q,
\]
which is \emph{not} $\widetilde h_2(x;q)$, while
$-V_2^{(-1)}(ix;q)=\widetilde h_2(x;q)$ as required. This checks both the rotation and its degree-dependent normalization.

\paragraph{Primary-source and duplicate check.}
The original 1998 Koekoek--Swarttouw report, Chapter 4, equation (4.25.1), lists this distinct-family arrow explicitly. The ordered family pair Al-Salam--Carlitz II $\to$ discrete $q$-Hermite II differs from PASS 51 (the type-I pair) and was not previously counted. The unrotated V31 formula has now been used with the required rotation, and its transformed coefficient matrix has been verified to tend to the identity.

\noindent\textbf{V109 verdict: NEW PASS 52---Al-Salam--Carlitz II $\to$ discrete $q$-Hermite II. RECORDED TOTAL $51\to52$.} This supersedes V108's pending classification for the type-II companion, but leaves V108's PASS 51 unchanged. Global completeness certification remains open; the V31 manuscript is unchanged.



\clearpage
\section{V110: duplicate-safe reassessment of two exact $q$-family embeddings}
\noindent\textbf{Purpose.} The revised eligibility rule permits exact specializations when they represent an independently recognized Koekoek--Swarttouw (K--S) limiting arrow. It does \emph{not} permit a second credit for an already counted ordered family pair, nor does a valid exact polynomial identity by itself establish a separately enumerated K--S arrow. We apply both distinctions to V83 and V82. This section supersedes their older blanket ``exact-specialization excluded'' rationale, without deleting their valid mathematical calculations.

\subsection*{A. Continuous $q$-Hahn $\to$ $q$-Meixner--Pollaczek (V83)}
In the notation of V83, $h=be^{i\phi}$ and the V31 reverse connection \texttt{eq:cqh-qmp23} is
\[
 \mathscr Q_n(x;a,b,0,a,\phi\mid q)
 =\sum_{k=0}^n {n\brack k}_q(q;q)_k
 (-h)^{n-k}q^{\binom{n-k}{2}}(a^2q^k;q)_{n-k}
 \mathscr P_k(x;a,\phi\mid q).
\]
At $b=0$, the coefficient for $k<n$ vanishes and the $k=n$ coefficient is $(q;q)_n$. Thus the \emph{normalized} coefficient matrix $D_{nk}(b)/(q;q)_n$ tends to $\delta_{nk}$, and
\[
 \lim_{b\to0}\frac{\mathscr Q_n(x;a,b,0,a,\phi\mid q)}{(q;q)_n}
 =\mathscr P_n(x;a,\phi\mid q).
\]
The result is mathematically sound and an exact endpoint is not a reason to exclude it. However, the ordered family pair continuous $q$-Hahn $\to$ $q$-Meixner--Pollaczek was already included among the four historical PASS entries explicitly retained in V99 (and identified in the original 43-entry ledger). The V83 identity is another realization of that \emph{same ordered family arrow}, not an additional arrow. Therefore the new incremental count is \textbf{zero}, irrespective of whether its parameterization is a useful independent proof. A second PASS would be a duplicate.

\subsection*{B. Al-Salam--Chihara $\leftrightarrow$ continuous $q$-Laguerre (V82)}
The V31 normalization is
\[
 L_n^{(\lambda)}(x\mid q)=\frac{a^n}{(q;q)_n}
 Q_n(x;a,aq^{1/2}\mid q),\qquad a=q^{\lambda/2+1/4}.
\]
V82's regularized reverse coefficient kernel is
\[
 \mathcal K_m(u,v;a,b)=\sum_{r=0}^m {m\brack r}_q
 a^r(u/a;q)_r b^{m-r}(v/b;q)_{m-r},\quad b=aq^{1/2}.
\]
At $(u,v)=(a,b)$, each term for $m>0$ contains $(1;q)_j=0$ for some $j>0$; hence $\mathcal K_m(a,b;a,b)=\delta_{m0}$. Consequently the reverse coefficients collapse to
\[
 D_{nk}(a,b)=\frac{(q;q)_k}{a^k}\delta_{nk},
 \qquad Q_n(x;a,aq^{1/2}\mid q)
 =\frac{(q;q)_n}{a^n}L_n^{(\lambda)}(x\mid q).
\]
This is a valid \emph{exact normalized embedding} and is not excluded merely for being exact. Yet V31's own literature context for \texttt{eq:lag-asc-hyper-rev51} identifies it as a parameter-specialized Al-Salam--Chihara basis change and does not cite a distinct K--S limiting arrow. The matched-pair identity alone cannot supply the missing independent K--S arrow: its parameter values do not approach a singular boundary, and no numbered K--S relation for this \emph{ordered pair and normalization} has been established in this audit. Its status is therefore \textbf{mathematically verified; independent K--S eligibility not established}, not a categorical rejection of exact specializations. No PASS is awarded at this stage.

\paragraph{V110 counting control.}
Starting from the V109 recorded total of \textbf{52}, the continuous $q$-Hahn/$q$-Meixner--Pollaczek realization adds \textbf{0} (duplicate); the Al-Salam--Chihara/continuous $q$-Laguerre matched embedding adds \textbf{0} (separate K--S arrow unverified). The \textbf{recorded PASS total remains 52}. The original V31 source is unchanged. The independent primary-source check and global completeness certification remain open.


\clearpage
\section{V111: independent coefficient checks and normalization controls for the Carlitz--Hermite reinstatements}
\noindent\textbf{Purpose.} The revised eligibility rule restored two exact specializations in V108--V109. This checkpoint independently compares the \emph{literal V31 reverse coefficients} with the claimed low-degree endpoints and isolates the indispensable rotation in the type-II case. It does not award a second PASS for either arrow.

\subsection*{A. Type I: literal coefficient collapse}
V31 equation \texttt{eq:dHI-carlitz-rev} has
\[
 C^{\rm I}_{nk}(a)=(-1)^{n-k}{n\brack k}_q(-a;q)_{n-k}.
\]
In degrees one and two this gives, directly,
\begin{align*}
 U_1^{(a)}(x;q)&=h_1(x;q)-(1+a)h_0(x;q),\\
 U_2^{(a)}(x;q)&=h_2(x;q)-(1+q)(1+a)h_1(x;q)
 +(1+a)(q+a)h_0(x;q).
\end{align*}
Both off-diagonal coefficients vanish at $a=-1$. In general, if $n>k$, the factor $(-a;q)_{n-k}$ contains $1+a$, whereas $C^{\rm I}_{nn}=1$. Thus the all-degree identity limit in V108 is confirmed directly by the \emph{source's displayed coefficient formula}, without any appeal to an indirect polynomial definition.

\subsection*{B. Type II: the unrotated reverse matrix does not collapse}
Set $a=-1$ in V31 equation \texttt{eq:dHII-carlitz-rev}. For $n=1$ the lower coefficient is
\[
 C^{\rm II}_{10}(a)=-\bigl[(a+i)+(1-i)\bigr]=-(a+1),
\]
so $C^{\rm II}_{10}(-1)=0$. This first-degree coincidence alone could conceal an invalid all-degree argument. In degree two, however, direct evaluation of the same V31 coefficient yields
\[
 C^{\rm II}_{22}(-1)=1,\qquad C^{\rm II}_{21}(-1)=0,\qquad
 C^{\rm II}_{20}(-1)=\frac{2(1-q)}q.
\]
Consequently the \emph{unrotated} matrix is not the identity. Since
$\widetilde h_2(x;q)=x^2-(1-q)/q$, the reverse connection gives
\[
 V_2^{(-1)}(x;q)=x^2+\frac{1-q}{q}.
\]
After rotation and normalization,
\[
 i^{-2}V_2^{(-1)}(ix;q)=-\left[-x^2+\frac{1-q}{q}\right]
 =\widetilde h_2(x;q).
\]
The all-degree result in V109 uses the transformed matrix
$F_{nj}(a)=i^{-n}\sum_{k=j}^n C^{\rm II}_{nk}(a)T_{kj}(i)$,
not the unrotated $C^{\rm II}_{nj}(a)$. This distinction is a \emph{necessary normalization control}, not an optional change of convention.

\subsection*{C. Primary-source equation cross-reference and ledger control}
The source manuscript explicitly states $h_n(x;q)=U_n^{(-1)}(x;q)$ in the text adjacent to \texttt{eq:dHI-carlitz-rev}, and explicitly states
$\widetilde h_n(x;q)=i^{-n}V_n^{(-1)}(ix;q)$ adjacent to
\texttt{eq:dHII-carlitz-rev}. The V108 and V109 claims therefore use the actual source normalization, not a silently substituted monic family. The K--S citations already recorded for these arrows are (4.24.1) and (4.25.1), respectively; this stage does not independently re-inspect the external report.

\noindent\textbf{V111 verdict:} The source-formula and normalization checks support the two \emph{existing} recorded PASS entries 51 and 52. No new distinct-family arrow was established, and no earlier PASS is rescinded. \textbf{Recorded total remains 52.} Global completeness and remaining K--S eligibility checks remain open. The original V31 manuscript has not been modified.



\clearpage
\section{V112: independent degree-three through degree-five check of the rotated type-II reverse connection}
\noindent\textbf{Scope and count.} This stage stress-tests PASS 52 beyond the degree-two check in V111. It derives both endpoint polynomial arrays \emph{independently} from their generating functions at the exact rational base $q=1/2$, and checks the necessary complex rotation through degree five. It is a finite-degree consistency check, not a replacement for the all-degree proof in V109, and awards no additional arrow.

\paragraph{Reproducible finite coefficient extraction.}
Write $[t^n]$ for coefficient extraction and $(q;q)_0=1$. At $a=-1$, expand the Al-Salam--Carlitz II generating function using
\[
 (xt;q)_\infty=\sum_{r\ge0}\frac{(-1)^r q^{\binom r2}x^r t^r}{(q;q)_r},\qquad
 \frac1{(t;q)_\infty(-t;q)_\infty}
 =\sum_{u,v\ge0}\frac{(-1)^v t^{u+v}}{(q;q)_u(q;q)_v}.
\]
For the discrete $q$-Hermite II endpoint, independently use
\[
 \frac1{(-t^2;q^2)_\infty}
 =\sum_{j\ge0}\frac{(-1)^j t^{2j}}{(q^2;q^2)_j}.
\]
In both cases multiply the extracted $t^n$ coefficient by
$(q;q)_n/[(-1)^n q^{\binom n2}]$ to obtain the polynomial in the normalization used in V109. These formulas involve only finite sums at any fixed $n$ and permit exact rational arithmetic.

\paragraph{Exact checks at $q=1/2$.}
Direct coefficient extraction gives the following two polynomial sequences (the common variable is $x$):
\[
\begin{array}{c|l|l}
 n & V_n^{(-1)}(x;1/2)&\widetilde h_n(x;1/2)\\\hline
 0&1&1\\
 1&x&x\\
 2&x^2+1&x^2-1\\
 3&x^3+7x&x^3-7x\\
 4&x^4+35x^2+28&x^4-35x^2+28\\
 5&x^5+155x^3+868x&x^5-155x^3+868x
\end{array}
\]
For each row $n=0,\ldots,5$, direct substitution verifies
\[
 \boxed{i^{-n}V_n^{(-1)}(ix;1/2)=\widetilde h_n(x;1/2).}
\]
In particular, degrees three, four and five provide new, nontrivial checks of the rotation: the $x^{n-2}$ terms change sign, whereas the $x^{n-4}$ terms do not. Omitting the rotation fails already in degree two and in all three newly checked degrees.

\paragraph{What this does and does not certify.}
Because V31's finite reverse connection is a polynomial identity and the $\widetilde h_j$ form a monic basis, the verified endpoint identities at these degrees imply that the corresponding transformed reverse coefficients $F_{nj}(-1)$ equal $\delta_{nj}$ for $n\le5$ at $q=1/2$. This independently checks the conclusion of V109 in several nontrivial degrees; the all-$q$, all-degree proof continues to rest on V109's formal generating-function identity. The original K--S equation (4.25.1) was cited in V109; this section makes no new claim of independent external source inspection.

\noindent\textbf{V112 verdict:} PASS 52 survives exact arithmetic spot checks through degree five. No new family transition has been certified. \textbf{Recorded PASS total remains 52.} The V31 manuscript remains unchanged; global completeness and outstanding K--S classifications remain open.



\clearpage
\section{V113: independent all-degree recurrence certification of the rotated Carlitz-II endpoint}
\noindent\textbf{Objective.} V109 proved the rotated Al-Salam--Carlitz II/discrete $q$-Hermite II endpoint using a generating-function substitution and uniqueness of the V31 reverse-basis expansion; V112 supplied exact checks at $q=1/2$ through degree five. This stage supplies a second, all-degree proof based on a polynomial recurrence extracted from a $q$-difference equation. It is a recertification of the existing PASS 52, not a new family arrow.

\paragraph{Common endpoint generating function.}
For $0<q<1$ define
\[
 G(x,t)=\frac{(xt;q)_\infty}{(-t^2;q^2)_\infty}
 =\sum_{n\ge0}c_n H_n(x)t^n,
 \qquad c_n=\frac{(-1)^nq^{n(n-1)/2}}{(q;q)_n}.
\]
The $q$-product shift laws give the exact formal power-series equation
\begin{equation}\label{eq:v113-qdiff}
 (1-xt)G(x,qt)=(1+t^2)G(x,t).
\end{equation}
The constant coefficient is $H_0=1$. For $n\ge1$, extracting $t^n$ from \eqref{eq:v113-qdiff} yields
\[
 (q^n-1)c_nH_n=xq^{n-1}c_{n-1}H_{n-1}+c_{n-2}H_{n-2},
\]
with the last term omitted when $n=1$. The elementary ratios
\[
 \frac{q^{n-1}c_{n-1}}{(q^n-1)c_n}=1,
 \qquad
 \frac{c_{n-2}}{(q^n-1)c_n}
 =-q^{3-2n}(1-q^{n-1})
\]
therefore imply the monic three-term recurrence
\begin{equation}\label{eq:v113-rec}
 H_0(x)=1,\qquad H_1(x)=x,\qquad
 H_n(x)=xH_{n-1}(x)-q^{3-2n}(1-q^{n-1})H_{n-2}(x)
 \quad(n\ge2).
\end{equation}
This derivation is valid formally over $\mathbb Q(q)[x][[t]]$ and hence does not require numerical evaluation or a limiting interchange.

\paragraph{Identification with both families.}
V31's Al-Salam--Carlitz II generating function at $a=-1$, after $y=ix$ and $t\mapsto t/i$, becomes $G(x,t)$ because
\[
 (t/i,-t/i;q)_\infty=(-t^2;q^2)_\infty.
\]
The coefficient of $t^n$ in that transformed series is
$c_n i^{-n}V_n^{(-1)}(ix;q)$. V31's discrete $q$-Hermite II generating function is also $G(x,t)$, with coefficient $c_n\widetilde h_n(x;q)$. Thus both polynomial sequences independently satisfy \eqref{eq:v113-rec} with the same initial data. Uniqueness of the recurrence gives, for every $n\ge0$,
\begin{equation}\label{eq:v113-endpoint}
 \boxed{i^{-n}V_n^{(-1)}(ix;q)=\widetilde h_n(x;q).}
\end{equation}
This is an independent recurrence-based all-degree certification of the endpoint identity; it does not assume the degree-five computations of V112.

\paragraph{Transfer to the literal V31 reverse matrix.}
V31 equation \texttt{eq:dHII-carlitz-rev} is a finite polynomial identity, with coefficients $C_{nk}(a)$ polynomial in $a$. Apply $y=ix$ and multiply by $i^{-n}$, then express $\widetilde h_k(ix;q)$ in the monic $\widetilde h_j(x;q)$ basis via $T_{kj}(i)$. The resulting coefficient matrix is precisely
\[
 F_{nj}(a)=i^{-n}\sum_{k=j}^{n}C_{nk}(a)T_{kj}(i)
\]
as defined in V109. The recurrence identity \eqref{eq:v113-endpoint} and uniqueness of polynomial-basis coordinates force
\[
 \boxed{F_{nj}(-1)=\delta_{nj}\quad(0\le j\le n),}
\]
and polynomial dependence on $a$ gives the same value as $a\to-1$. The unrotated $C_{nj}(-1)$ are \emph{not} claimed to equal $\delta_{nj}$.

\paragraph{Scope and ledger decision.}
The proof applies for $0<q<1$ and all fixed degrees, and explicitly handles the complex rotation and normalization. It strengthens the internal mathematical verification of the \emph{already counted} Al-Salam--Carlitz II $\to$ discrete $q$-Hermite II relation (PASS 52). The original K--S citation (4.25.1) remains the previously recorded external source; no new primary-source inspection is claimed here. No new distinct-family arrow is added. \textbf{Recorded PASS total remains 52.} Global completeness and outstanding K--S classifications remain open. The original V31 manuscript is unchanged.



\clearpage
\section{V114: symbolic parameter-dependent reverse-matrix check for the rotated type-II endpoint}
\noindent\textbf{Objective.} Earlier V111--V113 tests establish the endpoint at $a=-1$, including an all-degree argument. This stage performs a different check: retain the free Al-Salam--Carlitz II parameter $a$ and compute the first two nontrivial rows of the \emph{transformed} reverse-connection matrix as exact rational functions of $q,a$. This detects normalization or rotation mistakes that an endpoint-only identity could conceal. No new family arrow is claimed.

\paragraph{Derivation without assuming the endpoint.}
The source V31 gives, in its monic convention (adjacent to \texttt{eq:dHII-carlitz-rev}),
\[
 \sum_{n\ge0}\frac{(-1)^nq^{\binom n2}}{(q;q)_n}V_n^{(a)}(y;q)t^n
 =\frac{(yt;q)_\infty}{(t,at;q)_\infty}.
\]
The shift identity yields
\[
 (1-yt)G_a(y,qt)=(1-t)(1-at)G_a(y,t).
\]
Extracting the first three coefficients gives
\begin{align*}
 V_0^{(a)}(y;q)&=1,\\
 V_1^{(a)}(y;q)&=y-(1+a),\\
 V_2^{(a)}(y;q)&=y^2-\frac{(1+a)(1+q)}q\,y
 +\frac{a^2+a(1+q)+1}{q}.
\end{align*}
Independently, the discrete $q$-Hermite II basis starts with
$\widetilde h_0=1$, $\widetilde h_1=x$, and
$\widetilde h_2=x^2-(1-q)/q$.

\paragraph{Explicit rotated basis coefficients.}
Define $W_n^{(a)}(x;q)=i^{-n}V_n^{(a)}(ix;q)$, with $i^2=-1$.
Direct substitution and re-expansion in the $\widetilde h_j$ basis gives
\begin{align}
 W_0^{(a)}&=\widetilde h_0,\nonumber\\
 W_1^{(a)}&=\widetilde h_1+i(1+a)\widetilde h_0,\label{eq:v114-row1}\\
 W_2^{(a)}&=\widetilde h_2
 +i\frac{(1+a)(1+q)}q\widetilde h_1
 -\frac{(1+a)(a+q)}q\widetilde h_0.
 \label{eq:v114-row2}
\end{align}
Thus the literal transformed reverse-basis matrix, in degrees $0\le n\le2$, is
\[
 \boxed{F^{(\le2)}(a)=
 \begin{pmatrix}
 1&0&0\\
 i(1+a)&1&0\\
 -\dfrac{(1+a)(a+q)}q&i\dfrac{(1+a)(1+q)}q&1
 \end{pmatrix}.}
\]
Every off-diagonal entry contains $1+a$, and hence $F^{(\le2)}(-1)=I_3$ for $q\ne0$. This is a \emph{parameter-dependent} calculation, not merely a check of the already known endpoint. In particular the unrotated degree-two coefficients are not asserted to vanish at $a=-1$.

\paragraph{Consistency and limitations.}
Equations \eqref{eq:v114-row1}--\eqref{eq:v114-row2} follow from the same monic source normalization as V31 and can also be obtained by applying the V31 finite reverse connection followed by the dilation matrix $T_{kj}(i)$. Since polynomial-basis coordinates are unique, these are necessarily the corresponding entries of the transformed reverse matrix. This explicitly verifies its first three rows at \emph{general} $a$ and provides a concrete diagnostic for future symbolic or numerical implementations. The all-degree proof remains V113; this section does not independently inspect K--S (4.25.1).

\noindent\textbf{V114 verdict:} Exact parameter-dependent transformed reverse coefficients verified through degree two, with the required factor $1+a$ in all off-diagonal entries. Existing PASS 52 retained; \textbf{recorded total remains 52}. No additional arrow awarded. Global K--S eligibility and completeness checks remain outstanding; V31 is unchanged.



\clearpage
\section{V115: all-degree parameter-dependent factorization of the rotated Carlitz-II reverse matrix}
\noindent\textbf{Aim.} V114 checked the factor $1+a$ in the first three rows of the transformed reverse matrix. We now establish this factor for \emph{every} strictly lower-triangular entry, without extrapolating from low degrees. This is a structural consistency test of existing PASS 52, not a new K--S arrow.

\paragraph{Exact source recurrence at arbitrary $a$.}
Write the V31 generating function as
\[
 G_a(y,t)=\sum_{n\ge0}\frac{(-1)^nq^{n(n-1)/2}}{(q;q)_n}V_n^{(a)}(y;q)t^n
 =\frac{(yt;q)_\infty}{(t,at;q)_\infty}.
\]
Its $q$-shift equation is
\[
 (1-yt)G_a(y,qt)=(1-t)(1-at)G_a(y,t).
\]
Comparing coefficients of $t^n$ (with $V_{-1}=0$) yields, for $n\ge1$,
\begin{equation}
 \boxed{V_n^{(a)}(y;q)=
 \bigl(y-(1+a)q^{1-n}\bigr)V_{n-1}^{(a)}(y;q)
 -a q^{3-2n}(1-q^{n-1})V_{n-2}^{(a)}(y;q).}
 \label{eq:v115-carlitz-rec}
\end{equation}
In particular, the $n=1$ term is $y-(1+a)$, consistent with V114.

\paragraph{Rotated recurrence, retaining the free parameter.}
Set $W_n^{(a)}(x;q)=i^{-n}V_n^{(a)}(ix;q)$. Since $i^{-1}=-i$ and $i^{-2}=-1$, equation \eqref{eq:v115-carlitz-rec} becomes
\begin{equation}
 \boxed{W_n^{(a)}(x;q)=
 \bigl(x+i(1+a)q^{1-n}\bigr)W_{n-1}^{(a)}(x;q)
 +a q^{3-2n}(1-q^{n-1})W_{n-2}^{(a)}(x;q).}
 \label{eq:v115-rot-rec}
\end{equation}
At $a=-1$ this is exactly the monic discrete $q$-Hermite II recurrence
\[
 \widetilde h_n(x;q)=x\widetilde h_{n-1}(x;q)
 -q^{3-2n}(1-q^{n-1})\widetilde h_{n-2}(x;q),
\]
with matching initial values $W_0^{(-1)}=1$, $W_1^{(-1)}=x$. Thus $W_n^{(-1)}=\widetilde h_n$ for all $n$, as required.

\paragraph{All-degree divisibility of reverse coefficients.}
Fix $q\in(0,1)$, and expand the monic polynomials in the fixed monic basis:
\[
 W_n^{(a)}(x;q)=\sum_{j=0}^n F_{nj}(a;q)\widetilde h_j(x;q).
\]
By recurrence \eqref{eq:v115-rot-rec}, every coefficient of $W_n^{(a)}$ as a polynomial in $x$ belongs to $\mathbb C[a]$ (with $q$ fixed); there is no denominator involving $a$. Change of basis from $1,x,\ldots,x^n$ to the monic $\widetilde h_j$ basis is a unitriangular linear transformation independent of $a$. Hence every $F_{nj}(a;q)$ is a polynomial in $a$, and $F_{nn}(a;q)=1$ identically. The established endpoint identity $W_n^{(-1)}=\widetilde h_n$ gives $F_{nj}(-1;q)=0$ whenever $j<n$. The polynomial factor theorem therefore implies the \emph{all-degree} statement
\begin{equation}
 \boxed{F_{nj}(a;q)=\delta_{nj}+(1+a)G_{nj}(a;q),\qquad
 G_{nj}(a;q)\in\mathbb C[a],\quad G_{nn}=0.}
 \label{eq:v115-factor}
\end{equation}
For fixed $n$, this proves coefficientwise $F_{nj}(a;q)\to\delta_{nj}$ as $a\to-1$ with no numerical truncation and no assumption that the unrotated matrix is diagonal.

\paragraph{Relation to V31 and certification boundary.}
The coefficients $F_{nj}$ are exactly the coefficients obtained by composing the V31 reverse connection \texttt{eq:dHII-carlitz-rev} with the polynomial dilation $x\mapsto ix$ and multiplying degree $n$ by $i^{-n}$: uniqueness of expansion in the $\widetilde h_j$ basis identifies both constructions. Equation \eqref{eq:v115-factor} strengthens V114's degree-two observation to all degrees and establishes a useful test for implementations of the transformed reverse connection. This is \emph{not} a new independent limiting relation and does not recheck the original K--S source or the other historical PASS entries.

\noindent\textbf{V115 verdict:} All-degree factorization of the parameter-dependent transformed reverse matrix \textbf{proved}; PASS 52 retained; \textbf{recorded PASS total remains 52}. Independent K--S certification and global completeness remain open. The V31 source is unchanged.


\clearpage
\section{V116: exact Laguerre--Al-Salam--Chihara endpoint and independent-arrow test}
\noindent\textbf{Question.} V110 left the continuous $q$-Laguerre/Al-Salam--Chihara exact endpoint mathematically verified but without a separately identified Koekoek--Swarttouw (K--S) limiting arrow. This stage distinguishes the reverse-matrix assertion from the independent-arrow requirement. It does not infer K--S eligibility from a family identity.

\paragraph{V31 source and exact endpoint.}
The relevant source is V31 equation \texttt{eq:lag-asc-hyper-rev51}. Set
\[
 a=q^{\lambda/2+1/4},\qquad b=aq^{1/2},\qquad
 u=a,\quad v=b,\qquad 0<q<1.
\]
The normalization used in V31 is
\begin{equation}\label{eq:v116-asc-lag}
 L_n^{(\lambda)}(x\mid q)=\frac{a^n}{(q;q)_n}Q_n(x;a,b\mid q),
 \qquad Q_n(x;a,b\mid q)=\frac{(q;q)_n}{a^n}L_n^{(\lambda)}(x\mid q).
\end{equation}
These are exact polynomial identities, not asymptotic statements.

\paragraph{Reverse coefficients at the matching pair.}
Write the reverse expansion from V31 as
\[
 Q_n(x;u,v\mid q)=\sum_{k=0}^n D_{nk}(u,v;\lambda)L_k^{(\lambda)}(x\mid q).
\]
The V31 proof constructs $D_{nk}$ by first expanding $Q_n(x;u,v)$ in the monic Al-Salam--Chihara basis $Q_k(x;a,b)$ and then multiplying the $k$th coefficient by $(q;q)_k/a^k$. At $(u,v)=(a,b)$, the first basis change is the identity, so
\begin{equation}\label{eq:v116-diag}
 \boxed{D_{nk}(a,b;\lambda)=\frac{(q;q)_n}{a^n}\delta_{nk}.}
\end{equation}
This conclusion is independent of the apparent $0/0$ factors in the terminating ${}_2\phi_1$ representation: V31 explicitly supplies the finite polynomial coefficient-sum continuation at exceptional parameters. For fixed $n$, those finite coefficients are continuous at $(a,b)$; hence the same diagonal matrix is their limit as $(u,v)\to(a,b)$.

\paragraph{Degree-one normalization diagnostic.}
The monic Al-Salam--Chihara polynomial is $Q_1(x;u,v\mid q)=2x-u-v$. In the matched case, $Q_1(x;a,b)=2x-a-b$, whereas
\[
 L_1^{(\lambda)}(x\mid q)=\frac{a}{1-q}(2x-a-b).
\]
Thus $D_{11}=(1-q)/a$, in agreement with \eqref{eq:v116-diag}; an unnormalized claim $D_{11}=1$ would be false in general.

\paragraph{Counting decision.}
The exact endpoint and its correctly normalized reverse-matrix collapse are proved, but this alone is not evidence that the ordered transition Al-Salam--Chihara $\to$ continuous $q$-Laguerre appears as a \emph{separate} K--S limiting arrow. V31 itself describes the connection as a specialization of the Al-Salam--Chihara parameter-change construction, and does not give a numbered K--S arrow for this ordered pair. Therefore the status remains \textbf{mathematically verified; independent K--S eligibility unconfirmed}. This is not a rejection of exact specializations generally. No new PASS or duplicate is claimed.

\noindent\textbf{V116 verdict:} Endpoint reverse coefficients established with the exact diagonal normalization \eqref{eq:v116-diag}. Independent K--S arrow not certified. \textbf{Recorded PASS total remains 52.} The V31 manuscript is unchanged.


\clearpage
\section{V117: Full 52-entry ledger reconciliation and duplicate-screening certificate}
\noindent\textbf{Objective.} Previous checks counted 43 historical PASS entries (V96) and awarded nine additional entries (V100, V101, V106--V109). This stage reconciles these two blocks and tests the \emph{ordered family-pair keys} for duplicate entries. It is an accounting/source-traceability test, not an independent proof of all 52 limits or a new inspection of the original K--S publication.

\paragraph{Reconciliation.} The historical V96 ledger has 43 entries. The post-V99 additions are exactly the following nine, with their unchanged V31 source-equation labels:
\begin{center}\small
\begin{longtable}{@{}r p{0.48\linewidth} p{0.30\linewidth}@{}}
\toprule No. & Ordered source $\to$ target & V31 equation / award stage\\
\midrule\endfirsthead
\toprule No. & Ordered source $\to$ target & V31 equation / award stage\\
\midrule\endhead
\bottomrule\endfoot
44 & q-Meixner $\to$ ordinary q-Charlier & \texttt{eq:qC-qM-reverse}; V100 \\
45 & dual q-Hahn $\to$ affine q-Krawtchouk & \texttt{eq:finite-affdH-reverse}; V101 \\
46 & q-Hahn $\to$ affine q-Krawtchouk & \texttt{eq:finite-affH-reverse}; V106 (V102) \\
47 & q-Racah $\to$ dual q-Hahn & \texttt{eq:finite-dHR-reverse}; V106 (V103) \\
48 & q-Racah $\to$ ordinary q-Krawtchouk & \texttt{eq:finite-KR-reverse}; V106 (V104) \\
49 & q-Racah $\to$ q-Hahn & \texttt{eq:finite-HR-double-D}; V106 (V105) \\
50 & continuous q-Jacobi $\to$ continuous q-Laguerre & \texttt{eq:lag-jacobi-new-rev}; V107 \\
51 & Al-Salam--Carlitz I $\to$ discrete q-Hermite I & \texttt{eq:dHI-carlitz-rev}; V108 \\
52 & Al-Salam--Carlitz II $\to$ discrete q-Hermite II & \texttt{eq:dHII-carlitz-rev}; V109 \\
\end{longtable}\end{center}
\paragraph{Mechanical checks.} The nine newly listed V31 labels all occur literally as \texttt{\textbackslash label\{...\}} in the unchanged manuscript. Comparing the nine ordered source--target pairs against each other and against the 43 V96 ledger pairs gives \textbf{52 distinct ordered family-pair strings}. The 43 original entries were already unique as recorded in V96; none of the nine new entries repeats one of those keys. In particular, the continuous $q$-Hahn $\to$ $q$-Meixner--Pollaczek case rechecked in V110 is already historical entry 22, so V110 correctly added zero.

\paragraph{Scope limitation and caution.} Uniqueness of family-pair strings is weaker than uniqueness of a published K--S limiting arrow: alternative parameterizations can encode the same arrow, and source/target naming conventions may conceal equivalences. Conversely, two distinct paths between the same ordered families are intentionally counted at most once here. This check establishes \emph{no} independent validity of the nine published equation-number attributions and \emph{no} asymptotic or normalization recertification of the historical 43. The local file named \texttt{kls\_as98\_reference.pdf} is zero bytes in the present workspace and therefore cannot serve as an independently inspected primary source for this stage; prior external-source claims must be assessed from a valid copy before final certification.

\noindent\textbf{V117 verdict:} \textbf{52 recorded entries reconciled as 43+9; no duplicate ordered family pair detected; nine new V31 labels exist.} Newly awarded PASS: \textbf{0}. Recorded PASS total remains \textbf{52}. Primary-source and global mathematical certification remain open. The original V31 manuscript is unchanged.


\clearpage
\section{V118: Parameter-dependent reverse-coefficient audit of PASS 44}
\noindent\textbf{Objective.} Recheck the recorded $q$-Meixner $\to$ ordinary $q$-Charlier specialization (PASS 44, V100) at \emph{nonzero} $b$, including the first off-diagonal and the correct diagonal normalization. This tests the actual V31 reverse matrix beyond its endpoint value; it does not create a new limiting arrow or certify a K--S equation number.

\paragraph{Exact coefficients from the unchanged V31 source.}
Use V31 equation \texttt{eq:qC-qM-reverse}, with $c=a\ne0$ and $0<q<1$:
\begin{equation}\label{eq:v118-D}
 D_{nk}(b)={n\brack k}_q\frac{1}{(bq;q)_k}
 \sum_{r=0}^{n-k}{n-k\brack r}_q
 \frac{(-1)^r q^{\binom r2}}{(bq^{k+1};q)_r}.
\end{equation}
For each fixed $n$, choose $b$ sufficiently close to zero that all denominators are nonzero. The expression is then analytic in $b$ and has a coefficientwise limit.

\paragraph{Diagonal and first subdiagonal, with no endpoint substitution.}
Setting $k=n$ in \eqref{eq:v118-D} leaves only $r=0$, so
\begin{equation}\label{eq:v118-diag}
 \boxed{D_{nn}(b)=\frac{1}{(bq;q)_n}.}
\end{equation}
For $k=n-1$ the inner sum has precisely two terms:
\[
 1-\frac{1}{1-bq^n}=-\frac{bq^n}{1-bq^n}.
\]
Consequently,
\begin{equation}\label{eq:v118-subdiag}
 \boxed{D_{n,n-1}(b)=-{n\brack1}_q\frac{bq^n}{(bq;q)_n}.}
\end{equation}
These exact formulas show both that $D_{nn}(b)$ is \emph{not} identically one away from the endpoint and that $D_{n,n-1}(b)$ vanishes linearly as $b\to0$.

\paragraph{All lower entries and rate.}
At $b=0$ the finite sum in \eqref{eq:v118-D} is $(1;q)_{n-k}$ by the finite $q$-binomial theorem. Thus $D_{nk}(0)=\delta_{nk}$. Since the coefficients are analytic in $b$ near zero, for each fixed $n$ and each $k<n$,
\[
 D_{nk}(b)=O(b),\qquad
 D_{nn}(b)=1+b\sum_{j=1}^n q^j+O(b^2).
\]
Here the constants in the $O(b)$ statements may depend on $n,k,q$; no uniform-in-degree estimate is claimed. The finite expansion therefore gives $M_n(x;b,a;q)\to C_n(x;a;q)$ coefficientwise and locally uniformly in $x$ for each fixed $n$.

\paragraph{Degree-one independent consistency check.}
The two exact coefficients are
\[
 D_{11}(b)=\frac{1}{1-bq},\qquad
 D_{10}(b)=-\frac{bq}{1-bq}.
\]
Using the terminating definitions in V31 gives
\[
 C_1(x;a;q)=1+\frac{q}{a}(1-x),\qquad
 M_1(x;b,a;q)=1+\frac{q}{a(1-bq)}(1-x).
\]
Hence $M_1=D_{11}C_1+D_{10}C_0$ exactly, including the constant term. This provides a non-endpoint check of the V31 reverse coefficients and their normalization.

\paragraph{Certification boundary.}
The mathematical endpoint and its local analytic stability are verified directly from the unchanged V31 reverse formula. The reference to Koekoek--Swarttouw Eq.~(4.13.2) remains an attribution from V31, not a fresh primary-source inspection: the local file \texttt{kls\_as98\_reference.pdf} has zero bytes. Thus this stage neither upgrades the entire 52-entry ledger to externally certified status nor adds an entry.

\noindent\textbf{V118 verdict:} PASS 44 retained with an independent non-endpoint diagonal, subdiagonal, and degree-one check. \textbf{New PASS: 0; recorded total: 52.} External K--S and global completeness certification remain pending. V31 is unchanged.



\clearpage
\section{V119: Non-endpoint finite-grid verification of PASS 45}
\noindent\textbf{Objective.} Independently check the dual $q$-Hahn $\to$ affine $q$-Krawtchouk specialization (PASS 45) \emph{away from} $\delta=0$, using the exact degree-one polynomial and the actual finite-grid interpolation matrix. The result distinguishes the generally nontriangular grid connection from its identity endpoint. No new arrow is counted.

\paragraph{Setup and nondegeneracy.} Fix $0<q<1$, $N\ge1$, $p$ with $A=(1-pq)(1-q^{-N})\ne0$, and put $t=p\delta q$. Let $y_x=q^{-x}$ and $E_{xk}=K_k^{\mathrm{aff}}(y_x;p,N\mid q)$ for $0\le x,k\le N$. Assume the full affine $q$-Krawtchouk basis is nondegenerate, so $E$ is invertible. The source is unchanged V31 equation \texttt{eq:finite-affdH-reverse}.

\paragraph{An exact degree-one identity at arbitrary $t$.} The terminating ${}_3\phi_2$ definitions in V31 yield
\begin{align*}
 K_1^{\mathrm{aff}}(y;p,N\mid q)&=1-\frac{1-y}{A},\\
 R_1^{\mathrm{dH}}(y+t/y;p,\delta,N\mid q)&=1-\frac{(1-y)(1-t/y)}{A}.
\end{align*}
Thus on the grid, without an asymptotic expansion,
\begin{equation}\label{eq:v119-exact-difference}
 \boxed{R_1^{\mathrm{dH}}(y_x+t/y_x;p,\delta,N\mid q)
 =K_1^{\mathrm{aff}}(y_x;p,N\mid q)-\frac{t}{A}(1-q^x).}
\end{equation}
This proves a linear-in-$t$ deviation from the specialization and makes the normalization transparent.

\paragraph{Exact reverse-coefficient row.} Define $v_x=1-q^x$ and the unique interpolation coefficients $c_k=\sum_{x=0}^{N}(E^{-1})_{kx}v_x$. Equation \eqref{eq:v119-exact-difference} and uniqueness of the finite-grid basis expansion imply the \emph{actual} V31 reverse coefficients in row $n=1$ are
\begin{equation}\label{eq:v119-D-row}
 \boxed{D_{1k}(\delta;p)=\delta_{1k}-\frac{p\delta q}{A}\sum_{x=0}^{N}(E^{-1})_{kx}(1-q^x),\qquad 0\le k\le N.}
\end{equation}
In particular $D_{1k}(0;p)=\delta_{1k}$ and $D_{1k}(\delta;p)-\delta_{1k}=O(\delta)$ for fixed $q,p,N$; the formula is exact for all nonsingular $\delta$.

\paragraph{Why a full-grid connection is necessary.} For $N\ge2$ and $t\ne0$, the correction $1-1/y$ in \eqref{eq:v119-exact-difference} cannot be represented by a polynomial of degree at most one at all $N+1$ distinct nonzero nodes $y_x$. Indeed, if $1-1/y=a+by$ at three distinct nodes, multiplying by $y$ would give a quadratic polynomial with three distinct roots, hence identically zero, contradicting its constant term $-1$. Therefore at least one coefficient with $k\ge2$ is nonzero in row $n=1$ whenever $t\ne0$ (assuming the stated basis nondegeneracy). This is a concrete diagnostic that the general V31 reverse connection is \emph{not triangular}; treating it as triangular would be an error.

\paragraph{All-degree endpoint rate.} For any fixed $0\le n\le N$, the terminating dual $q$-Hahn evaluation at $y_x+t/y_x$ is a polynomial in $t$ (of degree at most $n$) with denominators independent of $t$. At $t=0$ it equals $E_{xn}$. Multiplication by the fixed inverse $E^{-1}$ gives, for all $k$,
\[
 D_{nk}(\delta;p)=\delta_{nk}+O(\delta)\qquad(\delta\to0),
\]
with constants allowed to depend on $n,k,N,q,p$. This is a direct finite-dimensional analytic proof for the entire reverse matrix; it does not assert any uniform $N\to\infty$ estimate.

\paragraph{Certification boundary.} V119 independently verifies the parameter-dependent degree-one reverse row, its nontriangular structure away from the endpoint, and all-degree fixed-grid $O(\delta)$ convergence. The V31 citation of Koekoek--Swarttouw Eq.~(4.7.1) is \emph{not} independently revalidated here. This does not change the historical ledger or settle global completeness.

\noindent\textbf{V119 verdict:} PASS 45 retained with an exact non-endpoint reverse-row check and all-degree endpoint-rate proof. \textbf{New PASS: 0; recorded total: 52.} The original V31 manuscript is unchanged.



\clearpage
\section{V120: Non-endpoint coefficient audit of PASS 46 ($q$-Hahn to affine $q$-Krawtchouk)}
\noindent\textbf{Objective.} Recheck the previously awarded PASS 46 directly from the unchanged V31 triangular reverse connection \texttt{eq:finite-affH-reverse}, keeping the parameter $B$ nonzero. This provides an independent normalization and first-degree check; it does not award a new limiting arrow or independently certify the published K--S equation number.

\paragraph{Setup.} Fix $0<q<1$, $N\ge1$, $p$ such that $(pq;q)_N\ne0$, and use the common lattice $y=q^{-x}$, $0\le x\le N$. Write $Q_n(y;p,B,N\mid q)$ for $q$-Hahn and $K_n^{\mathrm{aff}}(y;p,N\mid q)=Q_n(y;p,0,N\mid q)$ for affine $q$-Krawtchouk. All statements concern fixed finite $N$ and degrees $0\le n\le N$.

\paragraph{Exact diagonal at nonzero $B$.} In V31's displayed reverse coefficient formula, put $A=p$ and $k=n$. The terminating ${}_3\phi_2$ has only its zeroth term. Cancelling the identical $(q^{-n};q)_n$ and $(pq;q)_n$ factors in the prefactor yields
\begin{equation}\label{eq:v120-diagonal}
\boxed{D_{nn}(B)=(pBq^{n+1};q)_n.}
\end{equation}
In particular, $D_{nn}(B)$ is generally \emph{not} identically one away from $B=0$, but $D_{nn}(0)=1$. For $n=0$, the empty product is one.

\paragraph{Exact degree-one row.} The terminating defining series gives
\[
K_1^{\mathrm{aff}}(y;p,N\mid q)=1-\frac{1-y}{(1-pq)(1-q^{-N})},
\]
\[
Q_1(y;p,B,N\mid q)=1-\frac{(1-pBq^2)(1-y)}{(1-pq)(1-q^{-N})}.
\]
Hence, \emph{for arbitrary admissible $B$},
\begin{equation}\label{eq:v120-degree-one}
\boxed{Q_1(y;p,B,N\mid q)=(1-pBq^2)K_1^{\mathrm{aff}}(y;p,N\mid q)+pBq^2K_0^{\mathrm{aff}}(y;p,N\mid q).}
\end{equation}
Thus $D_{11}(B)=1-pBq^2$ and $D_{10}(B)=pBq^2$, in agreement with \eqref{eq:v120-diagonal}. This is a direct non-endpoint check of both the diagonal and constant coefficient.

\paragraph{All-degree coefficientwise endpoint.} For each fixed $n$, the defining terminating $q$-Hahn series is analytic in $B$ near zero because the denominators $(pq;q)_j$ and $(q^{-N};q)_j$ are independent of $B$ and nonzero for $j\le n$. At $B=0$ it equals the affine $q$-Krawtchouk defining series term by term. The latter polynomials have nonzero leading coefficients under the stated assumptions and form a degree-ordered basis. Therefore their unique triangular expansion coefficients satisfy
\begin{equation}\label{eq:v120-endpoint}
\boxed{D_{nk}(0)=\delta_{nk},\qquad D_{nk}(B)=\delta_{nk}+O(B)\quad(0\le k\le n\le N).}
\end{equation}
The $O(B)$ constants may depend on $n,k,N,p,q$. This conclusion applies to the \emph{same} reverse coefficients printed in V31, by uniqueness of the basis expansion. It is a triangular polynomial connection, in contrast to the generally nontriangular finite-grid dual $q$-Hahn connection discussed in V119.

\paragraph{Certification boundary.} The mathematical endpoint and local parameter dependence of PASS 46 are independently checked. Its earlier attribution to Koekoek--Swarttouw Eq.~(4.6.5) is not freshly authenticated from a readable primary-source PDF in this stage. No assertion about uniformity in $N$ or about global completeness is made.

\noindent\textbf{V120 verdict:} PASS 46 retained with an exact non-endpoint diagonal and degree-one reverse-row verification and an all-degree analytic endpoint argument. \textbf{New PASS: 0; recorded total: 52.} V31 remains unchanged.



\clearpage
\section{V121: Non-endpoint reverse-coefficient audit of PASS 47 ($q$-Racah to dual $q$-Hahn)}
\noindent\textbf{Objective.} Independently check the actual V31 reverse-connection coefficients in \texttt{eq:finite-dHR-reverse} at nonzero $\beta$, including an exact diagonal and the complete degree-one row. This tests the normalization of the previously awarded PASS 47 without counting a new arrow.

\paragraph{Setup.} Fix $0<q<1$, $N\ge1$, and $0\le n\le N$. The lattice $\mu(x)=q^{-x}+\gamma\delta q^{x+1}$ is common to both families. Assume $(\gamma q;q)_N\ne0$, and restrict $\beta$ to a neighborhood of zero in which $(\beta\delta q;q)_N\ne0$ and the relevant polynomial leading coefficients do not vanish. Write
\[
 R_n(\mu(x);\beta,\gamma,\delta,N\mid q)
 =\sum_{k=0}^n D_{nk}(\beta)R_k^{\mathrm{dH}}(\mu(x);\gamma,\delta,N\mid q).
\]

\paragraph{Exact diagonal.} In V31's coefficient formula, the terminating ${}_2\phi_1$ is $1$ when $k=n$. Using the printed Newton coefficients $T^R$ and $T^{\mathrm{dH}}$ gives
\begin{equation}\label{eq:v121-diagonal}
\boxed{D_{nn}(\beta)=\frac{(\beta q^{n-N};q)_n}{(\beta\delta q;q)_n}.}
\end{equation}
The diagonal is generally not identically $1$ away from the endpoint, but $D_{nn}(0)=1$. This explicitly tests both the numerator parameter and the denominator normalization.

\paragraph{Complete degree-one row.} Put
\[
 A=\frac{1-\beta q^{1-N}}{1-\beta\delta q}.
\]
For $n=k=1$, \eqref{eq:v121-diagonal} gives $D_{11}=A$. For $n=1,k=0$, the Newton-coefficient prefactor is $1$ and the terminating hypergeometric sum has two terms:
\begin{align*}
 D_{10}(\beta)
 &={}_2\phi_1\!\left(\begin{matrix}q^{-1},\beta q^{1-N}\\\beta\delta q\end{matrix};q,q\right)\\
 &=1+\frac{(1-q^{-1})(1-\beta q^{1-N})q}{(1-\beta\delta q)(1-q)}\\
 &=1-A=\frac{\beta(q^{1-N}-\delta q)}{1-\beta\delta q}.
\end{align*}
Consequently, for every admissible nonzero $\beta$,
\begin{equation}\label{eq:v121-degree-one}
\boxed{R_1(\mu(x);\beta,\gamma,\delta,N\mid q)
 =A\,R_1^{\mathrm{dH}}(\mu(x);\gamma,\delta,N\mid q)+(1-A).}
\end{equation}
This agrees independently with the terminating defining series: the coefficient of $(1-q^{-x})(1-\gamma\delta q^{x+1})$ in $R_1$ is exactly $A$ times that in $R_1^{\mathrm{dH}}$.

\paragraph{All-degree analytic endpoint.} For $k<n$, V31's terminating ${}_2\phi_1$ at $\beta=0$ is
\[
 \sum_{j=0}^{n-k}\frac{(q^{k-n};q)_j}{(q;q)_j}q^j
 =(q^{1-(n-k)};q)_{n-k}=0.
\]
The Newton-coefficient ratio is finite under the stated nondegeneracy assumptions. For $k=n$, \eqref{eq:v121-diagonal} tends to $1$. Since every coefficient is a rational function of $\beta$ regular near $0$, for fixed $n,k,N$,
\begin{equation}\label{eq:v121-endpoint}
\boxed{D_{nk}(\beta)=\delta_{nk}+O(\beta)\quad(0\le k\le n\le N).}
\end{equation}
This verifies the actual triangular reverse-connection matrix, not merely the equality of the endpoint polynomial definitions. No uniform-in-$N$ assertion is made.

\paragraph{Certification boundary.} The source equation and its parameter-dependent coefficients have been checked internally. The previously cited Koekoek--Swarttouw equation (4.2.5) is \emph{not} freshly authenticated from a readable primary-source copy in V121. This stage does not settle the global completeness of the ledger.

\noindent\textbf{V121 verdict:} PASS 47 retained after exact non-endpoint diagonal and degree-one coefficient checks and an all-degree local endpoint proof. \textbf{New PASS: 0; recorded total: 52.} The original V31 manuscript is unchanged.



\clearpage
\section{V122: Non-endpoint reverse-coefficient audit of PASS 48 ($q$-Racah to ordinary $q$-Krawtchouk)}
\noindent\textbf{Objective.} Recheck the previously counted PASS 48 using the actual finite-grid reverse connection \texttt{eq:finite-KR-reverse} of V31, including its exact diagonal and full degree-one row away from the matching specialization. No new limiting arrow is claimed.

\paragraph{Parameter path and nondegeneracy.} Fix $0<q<1$, $N\ge1$, and an admissible $p$ such that the ordinary $q$-Krawtchouk evaluation matrix is invertible. Set $\gamma=\delta=0$ and vary $\beta$ near $\beta_0=-pq^N$. The $q$-Racah lattice then reduces to $y=q^{-x}$. Assume the leading Newton coefficients of both families are nonzero in this neighborhood. The V31 reverse connection has the form
\[
 R_n(y;\beta,0,0,N\mid q)=\sum_{k=0}^{N}D_{nk}(\beta;p)K_k(y;p,N\mid q),\qquad y=q^{-x},\quad 0\le x\le N.
\]
Since $\gamma\delta=0$, both families are degree-$n$ polynomials in $y$, so uniqueness of the polynomial basis forces $D_{nk}=0$ for $k>n$ (even though the general V31 grid-interpolation formula is a full $N+1$ by $N+1$ matrix).

\paragraph{Exact diagonal away from the endpoint.} Write $P_j(y)=(y;q)_j$. Directly from V31's defining terminating series \texttt{eq:finite-racah-definition} and \texttt{eq:finite-q-krawtchouk-definition},
\begin{align*}
 R_n(y;\beta,0,0,N\mid q)
 &=\sum_{j=0}^n\frac{(q^{-n},\beta q^{n-N};q)_j q^j}{(q,q^{-N};q)_j}P_j(y),\\
 K_n(y;p,N\mid q)
 &=\sum_{j=0}^n\frac{(q^{-n},-pq^n;q)_j q^j}{(q,q^{-N};q)_j}P_j(y).
\end{align*}
Here the zero numerator and denominator parameters have $q$-Pochhammer factors $(0;q)_j=1$. Comparing the top Newton coefficient, which also determines the leading coefficient in $y$, gives the exact reverse-connection diagonal
\begin{equation}\label{eq:v122-diagonal}
 \boxed{D_{nn}(\beta;p)=\frac{(\beta q^{n-N};q)_n}{(-pq^n;q)_n}.}
\end{equation}
At $\beta=\beta_0=-pq^N$ the numerator and denominator agree, hence $D_{nn}(\beta_0;p)=1$. In particular the diagonal is generally \emph{not} identically one at nonmatching parameters.

\paragraph{Exact complete degree-one row.} For $n=1$ the terminating series have only two terms. Put
\[
 A(\beta)=\frac{1-\beta q^{1-N}}{1+pq}.
\]
Then, for all $y$,
\begin{equation}\label{eq:v122-degree-one}
 \boxed{R_1(y;\beta,0,0,N\mid q)=A(\beta)K_1(y;p,N\mid q)+\bigl(1-A(\beta)\bigr).}
\end{equation}
Thus the \emph{actual} reverse coefficients, by uniqueness of V31's finite-grid expansion, are
\begin{equation}\label{eq:v122-row}
 \boxed{D_{11}(\beta;p)=\frac{1-\beta q^{1-N}}{1+pq},\qquad
 D_{10}(\beta;p)=\frac{\beta q^{1-N}+pq}{1+pq}
 =\frac{q^{1-N}(\beta-\beta_0)}{1+pq}.}
\end{equation}
Both are exact at nonzero displacement $\beta-\beta_0$, and tend to $1$ and $0$, respectively. This tests the sign and the $q^N$ scaling in the specialization.

\paragraph{All-degree local endpoint and the interpolation matrix.} At $\beta=\beta_0$, the two terminating hypergeometric series above agree term by term for every $0\le n\le N$; in particular, $R_n(y;\beta_0,0,0,N\mid q)=K_n(y;p,N\mid q)$ as polynomials, not only on the grid. Let $E_{xk}=K_k(q^{-x};p,N\mid q)$ and $r_n(\beta)_x=R_n(q^{-x};\beta,0,0,N\mid q)$. The explicit V31 reverse coefficients satisfy $D_n(\beta)=E^{-1}r_n(\beta)$ in column-vector notation. The matrix $E$ is independent of $\beta$ and invertible by assumption; each $r_n(\beta)$ is polynomial in $\beta$ for fixed $n$. Hence
\begin{equation}\label{eq:v122-endpoint}
 \boxed{D_{nk}(\beta;p)=\delta_{nk}+O(\beta-\beta_0),\qquad 0\le n,k\le N,\quad \beta\to\beta_0.}
\end{equation}
This proves the collapse of the \emph{same full-grid reverse-connection coefficients} displayed in V31, not merely an informal substitution into the family names. The error constants can depend on $N,n,k,p,q$; no uniform-in-$N$ claim is made.

\paragraph{Certification boundary.} This stage checks the mathematical endpoint and its non-endpoint coefficient dependence. The previously reported Koekoek--Swarttouw equation (4.2.8) has \emph{not} been freshly authenticated against a readable primary-source copy in V122; that external bibliographic check and global completeness remain open.

\noindent\textbf{V122 verdict:} PASS 48 retained, with exact non-endpoint diagonal and degree-one reverse coefficients and a full-grid all-degree analytic endpoint proof. \textbf{New PASS: 0; recorded total: 52.} The original V31 manuscript is unchanged.



\clearpage
\section{V123: Independent non-endpoint verification of PASS 49 and completion-gate assessment}
\noindent\textbf{Scope.} Recheck the previously counted $q$-Racah $\to$ $q$-Hahn arrow (PASS 49) from the reverse finite-grid connection \texttt{eq:finite-HR-double-D} of V31. Distinguish this local mathematical certification from a global certification of all 52 entries.

\paragraph{Parameters.} Fix $0<q<1$, $N\ge1$, and admissible $A,B$ such that the $q$-Hahn evaluation basis is nondegenerate. Set $\gamma=A$, $\delta=0$, and let the $q$-Racah parameter $\beta_R$ vary near $\beta_0=ABq^{N+1}$. Both lattices then equal $y=q^{-x}$. In particular, the source $q$-Racah polynomial is a degree-$n$ polynomial in $y$, so its unique $q$-Hahn expansion is triangular, even though V31's generic finite-grid connection has $N+1$ terms.

\paragraph{Exact Newton coefficients.} With $P_j(y)=(y;q)_j$, the defining terminating series yield
\begin{align*}
R_n(y;\beta_R,A,0,N\mid q)
 &=\sum_{j=0}^n \frac{(q^{-n},\beta_R q^{n-N};q)_j q^j}{(q,q^{-N},Aq;q)_j}P_j(y),\\
Q_n(y;A,B,N\mid q)
 &=\sum_{j=0}^n \frac{(q^{-n},ABq^{n+1};q)_j q^j}{(q,q^{-N},Aq;q)_j}P_j(y).
\end{align*}
The top Newton coefficients are nonzero under the admissibility assumptions. Therefore the \emph{actual} reverse-connection diagonal is
\begin{equation}\label{eq:v123-diagonal}
\boxed{D_{nn}(\beta_R;A,B)=\frac{(\beta_R q^{n-N};q)_n}{(ABq^{n+1};q)_n}.}
\end{equation}
This is not identically one at nonmatching parameters, but it is exactly one at $\beta_R=\beta_0$.

\paragraph{Complete degree-one row.} Define
\[
T(\beta_R)=\frac{1-\beta_R q^{1-N}}{1-ABq^2}.
\]
Both degree-one terminating series have constant value $1$ at $y=1$. Comparing their linear Newton coefficients gives the \emph{exact} non-endpoint expansion
\begin{equation}\label{eq:v123-degree-one}
\boxed{R_1(y;\beta_R,A,0,N\mid q)=T(\beta_R)Q_1(y;A,B,N\mid q)+1-T(\beta_R).}
\end{equation}
Thus
\[
\boxed{D_{11}=\frac{1-\beta_R q^{1-N}}{1-ABq^2},\qquad
D_{10}=\frac{\beta_R q^{1-N}-ABq^2}{1-ABq^2}
=\frac{q^{1-N}(\beta_R-\beta_0)}{1-ABq^2}.}
\]
The diagonal tends to $1$ and the lower coefficient to $0$. This also tests the critical exponent $N+1$ in the specialization.

\paragraph{Full reverse-matrix endpoint.} At $\beta_R=\beta_0$, the two Newton series agree term by term, since $\beta_0q^{n-N}=ABq^{n+1}$. Let $E_{xk}=Q_k(q^{-x};A,B,N\mid q)$ and $r_n(\beta_R)_x=R_n(q^{-x};\beta_R,A,0,N\mid q)$. V31's explicit reverse coefficients are uniquely determined by the finite-grid identity $E D_n=r_n$. The matrix $E$ is independent of $\beta_R$ and invertible, and $r_n$ is polynomial in $\beta_R$ for fixed $n$. Hence
\begin{equation}\label{eq:v123-full}
\boxed{D_{nk}(\beta_R;A,B)=\delta_{nk}+O(\beta_R-ABq^{N+1}),\quad 0\le n,k\le N.}
\end{equation}
In particular, the full-grid coefficients with $k>n$ vanish identically in this specialized common-variable situation. The error bound is for fixed $N,n,k,A,B,q$; no uniform-in-$N$ claim is made.

\noindent\textbf{Verdict on PASS 49: RETAIN.} This independently confirms the mathematical endpoint and the non-endpoint normalization for the reverse connection \texttt{eq:finite-HR-double-D}. The previously reported K--S equation (4.2.4) has \emph{not} been freshly authenticated against a readable primary-source copy in V123.

\paragraph{Completion-gate assessment.} The cumulative ledger records \textbf{52} distinct ordered-family PASS entries, and V117 checked the arithmetic and absence of duplicate ordered pairs. Those bookkeeping checks do not constitute a global theorem: (i) a readable authoritative K--S source has not been systematically cross-checked against every claimed arrow; (ii) the mathematical and normalization proof of every historical PASS has not been independently revalidated; (iii) the previously listed uncounted candidates have not all received final K--S eligibility dispositions; and (iv) the 72 distinct-family V31 headings have been inventoried, but that inventory alone does not establish completeness of all K--S arrows derivable from the reverse connections. Therefore \emph{final global certification cannot responsibly be issued on the available evidence}. The audit must not label 52 as a proven exhaustive total.

\noindent\textbf{V123 status:} PASS 49 mathematically rechecked; \textbf{new PASS: 0; recorded total: 52}. \textbf{Global completeness and independent primary-source certification: OPEN}. The original V31 manuscript remains unchanged.


\clearpage
\section{V124: Independent normalization and local-analytic audit of PASS 50}
\noindent\textbf{Target.} Recheck continuous $q$-Jacobi $\to$ continuous $q$-Laguerre (PASS 50) against the \emph{explicit diagonal statement} in the unchanged V31 proposition \texttt{prop:lag-jacobi-new}, not merely the endpoint equality invoked in V107. This stage certifies a local mathematical property only; it does not independently authenticate the reported Koekoek--Swarttouw equation (4.10.1).

\paragraph{Parameter identification.} Fix $0<q<1$, $n\ge0$, and a real nondegenerate $\alpha$. In V31 put
\[
 a=A=q^{\alpha/2+1/4},\qquad b=aq^{1/2},\qquad B=q^{\beta/2+1/4}.
\]
Thus $\beta\to+\infty$ corresponds to $B\to0$. Let $D_{nk}(B)$ be the coefficients in the \emph{reverse} V31 equation \texttt{eq:lag-jacobi-new-rev}, namely
\[
 J_n(x;\alpha,\beta\mid q)=\sum_{k=0}^nD_{nk}(B)L_k^{(\alpha)}(x\mid q).
\]

\paragraph{Exact diagonal at nonzero $B$.} V31 states that the reverse diagonal is
\[
 \boxed{D_{nn}(B)=\frac{(q;q)_n}{a^n}(A^2B^2q^n;q)_n.}
\]
This expression is exact, not an asymptotic assertion. With $A=a$ and $B\to0$ it gives
\[
 D_{nn}(B)=\frac{(q;q)_n}{a^n}\left[1-a^2B^2q^n\frac{1-q^n}{1-q}+O(B^4)\right]
 \quad(n\ge1),
\]
where the first-order coefficient follows by expanding the finite product
$\prod_{j=0}^{n-1}(1-a^2B^2q^{n+j})$; the sum is $q^n(1-q^n)/(1-q)$. For $n=0$, the diagonal is identically $1$. This calculation verifies both the nontrivial $n$-dependent normalization and the absence of a linear-in-$B$ term on the diagonal.

\paragraph{Normalization to the published families.} In V107's normalization, $P_n^{(\alpha,\beta)}=a^nJ_n/(q;q)_n$ and $P_n^{(\alpha)}=L_n^{(\alpha)}$. The \emph{normalized} reverse diagonal therefore satisfies the stronger exact identity
\begin{equation}\label{eq:v124-normalized-diagonal}
 \boxed{\widehat D_{nn}(B):=\frac{a^n}{(q;q)_n}D_{nn}(B)=(a^2B^2q^n;q)_n=1+O(B^2).}
\end{equation}
An assertion that the \emph{unnormalized} reverse diagonal tends to $1$ would be incorrect unless $(q;q)_n/a^n=1$ accidentally.

\paragraph{Entire finite reverse matrix.} The terminating Askey--Wilson definition of $J_n$ has polynomial dependence on its two parameters $-B$ and $-Bq^{1/2}$ in a neighborhood of $B=0$, provided the fixed denominators are nonzero. The target Laguerre basis is independent of $B$ and has nonzero leading coefficients. Therefore the triangular change-of-basis coefficients $D_{nk}(B)$ are analytic near $0$. The exact specialization $J_n(0)=(q;q)_nL_n^{(\alpha)}/a^n$, previously established in V70 and V107, implies by uniqueness of basis expansion
\[
 D_{nk}(0)=\frac{(q;q)_n}{a^n}\delta_{nk}.
\]
Consequently, with $\widehat D_{nk}(B)=a^nD_{nk}(B)/(q;q)_n$,
\begin{equation}\label{eq:v124-matrix}
 \boxed{\widehat D_{nk}(B)=\delta_{nk}+O(B),\qquad 0\le k\le n,\quad B\to0.}
\end{equation}
The $O(B^2)$ improvement in \eqref{eq:v124-normalized-diagonal} applies to the diagonal only; it is \emph{not} asserted for every off-diagonal entry. All estimates are for fixed degree and fixed nondegenerate parameters.

\paragraph{Certification boundary and disposition.} This independent check uses the exact diagonal printed in V31, the source-to-KLS normalization established in V107, and a finite-dimensional analytic change-of-basis argument. It finds no normalization inconsistency in PASS 50. The external bibliographic identity of the claimed K--S arrow has not been freshly checked from a readable primary-source document here. The global completeness gate remains open.

\noindent\textbf{V124 verdict: PASS 50 retained. Newly awarded PASS: 0. Recorded total: 52.} The original V31 manuscript is unchanged.



\clearpage
\section{V125: Non-endpoint matrix certification of PASS 51 and exact error factorization}
\noindent\textbf{Target.} Recheck Al-Salam--Carlitz I $\to$ discrete $q$-Hermite I (PASS 51), using the actual V31 reverse formula \texttt{eq:dHI-carlitz-rev}. Unlike an endpoint-only argument, this stage determines the complete parameter-dependent reverse matrix and its first-order error at $a=-1$. No new arrow is awarded.

\paragraph{Exact coefficients for arbitrary $a$.} For fixed $0<q<1$, V31 states
\begin{equation}\label{eq:v125-matrix}
 U_n^{(a)}(x;q)=\sum_{k=0}^n D_{nk}(a;q)h_k(x;q),\qquad
 \boxed{D_{nk}(a;q)=(-1)^{n-k}{n\brack k}_q(-a;q)_{n-k}.}
\end{equation}
Set $m=n-k$. Since $(-a;q)_m=\prod_{j=0}^{m-1}(1+aq^j)$, for $m\ge1$ we have the \emph{exact} factorization
\begin{equation}\label{eq:v125-factor}
 \boxed{D_{nk}(a;q)=(1+a)(-1)^m{n\brack k}_q
 \prod_{j=1}^{m-1}(1+aq^j),\quad k<n.}
\end{equation}
The diagonal is $D_{nn}=1$ for every $a$, not just in the limit. Hence the full finite truncation $D^{[N]}(a)$ satisfies
\[
 \boxed{D^{[N]}(a)=I_{N+1}+(1+a)G^{[N]}(a),}
\]
where $G^{[N]}(a)$ is a lower-strictly-triangular polynomial matrix. In particular the limit is entrywise and, for each fixed $N$, holds in any matrix norm.

\paragraph{Exact first-order term.} Differentiating the factorization at $a=-1$ yields, for $k<n$,
\begin{equation}\label{eq:v125-derivative}
 \boxed{\left.\frac{\partial D_{nk}}{\partial a}\right|_{a=-1}
 =(-1)^{n-k}{n\brack k}_q(q;q)_{n-k-1},}
\end{equation}
while $\partial_aD_{nn}=0$. Thus, for fixed $n,k$,
\[
 D_{nk}(a;q)=\delta_{nk}+(1+a)\,(-1)^{n-k}{n\brack k}_q(q;q)_{n-k-1}+O((1+a)^2)
 \quad(k<n),
\]
with the diagonal separately equal to $1$. In particular
\[
 \boxed{D_{n,n-1}(a;q)=-{n\brack1}_q(1+a),\qquad
 D_{n,n-2}(a;q)={n\brack2}_q(1+a)(1+aq).}
\]
The second subdiagonal has slope ${n\brack2}_q(1-q)$ at $a=-1$; it is not generally of quadratic order in $1+a$.

\paragraph{Explicit low-degree matrix.} For degrees $0,1,2$, the actual reverse matrix is
\begin{equation}\label{eq:v125-small}
 \boxed{D^{[2]}(a)=\begin{pmatrix}
 1&0&0\\
 -(1+a)&1&0\\
 (1+a)(1+aq)&-(1+q)(1+a)&1
 \end{pmatrix}.}
\end{equation}
This agrees with the monic initial polynomial $U_1^{(a)}=x-(1+a)$ and gives an independent non-endpoint diagnostic. No complex rotation is needed for type I; the rotation in PASS 52 concerns type II only.

\paragraph{Inverse and conditioning near the endpoint.} V31 also gives the forward matrix $H(a)$ with
\[
 H_{nk}(a)={n\brack k}_q\prod_{j=0}^{n-k-1}(a+q^j).
\]
For $k<n$ its product contains $(a+1)$, and its diagonal is identically $1$. Hence $H^{[N]}(a)=I_{N+1}+O(1+a)$ for fixed $N$. The explicit V31 inverse identity $H^{[N]}(a)D^{[N]}(a)=I_{N+1}$ confirms the same endpoint from both directions, without assuming that a limit of inverse matrices exists in an infinite-dimensional topology.

\paragraph{Certification boundary.} The mathematical reverse-connection endpoint, its normalization, and its full finite-degree first-order behavior are independently checked. The claimed Koekoek--Swarttouw equation (4.24.1) is a previously recorded citation, \emph{not newly authenticated from a readable primary-source PDF in V125}. This local check does not resolve the other historical PASS entries, the uncounted K--S candidates, or global completeness.

\noindent\textbf{V125 verdict: PASS 51 retained. New PASS: 0. Recorded PASS total: 52. Global certification: OPEN.} The original V31 manuscript remains unchanged.


\clearpage
\section{V126: Independent differential-equation certification of historical PASS 43}
\noindent\textbf{Objective.} Audit the symmetric Jacobi $\to$ physicists' Hermite arrow recorded as historical PASS 43 in V71 by an independent method, with explicit control of every fixed-degree coefficient. The V31 reverse-connection source is \texttt{eq:v14-jac-herm}. This check tests the normalization and limit independently of the finite Gegenbauer resummation in V71; it does not purport to authenticate the external K--S equation number.

\paragraph{Scaled differential equation.} Let $t>0$ and, for fixed $n\ge0$, define
\[
 F_{n,t}(z)=\frac{2^n n!}{t^{n/2}}P_n^{(t,t)}\!\left(\frac{z}{\sqrt t}\right).
\]
The Jacobi differential equation in the symmetric case is
\[
 (1-x^2)y''(x)-(2t+2)xy'(x)+n(n+2t+1)y(x)=0.
\]
The change $x=z/\sqrt t$ and division by $t$ give the \emph{exact} equation
\begin{equation}\label{eq:v126-de}
 \boxed{\left(1-\frac{z^2}{t}\right)F_{n,t}''(z)
 -\left(2+\frac2t\right)zF_{n,t}'(z)
 +\left(2n+\frac{n(n+1)}t\right)F_{n,t}(z)=0.}
\end{equation}
Symmetry of $P_n^{(t,t)}$ implies $F_{n,t}(-z)=(-1)^nF_{n,t}(z)$.

\paragraph{Leading coefficient.} The coefficient of $x^n$ in the conventional Jacobi polynomial is $2^{-n}\binom{2n+2t}{n}$. Therefore the coefficient of $z^n$ in $F_{n,t}$ is exactly
\[
 c_n(t)=\frac{n!}{t^n}\binom{2n+2t}{n}
 =\prod_{j=0}^{n-1}\left(2+\frac{n+1+j}{t}\right)
 \longrightarrow 2^n.
\]
This step fixes the normalization without appealing to the Hermite limit itself.

\paragraph{Finite coefficient recurrence.} Write $F_{n,t}(z)=\sum_{j=0}^{\lfloor n/2\rfloor}c_{n-2j}(t)z^{n-2j}$. Equating the coefficient of $z^r$ in \eqref{eq:v126-de}, for $r=n-2,n-4,\ldots\ge0$, yields
\begin{equation}\label{eq:v126-rec}
 \boxed{c_r(t)=-\frac{(r+2)(r+1)}{(n-r)\bigl(2+(n+r+1)/t\bigr)}\,c_{r+2}(t).}
\end{equation}
For $t>0$ the denominator is nonzero. Since there are finitely many coefficients, induction from $c_n(t)\to2^n$ proves convergence of every coefficient. Iterating the limiting recurrence gives
\[
 \boxed{\lim_{t\to\infty}c_{n-2j}(t)
 =(-1)^j\frac{n!\,2^{n-2j}}{j!(n-2j)!}.}
\]
Thus the coefficientwise limit is exactly the physicists' Hermite polynomial,
\begin{equation}\label{eq:v126-limit}
 \boxed{\lim_{t\to\infty}\frac{2^n n!}{t^{n/2}}P_n^{(t,t)}\!\left(\frac z{\sqrt t}\right)
 =n!\sum_{j=0}^{\lfloor n/2\rfloor}\frac{(-1)^j(2z)^{n-2j}}{j!(n-2j)!}=H_n(z).}
\end{equation}
Because the degree is fixed, convergence is uniform for $z$ on compact subsets of $\mathbb C$.

\paragraph{Degree-one and degree-two diagnostics.} Directly,
\[
 F_{1,t}(z)=2\left(1+\frac1t\right)z\to2z=H_1(z),
\]
\[
 F_{2,t}(z)=\left(4+\frac{14}t+\frac{12}{t^2}\right)z^2-\left(2+\frac4t\right)
 \to4z^2-2=H_2(z).
\]
These are independent normalization checks, not substitutes for the all-degree recurrence.

\paragraph{Connection to the reverse-connection requirement.} V71 already establishes that resumming the finite V31 reverse identity \texttt{eq:v14-jac-herm} yields the same symmetric Jacobi polynomial before scaling. The present calculation independently certifies the \emph{limit and its normalization} using the Jacobi differential equation. It does not assert that the individual unscaled Hermite-basis coefficients have a termwise limit after substituting $z/\sqrt t$; cancellation among them is possible. The reverse-connection derivation remains V71, and the independent endpoint certification is now strengthened.

\paragraph{Verdict and global gate.} \textbf{Historical PASS 43 retained; no new PASS awarded.} The recorded total remains \textbf{52}. This independent mathematical check does not close the unresolved primary-source, remaining historical-proof, and candidate-eligibility gates. No exhaustive count is claimed. The original V31 manuscript remains unchanged.

\clearpage
\section{V127: Batch source-identity gate for all 52 recorded PASS entries}
\noindent\textbf{Purpose and standard.} In response to the request to finish the audit efficiently, this stage replaces serial one-proof-per-version work with a reproducible, whole-ledger gate. It checks the exact V31 source-label identity for \emph{every} recorded entry, distinguishes this mechanical result from proof and K--S certification, and identifies the remaining blocking issues. No mathematical PASS is awarded merely because a source label exists.

\paragraph{Reproducible extraction.} We parsed the 36 explicitly labeled entries in V96, the seven manually restored labels in V97, and the nine additions in V117, matching rows by ledger number. We then searched the \emph{unchanged} V31 \texttt{.tex} file for literal \texttt{\\label\{...\}} declarations. The test requires exactly one declaration for each ledger label; a zero or multiple occurrence would be a source-identity exception.

\begin{center}\small
\begin{longtable}{@{}r p{0.43\linewidth} p{0.39\linewidth}@{}}
\toprule PASS & Recorded transition & V31 equation label \\
\midrule\endfirsthead
\toprule PASS & Recorded transition & V31 equation label \\
\midrule\endhead
\bottomrule\endfoot
1 & Symmetric Jacobi $\ensuremath{\to}$ Hermite & \texttt{eq:v14-jac-herm} \\
2 & Normalized continuous $q$-Laguerre $\ensuremath{\to}$ continuous $q$-Hermite & \texttt{eq:qlag-hermite-fwd} \\
3 & Continuous $q$-Jacobi $\ensuremath{\to}$ continuous $q$-Hermite & \texttt{eq:JHforward} \\
4 & Hahn $\ensuremath{\to}$ Jacobi & \texttt{eq:jacobi-hahn-reverse} \\
5 & $q$-Hahn $\ensuremath{\to}$ little $q$-Jacobi & \texttt{eq:ljqh-reverse} \\
6 & $q$-Hahn $\ensuremath{\to}$ $q$-Meixner & \texttt{eq:168-D} \\
7 & Big $q$-Jacobi $\ensuremath{\to}$ $q$-Meixner & \texttt{eq:169-reverse} \\
8 & Big $q$-Jacobi $\ensuremath{\to}$ little $q$-Jacobi & \texttt{eq:ljbj166-reverse} \\
9 & Big $q$-Laguerre $\ensuremath{\to}$ little $q$-Laguerre (Wall) & \texttt{eq:wall-big-reverse} \\
10 & $q$-Hahn $\ensuremath{\to}$ quantum $q$-Krawtchouk & \texttt{eq:finite-qtmH-reverse} \\
11 & $q$-Hahn $\ensuremath{\to}$ ordinary $q$-Krawtchouk & \texttt{eq:finite-KH-reverse} \\
12 & Dual $q$-Hahn $\ensuremath{\to}$ dual $q$-Krawtchouk & \texttt{eq:finite-dKdH-reverse} \\
13 & Askey--Wilson $\ensuremath{\to}$ big $q$-Jacobi & \texttt{eq:aw-bigqj-reverse} \\
14 & Standard $q$-Krawtchouk $\ensuremath{\to}$ alternative $q$-Charlier ($q$-Bessel) & \texttt{eq:qBqK-reverse} \\
15 & Standard $q$-Krawtchouk $\ensuremath{\to}$ ordinary $q$-Charlier & \texttt{eq:qCqK-reverse} \\
16 & Little $q$-Jacobi $\ensuremath{\to}$ $q$-Bessel (alternative $q$-Charlier) & \texttt{eq:qBlqJ-reverse} \\
17 & $q$-Meixner $\ensuremath{\to}$ $q$-Laguerre & \texttt{eq:qLqM-D} \\
18 & Little $q$-Jacobi $\ensuremath{\to}$ $q$-Laguerre & \texttt{eq:qLlJ-D} \\
19 & Affine $q$-Krawtchouk $\ensuremath{\to}$ little $q$-Laguerre (Wall) & \texttt{eq:wall-affine-reverse} \\
20 & Little $q$-Jacobi $\ensuremath{\to}$ little $q$-Laguerre (Wall) & \texttt{eq:wall-lj-reverse} \\
21 & Big $q$-Jacobi $\ensuremath{\to}$ big $q$-Laguerre & \texttt{eq:blbj-reverse} \\
22 & Continuous $q$-Hahn $\ensuremath{\to}$ $q$-Meixner--Pollaczek & \texttt{eq:cqh-qmp23} \\
23 & Al-Salam--Chihara $\ensuremath{\to}$ continuous big $q$-Hermite & \texttt{eq:ASCBH} \\
24 & Continuous dual $q$-Hahn $\ensuremath{\to}$ Al-Salam--Chihara & \texttt{eq:CDASC} \\
25 & Continuous big $q$-Hermite $\ensuremath{\to}$ continuous $q$-Hermite & \texttt{eq:BHH} \\
26 & $q$-Laguerre $\ensuremath{\to}$ Stieltjes--Wigert & \texttt{eq:SWqL-rev} \\
27 & Quantum $q$-Krawtchouk $\ensuremath{\to}$ Al-Salam--Carlitz II & \texttt{eq:ASCII-qtm-reverse} \\
28 & $q$-Meixner $\ensuremath{\to}$ Al-Salam--Carlitz II & \texttt{eq:ASCII-qM-reverse} \\
29 & Askey--Wilson $\ensuremath{\to}$ continuous dual $q$-Hahn & \texttt{eq:AWCD} \\
30 & Racah $\ensuremath{\to}$ Hahn & \texttt{eq:racah-to-hahn-grid} \\
31 & Racah $\ensuremath{\to}$ dual Hahn & \texttt{eq:racah-dh-reverse} \\
32 & Krawtchouk $\ensuremath{\to}$ Hermite & \texttt{eq:v16-k-h} \\
33 & Hahn $\ensuremath{\to}$ Meixner & \texttt{eq:hahn-meixner-reverse} \\
34 & Hahn $\ensuremath{\to}$ Krawtchouk & \texttt{eq:hahn-kraw-reverse} \\
35 & Krawtchouk $\ensuremath{\to}$ Charlier & \texttt{eq:kraw-char-reverse} \\
36 & Meixner $\ensuremath{\to}$ Charlier & \texttt{eq:meix-char-reverse} \\
37 & Charlier $\to$ Hermite & \texttt{eq:v16-c-h} \\
38 & Wilson $\to$ continuous dual Hahn & \texttt{eq:consolidated-w-cdh} \\
39 & Meixner--Pollaczek $\to$ Hermite & \texttt{eq:cat-mp-to-hermite-bell} \\
40 & Continuous Hahn $\to$ Meixner--Pollaczek & \texttt{eq:consolidated-ch-mp} \\
41 & Meixner $\to$ generalized Laguerre & \texttt{eq:meixner-to-laguerre-v61} \\
42 & Jacobi $\to$ generalized Laguerre & \texttt{eq:v14-jac-lag} \\
43 & Wilson $\to$ Jacobi & \texttt{eq:wilson-to-jac} \\
44 & q-Meixner $\to$ ordinary q-Charlier & \texttt{eq:qC-qM-reverse} \\
45 & dual q-Hahn $\to$ affine q-Krawtchouk & \texttt{eq:finite-affdH-reverse} \\
46 & q-Hahn $\to$ affine q-Krawtchouk & \texttt{eq:finite-affH-reverse} \\
47 & q-Racah $\to$ dual q-Hahn & \texttt{eq:finite-dHR-reverse} \\
48 & q-Racah $\to$ ordinary q-Krawtchouk & \texttt{eq:finite-KR-reverse} \\
49 & q-Racah $\to$ q-Hahn & \texttt{eq:finite-HR-double-D} \\
50 & continuous q-Jacobi $\to$ continuous q-Laguerre & \texttt{eq:lag-jacobi-new-rev} \\
51 & Al-Salam--Carlitz I $\to$ discrete q-Hermite I & \texttt{eq:dHI-carlitz-rev} \\
52 & Al-Salam--Carlitz II $\to$ discrete q-Hermite II & \texttt{eq:dHII-carlitz-rev} \\
\end{longtable}\end{center}
\paragraph{Machine-check results.} The extracted ledger has \textbf{52 numbered entries} with no missing number in $1,\ldots,52$. All \textbf{52} corresponding V31 equation labels are present \textbf{exactly once} in the original source, and the 52 label strings are mutually distinct. Missing-label exceptions: \textbf{0}. Multiple-definition exceptions: \textbf{0}. This is a stronger, complete machine-readable source-identity gate than relying on the cumulative prose statements alone; it does \emph{not} prove that the formulas at those labels are mathematically correct.

\paragraph{Interpretation of the 52 count.} The ledger is \emph{provisionally recorded}, not independently certified as a complete K--S limiting scheme. Earlier V117 already checked duplicate ordered family pairs. Distinct equation labels do not exclude equivalent arrows under different parameterizations; conversely, a shared formula could support several legitimate transitions. This gate therefore neither adds to nor subtracts from the recorded 52.

\paragraph{Outstanding certification gates, grouped for completion.}
\begin{enumerate}
\item \textbf{Primary-source gate:} authenticate the claimed K--S arrows, equation numbers, parameter maps, and normalizations against a readable original reference. The local \texttt{kls\_as98\_reference.pdf} is empty, so no new primary-source authentication is asserted here. Particular historical checks already flagged are continuous Hahn $\to$ Meixner--Pollaczek and Wilson $\to$ Jacobi.
\item \textbf{Proof gate:} independently examine the remaining historical reverse-connection asymptotics in batches, prioritizing entries with unverified normalization, divergent or singular substitutions, or nontriangular finite-grid changes. A source-label match is not a substitute for this step.
\item \textbf{Candidate gate:} dispose of the eight historical K--S eligibility questions recorded in V96 and the Al-Salam--Chihara/continuous $q$-Laguerre candidate from V116. Exact specializations remain eligible if independently recognized by K--S; self-family connections and duplicates remain excluded.
\item \textbf{Exhaustiveness gate:} reconcile the decisions with all 72 distinct-family V31 headings, including multiple possible arrows per heading. The 72 headings are an inventory denominator, not a count of possible K--S arrows.
\end{enumerate}
\paragraph{Outcome and next batch.} \textbf{All 52 recorded entries pass the literal V31 source-identity test.} This does not close the independent K--S, remaining mathematical, candidate, or exhaustiveness gates. The next completion-focused batch should authenticate primary-source claims and resolve the outstanding candidate set rather than repeating the already certified Carlitz--Hermite computations. \textbf{New PASS: 0; recorded total: 52; final certification: OPEN.} V31 is unchanged.


\clearpage
\section{V128: Nine-candidate disposition docket and evidence thresholds}
\noindent\textbf{Scope.} This completion-oriented stage consolidates the eight unresolved candidate pairs from V96 and the additional Al-Salam--Chihara/continuous $q$-Laguerre pair from V116. The test is intentionally stricter than the existence of a finite connection: to award a new K--S arrow one must exhibit a distinct ordered family transition recognized in the original limiting scheme, a parameter/variable/normalization map, and the corresponding V31 reverse-matrix endpoint. In particular, a family relationship or finite expansion alone does not meet that test.

\paragraph{Candidate-by-candidate evidence matrix.}
\begingroup\small
\begin{longtable}{@{}p{0.045\linewidth}p{0.31\linewidth}p{0.31\linewidth}p{0.24\linewidth}@{}}
\toprule No. & Candidate ordered pair (orientation to be checked) & V31 evidence / current mathematical finding & Eligibility disposition\\
\midrule\endfirsthead
\toprule No. & Candidate & V31 evidence / finding & Disposition\\\midrule\endhead
\bottomrule\endfoot
1 & Big $q$-Laguerre / Al-Salam--Carlitz I & Section heading at original V31 line 8171; connection exists, but no independently authenticated K--S arrow here. & OPEN: original arrow and endpoint needed.\\
2 & Stieltjes--Wigert / alternative $q$-Charlier ($q$-Bessel) & V31 section \texttt{sec:SW-qBessel}, around line 8705; finite coefficients obtained through monomial projections. & OPEN: finite projection is not a limit proof.\\
3 & Stieltjes--Wigert / ordinary $q$-Charlier & V31 section \texttt{sec:SW-qCharlier}, around line 8807; source expressly distinguishes ordinary from alternative $q$-Charlier. & OPEN: separate K--S arrow not established.\\
4 & Jacobi / continuous Hahn & V31 heading around line 975; the cited Koelink Fourier-transform relation is explicitly distinguished from a polynomial-basis limiting arrow. & OPEN: transform relation does not certify an arrow.\\
5 & Meixner--Pollaczek / generalized Laguerre & V31 heading around line 1211. Earlier V96 records a genuine classical family limit but finds that the \emph{fixed-angle} reverse connection does not realize its required singular scaling. & BLOCKED for the displayed fixed-angle path; alternate compatible path not excluded.\\
6 & Meixner / dual Hahn & V31 heading around line 2027; no independent K--S map or coefficient endpoint certified. & OPEN: map and proof required.\\
7 & Pseudo-Jacobi / continuous Hahn & V31 heading around line 2505; existence of a finite connection alone is insufficient. & OPEN: K--S family status and limit required.\\
8 & Generalized Bessel / Jacobi & V31 heading around line 2659; no independently verified distinct K--S arrow. & OPEN: direction and limiting normalization required.\\
9 & Al-Salam--Chihara / continuous $q$-Laguerre & V116 proves the exact endpoint and nonunit diagonal normalization, but explicitly cannot authenticate a \emph{separate} K--S arrow. & MATHEMATICALLY VERIFIED; K--S ELIGIBILITY OPEN.\\
\end{longtable}\endgroup

\paragraph{Independent source-level cross-check.} The V31 descriptions of the Stieltjes--Wigert/$q$-Bessel and Stieltjes--Wigert/ordinary $q$-Charlier connections explicitly derive coefficients by composing monomial expansions with orthogonality projections; they do not assert a separately numbered K--S limit. These are two different target families, so combining them into one candidate would be an accounting error. The Jacobi/continuous Hahn source explicitly cites a Fourier-transform relation and distinguishes it from its finite polynomial-basis coefficients. None of these statements logically rules out a different K--S limiting path, but none proves one either.

\paragraph{What has and has not been closed.} All \textbf{nine} previously named candidates now have a single explicit evidence/disposition record. This closes the \emph{candidate inventory and triage} subtask, \textbf{not} the eligibility gate: eight require original-arrow evidence or a new endpoint analysis, while candidate 9 has an established endpoint but still lacks separate-arrow authentication. Candidate 5 is specifically obstructed for the fixed-angle path, not categorically rejected for every possible scaling. No candidate is silently promoted to PASS or permanently excluded without adequate evidence.

\paragraph{Required next evidence.} Obtain a readable original Koekoek--Swarttouw source and record for each candidate the exact numbered arrow (if present), source/target orientation, variable scaling, parameter substitutions, and normalization. Only then test the displayed V31 reverse coefficients along that same path. For the 52 recorded PASS entries, the same independent primary-source gate remains open; literal equation-label existence from V127 is not a substitute. The local \texttt{kls\_as98\_reference.pdf} is zero bytes and was not used as evidence.

\noindent\textbf{V128 verdict:} Nine-candidate triage documented; newly certified K--S arrows \textbf{0}; recorded PASS total \textbf{52}; global certification \textbf{OPEN}. Original V31 remains unchanged.

\clearpage
\section{V129: Big $q$-Laguerre / Al-Salam--Carlitz I: a failed naive endpoint}
\noindent\textbf{Purpose.} Resolve a concrete mathematical subquestion for candidate 1 of V128, using the \emph{actual reverse connection} \texttt{eq:ASCI-bL-reverse} in unchanged V31. This is an obstruction to one plausible endpoint, not a categorical verdict about other parameter or variable scalings and not an independently authenticated K--S arrow.

\paragraph{Exact diagonal at arbitrary parameters.}
V31 defines
\[
 P_n(x;A,B\mid q)=\sum_{j=0}^n
 \frac{(q^{-n};q)_j q^j}{(q,Aq,Bq;q)_j}(x;q)_j,
 \qquad
 \sum_{n\ge0}\frac{U_n^{(a)}(x;q)}{(q;q)_n}t^n
 =\frac{(t,at;q)_\infty}{(xt;q)_\infty}.
\]
The coefficient of $x^n$ in $P_n$ is $(Aq,Bq;q)_n^{-1}$: only $j=n$ contributes, and
\((q^{-n};q)_n=(-1)^n q^{-n(n+1)/2}(q;q)_n\), while the leading coefficient of $(x;q)_n$ is $(-1)^nq^{n(n-1)/2}$. The generating function shows $U_n^{(a)}$ is monic. Thus the reverse coefficient of $U_n^{(a)}$ in $P_n$ is exactly
\begin{equation}\label{eq:v129-diagonal}
 \boxed{D_{nn}(A,B;a)=\frac{1}{(Aq,Bq;q)_n}.}
\end{equation}
In particular $D_{nn}\to1$ when $A,B\to0$, but this alone does not imply convergence of the entire matrix.

\paragraph{The apparent $A=B=0$ endpoint fails in degree two.}
The terminating series at $A=B=0$ gives $P_1(x;0,0\mid q)=x$ and $P_2(x;0,0\mid q)=x^2$ (indeed $P_n(x;0,0\mid q)=x^n$ by the terminating $q$-binomial identity). The generating function gives
\[
 U_1^{(a)}(x;q)=x-(1+a),\qquad
 U_2^{(-1)}(x;q)=x^2-(1-q).
\]
To have $P_1(x;0,0\mid q)=U_1^{(a)}(x;q)$, necessarily $a=-1$. Yet for $0<q<1$,
\[
 P_2(x;0,0\mid q)=x^2
 =U_2^{(-1)}(x;q)+(1-q)U_0^{(-1)}(x;q),
\]
so the exact reverse matrix has
\begin{equation}\label{eq:v129-obstruction}
 \boxed{D_{20}(0,0;-1)=1-q\ne0.}
\end{equation}
Equivalently, in degree one the reverse coefficient is $D_{10}(0,0;a)=1+a$, which forces $a=-1$ for an identity endpoint, but in degree two the off-diagonal $D_{20}$ then survives. This obstruction is derived from the displayed V31 reverse relation, not just a comparison of family names.

\paragraph{Limits of the conclusion.} The simultaneous specialization $A=B=0$ with a fixed common variable and fixed Al-Salam--Carlitz parameter $a$ is \emph{not} an identity-matrix endpoint for all degrees. It does \emph{not} exclude a different K--S-recognized limiting arrow involving variable rescaling, a singular parameter path, or polynomial renormalization. An independently readable original K--S limiting diagram is still required for a final eligibility decision. Candidate 1 therefore remains \textbf{OPEN}, with this particular endpoint \textbf{REJECTED}. No other candidate status changes.

\noindent\textbf{V129 verdict:} One plausible candidate-1 endpoint rigorously excluded; new PASS $=0$; recorded PASS total $=52$; primary-source and global certification gates remain \textbf{OPEN}. Original V31 unchanged.



\section{V130: Stieltjes--Wigert / alternative $q$-Charlier candidate---exact diagonal and low-degree obstruction tests}
\noindent\textbf{Scope and evidence.} This stage examines the \emph{alternative} $q$-Charlier (also called $q$-Bessel) candidate, distinct from the ordinary $q$-Charlier family. The unchanged V31 source is \texttt{eq:SWqB-reverse} and the defining formula \texttt{eq:SWqB-def}. The computations below establish exact necessary conditions for an identity endpoint with the \emph{same polynomial variable and the unscaled V31 normalizations}. They do not certify or disprove a K--S arrow involving a variable transformation.

\paragraph{Leading coefficients and exact reverse diagonal.}
V31 defines
\[
 S_n(x;q)=\sum_{j=0}^n\frac{(-1)^j q^{j^2}}{(q;q)_j(q;q)_{n-j}}x^j,
 \qquad
 y_n(x;a;q)=\sum_{j=0}^n\frac{(q^{-n},-aq^n;q)_j q^j}{(q;q)_j}x^j.
\]
Using $(q^{-n};q)_n=(-1)^n q^{-n(n+1)/2}(q;q)_n$, the respective leading coefficients are
\[
 [x^n]S_n=\frac{(-1)^n q^{n^2}}{(q;q)_n},\qquad
 [x^n]y_n=(-1)^n q^{-n(n-1)/2}(-aq^n;q)_n.
\]
Therefore, in the V31 reverse expansion $y_n=\sum_{k\le n}D_{nk}(a;q)S_k$,
\begin{equation}\label{eq:v130-diagonal}
 \boxed{D_{nn}(a;q)=q^{-n(3n-1)/2}(q;q)_n(-aq^n;q)_n.}
\end{equation}
This is an exact consequence of the V31 defining polynomials, independent of the nested coefficient sum. In particular
\[
 D_{11}=\frac{(1-q)(1+aq)}q,\qquad
 D_{22}=\frac{(1-q)(1-q^2)(1+aq^2)(1+aq^3)}{q^5}.
\]
Thus a diagonal-identity endpoint at fixed $q$ would have to satisfy \emph{both} $D_{11}=1$ and $D_{22}=1$, not merely one leading-coefficient match.

\paragraph{A stronger degree-one necessary condition.}
Direct expansion gives
\[
 S_1(x;q)=\frac{1-qx}{1-q},\qquad
 y_1(x;a;q)=1-(1+aq)x.
\]
Consequently, comparing the constant terms and leading terms in $y_1=D_{11}S_1+D_{10}S_0$ yields
\begin{equation}\label{eq:v130-degree-one}
 \boxed{D_{11}=\frac{(1-q)(1+aq)}q,\qquad D_{10}=1-D_{11}.}
\end{equation}
The degree-one identity endpoint forces $a=(2q-1)/(q(1-q))$, which is admissible for positive $a$ only if $q>1/2$. This eliminates the possibility of an unscaled identity endpoint for $0<q\le1/2$ with $a>0$; it does \emph{not} settle the remaining range $1/2<q<1$, nor paths involving a variable or normalization change.

\paragraph{Independent numerical countercheck.} At $q=3/4$, the unique degree-one-matching parameter is $a=8/3$. The exact diagonal formula then gives
\[
 D_{22}=\frac{(1/4)(7/16)(1+(8/3)(9/16))(1+(8/3)(27/64))}{(3/4)^5}
 =\frac{(7/64)(5/2)(17/8)}{243/1024}
 =\frac{595}{243}\ne1.
\]
Hence even though the first-degree reverse matrix becomes the identity at this point, the degree-two diagonal does not. This is an explicit counterexample to inferring a full identity endpoint from degree one.

\paragraph{Certification boundary.} The tests above concern a fixed $q$, fixed common variable, and the exact V31 normalizations. A genuine K--S limiting relation could use a varying parameter, variable rescaling, or family-dependent renormalization. No primary-source K--S arrow has been authenticated here, so the candidate remains \textbf{OPEN}; no PASS is added. The recorded total stays \textbf{52}. The original V31 file is unchanged.



\clearpage
\section{V131: uniform degree-two obstruction for the Stieltjes--Wigert/alternative $q$-Charlier fixed-variable endpoint}
\label{sec:v131-uniform-obstruction}

\paragraph{Scope and correction of the earlier evidence.}
V130 supplied one counterexample at $q=3/4$, but a single failed point cannot exclude an endpoint at a different $q$. Here we strengthen that calculation to an \emph{all-$q$ impossibility theorem} for the precise fixed-variable, unrenormalized identity specialization. The result does not exclude a genuine Koekoek--Swarttouw arrow with variable or polynomial renormalization, and it does not settle K--S eligibility.

\paragraph{Source and exact leading coefficients.}
The unchanged V31 equation \texttt{eq:SWqB-reverse} expresses the alternative $q$-Charlier polynomial $y_n(x;a;q)$ in the Stieltjes--Wigert basis $S_k(x;q)$. Their defining monomial coefficients give
\[
 [x^n]y_n(x;a;q)=\frac{(q^{-n},-aq^n;q)_n q^n}{(q;q)_n},\qquad
 [x^n]S_n(x;q)=\frac{(-1)^nq^{n^2}}{(q;q)_n}.
\]
Thus the reverse diagonal, independently of the nested summation, is
\begin{equation}\label{eq:v131-diagonal}
 D_{nn}(a;q)=q^{-n(3n-1)/2}(q;q)_n(-aq^n;q)_n.
\end{equation}
At degree one, the exact expansions $y_1=1-(1+aq)x$ and $S_1=(1-qx)/(1-q)$ give
\[
 D_{11}=\frac{(1-q)(1+aq)}q,\qquad D_{10}=1-D_{11}.
\]
For a positive parameter $a$, the degree-one identity requirement forces
\begin{equation}\label{eq:v131-a-star}
 a=a_*(q):=\frac{2q-1}{q(1-q)},\qquad \tfrac12<q<1.
\end{equation}
For $0<q\le\tfrac12$, no positive $a$ satisfies even the degree-one requirement.

\paragraph{Exact degree-two contradiction for every remaining $q$.}
Substituting \eqref{eq:v131-a-star} into \eqref{eq:v131-diagonal} at $n=2$ and simplifying yields
\begin{align}
 D_{22}(a_*(q);q)
 &=\frac{(1+q)(2q^2-2q+1)(2q^3-q^2-q+1)}{q^5},\\
 D_{22}(a_*(q);q)-1
 &=\frac{P(q)}{q^5},\qquad
 P(q)=4q^6-3q^5-4q^4+5q^3-2q+1.
 \label{eq:v131-polynomial}
\end{align}
We prove $P(q)>0$ throughout $\tfrac12\le q<1$ without numerical root finding. Make the bijective substitution
\[
 q=\frac{1+2t}{2(1+t)},\qquad t\ge0.
\]
A direct polynomial expansion gives the \emph{positive-coefficient certificate}
\begin{equation}\label{eq:v131-positive-certificate}
 32(1+t)^6 P\!\left(\frac{1+2t}{2(1+t)}\right)
 =32t^6+96t^5+168t^4+196t^3+142t^2+59t+11>0.
\end{equation}
Every coefficient on the right is strictly positive. Hence $D_{22}(a_*(q);q)>1$ for \emph{every} $\tfrac12<q<1$. Combining the two $q$ ranges proves
\begin{equation}\label{eq:v131-impossibility}
 \boxed{\text{There is no }(q,a)\in(0,1)\times(0,\infty)
 \text{ for which }D_{nk}(a;q)=\delta_{nk}\text{ for all }n,k.}
\end{equation}
Indeed, the identity already fails by degree two. This is a uniform theorem, not merely the example from V130.

\paragraph{Precise verdict and ledger.}
\textbf{REJECTED:} the fixed-$q$, fixed-variable, unrenormalized identity specialization $y_n(x;a;q)=S_n(x;q)$ with $a>0$. \textbf{OPEN:} any distinct limiting arrow allowing $q$ to vary, rescaling of $x$, degree-dependent renormalization, or other K--S parameter paths. The proof does not authenticate an original K--S arrow. No new PASS is awarded; the cumulative \emph{recorded} PASS total remains \textbf{52}. The original V31 manuscript has not been modified.



\clearpage
\section{V132: Stieltjes--Wigert versus ordinary $q$-Charlier---an all-degree diagonal obstruction}
\label{sec:v132-SW-ordinary-qCharlier}

\paragraph{Scope and exact V31 source.}
This stage tests a \emph{fixed-variable, fixed-$q$, unrenormalized identity specialization} of the distinct candidate Stieltjes--Wigert $\to$ \emph{ordinary} $q$-Charlier, not the alternative $q$-Charlier candidate excluded in V131.  The unchanged V31 source is equation \texttt{eq:SWqC-rev} in the section \texttt{sec:SW-qCharlier}.  The reverse direction in that equation is
\begin{equation}\label{eq:v132-reverse}
 C_n(x;a;q)=\sum_{k=0}^n D_{nk}(a;q)S_k(x;q),\qquad 0<q<1,\quad a>0.
\end{equation}
The V31 definitions give
\begin{align}
 S_n(x;q)&=\sum_{j=0}^n\frac{(-1)^j q^{j^2}}{(q;q)_j(q;q)_{n-j}}x^j,
 & [x^n]S_n&=\frac{(-1)^nq^{n^2}}{(q;q)_n},\\
 C_n(x;a;q)&={}_2\phi_1\!\left(\begin{matrix}q^{-n},x\\0\end{matrix};q,-q^{n+1}/a\right),
 & [x^n]C_n&=\frac{(-1)^nq^{n^2}}{a^n}.
\end{align}
The $q$-Charlier leading coefficient also follows directly from the $r=n$ summand: $[x^n](x;q)_n=(-1)^nq^{n(n-1)/2}$ and $(q^{-n};q)_n=(-1)^nq^{-n(n+1)/2}(q;q)_n$.

\paragraph{Exact diagonal of the actual reverse connection.}
Comparing the $x^n$ coefficients in \eqref{eq:v132-reverse} gives, for \emph{every} $n\ge0$,
\begin{equation}\label{eq:v132-diagonal}
 \boxed{D_{nn}(a;q)=\frac{(q;q)_n}{a^n}.}
\end{equation}
This identity is exact for arbitrary $a>0$; it is not an endpoint-only formula.  In particular,
\begin{equation}\label{eq:v132-first-two}
 D_{11}=\frac{1-q}{a},\qquad
 D_{22}=\frac{(1-q)(1-q^2)}{a^2}.
\end{equation}

\paragraph{Uniform impossibility of an unrenormalized identity endpoint.}
If the reverse matrix were the identity at some fixed $a$, then $D_{11}=1$ would force $a=1-q$.  Substitution into the degree-two diagonal yields
\begin{equation}\label{eq:v132-obstruction}
 D_{22}(1-q;q)=\frac{1-q^2}{1-q}=1+q>1
 \qquad(0<q<1).
\end{equation}
Thus no $a>0$ can simultaneously make the first two diagonals equal to one, regardless of any off-diagonal behavior.  More generally, if a parameter path $a(t)>0$ with \emph{fixed} $q\in(0,1)$ made the full reverse matrix tend to the identity, then $D_{11}(a(t);q)\to1$ would imply $a(t)\to1-q$, contradicting \eqref{eq:v132-obstruction}.  Therefore
\begin{equation}\label{eq:v132-no-path}
 \boxed{\text{For fixed }0<q<1,\text{ no positive-}a\text{ path yields }D_{nk}\to\delta_{nk}
 \text{ with these fixed-variable, unrenormalized bases.}}
\end{equation}
This rules out an entire class of fixed-$q$ paths, not just a single parameter value.  It does \emph{not} rule out a rescaled-variable, renormalized, or $q$-varying limit.

\paragraph{Certification and count.}
\textbf{REJECTED:} the fixed-$q$, fixed-variable, unrenormalized identity specialization for Stieltjes--Wigert $\to$ ordinary $q$-Charlier.  \textbf{OPEN:} recognition of any other distinct K--S limiting arrow and any path requiring scaling or renormalization.  No primary K--S equation is authenticated by this computation.  \textbf{New PASS: 0; cumulative recorded PASS: 52.}  This is separate from the alternative $q$-Charlier obstruction of V131.  The original V31 file is unchanged.


\clearpage
\section{V133: Generalised Bessel / Jacobi --- a constructive confluent reverse limit}
\noindent\textbf{Scope and decision.} We revisit the unresolved V31 pair
\texttt{eq:jacobi-gb-reverse} (Jacobi expanded in the Generalised Bessel basis).
Unlike an obstruction for a fixed-variable identity endpoint, this section
constructs an explicit parameter-and-variable scaling for which the \emph{entire}
normalised reverse-connection matrix tends to the identity. This proves a
mathematical limiting relation, but \emph{does not yet certify} that the
original Koekoek--Swarttouw chart recognises this as a separate directed arrow.

\paragraph{Source definitions and confluent parameter path.}
The V31 Generalised Bessel convention is
\begin{equation}\label{eq:v133-gb-def}
 \mathcal{B}_n(x;\gamma)=\sum_{j=0}^n
 \frac{(-n)_j(n+\gamma+1)_j}{j!}\left(-\frac{x}{2}\right)^j.
\end{equation}
For Jacobi polynomials in the usual normalisation,
\begin{equation}\label{eq:v133-jac-def}
 P_n^{(\alpha,\beta)}(y)=\frac{(\alpha+1)_n}{n!}
 {}_2F_1\!\left(\begin{matrix}-n,n+\alpha+\beta+1\\\alpha+1\end{matrix};\frac{1-y}{2}\right).
\end{equation}
Fix a generic complex $\gamma$ such that $\mathcal{B}_0,\ldots,\mathcal{B}_N$
form a nondegenerate degree-graded basis for each fixed $N$ (in particular,
$(n+\gamma+1)_n\ne0$ for the degrees under consideration).  Let
\begin{equation}\label{eq:v133-path}
 \alpha=c-1,\qquad \beta=\gamma-c+1,\qquad y=1+cx,
 \qquad c\longrightarrow+\infty.
\end{equation}
The Jacobi parameters are understood by algebraic continuation, not by
simultaneous positivity of both orthogonality weights. Define
\begin{equation}\label{eq:v133-normalization}
 \widehat P_{n,c}(x):=\frac{n!}{(c)_n}
 P_n^{(c-1,\gamma-c+1)}(1+cx).
\end{equation}
Then the terminating series gives the \emph{exact} polynomial formula
\begin{equation}\label{eq:v133-exact}
 \widehat P_{n,c}(x)
 =\sum_{j=0}^n \frac{(-n)_j(n+\gamma+1)_j}{j!}
 \left(-\frac{x}{2}\right)^j\frac{c^j}{(c)_j}.
\end{equation}
For each fixed $j$, $c^j/(c)_j=\prod_{r=0}^{j-1}(1+r/c)^{-1}
=1+O(c^{-1})$. Consequently, for each fixed $n$,
\begin{equation}\label{eq:v133-limit}
 \boxed{\widehat P_{n,c}(x)=\mathcal{B}_n(x;\gamma)+O(c^{-1})}
\end{equation}
coefficientwise and locally uniformly in $x$. This is an all-fixed-degree
confluent limit, not an assertion about degree growing with $c$.

\paragraph{Verification of the actual reverse-connection matrix.}
Write the V31 reverse connection
$P_n^{(\alpha,\beta)}(y)=\sum_{k=0}^n
D_{nk}^{\mathrm{J,GB}}(\alpha,\beta;\gamma)\mathcal{B}_k(y;\gamma)$.
\emph{Substituting parameters alone into that fixed-variable matrix is not
sufficient}, because $y=1+cx$ also transforms the target basis.
Define the induced, rescaled reverse coefficients $\widehat D_{nk}(c)$ by
\begin{equation}\label{eq:v133-induced}
 \widehat P_{n,c}(x)=\sum_{k=0}^n
 \widehat D_{nk}(c)\mathcal{B}_k(x;\gamma).
\end{equation}
These are the coefficients obtained from the V31 identity by substitution,
normalisation and a finite change of polynomial basis, uniquely determined
for generic parameters.  Since the basis $\mathcal{B}_k(x;\gamma)$ is independent
of $c$ and degree-graded, the change-of-basis map from monomials is a fixed
invertible triangular matrix. Equation~\eqref{eq:v133-exact} therefore
implies the \emph{full} matrix estimate
\begin{equation}\label{eq:v133-matrix}
 \boxed{\widehat D_{nk}(c)=\delta_{nk}+O(c^{-1})
 \qquad(0\le k\le n\le N),}
\end{equation}
for every fixed $N$.

For an independent diagonal check, the top coefficient of
\eqref{eq:v133-exact} is
$[x^n]\widehat P_{n,c}=(-n)_n(n+\gamma+1)_n(-1/2)^n c^n/[n!(c)_n]$,
while $[x^n]\mathcal{B}_n=(-n)_n(n+\gamma+1)_n(-1/2)^n/n!$.
Thus
\begin{equation}\label{eq:v133-diagonal}
 \boxed{\widehat D_{nn}(c)=\frac{c^n}{(c)_n}
 =1-\frac{n(n-1)}{2c}+O(c^{-2}).}
\end{equation}
At degree one the factor is exactly one, and
$\widehat P_{1,c}=\mathcal{B}_1$ for all $c$. At degree two, the $j=2$
coefficient differs by the factor $c/(c+1)$, consistent with
\eqref{eq:v133-matrix} and \eqref{eq:v133-diagonal}.

\paragraph{Audit verdict and limits of the conclusion.}
\textbf{MATHEMATICALLY VERIFIED:} a Jacobi $\to$ Generalised Bessel
confluent limit with explicit singular parameter path, affine variable
scaling, degree-dependent normalisation and convergence of the full induced
reverse matrix. \textbf{K--S ELIGIBILITY OPEN:} no independently authenticated
primary-source chart or equation has been provided establishing that this
is a \emph{distinct counted} K--S arrow; the source V31 subsection itself
warns that its fixed-parameter connection is not automatically a K--S
limit. We therefore \emph{do not} promote this candidate to PASS yet.
\textbf{New PASS: 0; cumulative recorded PASS: 52.} The original V31
manuscript has not been changed.


\clearpage
\section{V134: Meixner / dual Hahn candidate---degree filtration and grid-versus-polynomial audit}
\label{sec:v134-meixner-dual-hahn}
\paragraph{Source and scope.}
We inspect the V31 equations \texttt{eq:meixner-dh-forward} and
\texttt{eq:dh-meixner-reverse}, without assuming that the existence of a
finite-grid connection constitutes a Koekoek--Swarttouw limiting arrow.
Write $\sigma=\gamma+\delta+1$ and
$\Lambda(x)=x(x+\sigma)$.  The source expressly restricts the forward
Meixner-to-dual-Hahn expansion to $x=0,\ldots,N$, whereas the reverse
expansion is an unrestricted polynomial identity in $x$.

\paragraph{A rigorous degree obstruction.}
The dual Hahn polynomial $\mathrm{DH}_n(\Lambda;\gamma,\delta,N)$ is,
under its nondegeneracy assumptions, of degree exactly $n$ in $\Lambda$.
Consequently
\begin{equation}\label{eq:v134-degree}
 \deg_x\mathrm{DH}_n(\Lambda(x);\gamma,\delta,N)=2n
 \qquad(n\ge1),
\end{equation}
because $\Lambda(x)=x^2+\sigma x$ is monic quadratic.  In contrast,
$\mathrm M_j(x;\beta,c)$ has degree $j$ in $x$ for admissible $c\ne0,1$.
Thus any unrestricted polynomial expansion of the dual Hahn member in
Meixner polynomials must include the degree-$2n$ basis member with
nonzero coefficient.  More precisely, denoting leading coefficients by
$\operatorname{lc}$,
\begin{equation}\label{eq:v134-top}
 D_{n,0}=
 \frac{\operatorname{lc}_{\Lambda}\mathrm{DH}_n(\Lambda;\gamma,\delta,N)}
 {\operatorname{lc}_{x}\mathrm M_{2n}(x;\beta,c)}\ne0
 \quad\text{for all nondegenerate finite parameters}.
\end{equation}
Here $D_{n,0}$ is the coefficient of $\mathrm M_{2n}$ in the indexing
of \texttt{eq:dh-meixner-reverse}.  This is a leading-coefficient identity,
not an asymptotic assertion: it follows by comparing the $x^{2n}$ terms
on the two sides.  In particular, the $n=1$ reverse relation requires
$\mathrm M_2$; it cannot be an identity of the form
$\mathrm{DH}_1(\Lambda(x))=\mathrm M_1(x)$ on the unrestricted variable.

\paragraph{Why the forward grid identity does not remove the obstruction.}
On the grid $\{0,\ldots,N\}$ every function has an expansion in the
$N+1$ dual Hahn grid basis, and the polynomial $\Lambda(x)$ separates
nodes by the V31 hypothesis.  Hence the grid connection does not imply
a degree-preserving polynomial identity away from the grid.  Indeed,
for any polynomial $F(x)$, adding a multiple of
$\prod_{j=0}^N(x-j)$ leaves all its grid values unchanged.  The
forward coefficient matrix therefore encodes interpolation on a finite
set, not by itself a polynomial-family limit.

\paragraph{Precisely delimited verdict.}
\textbf{VERIFIED:} V31's reverse expansion has an intrinsic degree-doubling
structure and its forward expansion is grid-only.  A fixed-variable,
degree-preserving \emph{identity specialization} between
$\mathrm{DH}_n(\Lambda(x))$ and $\mathrm M_n(x)$ is ruled out at
nondegenerate finite parameters by \eqref{eq:v134-degree}.  This does
\emph{not} rule out a singular limiting regime with simultaneous variable
and parameter scaling, in which top coefficients may vanish after
normalisation, nor does it establish or refute recognition in the original
K--S chart.  \textbf{K--S eligibility remains open.}
\textbf{New PASS: 0; cumulative recorded PASS: 52.}
The original V31 manuscript is unchanged.


\clearpage
\section{V135: Jacobi to Generalised Bessel---first-order confluent asymptotics and an independent residual test}
\label{sec:v135-jacobi-bessel-first-order}

\paragraph{Purpose and source discipline.}
This section strengthens the constructive V133 candidate without adding a
K--S PASS.  It independently checks the \emph{whole polynomial}, rather than
only its top coefficient, by deriving the exact first asymptotic correction
from the terminating Jacobi series.  The unchanged V31 source is
\texttt{eq:jacobi-gb-reverse}; the parameter and variable transformation
is exactly that in \eqref{eq:v133-path}.  This calculation is a mathematical
cross-check, not an independent authentication of the K--S limiting chart.

\paragraph{Uniform expansion of the terminating factors.}
For each fixed nonnegative integer $j$, with $c>0$,
\begin{equation}\label{eq:v135-factor}
 \frac{c^j}{(c)_j}=\prod_{r=0}^{j-1}(1+r/c)^{-1}
 =1-\frac{j(j-1)}{2c}+O_j(c^{-2}).
\end{equation}
Indeed, the coefficient of $1/c$ in the product is
$-\sum_{r=0}^{j-1}r=-j(j-1)/2$.  For any fixed maximum degree $N$,
the remainder is bounded uniformly for $0\le j\le N$ as $c\to\infty$.
This uniformity is only for fixed $N$, and no growing-degree claim is made.

\paragraph{An exact first-order differential-operator correction.}
Write $\mathcal B_n(x;\gamma)$ for the V31 Generalised Bessel polynomial
as defined in \eqref{eq:v133-gb-def}.  Applying
\eqref{eq:v135-factor} term by term to the \emph{exact} terminating
formula \eqref{eq:v133-exact} yields
\begin{equation}\label{eq:v135-operator}
 \boxed{\widehat P_{n,c}(x)=\mathcal B_n(x;\gamma)
 -\frac{x^2}{2c}\,\mathcal B_n''(x;\gamma)
 +O_{n,\gamma}(c^{-2}).}
\end{equation}
The error is coefficientwise and locally uniform in $x$.  The operator
$x^2\partial_x^2$ multiplies the coefficient of $x^j$ by $j(j-1)$,
which proves the claimed correction without invoking any unverified
hypergeometric transformation.  In particular, the correction is zero
for degrees $0$ and $1$, and generally nonzero for degree $2$ and above.

\paragraph{Exact degree-two residual.}
From the defining terminating series,
\begin{align}
 \mathcal B_2(x;\gamma)
 &=1+(\gamma+3)x+\frac{(\gamma+3)(\gamma+4)}4x^2,\\
 \widehat P_{2,c}(x)
 &=1+(\gamma+3)x
 +\frac{(\gamma+3)(\gamma+4)}4\frac{c}{c+1}x^2.
\end{align}
Consequently, for every admissible $c$,
\begin{equation}\label{eq:v135-degree-two}
 \boxed{\widehat P_{2,c}(x)-\mathcal B_2(x;\gamma)
 =-\frac{(\gamma+3)(\gamma+4)}{4(c+1)}x^2.}
\end{equation}
Expanding $1/(c+1)=1/c+O(c^{-2})$ agrees exactly with
\eqref{eq:v135-operator}.  This is a nontrivial lower-degree as well as
leading-coefficient check: the constant and linear coefficients agree
\emph{identically} in $c$ while the quadratic coefficient has the
predicted first-order discrepancy.

\paragraph{Consequences for the induced reverse matrix.}
Let $\mathsf C_{kj}(\gamma)$ be the unique coefficients in the finite
expansion
\begin{equation}\label{eq:v135-C}
 x^2\mathcal B_n''(x;\gamma)
 =\sum_{k=0}^n \mathsf C_{nk}(\gamma)\mathcal B_k(x;\gamma).
\end{equation}
For generic $\gamma$, the basis is degree graded, so these constants
exist and are independent of $c$.  Substituting into
\eqref{eq:v135-operator} proves the sharper matrix expansion
\begin{equation}\label{eq:v135-matrix}
 \boxed{\widehat D_{nk}(c)
 =\delta_{nk}-\frac{\mathsf C_{nk}(\gamma)}{2c}
 +O_{N,\gamma}(c^{-2}),\qquad 0\le k\le n\le N.}
\end{equation}
The leading coefficient of $x^2\mathcal B_n''$ is $n(n-1)$ times
that of $\mathcal B_n$, hence $\mathsf C_{nn}=n(n-1)$.
This recovers the independent diagonal expansion
$\widehat D_{nn}=1-n(n-1)/(2c)+O(c^{-2})$ from V133.
No assertion that every off-diagonal first-order term is nonzero is
needed or made.

\paragraph{Certification boundary and tally.}
\textbf{VERIFIED:} the first-order correction for every fixed degree,
the exact degree-two residual, and the first-order expansion of the
entire \emph{induced, rescaled} reverse matrix.  This is not a claim
that the untransformed V31 reverse matrix itself tends to the identity:
the affine variable change and normalization are essential.
\textbf{STILL OPEN:} whether this ordered family pair is a separate
arrow in the original Koekoek--Swarttouw chart.  No new PASS is awarded.
\textbf{Cumulative recorded PASS: 52.} The original V31 source remains
unchanged.


\clearpage
\section{V136: Explicit first subdiagonal of the Jacobi--generalised Bessel limiting matrix}
\label{sec:v136-first-subdiagonal}
\paragraph{Scope.}
This section sharpens V135 by computing an \emph{off-diagonal} first-order
coefficient of the induced, rescaled Jacobi-to-Generalised-Bessel reverse
matrix. It is an internal mathematical check, not independent authentication
of an arrow in the Koekoek--Swarttouw chart. The original V31 manuscript is
not changed.

\paragraph{Polynomial normalization.}
Let
\[
 \mathcal B_n(x;\gamma)=\sum_{j=0}^n b_{n,j}x^j,
 \qquad b_{n,j}=\binom nj\frac{(n+\gamma+1)_j}{2^j}.
\]
These coefficients follow directly from the terminating representation
of the Generalised Bessel polynomials in V31. Assume all leading
coefficients used below are nonzero. The first-order expansion in V135 is
\[
 \widehat P_{n,c}(x)=\mathcal B_n(x;\gamma)
 -\frac{x^2}{2c}\mathcal B_n''(x;\gamma)+O(c^{-2}).
\]
Define the unique triangular coefficients $\mathsf C_{nk}$ by
\[
 x^2\mathcal B_n''(x;\gamma)
 =\sum_{k=0}^n\mathsf C_{nk}(\gamma)\mathcal B_k(x;\gamma).
\]

\paragraph{Exact diagonal and first subdiagonal.}
Comparison of the $x^n$ coefficient gives
$\mathsf C_{nn}=n(n-1)$. Comparing $x^{n-1}$ gives
\[
 (n-1)(n-2)b_{n,n-1}
 =n(n-1)b_{n,n-1}+\mathsf C_{n,n-1}b_{n-1,n-1}.
\]
Hence for $n\ge2$,
\begin{equation}\label{eq:v136-C-subdiag}
 \boxed{\mathsf C_{n,n-1}
 =-2(n-1)\frac{b_{n,n-1}}{b_{n-1,n-1}}
 =-\frac{2n(n-1)(\gamma+2n-1)}{\gamma+n}.}
\end{equation}
The ratio uses
$\,b_{n,n-1}=n(n+\gamma+1)_{n-1}/2^{n-1}\,$ and
$\,b_{n-1,n-1}=(n+\gamma)_{n-1}/2^{n-1}\,$.
At removable exceptional parameter values, the formula should be read
through the polynomial coefficient equations rather than dividing by a
vanishing leading coefficient.

\paragraph{Consequences for the actual transformed connection.}
Combining \eqref{eq:v136-C-subdiag} with \eqref{eq:v135-matrix} yields
\begin{equation}\label{eq:v136-D-subdiag}
 \boxed{\widehat D_{n,n-1}(c)
 =\frac{n(n-1)(\gamma+2n-1)}{c(\gamma+n)}
 +O_{n,\gamma}(c^{-2}),\qquad n\ge2.}
\end{equation}
In particular, the off-diagonal coefficients are generally of order
$c^{-1}$, so V135's $O(c^{-1})$ estimate for the full matrix cannot in
general be improved to $O(c^{-2})$. The degree-one row is exceptional:
$\widehat P_{1,c}=\mathcal B_1$ exactly, so its off-diagonal entry is zero.

\paragraph{Independent degree-two check.}
Since
$\mathcal B_1=1+(\gamma+2)x/2$ and
$\mathcal B_2=1+(\gamma+3)x+(\gamma+3)(\gamma+4)x^2/4$,
\[
 x^2\mathcal B_2''
 =2\mathcal B_2
 -\frac{4(\gamma+3)}{\gamma+2}\mathcal B_1
 +\frac{2(\gamma+4)}{\gamma+2}\mathcal B_0.
\]
Thus
\[
 \widehat D_{21}(c)=\frac{2(\gamma+3)}{c(\gamma+2)}+O(c^{-2}),
 \quad
 \widehat D_{20}(c)=-\frac{\gamma+4}{c(\gamma+2)}+O(c^{-2}).
\]
These two terms cancel the constant and linear contributions introduced
by the diagonal first-order term, leaving precisely the quadratic
correction derived in V135. This is an explicit nontrivial consistency
check across the \emph{entire degree-two row}.

\paragraph{Certification and count.}
\textbf{Verified:} the generic first subdiagonal and the complete
first-order degree-two row of the induced reverse matrix.
\textbf{Not verified here:} recognition of Jacobi to Generalised Bessel
as an independent Koekoek--Swarttouw limiting arrow, or global
completeness of the 52-entry ledger.
\textbf{New PASS: 0. Cumulative recorded PASS: 52.}


\clearpage
\section{V137: first direct authentication against the 1998 Koekoek--Swarttouw primary source}
\label{sec:v137-primary-kls}
\paragraph{Primary-source discovery and scope.}
The authors' 1998 Delft report, \emph{The Askey-scheme of hypergeometric orthogonal polynomials and its $q$-analogue}, Report 98--17, is accessible from the Delft-hosted PDF
\texttt{https://tudelft.roelofkoekoek.nl/documents/as98.pdf}.
The following findings are based on the text and rendered pages of that primary report, rather than on a secondary diagram.  This is a \emph{partial} primary-source authentication, not a check of all 52 entries.

\paragraph{Finding A: a previously pending arrow is explicitly present.}
In section 4.11, printed pp.~128--129, the authors explicitly list
\emph{Big $q$-Laguerre $\longrightarrow$ Al-Salam--Carlitz I} and give equation (4.11.2):
\[
 \lim_{a\to0}\frac{P_n(aqx;a,ab;q)}{a^n}=U_n^{(b)}(x;q)
 \qquad\text{(as printed in the 1998 report).}
\]
Thus the V128--V129 description ``KLS eligibility unresolved'' for this \emph{ordered family pair} is now obsolete: \textbf{the KLS arrow is independently authenticated}.  The obstruction in V129 applies only to the distinct fixed-variable endpoint $A=B=0$, and cannot disprove the scaled KLS limit.

\paragraph{A critical normalization warning requiring a further check.}
The V31 normalization, recorded beside \texttt{eq:ASCI-bL-reverse}, is
\[
 P_n(y;A,B\mid q)={}_3\phi_2(q^{-n},0,y;Aq,Bq;q,q),
 \qquad \operatorname{lc}_yP_n=\frac1{(Aq,Bq;q)_n},
\]
and $U_n^{(b)}(x;q)$ is monic by its generating function. Consequently the leading coefficient in $x$ of the \emph{V31-normalized} expression
\[
 a^{-n}P_n(aqx;a,ab\mid q)
\]
is exactly $q^n/[(aq;q)_n(abq;q)_n]$, which tends to $q^n$, not $1$.
Therefore the printed KLS formula must \emph{not} be transplanted verbatim into the V31 normalization without checking the source's polynomial argument and normalization conventions.  In particular, an identity-matrix endpoint for the V31 reverse coefficients does not follow merely from the primary-source citation.  \textbf{Decision: KLS recognition YES; V31-normalized transformed reverse-matrix verification PENDING; no new PASS in V137.}

\paragraph{Finding B: a second previously pending exact arrow is explicit.}
Section 4.8, printed p.~127, explicitly lists
\emph{Al-Salam--Chihara $\longrightarrow$ continuous $q$-Laguerre}, equation (4.8.2), with the special parameters
\[
 a=q^{\alpha/2+1/4},\qquad b=q^{\alpha/2+3/4}
\]
and a degree-dependent proportionality factor.  This is \emph{explicit recognition of the ordered arrow}, including an exact specialization rather than a singular limit.  Under the user's counting rule, an exact specialization is eligible once its V31 normalization and full reverse coefficients are verified.  V116 established an endpoint relation, but this stage does not independently reconcile its exact factor with the printed primary-source formula.  \textbf{Decision: KLS recognition YES; final V31 matrix and duplicate check PENDING; no new PASS in V137.}

\paragraph{Finding C: the Stieltjes--Wigert direction must be respected.}
Section 4.23, printed p.~134, gives equation (4.23.1) under the heading
\emph{$q$-Charlier $\longrightarrow$ Stieltjes--Wigert}:
\[
 \lim_{a\to\infty}C_n(ax;a;q)=(q;q)_nS_n(x;q).
\]
This is the \emph{opposite direction} to the candidate
Stieltjes--Wigert $\longrightarrow$ ordinary $q$-Charlier investigated in V132.  The published arrow therefore does not establish the reverse ordered pair.  The negative fixed-variable calculation in V132 remains valid for its stated assumptions; it should not be read as excluding the KLS forward arrow.

\paragraph{Finding D: an existing PASS has primary-source confirmation.}
Section 4.24, printed p.~135, gives
$U_n^{(-1)}(x;q)=h_n(x;q)$ as equation (4.24.1), confirming the KLS recognition of the Al-Salam--Carlitz I $\longrightarrow$ discrete $q$-Hermite I specialization (PASS 51). Section 4.25, printed p.~135, also records the Al-Salam--Carlitz II $\longrightarrow$ discrete $q$-Hermite II complex-variable specialization as (4.25.1), corresponding to PASS 52.

\paragraph{Status and count.}
This stage resolves \emph{source-recognition} for two previously open ordered pairs, authenticates two recorded arrows, and identifies a directionality trap. It does \emph{not} resolve the normalization discrepancy for the first pair, verify all 52 KLS entries, or certify completeness. \textbf{New PASS: 0; cumulative recorded PASS: 52.}


\clearpage
\section{V138: primary-source and exact reverse-matrix certification of the Al-Salam--Chihara to continuous $q$-Laguerre arrow}
\label{sec:v138-asc-lag-certified}
\paragraph{Primary-source gate.}
The original Koekoek--Swarttouw 1998 report, Section 4.8, equation (4.8.2), explicitly lists the ordered arrow
Al-Salam--Chihara $\longrightarrow$ continuous $q$-Laguerre, with
\[
 a=q^{\alpha/2+1/4},\qquad b=q^{\alpha/2+3/4},\qquad
 Q_n(x;a,b\mid q)=\frac{(q;q)_n}{a^n}P_n^{(\alpha)}(x\mid q).
\]
The cited primary-source URL is
\texttt{https://tudelft.roelofkoekoek.nl/documents/as98.pdf},
PDF page 128 (printed page 127). This is an explicit \emph{distinct ordered-family entry}, not merely an inference from a general Askey--Wilson specialization.

\paragraph{V31 normalization and exact reverse coefficients.}
V31 defines its continuous $q$-Laguerre polynomial $L_n^{(\alpha)}$ by
\texttt{eq:qlag-definition} and records the exact identity
\texttt{eq:qlag-asc}:
\[
 L_n^{(\alpha)}(x\mid q)=\frac{a^n}{(q;q)_n}Q_n(x;a,b\mid q).
\]
This agrees with KLS (4.8.2), with $L_n^{(\alpha)}$ denoting the KLS polynomial $P_n^{(\alpha)}$ in V31 notation.
The actual V31 reverse connection \texttt{eq:lag-asc-hyper-rev51} is
\[
 Q_n(x;u,v\mid q)=\sum_{k=0}^n D_{nk}(u,v;a,b) L_k^{(\alpha)}(x\mid q),
\]
where, putting $m=n-k$, the equivalent \emph{finite polynomial kernel} form is
\[
 D_{nk}(u,v;a,b)=\begin{bmatrix}n\\k\end{bmatrix}_q
 \frac{(q;q)_k}{a^k}
 \sum_{r=0}^{m}\begin{bmatrix}m\\r\end{bmatrix}_q
 a^r(u/a;q)_r\,b^{m-r}(v/b;q)_{m-r}.
\]
The finite sum, unlike a quotient-based ${}_2\phi_1$ representation, is well-defined at the matching parameters. It follows by taking the ratio of the two Al-Salam--Chihara generating functions, as in V31's proof of the cited connection.
At $u=a$, $v=b$, each term for $m>0$ vanishes: for $r>0$, $(u/a;q)_r=(1;q)_r=0$, while for $r=0$ the factor $(v/b;q)_m=(1;q)_m=0$. For $m=0$ the sum equals one. Thus
\[
 \boxed{D_{nk}(a,b;a,b)=\delta_{nk}\frac{(q;q)_n}{a^n}.}
\]
This diagonal is \emph{not} one in V31's raw reverse connection; the degree-dependent KLS factor is essential.

\paragraph{Normalized complete-matrix endpoint.}
Multiply the source polynomial by $a^n/(q;q)_n$ and define
\[
 \widehat D_{nk}(u,v):=\frac{a^n}{(q;q)_n}D_{nk}(u,v;a,b).
\]
The preceding exact evaluation gives
\[
 \boxed{\widehat D_{nk}(a,b)=\delta_{nk}\qquad(0\le k\le n).}
\]
Moreover the kernel is polynomial in $(u,v)$, since
$a^r(u/a;q)_r=\prod_{j=0}^{r-1}(a-uq^j)$ and similarly for $b$.
Hence for fixed degree and $a,b\ne0$,
\[
 \widehat D_{nk}(u,v)=\delta_{nk}+O(|u-a|+|v-b|)
 \quad\text{as }(u,v)\to(a,b).
\]
This verifies the \emph{actual V31 reverse coefficient matrix} at every degree, including the nonunit raw diagonal and the complete off-diagonal collapse.

\paragraph{Ledger and decision.}
V116 established the exact endpoint but withheld a PASS because the separate KLS arrow had not been authenticated. V137 supplied the missing primary-source recognition. This stage combines both gates and checks the full V31 matrix explicitly. The pair is not among the previous 52 counted ordered-family PASS entries (the continuous $q$-Laguerre $\to$ Al-Salam--Chihara direction in historical discussions is distinct). Under the user's rule admitting independently recognized exact specializations, this is a \textbf{new PASS 53}.
\[
 \boxed{\text{Previously recorded PASS}=52,\quad\text{new PASS}=1,\quad\text{updated recorded PASS}=53.}
\]
This is \emph{not} global certification: other historical KLS citations, the big $q$-Laguerre scaling discrepancy, and remaining candidate decisions still require checking.

\paragraph{Separate caution on KLS (4.11.2).}
For the other KLS-recognized candidate, the printed expression $a^{-n}P_n(aqx;a,ab\mid q)$ has leading coefficient in $x$ equal to
$q^n/[(aq,abq;q)_n]$, hence limit $q^n$, whereas $U_n^{(b)}$ is monic. In fact the exact degree-one formula is
\[
 a^{-1}P_1(aqx;a,ab\mid q)=
 \frac{q(x-1-b+abq)}{(1-aq)(1-abq)}\longrightarrow q(x-1-b)=qU_1^{(b)}(x;q).
\]
Thus the printed KLS equation (4.11.2), as extracted, disagrees with the V31 normalization already in degree one; a missing factor or convention issue must be resolved before awarding that separate arrow. No PASS is granted for it here.


\clearpage
\section{V139: resolving the big $q$-Laguerre to Al-Salam--Carlitz I normalization and full reverse limit}
\label{sec:v139-biglag-asci}
\paragraph{Primary-source versus V31 normalization.}
The 1998 Koekoek--Swarttouw report, section 4.11, equation (4.11.2), explicitly labels the arrow
\emph{big $q$-Laguerre $\longrightarrow$ Al-Salam--Carlitz I} and prints
$a^{-n}P_n(aqx;a,ab;q)\to U_n^{(b)}(x;q)$.
The displayed definitions in its sections 3.11 and 3.24 and V31's
\texttt{eq:ASCI-bL-reverse} imply, however, that the left-hand leading coefficient tends to $q^n$ while the right-hand polynomial is monic. Thus the formula \emph{as printed} fails already for $n=1$ when $0<q<1$; this is a normalization discrepancy, not a failure of the recognized directed arrow.
The corrected normalization for the exact V31 polynomial definition is
\begin{equation}\label{eq:v139-corrected-limit}
 \boxed{\lim_{a\to0}(aq)^{-n}P_n(aqx;a,ab\mid q)=U_n^{(b)}(x;q).}
\end{equation}
Here $b$ and $q$ are fixed, $0<q<1$, and $n$ is fixed. The factor $q^{-n}$ relative to the printed KLS formula is necessary.

\paragraph{Independent proof for every degree using monic recurrences.}
Set $A=a$, $B=ab$ and define the monic big $q$-Laguerre polynomial
$p_n(y;a,ab)=(aq,abq;q)_nP_n(y;a,ab\mid q)$.
The monic recurrence given by the big $q$-Laguerre recurrence in KLS section 3.11 is
\[
 y p_n=p_{n+1}+[1-(\mathcal A_n+\mathcal C_n)]p_n
       +\mathcal A_{n-1}\mathcal C_n p_{n-1},
\]
where
\[
 \mathcal A_n=(1-aq^{n+1})(1-abq^{n+1}),\qquad
 \mathcal C_n=-a^2bq^{n+1}(1-q^n),\quad p_0=1.
\]
Define $T_n(x;a)=(aq)^{-n}p_n(aqx;a,ab)$, with $T_{-1}=0$ and $T_0=1$.
Dividing the recurrence by $(aq)^{n+1}$ yields
\[
 xT_n=T_{n+1}+B_n(a)T_n+K_n(a)T_{n-1},
\]
\[
 B_n(a)=\frac{1-\mathcal A_n-\mathcal C_n}{aq}
 =(1+b)q^n-abq^{2n+1}+abq^n(1-q^n),
\]
\[
 K_n(a)=\frac{\mathcal A_{n-1}\mathcal C_n}{a^2q^2}
 =-bq^{n-1}(1-q^n)(1-aq^n)(1-abq^n).
\]
In particular, $B_n(a)\to(1+b)q^n$ and $K_n(a)\to-bq^{n-1}(1-q^n)$.
The monic Al-Salam--Carlitz I family defined by V31's generating function obeys
\[
 xU_n^{(b)}=U_{n+1}^{(b)}+(1+b)q^nU_n^{(b)}
                    -bq^{n-1}(1-q^n)U_{n-1}^{(b)}.
\]
Induction on $n$ therefore gives coefficientwise polynomial convergence
$T_n(x;a)\to U_n^{(b)}(x;q)$, uniformly on compact $x$-sets.
Since $(aq,abq;q)_n\to1$, this proves equation \eqref{eq:v139-corrected-limit}.

\paragraph{Exact first-degree check and normalization necessity.}
From the terminating definition, the exact formula is
\[
 P_1(y;a,ab\mid q)
 =\frac{y-q(a+ab)+a^2bq^2}{(1-aq)(1-abq)}.
\]
Thus
\[
 (aq)^{-1}P_1(aqx;a,ab\mid q)
 =\frac{x-1-b+abq}{(1-aq)(1-abq)}\to x-1-b=U_1^{(b)}(x;q).
\]
This verifies the extra factor $q^{-1}$ required in the KLS printed expression at degree one.

\paragraph{Complete transformed V31 reverse matrix.}
The reverse connection \texttt{eq:ASCI-bL-reverse} is an exact finite polynomial identity for all admissible nonzero $a$:
$P_n(y;a,ab)=\sum_{k=0}^n D_{nk}(a,ab;b)U_k^{(b)}(y;q)$.
After replacing $y$ by $aqx$, the functions $U_k^{(b)}(aqx)$ are \emph{not} the fixed target basis $U_k^{(b)}(x)$; therefore a naive row rescaling of $D_{nk}$ is invalid.
Define the induced transformed reverse coefficients uniquely by
\[
 (aq)^{-n}P_n(aqx;a,ab\mid q)
 =\sum_{k=0}^n\widehat D_{nk}(a;b,q)U_k^{(b)}(x;q).
\]
This is the exact composition of the V31 reverse matrix with the affine-variable change matrix and the degree normalization. The target basis is monic, so its coefficient-to-polynomial map is an invertible triangular matrix independent of $a$. Coefficientwise convergence proved above therefore gives, for every fixed $n$,
\[
 \boxed{\widehat D_{nk}(a;b,q)=\delta_{nk}+O(a),\qquad 0\le k\le n.}
\]
More explicitly, the exact diagonal is
\[
 \boxed{\widehat D_{nn}(a;b,q)=\frac{1}{(aq,abq;q)_n}=1+O(a).}
\]
The error is coefficientwise analytic at $a=0$, and all lower entries vanish there. This proves the full transformed reverse-matrix limit, not merely the polynomial endpoint.

\paragraph{Eligibility and ledger.}
The ordered pair is explicitly a distinct arrow in KLS section 4.11 and was not included among the 53 recorded PASS entries through V138. The V31 reverse identity, necessary variable transformation, corrected normalization, and all-degree transformed-matrix convergence are now verified. Under the user's rule permitting recognized KLS limits with explicitly reconciled normalizations, award \textbf{PASS 54: big $q$-Laguerre $\to$ Al-Salam--Carlitz I}. The printed KLS equation (4.11.2) should be cited with the above normalization caveat, not quoted as a literally correct V31-normalized formula.
\[
 \boxed{\text{Previous recorded PASS}=53,\quad\text{new PASS}=1,\quad\text{current recorded PASS}=54.}
\]
Global completeness and independent authentication of all other historical arrows remain outstanding.


\section{V140: primary-source directionality audit and definitive candidate dispositions}
\label{sec:v140}
\paragraph{Scope and source.}
This stage checks the actual \emph{ordered} arrows against Koekoek--Swarttouw,
\emph{The Askey-scheme of hypergeometric orthogonal polynomials and its
$q$-analogue} (1998), available at
\texttt{https://tudelft.roelofkoekoek.nl/documents/as98.pdf}.
The primary source is structured into the classical limit chapter~2 and
$q$-limit chapter~4.  A mathematical connection, an inverse connection, or a
valid extra confluent limit is not automatically a separately listed arrow
in those chapters.  This distinction is required by the audit counting rule.

\paragraph{Jacobi versus generalized Bessel: a mathematically valid limit
that is not a listed KLS arrow.}
V133--V136 establish the genuine coefficientwise confluent limit
\[
 \lim_{c\to\infty}\frac{n!}{(c)_n}
 P_n^{(c-1,\gamma-c+1)}(1+cx)
 =\mathcal B_n(x;\gamma).
\]
Nevertheless, the 1998 report's chapter~1 family list contains Jacobi
(section~1.8) but no family of generalized Bessel \emph{polynomials}.
Its section~2.8 lists the arrows Wilson $\to$ Jacobi, continuous Hahn
$\to$ Jacobi, Hahn $\to$ Jacobi, Jacobi $\to$ Laguerre (2.8.1), and
Jacobi $\to$ Hermite (2.8.2), with a subsequent ultraspherical subsection.
It does \emph{not} list Jacobi $\to$ generalized Bessel as a separate
arrow.  The mentions of Bessel \emph{functions} in the preliminary
special-function discussion are not an entry for generalized Bessel
polynomials in the classical Askey scheme.
\emph{Disposition: mathematically proved auxiliary limit, but EXCLUDED
from the KLS-arrow PASS ledger.}  This does not contradict V133--V136;
it resolves their explicitly pending eligibility question.

\paragraph{Meixner versus dual Hahn: the primary source gives the
opposite orientation.}
KLS section~2.9 expressly lists \emph{dual Hahn $\to$ Meixner}, with
\[
 \lim_{N\to\infty}
 R_n\!\left(\lambda(x);\beta-1,
             \frac{N(1-c)}{c},N\right)
 =M_n(x;\beta,c),
\]
under the conventions of its definitions (section~2.9, displayed relation
following ``Dual Hahn $\to$ Meixner'').
This is not a KLS citation for the reverse ordered arrow
Meixner $\to$ dual Hahn.  V134 independently found a degree-filtration
obstruction to the naive unscaled reverse identity.
\emph{Disposition: EXCLUDE Meixner $\to$ dual Hahn from the
KLS-directed PASS ledger; the opposite KLS arrow must be assessed
separately if present among V31's headings.}

\paragraph{Jacobi versus continuous Hahn: another directionality check.}
KLS section~2.8 explicitly lists \emph{continuous Hahn $\to$ Jacobi},
using a large-parameter scaling of the continuous Hahn variable and
parameters.  It does not thereby authorize a reverse ordered
Jacobi $\to$ continuous Hahn arrow.  The Fourier-transform connection
noted in V128 is not a polynomial limiting arrow.
\emph{Disposition: EXCLUDE Jacobi $\to$ continuous Hahn under the
specified KLS-directed counting criterion.}

\paragraph{Methodological boundary and ledger.}
These are \emph{primary-source classification} decisions, not claims
that every conceivable analytic continuation or scaled limit is
impossible.  In particular, the Jacobi--Bessel confluent formula
remains mathematically valid and useful outside the KLS ledger.
No historical PASS entry is deleted by these three candidate
classifications, and no new PASS is awarded:
\[
 \boxed{\text{Recorded PASS}=54,\qquad
        \text{new PASS in V140}=0.}
\]
The outstanding candidates still require their own exact KLS
section-and-direction checks, and the complete 54-entry ledger has
not yet undergone an exhaustive independent primary-source audit.
Accordingly, \emph{global completeness remains unproved}.

\clearpage
\section{V141: Primary-source direction audit of Stieltjes--Wigert and classical borderline candidates}
\label{sec:v141-kls-directions}

\paragraph{Evidence and scope.}
This stage checks the original Koekoek--Swarttouw (1998) report,
\emph{The Askey-scheme of hypergeometric orthogonal polynomials and its
$q$-analogue}, directly at
\texttt{https://tudelft.roelofkoekoek.nl/documents/as98.pdf},
PDF pages 135--137 (sections 4.22, 4.23 and 4.27) and PDF page 57
(section 2.7).  The conclusions concern \emph{ordered arrows};
reversing a connection identity does not reverse a limiting arrow.

\paragraph{Stieltjes--Wigert to alternative $q$-Charlier: EXCLUDE.}
KLS section 4.22 explicitly identifies the arrow
\emph{alternative $q$-Charlier $\to$ Stieltjes--Wigert}, equation
(4.22.1), with the limit
\[
 \lim_{a\to\infty}K_n(x/a;a;q)=(q;q)_n S_n(x;q).
\]
Section 4.27 repeats the same ordered arrow.  Neither occurrence
supports the reverse arrow Stieltjes--Wigert $\to$ alternative
$q$-Charlier.  The obstruction proved in V130--V131 concerns the
fixed-variable unrenormalized identity specialization only; the
\emph{primary-source direction check} is the reason for exclusion
under the audit's KLS-directed counting rule.  This disposition does
not deny that V31 contains a reverse connection between the families.

\paragraph{Stieltjes--Wigert to ordinary $q$-Charlier: EXCLUDE.}
KLS section 4.23, equation (4.23.1), explicitly gives the opposite
ordered arrow:
\[
 \lim_{a\to\infty} C_n(ax;a;q)=(q;q)_n S_n(x;q).
\]
Section 4.27 repeats $q$-Charlier $\to$ Stieltjes--Wigert.
There is no independent KLS arrow in the reverse direction established
by either citation.  V132's diagonal obstruction independently rules
out the fixed-$q$, fixed-variable unnormalized identity path, but is
not used to assert impossibility of every conceivable rescaled limit.
\emph{Disposition: EXCLUDE the reverse ordered candidate from the
KLS-directed PASS ledger.}

\paragraph{Pseudo-Jacobi to continuous Hahn: EXCLUDE from this KLS ledger.}
The primary-source table of contents and classical scheme enumerate
continuous Hahn, Jacobi, Meixner--Pollaczek, Laguerre and related
families, but do not include pseudo-Jacobi as an independent family.
A search of the searchable report for the name ``Pseudo-Jacobi'' yields
no occurrence.  Thus no independent KLS arrow involving that named
family is authenticated.  This is a \emph{classification exclusion},
not a mathematical assertion that pseudo-Jacobi analytic
continuations or connection formulas are invalid.  V76's generic
basis expansion does not substitute for a KLS arrow.

\paragraph{Meixner--Pollaczek to generalized Laguerre: KLS YES;
V31 reverse certification still required.}
KLS section 2.7, equation (2.7.1), \emph{does} explicitly list
Meixner--Pollaczek $\to$ Laguerre, with
$\lambda=(\alpha+1)/2$, $x\mapsto-x/(2\phi)$ and $\phi\to0$:
\[
 \lim_{\phi\to0}
 P_n^{((\alpha+1)/2)}\!\left(-\frac{x}{2\phi};\phi\right)
 =L_n^{(\alpha)}(x).
\]
Here generalized Laguerre means the ordinary parameterized Laguerre
family $L_n^{(\alpha)}$, not a distinct additional KLS family.
The fixed-angle reverse-connection argument previously examined in
V128 does not verify this \emph{singular, variable-scaled} endpoint.
Before awarding any PASS, one must identify the relevant V31 ordered
heading, check whether this ordered pair is already in the 54-entry
ledger, and prove convergence of the correctly transformed reverse
coefficients.  This is now a \textbf{confirmed KLS arrow with a pending
V31 mathematical/duplicate audit}, rather than a KLS-eligibility unknown.

\paragraph{Cumulative decision ledger.}
This stage closes three previously open KLS classification questions:
two Stieltjes--Wigert reverse directions and pseudo-Jacobi to
continuous Hahn.  It also positively authenticates the classical
Meixner--Pollaczek $\to$ Laguerre direction without awarding a PASS.
\[
 \boxed{\text{Recorded PASS}=54,\qquad
 \text{new PASS in V141}=0.}
\]
The remaining tasks include the transformed V31 Meixner--Pollaczek
limit, historical 54-entry citation authentication, and a final
72-heading global completeness and duplicate reconciliation.
\emph{No global completeness claim is made in this stage.}


\clearpage
\section{V142: Meixner--Pollaczek to Laguerre---primary-source verification and scope of the V31 reverse connection}
\noindent\textbf{Audit question.} KLS (1998), Section 2.7, equation (2.7.1), recognizes the ordered arrow Meixner--Pollaczek $\to$ Laguerre. Does the actual \emph{angle-fixed} V31 reverse formula prove this singular arrow? We distinguish the polynomial limit from the particular connection formula; they are not interchangeable.

\paragraph{Exact source inventory.}
The unchanged V31 contains the heading \texttt{sec:ordinary-laguerre-mp-bell}, the two finite-connection labels \texttt{eq:ordinary-lag-to-mp} and \texttt{eq:ordinary-mp-to-lag}, and the independent direct-limit label \texttt{eq:mp-laguerre-direct-limit}. In the first two formulas V31 explicitly fixes the Meixner--Pollaczek angle at $\phi=\pi/2$ and uses $L_n^{(\lambda)}(x)$ with the \emph{same} parameter $\lambda$ as $P_n^{(\lambda)}(x;\pi/2)$. In contrast, the KLS arrow requires $\phi\to0$, $\lambda=(\alpha+1)/2$, and $x_{\mathrm{MP}}=-X/(2\phi)$; its target is $L_n^{(\alpha)}(X)=L_n^{(2\lambda-1)}(X)$. Thus the printed Bell-polynomial reverse coefficients cannot be evaluated at $\phi\to0$; that angle is not a free parameter in those coefficients. The V31 direct-limit equation does explicitly state the appropriate arrow, but it is not a finite reverse-connection coefficient formula.

\paragraph{Independent all-degree verification of the KLS polynomial limit.}
For $\lambda>0$, the standard Meixner--Pollaczek generating function is
\[
 \sum_{n\ge0}P_n^{(\lambda)}(u;\phi)t^n
 =(1-e^{i\phi}t)^{-\lambda+iu}
  (1-e^{-i\phi}t)^{-\lambda-iu},\qquad |t|<1.
\]
Set $u=-X/(2\phi)$, where $X$ is fixed. Taking the logarithm and expanding about $\phi=0$ gives, locally uniformly for $|t|\le r<1$,
\[
\begin{aligned}
 \log G_{\phi}(t)
 &=-\lambda\log(1-e^{i\phi}t)
   -\lambda\log(1-e^{-i\phi}t)\\
 &\quad-\frac{iX}{2\phi}
  \bigl[\log(1-e^{i\phi}t)-\log(1-e^{-i\phi}t)\bigr]\\
 &=-2\lambda\log(1-t)-\frac{Xt}{1-t}+O(\phi^2).
\end{aligned}
\]
The sign follows from $\log(1-e^{i\phi}t)-\log(1-e^{-i\phi}t)=-2i\phi t/(1-t)+O(\phi^3)$. Consequently
\[
 G_{\phi}(t)\longrightarrow
 (1-t)^{-2\lambda}\exp\!\left(-\frac{Xt}{1-t}\right)
 =\sum_{n\ge0}L_n^{(2\lambda-1)}(X)t^n.
\]
Uniform convergence on a circle $|t|=r<1$ and Cauchy's coefficient formula prove, for each fixed $n$,
\[
\boxed{\lim_{\phi\to0^+}P_n^{(\lambda)}\!\left(-\frac{X}{2\phi};\phi\right)
 =L_n^{(2\lambda-1)}(X).}
\]
This proves the KLS relation in every degree, independently of the V31 Bell connection. The error is $O(\phi^2)$ coefficientwise for fixed $n$, on compact $X$-sets, as long as the stated analytic branches are used near $t=0$.

\paragraph{Low-degree consistency.}
The exact degree-one polynomial is
$P_1^{(\lambda)}(u;\phi)=2u\sin\phi+2\lambda\cos\phi$.
After $u=-X/(2\phi)$, its limit is $2\lambda-X=L_1^{(2\lambda-1)}(X)$. This also checks the variable sign and the target parameter.

\paragraph{Reverse-connection certification boundary.}
The generating-function argument supplies an all-degree \emph{polynomial} limit and verifies KLS eligibility. It does \emph{not} turn V31's fixed-$\pi/2$, same-$\lambda$ Bell coefficients into a $\phi$-dependent reverse matrix. In particular, the target Laguerre parameter must change from $\lambda$ to $2\lambda-1$ and the variable must be rescaled. One could construct a new angle-dependent, transformed Laguerre-basis expansion and prove its matrix tends to the identity, but that would be an additional derivation, not a verification of the coefficients printed in \texttt{eq:ordinary-mp-to-lag}. Under the audit's requirement that the actual V31 reverse coefficients realize the arrow, \textbf{no new PASS is awarded at V142}. This is a precise evidence limitation, not a rejection of the KLS arrow itself.

\paragraph{Ledger decision.}
\[\boxed{\text{KLS arrow verified; polynomial limit verified; V31 printed reverse realization unresolved.}}\]
\[\boxed{\text{Recorded PASS}=54,\qquad\text{new PASS in V142}=0.}\]
The next step is either to certify an admissible transformed reverse connection from V31's definitions with a clear distinction from its printed fixed-angle Bell formula, or to mark this arrow as outside the finite-coefficient scope. Global 72-heading completeness and remaining primary-source checks are still pending. The original V31 is unchanged.



\section{V143: Explicit angle-dependent transformed reverse rows and the scope of the printed V31 connection}
\label{sec:v143-mp-lag-angle}

\paragraph{Purpose and distinction.}
V142 proved the KLS Meixner--Pollaczek-to-Laguerre polynomial limit but observed that V31's printed Bell-polynomial connection \texttt{eq:ordinary-mp-to-lag} is restricted to angle $\phi=\pi/2$ and uses the same parameter in both families. We now calculate the first two \emph{transformed} reverse rows exactly at arbitrary $\phi$, with the KLS parameter matching. These rows are independently derived from the Meixner--Pollaczek recurrence, \emph{not} asserted to be coefficients of V31's printed fixed-angle formula.

\paragraph{Definitions.}
Fix $\lambda$ with $2\lambda\notin\{0,-1,-2,\ldots\}$, set $\alpha=2\lambda-1$, and write
\[
 F_n(X;\phi)=P_n^{(\lambda)}\!\left(-\frac{X}{2\phi};\phi\right),\qquad
 t=\frac{\sin\phi}{\phi},\quad c=\cos\phi,
 \qquad \ell_j(X)=L_j^{(2\lambda-1)}(X).
\]
The degree-one formula and three-term recurrence give the \emph{exact} identities
\[
 F_0=1,\qquad F_1=2\lambda c-tX,
\]
\[
 F_2=\frac{t^2X^2}{2}-(2\lambda+1)ctX
 +\lambda\bigl(2(\lambda+1)c^2-1\bigr).
\]
In the target Laguerre basis,
\[
 \ell_0=1,\quad \ell_1=2\lambda-X,\quad
 \ell_2=\frac{X^2}{2}-(2\lambda+1)X+\lambda(2\lambda+1).
\]
Coefficient comparison therefore yields the following complete rows of the uniquely defined transformed connection $F_n=\sum_{k=0}^n E_{nk}(\phi)\ell_k$:
\begin{align*}
&E_{00}=1,\\
&E_{11}=t,\qquad E_{10}=2\lambda(c-t),\\
&E_{22}=t^2,\qquad E_{21}=(2\lambda+1)(ct-t^2),\\
&E_{20}=\lambda\bigl(2(\lambda+1)c^2-1\bigr)
       -\lambda(2\lambda+1)t^2-2\lambda(2\lambda+1)(ct-t^2).
\end{align*}
All these expressions are exact for nonzero $\phi$ near zero. Since
\[
 c=1-\frac{\phi^2}{2}+O(\phi^4),\qquad
 t=1-\frac{\phi^2}{6}+O(\phi^4),
\]
we obtain $E_{nk}(\phi)=\delta_{nk}+O(\phi^2)$ for $n\le2$. In particular, the first off-diagonal entry is
\[
 E_{10}= -\frac{2\lambda}{3}\phi^2+O(\phi^4),
 \qquad E_{21}=-\frac{2(2\lambda+1)}{3}\phi^2+O(\phi^4).
\]
This gives a nontrivial check of the KLS parameter map and the \emph{full} low-degree transformed matrix, rather than just its diagonal.

\paragraph{All-degree analyticity and matrix convergence.}
For each fixed $n$, the recurrence
\[
 (n+1)P_{n+1}^{(\lambda)}(u;\phi)
 =2\bigl((n+\lambda)\cos\phi+u\sin\phi\bigr)P_n^{(\lambda)}(u;\phi)
 -(n+2\lambda-1)P_{n-1}^{(\lambda)}(u;\phi)
\]
shows after substituting $u=-X/(2\phi)$ that $F_n(X;\phi)$ is polynomial in $X$ with coefficients analytic and \emph{even} in $\phi$ near zero. The limit recurrence is the Laguerre recurrence with parameter $2\lambda-1$, so $F_n\to\ell_n$ coefficientwise. Because the fixed Laguerre basis has nonzero leading coefficients $(-1)^k/k!$, the change-of-basis matrix is triangular and invertible in every finite degree. Its coefficients $E_{nk}(\phi)$ are consequently analytic even functions of $\phi$ and satisfy
\[
 \boxed{E_{nk}(\phi)=\delta_{nk}+O(\phi^2)\quad(0\le k\le n,\ n\text{ fixed}).}
\]
This establishes an \emph{independently constructed} transformed reverse connection at every degree, but it does not alter the angle restriction of the coefficients actually printed in V31.

\paragraph{Counting verdict and remaining evidence.}
The original KLS arrow and its transformed polynomial and matrix limits are mathematically established. Nevertheless, the V31 Bell coefficients at $\phi=\pi/2$ cannot be varied to $\phi\to0$; identifying the newly derived $E_{nk}(\phi)$ with those printed coefficients would be incorrect. Under the strict ``printed V31 reverse-coefficient realization'' criterion retained in V142, this arrow remains \textbf{unawarded}; it may be reconsidered only if the audit formally permits reverse matrices newly constructed from V31's polynomial definitions. No new PASS is awarded here, and the recorded count remains \textbf{54}. Primary-source cross-checking and global completeness certification remain outstanding. The original V31 file has not been changed.



\clearpage
\section{V144: Completion-gate decision on the Meixner--Pollaczek arrow and candidate-disposition reconciliation}
\noindent\textbf{Purpose and counting rule.} This stage settles the outstanding \emph{classification of the evidence} for Meixner--Pollaczek $\to$ Laguerre and consolidates the previously adjudicated candidate directions. The cumulative audit distinguishes (i) an original KLS directed arrow, (ii) a proved polynomial limit, and (iii) a limit of the \emph{printed V31 finite reverse coefficients}. These are separate predicates. The standing rule for a PASS in this reverse-connection audit requires all three, including any justified change of variables and normalization; constructing a new connection from the defining polynomials alone does not automatically establish (iii).

\paragraph{Direct inspection of the unchanged V31 source.}
The printed reverse formula \texttt{eq:ordinary-mp-to-lag} is, literally,
\[
 P_n^{(\lambda)}(x;\pi/2)=\sum_{m=0}^n B_{nm}(\lambda)
 L_m^{(\lambda)}(x),
\]
where $B_{nm}(\lambda)$ is a finite Bell-polynomial expression independent of any variable angle $\phi$. The companion \texttt{eq:ordinary-lag-to-mp} is likewise fixed at $\phi=\pi/2$. Immediately after the proposition, V31 separately prints \texttt{eq:mp-laguerre-direct-limit}:
\[
 \lim_{\phi\to0^+}P_n^{((\alpha+1)/2)}
 \left(-\frac{x}{2\phi};\phi\right)=L_n^{(\alpha)}(x).
\]
The V31 manuscript itself explicitly calls this a ``separate, non-finite-parameter specialization'' and says it ``is not an additional finite connection formula.'' The first two formulas therefore have \emph{no $\phi\to0$ parameter to take}. A change of the polynomial argument cannot alter their fixed angle or automatically change the target parameter from $\lambda$ to $2\lambda-1$.

\paragraph{What V142--V143 do and do not prove.}
Let $F_n(X;\phi)=P_n^{(\lambda)}(-X/(2\phi);\phi)$ and $\alpha=2\lambda-1$. The independent expansion
\[
 F_n(X;\phi)=\sum_{k=0}^n E_{nk}(\phi)L_k^{(2\lambda-1)}(X)
\]
exists uniquely for each finite degree. V143 proved $E_{nk}(\phi)=\delta_{nk}+O(\phi^2)$, with exact low-degree checks. This verifies the KLS polynomial limit and a \emph{newly constructed} transformed connection. It does not show that $E_{nk}(\phi)$ is obtained by specializing the printed $B_{nm}(\lambda)$, because those are different coefficient arrays with different parameters and angles. Consequently, this ordered pair is marked \textbf{KLS-confirmed / polynomial-limit verified / printed-reverse-realization NOT ESTABLISHED}. It is \textbf{not a PASS} under the standing strict criterion. This is an evidentiary disposition, not a claim that the KLS arrow fails mathematically.

\paragraph{Reconciled status of the nine V128 candidate directions.}
The following table reports the \emph{ordered} candidates and the strongest decision established in V128--V143; ``excluded'' means excluded from this KLS-directed count, not that every conceivable polynomial identity is impossible.
\begin{center}\small
\begin{tabular}{p{0.40\linewidth}p{0.49\linewidth}}
\hline
Candidate direction & Current disposition and evidence \\
\hline
Big $q$-Laguerre $\to$ Al-Salam--Carlitz I & \textbf{PASS 54}; KLS recognition and corrected scaled V31 reverse matrix, V137--V139.\\
Stieltjes--Wigert $\to$ alternative $q$-Charlier & \textbf{EXCLUDE}; original KLS arrow runs in the opposite direction; fixed-variable obstruction, V130--V131 and V141.\\
Stieltjes--Wigert $\to$ ordinary $q$-Charlier & \textbf{EXCLUDE}; original KLS arrow runs in the opposite direction; diagonal obstruction, V132 and V141.\\
Jacobi $\to$ continuous Hahn & \textbf{EXCLUDE}; KLS recognizes continuous Hahn $\to$ Jacobi, V140.\\
Meixner--Pollaczek $\to$ Laguerre & \textbf{UNAWARDED / DIRECT-ONLY}; KLS and transformed polynomial limit verified, but printed V31 reverse coefficients are fixed-angle, V142--V144.\\
Meixner $\to$ dual Hahn & \textbf{EXCLUDE}; KLS recognizes dual Hahn $\to$ Meixner; degree and finite-grid caveats, V134 and V140.\\
Pseudo-Jacobi $\to$ continuous Hahn & \textbf{EXCLUDE}; pseudo-Jacobi not an independent original KLS family, V141.\\
Generalized Bessel / Jacobi candidate & \textbf{EXCLUDE from KLS count}; a Jacobi $\to$ generalized Bessel limit is proved, but the target is outside the original KLS family list, V133--V136 and V140.\\
Al-Salam--Chihara $\to$ continuous $q$-Laguerre & \textbf{PASS 53}; exact normalized specialization and independent KLS recognition, V137--V138.\\
\hline
\end{tabular}
\end{center}
\noindent The V128 list contained a fixed-angle Meixner--Pollaczek/Laguerre candidate; its present status must not be mistaken for a negative KLS finding. Likewise, a mathematical Jacobi/generalized-Bessel limit is not counted merely because it is correct. The nine rows are \emph{candidate-direction dispositions}, not a partition of the 72 V31 headings or of the 54 PASS arrows.

\paragraph{Next certification gates.}
\begin{enumerate}
\item Reconstruct the complete \textbf{54-row ordered PASS ledger}, including the two V138--V139 additions, and authenticate the specific original KLS location and direction for \emph{each} row. V127's 52/52 unique-equation-label check proves source traceability, not KLS authentication; the two newer rows require inclusion.
\item Compare that ledger with the complete \textbf{72-heading V31 inventory} and document every eligible KLS arrow not yet counted, including whether the actual V31 printed reverse coefficient can realize it. Heading count and PASS count are different quantities, so $72-54$ is not a justified number of missing arrows.
\item Audit the proof obligations and normalization conventions for the surviving 54 entries, and only then issue an exhaustive certified count. If the strict printed-coefficient criterion is to be broadened, that is a \emph{change of audit scope} requiring an explicit separate ledger, not an automatic extra PASS.
\end{enumerate}
\paragraph{V144 verdict.} No new PASS is awarded. The \textbf{recorded total remains 54}, with one confirmed KLS/direct-limit relation deliberately unawarded under the printed-V31-reverse criterion. This section closes the candidate \emph{triage} but does \emph{not} certify the global completeness or every historical KLS reference. The original V31 manuscript remains unchanged.


\clearpage
\section{V145: Complete 54-entry source ledger and primary-source authentication boundary}
\label{sec:v145-full-ledger}
\noindent\textbf{Scope.} This stage reconstructs the complete numbered PASS ledger after the V138--V139 awards and performs a fresh literal check of each cited equation label against the \emph{unchanged} V31 source. It is an independent mechanical source-identity test, \textbf{not} a claim to have verified every arrow against the original KLS publication. That stronger gate requires a row-by-row comparison of the published direction, parameter map, variable transformation, and normalization; a unique V31 label alone does not establish it.

\paragraph{Full ledger and fresh source checks.} The fourth column is the number of exact \texttt{\\label\{...\}} definitions in V31. Every row is an ordered source-to-target relation; the final two rows were not part of the earlier V127 52-row table.
\begingroup\scriptsize
\begin{longtable}{@{}r p{0.43\linewidth} p{0.39\linewidth} r@{}}
\toprule No. & Recorded ordered transition & V31 reverse/connection equation label & Occ.\\
\midrule\endfirsthead
\toprule No. & Recorded transition & V31 equation label & Occ.\\
\midrule\endhead
\bottomrule\endfoot
1 & Symmetric Jacobi $\ensuremath{\to}$ Hermite & \texttt{eq:v14-jac-herm} & 0 \\
2 & Normalized continuous $q$-Laguerre $\ensuremath{\to}$ continuous $q$-Hermite & \texttt{eq:qlag-hermite-fwd} & 0 \\
3 & Continuous $q$-Jacobi $\ensuremath{\to}$ continuous $q$-Hermite & \texttt{eq:JHforward} & 0 \\
4 & Hahn $\ensuremath{\to}$ Jacobi & \texttt{eq:jacobi-hahn-reverse} & 0 \\
5 & $q$-Hahn $\ensuremath{\to}$ little $q$-Jacobi & \texttt{eq:ljqh-reverse} & 0 \\
6 & $q$-Hahn $\ensuremath{\to}$ $q$-Meixner & \texttt{eq:168-D} & 0 \\
7 & Big $q$-Jacobi $\ensuremath{\to}$ $q$-Meixner & \texttt{eq:169-reverse} & 0 \\
8 & Big $q$-Jacobi $\ensuremath{\to}$ little $q$-Jacobi & \texttt{eq:ljbj166-reverse} & 0 \\
9 & Big $q$-Laguerre $\ensuremath{\to}$ little $q$-Laguerre (Wall) & \texttt{eq:wall-big-reverse} & 0 \\
10 & $q$-Hahn $\ensuremath{\to}$ quantum $q$-Krawtchouk & \texttt{eq:finite-qtmH-reverse} & 0 \\
11 & $q$-Hahn $\ensuremath{\to}$ ordinary $q$-Krawtchouk & \texttt{eq:finite-KH-reverse} & 0 \\
12 & Dual $q$-Hahn $\ensuremath{\to}$ dual $q$-Krawtchouk & \texttt{eq:finite-dKdH-reverse} & 0 \\
13 & Askey--Wilson $\ensuremath{\to}$ big $q$-Jacobi & \texttt{eq:aw-bigqj-reverse} & 0 \\
14 & Standard $q$-Krawtchouk $\ensuremath{\to}$ alternative $q$-Charlier ($q$-Bessel) & \texttt{eq:qBqK-reverse} & 0 \\
15 & Standard $q$-Krawtchouk $\ensuremath{\to}$ ordinary $q$-Charlier & \texttt{eq:qCqK-reverse} & 0 \\
16 & Little $q$-Jacobi $\ensuremath{\to}$ $q$-Bessel (alternative $q$-Charlier) & \texttt{eq:qBlqJ-reverse} & 0 \\
17 & $q$-Meixner $\ensuremath{\to}$ $q$-Laguerre & \texttt{eq:qLqM-D} & 0 \\
18 & Little $q$-Jacobi $\ensuremath{\to}$ $q$-Laguerre & \texttt{eq:qLlJ-D} & 0 \\
19 & Affine $q$-Krawtchouk $\ensuremath{\to}$ little $q$-Laguerre (Wall) & \texttt{eq:wall-affine-reverse} & 0 \\
20 & Little $q$-Jacobi $\ensuremath{\to}$ little $q$-Laguerre (Wall) & \texttt{eq:wall-lj-reverse} & 0 \\
21 & Big $q$-Jacobi $\ensuremath{\to}$ big $q$-Laguerre & \texttt{eq:blbj-reverse} & 0 \\
22 & Continuous $q$-Hahn $\ensuremath{\to}$ $q$-Meixner--Pollaczek & \texttt{eq:cqh-qmp23} & 0 \\
23 & Al-Salam--Chihara $\ensuremath{\to}$ continuous big $q$-Hermite & \texttt{eq:ASCBH} & 0 \\
24 & Continuous dual $q$-Hahn $\ensuremath{\to}$ Al-Salam--Chihara & \texttt{eq:CDASC} & 0 \\
25 & Continuous big $q$-Hermite $\ensuremath{\to}$ continuous $q$-Hermite & \texttt{eq:BHH} & 0 \\
26 & $q$-Laguerre $\ensuremath{\to}$ Stieltjes--Wigert & \texttt{eq:SWqL-rev} & 0 \\
27 & Quantum $q$-Krawtchouk $\ensuremath{\to}$ Al-Salam--Carlitz II & \texttt{eq:ASCII-qtm-reverse} & 0 \\
28 & $q$-Meixner $\ensuremath{\to}$ Al-Salam--Carlitz II & \texttt{eq:ASCII-qM-reverse} & 0 \\
29 & Askey--Wilson $\ensuremath{\to}$ continuous dual $q$-Hahn & \texttt{eq:AWCD} & 0 \\
30 & Racah $\ensuremath{\to}$ Hahn & \texttt{eq:racah-to-hahn-grid} & 0 \\
31 & Racah $\ensuremath{\to}$ dual Hahn & \texttt{eq:racah-dh-reverse} & 0 \\
32 & Krawtchouk $\ensuremath{\to}$ Hermite & \texttt{eq:v16-k-h} & 0 \\
33 & Hahn $\ensuremath{\to}$ Meixner & \texttt{eq:hahn-meixner-reverse} & 0 \\
34 & Hahn $\ensuremath{\to}$ Krawtchouk & \texttt{eq:hahn-kraw-reverse} & 0 \\
35 & Krawtchouk $\ensuremath{\to}$ Charlier & \texttt{eq:kraw-char-reverse} & 0 \\
36 & Meixner $\ensuremath{\to}$ Charlier & \texttt{eq:meix-char-reverse} & 0 \\
37 & Charlier $\to$ Hermite & \texttt{eq:v16-c-h} & 0 \\
38 & Wilson $\to$ continuous dual Hahn & \texttt{eq:consolidated-w-cdh} & 0 \\
39 & Meixner--Pollaczek $\to$ Hermite & \texttt{eq:cat-mp-to-hermite-bell} & 0 \\
40 & Continuous Hahn $\to$ Meixner--Pollaczek & \texttt{eq:consolidated-ch-mp} & 0 \\
41 & Meixner $\to$ generalized Laguerre & \texttt{eq:meixner-to-laguerre-v61} & 0 \\
42 & Jacobi $\to$ generalized Laguerre & \texttt{eq:v14-jac-lag} & 0 \\
43 & Wilson $\to$ Jacobi & \texttt{eq:wilson-to-jac} & 0 \\
44 & q-Meixner $\to$ ordinary q-Charlier & \texttt{eq:qC-qM-reverse} & 0 \\
45 & dual q-Hahn $\to$ affine q-Krawtchouk & \texttt{eq:finite-affdH-reverse} & 0 \\
46 & q-Hahn $\to$ affine q-Krawtchouk & \texttt{eq:finite-affH-reverse} & 0 \\
47 & q-Racah $\to$ dual q-Hahn & \texttt{eq:finite-dHR-reverse} & 0 \\
48 & q-Racah $\to$ ordinary q-Krawtchouk & \texttt{eq:finite-KR-reverse} & 0 \\
49 & q-Racah $\to$ q-Hahn & \texttt{eq:finite-HR-double-D} & 0 \\
50 & continuous q-Jacobi $\to$ continuous q-Laguerre & \texttt{eq:lag-jacobi-new-rev} & 0 \\
51 & Al-Salam--Carlitz I $\to$ discrete q-Hermite I & \texttt{eq:dHI-carlitz-rev} & 0 \\
52 & Al-Salam--Carlitz II $\to$ discrete q-Hermite II & \texttt{eq:dHII-carlitz-rev} & 0 \\
53 & Al-Salam--Chihara $\to$ continuous $q$-Laguerre & \texttt{eq:lag-asc-hyper-rev51} & 0 \\
54 & Big $q$-Laguerre $\to$ Al-Salam--Carlitz I & \texttt{eq:ASCI-bL-reverse} & 0 \\
\end{longtable}\endgroup

\paragraph{Mechanical verdict.} The ledger contains \textbf{54 consecutive numbers}, 54 distinct label strings, and \textbf{54/54} labels each defined exactly once in V31. There are zero missing or multiply defined labels. This extends the V127 52-entry source check to the two later awards; it does not recertify the original KLS citation or the analytic correctness of the coefficients at those labels.

\paragraph{A material numbering inconsistency in the cumulative prose.} V127's explicit reconstructed ledger identifies \textbf{PASS 1} as symmetric Jacobi $\to$ Hermite at \texttt{eq:v14-jac-herm}, and \textbf{PASS 43} as Wilson $\to$ Jacobi at \texttt{eq:wilson-to-jac}. However, V126's heading and verdict describe its independent Jacobi $\to$ Hermite proof as ``historical PASS 43.'' These statements \emph{cannot both be valid ledger identifiers}. The differential-equation proof in V126 can support the Jacobi--Hermite \emph{relation}, but under the explicit V127 numbering it is an independent check of \textbf{PASS 1}, not PASS 43. No new PASS or deletion follows from this editorial-numbering correction. A final certified report must normalize all references to the 54-row ledger above and distinguish the original V71 historical chronological numbering from the later reconstructed ledger numbering, if they differ.

\paragraph{Primary-source authentication still required.} The original KLS 1998 report was consulted for selected arrows in V137--V144, but those selective checks do not establish 54 independent primary-source matches. For a valid exhaustive KLS certificate, each row needs an independently checked published location, directed arrow, parameter map, normalization, and any exceptional convention. In particular, the V139 corrected factor in the Big $q$-Laguerre limit must remain explicitly distinguished from the printed KLS normalization. No unverifiable equation numbers are manufactured here.

\paragraph{V145 decision.} \textbf{54/54 source labels verified; 54 recorded PASS entries retained; new PASS 0.} A numbering inconsistency has been identified for correction in the final report. \textbf{KLS authentication for all 54, mathematical recertification of historical entries, and global 72-heading exhaustiveness remain open.} The original V31 manuscript has not been modified.


\clearpage
\section{V146: Global source-heading crosswalk and correction of the V145 occurrence column}
\label{sec:v146-crosswalk}
\noindent\textbf{Scope and reproducible method.} This stage independently parses the unchanged V31 source and the 54-row ledger printed in V145. For each ledger label, it finds the exact \verb|\label{...}| occurrence, records its physical V31 line number, and assigns it to the most recent V31 \verb|\subsubsection| heading. This is a \emph{source-location crosswalk}, not a mathematical or KLS arrow certificate. A heading may support more than one ordered limiting arrow, and the source heading's displayed family order need not equal the limiting direction.

\paragraph{Material erratum to V145.} The fourth column of V145's printed 54-row table says ``0'' in \emph{every row}, notwithstanding the accompanying claim that each label appears once. This is a table-generation error, not evidence that the labels are missing. A fresh literal source search in this stage returns \textbf{exactly one occurrence for each of the 54 labels}. Thus the V145 narrative assertion is supported by a new check, but \textbf{its displayed occurrence column is wrong and must not be cited as accurate}. The following table supplies the independently recomputed line numbers and actual occurrence counts.

\begingroup\scriptsize\setlength{\tabcolsep}{4pt}
\begin{longtable}{@{}r l r r p{0.35\linewidth}@{}}
\toprule PASS & V31 label & V31 line & Occ. & Containing heading (V31 line)\\
\midrule\endfirsthead
\toprule PASS & V31 label & V31 line & Occ. & Containing heading (V31 line)\\
\midrule\endhead\bottomrule\endfoot
38 & \texttt{eq:consolidated-w-cdh} & 188 & 1 & heading at line 176 \\
40 & \texttt{eq:consolidated-ch-mp} & 443 & 1 & heading at line 426 \\
43 & \texttt{eq:wilson-to-jac} & 876 & 1 & heading at line 822 \\
42 & \texttt{eq:v14-jac-lag} & 1320 & 1 & heading at line 1296 \\
39 & \texttt{eq:cat-mp-to-hermite-bell} & 1483 & 1 & heading at line 1454 \\
1 & \texttt{eq:v14-jac-herm} & 1549 & 1 & heading at line 1527 \\
30 & \texttt{eq:racah-to-hahn-grid} & 1765 & 1 & heading at line 1731 \\
31 & \texttt{eq:racah-dh-reverse} & 1863 & 1 & heading at line 1837 \\
33 & \texttt{eq:hahn-meixner-reverse} & 1998 & 1 & heading at line 1978 \\
34 & \texttt{eq:hahn-kraw-reverse} & 2161 & 1 & heading at line 2154 \\
36 & \texttt{eq:meix-char-reverse} & 2373 & 1 & heading at line 2362 \\
35 & \texttt{eq:kraw-char-reverse} & 2408 & 1 & heading at line 2397 \\
4 & \texttt{eq:jacobi-hahn-reverse} & 2610 & 1 & heading at line 2579 \\
41 & \texttt{eq:meixner-to-laguerre-v61} & 2769 & 1 & heading at line 2745 \\
32 & \texttt{eq:v16-k-h} & 2838 & 1 & heading at line 2817 \\
37 & \texttt{eq:v16-c-h} & 2891 & 1 & heading at line 2874 \\
29 & \texttt{eq:AWCD} & 3171 & 1 & heading at line 3153 \\
24 & \texttt{eq:CDASC} & 3463 & 1 & heading at line 3441 \\
22 & \texttt{eq:cqh-qmp23} & 3606 & 1 & heading at line 3582 \\
23 & \texttt{eq:ASCBH} & 3843 & 1 & heading at line 3829 \\
53 & \texttt{eq:lag-asc-hyper-rev51} & 4013 & 1 & heading at line 3991 \\
50 & \texttt{eq:lag-jacobi-new-rev} & 4168 & 1 & heading at line 4145 \\
25 & \texttt{eq:BHH} & 4332 & 1 & heading at line 4315 \\
2 & \texttt{eq:qlag-hermite-fwd} & 4392 & 1 & heading at line 4366 \\
3 & \texttt{eq:JHforward} & 4486 & 1 & heading at line 4469 \\
49 & \texttt{eq:finite-HR-double-D} & 5429 & 1 & heading at line 5411 \\
47 & \texttt{eq:finite-dHR-reverse} & 5639 & 1 & heading at line 5620 \\
21 & \texttt{eq:blbj-reverse} & 5795 & 1 & heading at line 5783 \\
8 & \texttt{eq:ljbj166-reverse} & 5941 & 1 & heading at line 5923 \\
5 & \texttt{eq:ljqh-reverse} & 6030 & 1 & heading at line 6012 \\
7 & \texttt{eq:169-reverse} & 6161 & 1 & heading at line 6144 \\
6 & \texttt{eq:168-D} & 6242 & 1 & heading at line 6226 \\
10 & \texttt{eq:finite-qtmH-reverse} & 6405 & 1 & heading at line 6385 \\
11 & \texttt{eq:finite-KH-reverse} & 6595 & 1 & heading at line 6560 \\
48 & \texttt{eq:finite-KR-reverse} & 6691 & 1 & heading at line 6654 \\
46 & \texttt{eq:finite-affH-reverse} & 6832 & 1 & heading at line 6811 \\
45 & \texttt{eq:finite-affdH-reverse} & 6900 & 1 & heading at line 6878 \\
12 & \texttt{eq:finite-dKdH-reverse} & 7034 & 1 & heading at line 7011 \\
9 & \texttt{eq:wall-big-reverse} & 7180 & 1 & heading at line 7149 \\
20 & \texttt{eq:wall-lj-reverse} & 7262 & 1 & heading at line 7243 \\
19 & \texttt{eq:wall-affine-reverse} & 7362 & 1 & heading at line 7337 \\
18 & \texttt{eq:qLlJ-D} & 7607 & 1 & heading at line 7584 \\
17 & \texttt{eq:qLqM-D} & 7669 & 1 & heading at line 7648 \\
16 & \texttt{eq:qBlqJ-reverse} & 7799 & 1 & heading at line 7787 \\
14 & \texttt{eq:qBqK-reverse} & 7864 & 1 & heading at line 7850 \\
44 & \texttt{eq:qC-qM-reverse} & 7973 & 1 & heading at line 7962 \\
15 & \texttt{eq:qCqK-reverse} & 8055 & 1 & heading at line 8039 \\
54 & \texttt{eq:ASCI-bL-reverse} & 8186 & 1 & heading at line 8171 \\
28 & \texttt{eq:ASCII-qM-reverse} & 8406 & 1 & heading at line 8389 \\
27 & \texttt{eq:ASCII-qtm-reverse} & 8503 & 1 & heading at line 8491 \\
26 & \texttt{eq:SWqL-rev} & 8640 & 1 & heading at line 8623 \\
51 & \texttt{eq:dHI-carlitz-rev} & 8997 & 1 & heading at line 8985 \\
52 & \texttt{eq:dHII-carlitz-rev} & 9113 & 1 & heading at line 9093 \\
13 & \texttt{eq:aw-bigqj-reverse} & 9513 & 1 & heading at line 9479 \\
\end{longtable}\endgroup
\paragraph{Heading coverage: what the arithmetic can and cannot prove.} The source contains 117 literal \verb|\subsubsection| headings. Exactly 54 distinct headings contain at least one of the 54 ledger equation labels, while 63 headings contain none. These are \emph{physical heading} counts, including self-family and auxiliary headings. They must not be equated with the previously reported 72 distinct-family inventory entries. In particular, a ledger label may lie in a section different from the intuitive named pair, and a heading without a ledger label is not automatically an omitted eligible arrow.

\paragraph{Previously identified inventory ambiguity.} The V90 census includes an auxiliary heading ``Parameter cancellation: one terminating 2F1 in each direction,'' which is not itself a named pair of distinct polynomial families. The V90 table actually marks 73 entries as ``KLS/record check required, including this auxiliary heading; excluding that one leaves 72 candidate distinct-family headings. The 72-heading denominator is therefore not a proof that there are 72 independent eligible KLS arrows. Counting heading occurrences, distinct family pairs, and directed published limits are three separate operations.

\paragraph{Noncovered source-heading locations.} The following 63 heading line numbers contain no label from the 54-entry ledger. They include same-family and auxiliary headings and do not by themselves establish missing eligible arrows.\par
\begingroup\small\noindent 106, 229, 267, 355, 509, 609, 701, 975, 1074, 1211, 1406, 1582, 1622, 1700, 1786, 1908, 1927, 2027, 2123, 2236, 2266, 2298, 2432, 2459, 2505, 2659, 2977, 3204, 3324, 3386, 3493, 3637, 3701, 3773, 3943, 4060, 4230, 5014, 5078, 5230, 5345, 5527, 5683, 5856, 6076, 6291, 6488, 6745, 6951, 7084, 7500, 7719, 7925, 8117, 8247, 8336, 8566, 8705, 8807, 8925, 9044, 9167, 9287.\par\endgroup

\paragraph{Certification verdict.} \textbf{54/54 ledger labels occur exactly once in V31; V145's displayed zero-occurrence column is corrected; no new PASS and no deletion.} This source-heading crosswalk does not establish that the 54 entries are KLS-recognized, that all analytic limits are valid, or that no additional eligible arrows exist. A definitive global certificate still requires a verified directed KLS-to-V31 crosswalk, including parameter and normalization checks for each arrow. The original V31 manuscript is unchanged.


\clearpage
\section{V147: Certification decision, evidence ledger, and unresolved obligations}
\label{sec:v147-certification-boundary}
\noindent\textbf{Purpose.} This stage answers whether the cumulative 54-entry ledger can presently be signed as a \emph{complete, independently verified KLS-directed reverse-connection audit}. The answer is \textbf{no}. We distinguish what has actually been checked from what remains to be proved, rather than promoting source-label traceability to primary-source or mathematical certification.

\paragraph{Reproducible independent source-identity test.} A fresh extraction of the V145 numbered ledger (not the V145 occurrence column) yields exactly 54 rows, numbered 1--54. All 54 equation-label strings are distinct, and an independent literal search of the unchanged V31 source finds \emph{exactly one} definition for each label. This confirms the V146 correction to V145's erroneous displayed zero counts. These tests establish source traceability only: a correct equation label does not demonstrate a correct KLS direction, scaling, or limiting proof.

\paragraph{Three distinct logical certification gates.} For each recorded directed pair $(F,G)$, a complete certificate requires (K) an authenticated original KLS source that explicitly recognizes the directed arrow $F\to G$ with its conventions; (M) a correct all-degree polynomial limit, including parameter path, argument scaling, and normalization; and (R) a demonstrated limiting realization of the \emph{printed V31 reverse connection} after any justified transformations. Source-identity verification (S) is a separate mechanical check. A row is globally certified only when $K\wedge M\wedge R\wedge S$ is established. The current ledger cannot be described as 54 fully certified rows merely because S holds 54 times.

\paragraph{A concrete unsettled arrow and the scope of the rule.} KLS recognizes Meixner--Pollaczek $\to$ Laguerre and V142--V143 prove the transformed polynomial limit and the convergence of a newly constructed angle-dependent connection. The coefficients \emph{printed} in V31 fix $\phi=\pi/2$, whereas the KLS arrow requires $\phi\to0$. Under the established strict R gate this is \textbf{not a new PASS}. This is an evidentiary boundary, not a counterexample to the KLS polynomial limit. Relaxing R would change the question being audited and must not be done silently.

\paragraph{A concrete source-convention warning.} V137--V139 identify a normalization mismatch in the printed KLS Big $q$-Laguerre $\to$ Al-Salam--Carlitz I limit. The corrected V31-normalized expression includes the factor $(aq)^{-n}$, not merely $a^{-n}$. PASS 54 relies on this explicitly documented correction. Any final primary-source crosswalk must record such convention differences, not mark the original printed formula as literally identical to V31.

\paragraph{Numbering and retrospective claims.} The authoritative reconstructed ledger calls Jacobi $\to$ Hermite PASS~1 and Wilson $\to$ Jacobi PASS~43. The V126 prose describing Jacobi $\to$ Hermite as PASS~43 is a historical \emph{indexing error}; the V126 proof is evidence about PASS~1. Likewise the V145 table's zero occurrence values were a reporting defect corrected in V146. Neither issue creates or removes an ordered limiting arrow.

\paragraph{What a final global certificate must still supply.} (i) An explicit row-by-row primary-source KLS crosswalk for all 54 recorded arrows, including exact original locations and direction, with no inferred citations; (ii) a mathematical proof/evidence pointer for each row with its variable and polynomial normalization; (iii) a complete directed-pair review of the 72 candidate distinct-family V31 headings, accounting for pairs with multiple possible arrows and for connections supporting none; and (iv) final reconciliation of all exclusions and the strict printed-reverse criterion. The 117 physical headings and 72 candidate distinct-family headings are not counts of eligible KLS arrows. The nine candidate decisions summarized in V144 are not a proof that no further arrows were overlooked.

\paragraph{Decision and release status.} \textbf{V147 is a final-status assessment, not an unconditional completion certificate.} The source-backed recorded count remains \textbf{54}; no new PASS is awarded and none is rescinded. A definitive exhaustive count and full primary-source/mathematical certification remain \textbf{open}. A report claiming that every one of the 54 arrows has been independently authenticated against KLS would exceed the evidence currently in this cumulative audit. The original V31 manuscript has not been modified.

\end{document}
